\section{An explicit secant plane in every dimension}\label{sec:construction}

\Cref{q:higher} asks for four things: an abelian $n$-fold $X$, a plane
$P\subset H^{\mathrm{ev}}(X,\QQ)$ of Hodge classes meeting the spinor variety
in two points conjugate over $K$, a pair of secant sheaves realising the
plane, and a deformation argument. This section supplies the first two in
closed form, for every $n\ge2$ and every imaginary quadratic $K$, and reduces
the third to an explicit equation on a Chern character. The fourth is
untouched and is the hardest of the four; we say so at the end.

The framework is Markman's: the Weil structure on $X\times\widehat{X}$
obtained by embedding $K$ into $\End^{0}(X\times\widehat{X})$, and the
notion of a $K$-secant sheaf whose Chern character lies in a plane meeting
the spinor variety at two conjugate points, are introduced in
\cite{Mar25a,Mar25b}, and the deformation step that finishes the argument
at $n=3$ is his. What we do here is write the first two ingredients out in
closed form, uniformly in $n$, $K$ and the discriminant, so that the third
becomes a finite system; we claim nothing beyond that.
What follows is elementary once set up, and we were surprised that it comes
out in closed form. The reason it does is that the two conjugate isotropic
subspaces demanded by \Cref{cor:orbitcriterion} can be produced by a single
element of the N\'eron--Severi group, and that the pure spinor of the graph of
a two-form is the exponential of that form.

\subsection{The Weil structure}

\begin{setupenv}\label{setup:construction}
Let $X$ be a complex abelian variety of dimension $n$ with
$H_{1}(X,\QQ)=U$ of dimension $2n$ and complex structure $J$ on
$U\otimes\RR$. Let $d$ be a squarefree positive integer, $s=\sqrt{-d}$ and
$K=\QQ(s)$. Let $\beta\in\NS(X)\otimes\QQ$ be nondegenerate, regarded as an
alternating form on $U$, so that the Riemann relation $J^{\top}\beta J=\beta$
holds, and write $\phi_{\beta}\colon X\to\widehat{X}$ for the associated
isogeny. Put
\[
  A=X\times\widehat{X},\qquad \dim A=2n,\qquad
  H_{1}(A,\QQ)=U\oplus U^{\vee},
\]
and let
\begin{equation}\label{eq:Mdef}
  M=\begin{pmatrix}0 & -\beta^{-1}\\[2pt] d\,\beta & 0\end{pmatrix}
  \in\End^{0}(A),
\end{equation}
acting on $H_{1}(A,\QQ)=U\oplus U^{\vee}$ by
$(u,\phi)\mapsto(-\beta^{-1}\phi,\;d\,\beta u)$.%
\footnote{The placement of the factor $d$ is a convention, and it has to be
made once and kept. The endomorphism $\begin{psmallmatrix}0&-d\beta^{-1}\\
\beta&0\end{psmallmatrix}$ also squares to $-d$ and also commutes with the
complex structure, and it defines a Weil structure on the same $A$; but it
is a different embedding of $K$, its eigenspaces are the graphs of
$\mp s\beta/d$ rather than of $\mp s\beta$, its polarisation is
$\beta+d\widehat\beta$ rather than $d\beta+\widehat\beta$, and the class in
$\bigwedge^{2}V_{+}\oplus\bigwedge^{2}V_{-}$ it produces is
$\beta-d\widehat\beta$ rather than $\gamma=d\beta-\widehat\beta$. The
convention of \eqref{eq:Mdef} is the one whose action on cohomology is
$x\mapsto-Bx$, $\xi\mapsto dB^{-1}\xi$ in the proof of
\Cref{thm:splitclosed}, and it is the one every script uses; item (XI) of
\Cref{sec:verification} checks the three eigenvalues that pin it down, and
item (XXV) shows what goes wrong if the two are mixed: the space of
polarised deformations killing the Weil line then has dimension
$n(n+1)/2$ instead of $n^{2}$. For $d=1$ the two conventions coincide.}
\end{setupenv}

\begin{theorem}[An explicit Weil structure]\label{thm:explicitweil}
In \Cref{setup:construction}:
\begin{enumerate}[label=\textup{(\roman*)}]
\item $M^{2}=-d$, so $\iota(s)=M$ defines an embedding
  $\iota\colon K\hookrightarrow\End^{0}(A)$;
\item $M$ commutes with the complex structure of $A$, and this holds exactly
  because $\beta$ satisfies the Riemann relation, that is because
  $\beta\in\NS(X)\otimes\QQ$;
\item $M$ acts on $H^{1,0}(A)$ with eigenvalues $s$ and $-s$, each of
  multiplicity $n$, so $(A,\iota,\eta)$ is of $(K,-1,n)$-Weil type for a
  suitable polarisation $\eta$;
\item the eigenspaces $W_{\pm}\subset H_{1}(A,\CC)$ of $M$ for $\pm s$ are
  maximal isotropic for the canonical pairing on $U\oplus U^{\vee}$, are
  interchanged by complex conjugation, and are the graphs of the two-forms
  $\mp s\beta$.
\end{enumerate}
\end{theorem}

\begin{proof}
(i) Block multiplication gives
$M^{2}=\mathrm{diag}(-\beta^{-1}\,d\beta,\;-d\beta\,\beta^{-1})=-d\cdot\mathrm{id}$.

(ii) The complex structure of $A=X\times\widehat{X}$ is
$J_{A}=\mathrm{diag}(J,-J^{\top})$. The relation $MJ_{A}=J_{A}M$ amounts to
the two identities $d\beta J=-dJ^{\top}\beta$ and
$-\beta^{-1}(-J^{\top})=J(-\beta^{-1})$, and each is equivalent to
$\beta J=-J^{\top}\beta$. That is exactly the Riemann relation: from
$J^{\top}\beta J=\beta$ and $J^{2}=-1$ one gets
$J^{\top}\beta=\beta J^{-1}=-\beta J$, as required. Conversely
$\beta J=-J^{\top}\beta$ gives $J^{\top}\beta J=-\beta J J=\beta$.

(iii) $M$ is block anti-diagonal and $H^{1,0}(A)=H^{1,0}(X)\oplus
H^{1,0}(\widehat{X})$ with both summands of dimension $n$, and $M$ exchanges
them. If $M(x,y)=\lambda(x,y)$ then $-\beta^{-1}y=\lambda x$ and
$d\beta x=\lambda y$, so $y=d\beta x/\lambda$ and, substituting,
$-dx/\lambda=\lambda x$, that is $\lambda^{2}=-d$. For each of
the two values of $\lambda$ the assignment $x\mapsto(x,d\beta x/\lambda)$ is an
isomorphism from $H^{1,0}(X)$ onto the $\lambda$-eigenspace, which is
therefore of dimension $n$. That is the balanced signature of
\Cref{def:weiltype}(ii). For (iii) of that definition one takes for $\eta$ the
polarisation of $A$ obtained from the canonical principal polarisation of
$X\times\widehat{X}$ twisted so that the Rosati involution sends $M$ to $-M$;
this is possible because $\beta$ is symmetric under the Rosati involution of
$X$, being a class in $\NS$.

(iv) The same computation identifies $W_{\lambda}=\{(x,d\beta x/\lambda)\}$,
the graph of the form $d\beta/\lambda$, and
$d\beta/\lambda=d\beta/(\pm s)=\mp s\beta$ since $d/s=-s$. The canonical
pairing on $U\oplus U^{\vee}$ is
$\langle(x,\xi),(x',\xi')\rangle=\xi'(x)+\xi(x')$, and on the graph it gives
$d\beta(x',x)/\lambda+d\beta(x,x')/\lambda=0$ because $\beta$ is alternating; so
$W_{\lambda}$ is isotropic, and it has dimension $2n$, half of
$\dim(U\oplus U^{\vee})=4n$, hence is maximal. Conjugation sends $s$ to $-s$
and so interchanges the two.
\end{proof}

\begin{remark}\label{rem:abundant}
\Cref{thm:explicitweil} says that abelian $2n$-folds of Weil type of the split
form $X\times\widehat{X}$ exist for every $n$, every imaginary quadratic $K$
and every polarised abelian $n$-fold $X$, and that the Weil structure is
written down by one element of $\NS(X)$. These are not general members of the
moduli of Weil type, and the Weil classes on them are of interest precisely
because Markman's deformation argument propagates algebraicity from them to
the general member; see \Cref{rem:deformation} below.
\end{remark}

\subsection{The plane}

Write $t=\beta\in\NS(X)\otimes\QQ$, so that by
\Cref{thm:explicitweil}(iv) the two conjugate maximal isotropic subspaces are
the graphs of $\mp s\,t$; the letter $t$ is kept because everything below
holds for any nondegenerate $t\in\NS(X)\otimes\QQ$. The pure spinor attached to the graph of a two-form
$\omega$ is $e^{\omega}$, so the two conjugate points of the even spinor
variety determined by the Weil structure are $[e^{-st}]$ and $[e^{st}]$.

\begin{definition}[The secant plane]\label{def:secantplane}
With $t$ and $d$ as above, put
\begin{equation}\label{eq:uv}
  u_{t}=\sum_{j\ge0}\frac{(-d)^{j}}{(2j)!}\,t^{2j},
  \qquad
  v_{t}=\sum_{j\ge0}\frac{(-d)^{j}}{(2j+1)!}\,t^{2j+1},
\end{equation}
both sums being finite because $t^{k}=0$ for $k>n$, and let
$P_{t}=\QQ u_{t}\oplus\QQ v_{t}\subset H^{\mathrm{ev}}(X,\QQ)$.
\end{definition}

\begin{theorem}[The plane is an explicit rational secant]\label{thm:plane}
With notation as above:
\begin{enumerate}[label=\textup{(\roman*)}]
\item $u_{t}$ and $v_{t}$ are rational Hodge classes, and
  $e^{\pm st}=u_{t}\pm s\,v_{t}$;
\item $P_{t}$ is a two-dimensional $\QQ$-subspace of $H^{\mathrm{ev}}(X,\QQ)$
  spanned by Hodge classes, and $\PP(P_{t})$ is a line defined over $\QQ$
  passing through the two conjugate points $[e^{\pm st}]$ of the even spinor
  variety;
\item that line is not contained in the spinor variety, so it meets it in
  exactly those two points.
\end{enumerate}
\end{theorem}

\begin{proof}
(i) Each $t^{k}$ is a rational Hodge class because $t\in\NS(X)\otimes\QQ$ is
of type $(1,1)$ and the algebraic classes form a subring. Separating
$e^{st}=\sum_{k}s^{k}t^{k}/k!$ into even and odd powers of $s$ and using
$s^{2j}=(-d)^{j}$, $s^{2j+1}=(-d)^{j}s$ gives $e^{st}=u_{t}+sv_{t}$, and
conjugating gives $e^{-st}=u_{t}-sv_{t}$.

(ii) The two elements are linearly independent over $\QQ$ because $u_{t}$ has
nonzero component in $H^{0}$ and $v_{t}$ does not, while $v_{t}\ne0$ since
$t\ne0$. From (i), both $e^{st}$ and $e^{-st}$ lie in $P_{t}\otimes K$, so
$\PP(P_{t})$ passes through the two points.

(iii) The even spinor variety is cut out by quadrics, by Cartan's description
of pure spinors, so a line either lies inside it or meets it in a subscheme of
length at most two. The point $[u_{t}]$ of $\PP(P_{t})$ is not a pure spinor:
a pure spinor with nonzero component in $H^{0}$ is a scalar multiple of
$e^{\omega}$ for a two-form $\omega$, and $u_{t}$ has zero component in
$H^{2}$ while its component in $H^{4}$ is $-\tfrac{d}{2}t^{2}\ne0$, whereas
$e^{\omega}$ with vanishing $H^{2}$ part is $1$. So the line is not contained
in the spinor variety, and the two points of (ii) exhaust the intersection.
\end{proof}

\begin{example}\label{ex:planes}
For $n=4$ and $d=1$,
$u_{t}=1-\tfrac{1}{2}t^{2}+\tfrac{1}{24}t^{4}$ and
$v_{t}=t-\tfrac{1}{6}t^{3}$. For $n=4$ and $d=3$,
$u_{t}=1-\tfrac{3}{2}t^{2}+\tfrac{3}{8}t^{4}$ and
$v_{t}=t-\tfrac{1}{2}t^{3}$. Both are visibly rational and visibly in the
subring generated by $t$.
\end{example}

\subsection{The Chern character equation}

We can now say exactly what a secant sheaf has to be. The following is the
substance of this section: the four demands of \Cref{q:higher} collapse, on
the sheaf side, to one equation with all invariants computed.

\begin{theorem}[The invariants of a secant sheaf]\label{thm:secantinvariants}
Let $X$, $t$, $d$ be as above and let $\cF$ be a coherent sheaf on $X$ with
$\ch(\cF)\in P_{t}$, say $\ch(\cF)=a\,u_{t}+b\,v_{t}$ with $a,b\in\QQ$. Then
\[
  \rk\cF=a,\qquad
  c_{1}(\cF)=b\,t,\qquad
  \ch_{2}(\cF)=-\tfrac{ad}{2}\,t^{2},\qquad
  c_{2}(\cF)=\tfrac{ad+b^{2}}{2}\,t^{2},
\]
and more generally $\ch_{k}(\cF)$ is the degree $k$ term of \eqref{eq:uv}
scaled by $a$ or $b$ according to the parity of $k$. The Bogomolov
discriminant is
\begin{equation}\label{eq:discnorm}
  \Delta(\cF)=2a\,c_{2}(\cF)-(a-1)c_{1}(\cF)^{2}
  =\Nm_{K/\QQ}\bigl(b+a s\bigr)\cdot t^{2}.
\end{equation}
In particular $\Delta(\cF)\cdot H^{n-2}>0$ for every ample $H$, so the
Bogomolov inequality places no obstruction on the existence of $\cF$.
\end{theorem}

\begin{proof}
The first four identities are read off from \eqref{eq:uv} degree by degree,
using $\ch_{2}=\tfrac12(c_{1}^{2}-2c_{2})$ to convert $\ch_{2}$ into $c_{2}$.
For \eqref{eq:discnorm},
\[
  2a\cdot\tfrac{ad+b^{2}}{2}t^{2}-(a-1)b^{2}t^{2}
  =\bigl(a^{2}d+ab^{2}-ab^{2}+b^{2}\bigr)t^{2}
  =\bigl(a^{2}d+b^{2}\bigr)t^{2},
\]
and $\Nm_{K/\QQ}(b+as)=(b+as)(b-as)=b^{2}+a^{2}d$. Positivity of the norm on
$\Kx$ and of $t^{2}H^{n-2}$ for $t$ in the positive cone give the last
assertion.
\end{proof}

\begin{remark}[The norm form as the discriminant]\label{rem:normdisc}
Identity \eqref{eq:discnorm} is the reason the construction has a chance at
all. The discriminant of the sheaf is the norm form of $K$ evaluated at the
coordinates of its Chern character in the plane. The two conjugate points of
the spinor variety are the two isotropic directions of that norm form after
extension to $K$, and the rational points of the plane are exactly the
non-isotropic ones. A sheaf therefore sits at a rational point of a conic
whose isotropic directions carry the Weil structure, which is the precise
sense in which the conjugate pair of \Cref{cor:orbitcriterion} is built into
the invariants rather than imposed on them.
\end{remark}

\begin{corollary}[The case of an abelian surface]\label{cor:surfacecase}
Let $n=2$, so $X$ is an abelian surface and $A=X\times\widehat{X}$ an abelian
fourfold of Weil type. A sheaf with $\ch(\cF)\in P_{t}$ has Mukai vector
$v(\cF)=(a,\,bt,\,-\tfrac{ad}{2}t^{2})$ and Mukai self-pairing
\[
  \langle v(\cF),v(\cF)\rangle=\Nm_{K/\QQ}(b+as)\cdot t^{2}>0 .
\]
\end{corollary}

\begin{proof}
On an abelian surface the Mukai pairing of $v=(r,c,e)$ with itself is
$c^{2}-2re$, which here is $b^{2}t^{2}-2a(-\tfrac{ad}{2}t^{2})=(b^{2}+a^{2}d)t^{2}$.
\end{proof}

\Cref{cor:surfacecase} is consistent with the known range: the positivity of
the Mukai self-pairing is exactly the numerical condition under which moduli
of stable sheaves on an abelian surface with a given Mukai vector are
non-empty, and $n=2$ is settled.

\subsection{The genus three case, and what the field predicts}

At $n=3$ the equation of \Cref{thm:secantinvariants} pins down the secant
sheaf far more tightly, and it reproduces the shape of the object Markman
uses. We record the computation because it is a check on the whole framework.

\begin{proposition}[Rank one at $n=3$]\label{prop:genus3}
Let $X$ be the Jacobian of a smooth curve of genus three, let $\Theta$ be the
theta divisor and $C\subset X$ the Abel--Jacobi curve, so that
$\Theta^{2}/2=[C]$. Take $t=\Theta$ and $a=1$ in
\Cref{thm:secantinvariants}. Then a rank one torsion-free sheaf $\cF$ with
$\ch(\cF)\in P_{\Theta}$ and $c_{1}(\cF)=b\Theta$ is of the form
$\cF=\mathcal{I}_{Z}(b\Theta)$ with
\[
  [Z]=\frac{b^{2}+d}{2}\,\Theta^{2}
     =\frac{b^{2}+d}{b^{2}}\cdot b^{2}[C].
\]
If $Z$ is a disjoint union of translates of $C$, the number of translates is
$m=(b^{2}+d)/b^{2}$, which is a positive integer precisely when $b^{2}\mid d$;
for $b=1$,
\[
  m=d+1 .
\]
\end{proposition}

\begin{proof}
A rank one torsion-free sheaf with $c_{1}=b\Theta$ on a smooth projective
threefold is $\mathcal{I}_{Z}(b\Theta)$ for a subscheme $Z$ of codimension at
least two, and
$\ch(\mathcal{I}_{Z}(b\Theta))=e^{b\Theta}-[Z]-(\text{terms in }H^{6})$.
Comparing the components in $H^{4}$ with \Cref{thm:secantinvariants} at $a=1$,
which gives $\ch_{2}=-\tfrac{d}{2}\Theta^{2}$, yields
$\tfrac{b^{2}}{2}\Theta^{2}-[Z]=-\tfrac{d}{2}\Theta^{2}$, that is
$[Z]=\tfrac{b^{2}+d}{2}\Theta^{2}$. The Poincar\'e formula
$\Theta^{g-1}/(g-1)!=[C]$ at $g=3$ reads $\Theta^{2}/2=[C]$, whence
$[Z]=(b^{2}+d)[C]$, and if $Z$ is $m$ translates of $C$ then $m=b^{2}+d$ when
$b=1$ and $m=(b^{2}+d)/b^{2}$ after dividing by the class of $b^{2}$
translates in general.
\end{proof}

\begin{table}[htbp]
\centering
\renewcommand{\arraystretch}{1.15}
\begin{tabular}{@{}rrrrrrr@{}}
\toprule
$d$ & $1$ & $2$ & $3$ & $5$ & $7$ & $11$ \\
\midrule
$m=d+1$ & $2$ & $3$ & $4$ & $6$ & $8$ & $12$ \\
\bottomrule
\end{tabular}
\caption{The number of translates of the Abel--Jacobi curve required at
$n=3$, $b=1$, as a function of the field $K=\QQ(\sqrt{-d})$, by
\Cref{prop:genus3}.}
\label{tab:translates}
\end{table}

\begin{remark}\label{rem:markmancheck}
The sheaf used in \cite{Mar25b} on a genus three Jacobian is
$\mathcal{I}_{\cup C_{i}}(\Theta)$, the twist by the theta bundle of the ideal
sheaf of a union of translates of the Abel--Jacobi curve. That is exactly the
shape \Cref{prop:genus3} forces, with the number of translates determined by
the field. We take this as evidence that the framework of this section is the
right one and not merely a consistent one.
\end{remark}

\subsection{The third ingredient, in every dimension}\label{ssec:thirdexists}

\Cref{thm:secantinvariants} says what a secant sheaf must look like. It does
not say that one exists, and at $n\ge4$ that is the question the section has so
far left open. We settle it here, and the answer is that the object exists on
every polarised abelian variety, in every dimension, for every imaginary
quadratic field, and that it can be assembled out of powers of a single line
bundle.

Two observations do the work. The first is that the object wanted is a
rational point of a plane, and a plane carries no preferred scale. If a
positive multiple of a rational point of $P_{t}$ is realised by a sheaf, then a
rational point of $P_{t}$ is realised by a sheaf, and since the Hodge
conjecture is a statement about classes with rational coefficients, replacing
$\omega$ by $M\omega$ costs nothing. The second is that once one is free to
clear denominators, the Chern characters of the powers of one line bundle
already span the subring in which the secant plane lives, because the matrix
expressing them in the basis $t^{k}/k!$ is a Vandermonde matrix on distinct
nodes.

What is missing is only the passage from a class in $K_{0}$ to a sheaf, and
for that the sign of the rank is the whole of the hypothesis.

\begin{lemma}[Positive rank is enough]\label{lem:posrank}
Let $Y$ be a smooth projective variety over a field and let $\alpha\in
K_{0}(Y)$ satisfy $\rk\alpha>0$. Then $\alpha=[\cF]$ for some coherent sheaf
$\cF$ on $Y$.
\end{lemma}

\begin{proof}
Write $m=\dim Y$ and $a=\rk\alpha\ge1$, and let $F^{c}K_{0}(Y)$ be the
topological filtration, the subgroup generated by the classes $[\cO_{Z}]$ of
closed integral subschemes $Z\subseteq Y$ of codimension at least $c$. Since
$Y$ is smooth, every coherent sheaf has a finite locally free resolution, so
$K_{0}(Y)=K_{0}'(Y)$ and the filtration is exhaustive with $F^{m+1}=0$;
moreover $F^{c}/F^{c+1}$ maps onto $\CH^{c}(Y)$ by sending the class of a sheaf
supported in codimension $c$ to its associated cycle, and $F^{1}$ is the kernel
of the rank while $F^{1}/F^{2}$ is $\Pic(Y)$ by the first Chern class
\cite[\S 15.1 and Ex.~15.1.5]{Ful98}. We construct coherent
sheaves $\cF_{c}$ of rank $a$ with $\alpha-[\cF_{c}]\in F^{c}$, by induction on
$c$.

Start at $c=2$. Let $L$ be a line bundle with $c_{1}(L)=c_{1}(\alpha)$ and set
$\cF_{2}=L\oplus\cO_{Y}^{\oplus(a-1)}$. Then $\alpha-[\cF_{2}]$ has rank zero
and first Chern class zero, hence lies in $F^{2}$.

For the inductive step let $\delta=\alpha-[\cF_{c}]\in F^{c}$ with $c\ge2$.
Modulo $F^{c+1}$ the class $\delta$ is a finite sum $\sum_{j}m_{j}
[\cO_{Z_{j}}]$ with $m_{j}\in\ZZ$ and $\codim Z_{j}=c$. The terms with
$m_{j}>0$ are removed by replacing $\cF_{c}$ with $\cF_{c}\oplus\bigoplus_{
m_{j}>0}\cO_{Z_{j}}^{\oplus m_{j}}$, which changes neither the rank nor the
codimension of the remaining difference. For a term with $m_{j}<0$ we must
subtract, and here the rank is used. The restriction $\cF_{c}\otimes
\cO_{Z_{j}}$ is a quotient of $\cF_{c}$ and has rank $a\ge1$ at the generic
point of the integral scheme $Z_{j}$; choosing a subsheaf of generic corank
one and dividing by it gives a surjection $\cF_{c}\twoheadrightarrow\cQ$ with
$\cQ$ of rank one on $Z_{j}$. Then $[\cQ]\in F^{c}$ with associated cycle
$[Z_{j}]$, so $[\cQ]-[\cO_{Z_{j}}]\in F^{c+1}$. The kernel of that surjection
is a coherent sheaf, still of rank $a$ because $\cQ$ is supported in positive
codimension, and its class is $[\cF_{c}]-[\cQ]$. Performing this $|m_{j}|$
times, and once for each $j$, produces $\cF_{c+1}$ with
$\alpha-[\cF_{c+1}]\in F^{c+1}$ and $\rk\cF_{c+1}=a$. At $c=m+1$ the
difference lies in $F^{m+1}=0$, so $\alpha=[\cF_{m+1}]$.
\end{proof}

\begin{theorem}[Secant sheaves exist in every dimension]\label{thm:secantexists}
Let $X$ be an abelian variety of dimension $n$, let $t\in\NS(X)\otimes\QQ$ be a
polarisation and let $d>0$. Let $p$ be a rational point of the secant plane
$P_{t}$ whose $u_{t}$ coordinate is positive. Then there are a positive
integer $M$, an integer $D>0$ with $Dt\in\NS(X)$, integers $n_{0},\dots,n_{n}$
and a coherent sheaf $\cF$ on $X$ such that, setting
\begin{equation}\label{eq:lbwitness}
  \alpha=\sum_{j=0}^{n}n_{j}\bigl[\cO_{X}(jDt)\bigr]\in K_{0}(X),
\end{equation}
one has $\alpha=[\cF]$ in $K_{0}(X)$, $\ch(\cF)=M\,p$ and
$\rk\cF=M\,a>0$, where $a$ is the $u_{t}$ coordinate of $p$. The integers $M$
and $n_{j}$ depend only on $n$, $d$ and the coordinates of $p$, and are
obtained by inverting one Vandermonde matrix. Moreover $Mp$ is again a
rational point of $P_{t}$; applying the statement to two linearly independent
rational points of $P_{t}$ produces a pair of coherent sheaves whose Chern
characters span $P_{t}$.
\end{theorem}

\begin{proof}
Scaling $p$ by a positive integer changes neither the hypothesis nor the
conclusion, so we may assume $p=au_{t}+bv_{t}$ with $a,b\in\ZZ$ and $a>0$.
Choose $D>0$ with $Dt\in\NS(X)$, so that $\cO_{X}(jDt)$ is a line bundle for
every $j\in\ZZ$.

Since $t$ is a polarisation, $t^{n}\ne0$ and $t^{n+1}=0$ in
$H^{\mathrm{ev}}(X,\QQ)$, so the $n+1$ classes $t^{k}/k!$ with $0\le k\le n$
are a basis of the subring $\QQ[t]$. In that basis
\[
  p=\sum_{k=0}^{n}c_{k}\,\frac{t^{k}}{k!},\qquad
  c_{k}=\begin{cases}a(-d)^{k/2}, & k\ \text{even},\\[2pt]
                     b(-d)^{(k-1)/2}, & k\ \text{odd},\end{cases}
\]
by \Cref{prop:notforbidden}(i), and all the $c_{k}$ are integers. On the other
side,
\[
  \ch\bigl(\cO_{X}(jDt)\bigr)=e^{jDt}=\sum_{k=0}^{n}(jD)^{k}\,\frac{t^{k}}{k!},
\]
so in the same basis the $n+1$ vectors $\ch(\cO_{X}(jDt))$ for $j=0,\dots,n$
are the columns of the Vandermonde matrix $\bigl((jD)^{k}\bigr)_{0\le j,k\le n}$
on the $n+1$ distinct nodes $0,D,2D,\dots,nD$. Its determinant is
$\prod_{0\le i<j\le n}(jD-iD)=D^{n(n+1)/2}\prod_{0\le i<j\le n}(j-i)\ne0$, so
the system
\[
  \sum_{j=0}^{n}x_{j}\,(jD)^{k}=c_{k}\qquad(0\le k\le n)
\]
has a unique solution $x\in\QQ^{n+1}$. Let $M$ be the least common multiple of
the denominators of the $x_{j}$ and put $n_{j}=Mx_{j}\in\ZZ$. With $\alpha$ as
in \eqref{eq:lbwitness},
\[
  \ch(\alpha)=\sum_{j}n_{j}e^{jDt}
  =\sum_{k=0}^{n}\Bigl(\sum_{j}Mx_{j}(jD)^{k}\Bigr)\frac{t^{k}}{k!}
  =M\sum_{k=0}^{n}c_{k}\frac{t^{k}}{k!}=Mp .
\]
In particular $\rk\alpha=\ch_{0}(\alpha)=Mc_{0}=Ma>0$, so \Cref{lem:posrank}
gives a coherent sheaf $\cF$ with $[\cF]=\alpha$. The Chern character is a
function of the class in $K_{0}(X)$, whence $\ch(\cF)=Mp$ and $\rk\cF=Ma$.
Finally $Mp=(Ma)u_{t}+(Mb)v_{t}$ is a rational point of $P_{t}$, and two
linearly independent rational points stay linearly independent after scaling,
which gives the last assertion.
\end{proof}

\begin{remark}[The multiple is real, and it is the factorials]%
\label{rem:themultiple}
The integer $M$ is not an artifact of the proof. The lattice spanned by the
Chern characters of the powers of a single line bundle is spanned equally well
by the classes $(e^{t}-1)^{k}$ for $0\le k\le n$, because the $(n+1)$st finite
difference of a polynomial of degree $n$ vanishes; and $(e^{t}-1)^{k}=t^{k}+
\cdots$ has leading term $t^{k}$, not $t^{k}/k!$. The secant plane, by
contrast, is spanned by classes whose coordinates in the basis $t^{k}/k!$ are
integers. The index between the two lattices is what $M$ measures, and it is
created by the factorials in $\ch$. Item (XXI) of \Cref{sec:verification}
computes the least $M$ exactly for $1\le n\le10$, every squarefree $d\le11$ and
a range of $a$ and $b$, together with the witness \eqref{eq:lbwitness} in each
case. At $n=4$, $d=3$, $a=b=1$ the least multiple is $M=6$ and the witness is
\[
  \alpha=-23\,[\cO_{X}]+66\,[\cO_{X}(t)]-54\,[\cO_{X}(2t)]
         +20\,[\cO_{X}(3t)]-3\,[\cO_{X}(4t)],
\]
of rank $6=Ma$; at $n=8$ the largest multiple over the range tested is $2016$,
and $M=1$ occurs in none of the seventy two cases tested once $n\ge7$. So the
freedom to rescale is used, and used essentially.

The same computation checks, independently and in the exterior algebra of a
symplectic lattice, that $\theta^{k}/k!$ has integer coefficients for a
principal polarisation. That is the reason the coordinates $c_{k}$ are
integers to begin with, and it is Poincar\'e's formula in the form that does not
need a curve.
\end{remark}

\begin{figure}[htbp]
\centering
  \includegraphics[width=0.86\textwidth]{fig_moment}
\caption{\Cref{thm:secantexists} at $n=2$ and $d=2$, in the coordinates
$(\ch_{0},\ch_{1},\ch_{2})$ of the basis $1$, $t$, $t^{2}/2$, the third axis at
two fifths scale. The characters of the powers of one line bundle lie on the
moment curve $j\mapsto(1,j,j^{2})$, and the Vandermonde determinant says that
any three of its points are affinely independent. The integral points of
$P_{t}=\QQ u_{t}\oplus\QQ v_{t}$ split by the criterion
$\ch_{2}\equiv\ch_{1}\bmod2$ into those the line bundles reach, drawn filled,
and those they miss, drawn open. Doubling $p=u_{t}+v_{t}$ lands on the
lattice, at the witness $-3[\cO_{X}]+8[\cO_{X}(t)]-3[\cO_{X}(2t)]$ of the
three dashed leaders.}
\label{fig:moment}
\end{figure}

\begin{remark}[What this closes, and what it does not]\label{rem:thirdclosed}
\Cref{thm:secantexists} settles the third of the four demands of
\Cref{q:higher}, for every $n\ge2$, every imaginary quadratic field and every
polarised abelian $n$-fold, and it does so with line bundles alone. Three
qualifications belong with it, and they are not small.

First, the theorem produces a sheaf and nothing more. It says nothing about
stability, about the rank being one, about the Ext groups, or about
semiregularity, and the construction of \Cref{lem:posrank} is not designed to
give any of them: the sheaf is assembled from line bundles and from quotients
supported on subvarieties, and its only controlled invariant is its class.
\Cref{prop:objectsize} applies to it as to any other object carrying the class,
and bounds $\sum_{i}\dim\Ext^{i}(\cF,\cF)$ from below.

This does not collide with \Cref{thm:positivity}, for the following reason. That
theorem concerns a perfect complex on the $2n$-fold $A$ whose Chern character
is an effective combination $\sum_{i}m_{i}e^{\lambda_{i}}$ with every $m_{i}$
positive, and it rules such a complex out because the norm form of $K$ is
positive definite. The class \eqref{eq:lbwitness} lives on the $n$-fold $X$,
not on $A$, and its coefficients alternate in sign; the witness at $n=4$ has
three negative ones. Positivity is exactly what is given up, and
\Cref{lem:posrank} is what replaces it, by requiring only that the total rank
be positive.

Second, the rank is large. \Cref{prop:genus3} at $n=3$ produces a rank one
sheaf, which is what Markman's argument uses and what \Cref{rem:markmancheck}
matches against the literature. The sheaf of \Cref{thm:secantexists} has rank
$Ma$, and $M$ grows quickly with $n$. Whether the rank can be brought down to
one at $n\ge4$, equivalently whether the lattice generated by the classes
$t^{k}/k!$ is realised by algebraic cycles on a general abelian $n$-fold rather
than only after clearing denominators, is a separate and open
question; on a Jacobian it is answered by the Abel--Jacobi cycles and
Poincar\'e's formula, and $M=1$ can be taken there.

Third, the fourth ingredient is untouched, and it is the one that matters. A
pair of secant sheaves gives a class on $X\times\widehat{X}$; making that class
remain of Hodge type across the moduli of polarised abelian $2n$-folds of Weil
type is a variational statement, and nothing here bears on it.
\Cref{rem:deformation} says what that step requires and stands unchanged.
We therefore do not claim the Hodge conjecture for abelian $2n$-folds of Weil
type at any $n\ge4$; what has changed is that three of the four ingredients of
\Cref{q:higher} are now theorems rather than questions.
\end{remark}

\subsection{What is supplied and what is not}

\begin{theorem}[Three of the four ingredients, in general]\label{thm:twoofthefour}
For every $n\ge2$, every squarefree $d>0$ and every polarised abelian $n$-fold
$X$ carrying a nondegenerate $\beta\in\NS(X)\otimes\QQ$, the data required by
\Cref{q:higher} exist and are explicit in the first three items:
\begin{enumerate}[label=\textup{(\alph*)}]
\item the abelian $2n$-fold of Weil type is $A=X\times\widehat{X}$ with the
  $K$-action \eqref{eq:Mdef}, by \Cref{thm:explicitweil};
\item the plane is $P_{t}$ of \Cref{def:secantplane} with $t=\beta$, it is
  spanned by Hodge classes and it meets the spinor variety in exactly the two
  conjugate points $[e^{\pm st}]$, by \Cref{thm:plane};
\item a pair of coherent sheaves $\cF_{1},\cF_{2}$ on $X$ whose Chern
  characters are linearly independent rational points of $P_{t}$ exists, by
  \Cref{thm:secantexists}, and each may be taken to have the class of an
  explicit integer combination of powers of one line bundle.
\end{enumerate}
The invariants of any such sheaf are forced by the explicit system of
\Cref{thm:secantinvariants}, which imposes $n-1$ linear conditions on a Chern
character already constrained to $\QQ[t]$, and are unobstructed numerically by
\eqref{eq:discnorm}.
\end{theorem}

\begin{proof}
(a), (b) and (c) are the cited theorems. For the count in the last sentence, a
Chern character in the subring $\QQ[t]\subset H^{\mathrm{ev}}(X,\QQ)$ is a
polynomial of degree at most $n$ in $t$, so lies in a space of dimension
$n+1$, while $P_{t}$ has dimension two.
\end{proof}

\begin{remark}[What is not supplied]\label{rem:deformation}
The fourth ingredient of \Cref{q:higher} is untouched by anything in this
section, and it is the hardest. Granting the pair of secant sheaves
$\cF_{1},\cF_{2}$ on $X$ that \Cref{thm:secantexists} now supplies, one must
still show that the characteristic class of
$\Phi(\cF_{1}^{\vee}\boxtimes\cF_{2})$ remains of Hodge type under every
deformation of $X\times\widehat{X}$ as a polarised abelian $2n$-fold of Weil
type, and that the object deforms with the variety, as a twisted sheaf to
first order. That step is carried out by Markman at $n=3$ \cite{Mar25b}, and it
rests on a semiregularity analysis. What we add to it is negative:
\Cref{thm:factor} identifies exactly which Hochschild classes the weakened
form of that analysis must kill for a secant object, and
\Cref{thm:candidatefails} shows that the only candidate proposed above
$n=3$ does not kill them. Without the step one obtains algebraicity only on
the split locus $X\times\widehat{X}$, where the Weil classes are of much
less interest, since the general abelian $2n$-fold of Weil type is simple and
is not of that form.

We therefore do not claim a proof of the Hodge conjecture for abelian
$2n$-folds of Weil type at any $n\ge4$, and nothing in this section should be
read as one. What we claim is that three of the four ingredients are now
written down for every $n$, the third of them with an explicit witness in
line bundles, and that the fourth is isolated.
\end{remark}

\begin{figure}[htbp]
\centering
  \includegraphics[width=0.88\textwidth]{fig_quadric}
\caption{The geometry of \Cref{thm:plane}, drawn in the smallest case in which
the even spinor variety is a surface, where it is a smooth quadric with its two
rulings. The line $P_{t}$ of \Cref{def:secantplane} is rational, and it meets
the variety in exactly the two points $[e^{\pm st}]$, which are conjugate over
$K$ and not individually rational. That is the configuration the Weil line
demands: a class defined over $\QQ$ cannot pick out one of the two points, so
the construction must carry the pair, and the pair spans a plane rather than a
line. \Cref{prop:orbitplane} is the same statement on the cohomological side,
and \Cref{fig:escape} places it among the hypotheses a construction can
satisfy.}
\label{fig:spinor}
\end{figure}

\subsection{The split member, in closed form}\label{ssec:split}

Before turning to what is missing, we record what the family gives outright.%
\footnote{The split locus is thin: a general abelian $2n$-fold of Weil type is simple, so it is not isogenous to a product at all, and the members $X\times\widehat{X}$ form a subvariety of $\cD_{n,n}$ of positive codimension. That is exactly why \Cref{cor:splitalgebraic} is a base case and not a solution, and why \Cref{q:residual} is about propagation.}
On the split member the Weil classes are not merely known to be algebraic;
they can be written down as polynomials in three divisor classes, and the
polynomial is the same for every $n$ and every $d$.

\begin{notation}\label{not:threeclasses}
On $A=X\times\widehat{X}$ write
\[
  \beta\in H^{2}(A,\QQ),
  \qquad
  \widehat{\beta}\in H^{2}(A,\QQ),
  \qquad
  \ell=c_{1}(\cP)\in H^{2}(A,\QQ)
\]
for the pullback of the polarisation of $X$, the pullback of the dual
polarisation of $\widehat{X}$, and the class of the Poincar\'e bundle. All
three are divisor classes, so every polynomial in them is algebraic. Put
\begin{equation}\label{eq:gammadef}
  \gamma=d\,\beta-\widehat{\beta}.
\end{equation}
\end{notation}

\begin{theorem}[The Weil line of the split member]\label{thm:splitclosed}
Let $A=X\times\widehat{X}$ carry the $K$-action \eqref{eq:Mdef} determined by a
principal $\beta$, and let $V_{\pm}\subset H^{1}(A,\CC)$ be the eigenspaces of
$\iota(\sqrt{-d})$. Then
\begin{enumerate}[label=\textup{(\roman*)},nosep]
\item $\gamma\mp\sqrt{-d}\,\ell$ lies in $\bigwedge^{2}V_{\pm}$, so it is an
  eigenclass for the character $\tau\mapsto\tau^{2}$ on $H^{2}(A,\CC)$
  respectively for its conjugate;
\item the generators of $\bigwedge^{2n}V_{\pm}$ are
  \begin{equation}\label{eq:splitclosed}
    \omega_{\pm}
    =\frac{(-1)^{n(n-1)/2}}{n!\,d^{\,n}}
      \bigl(\gamma\mp\sqrt{-d}\,\ell\bigr)^{n};
  \end{equation}
\item consequently the Weil line of $A$ is
  \[
    \HW(A)=\Bigl\langle\;
      \Re\bigl(\gamma-\sqrt{-d}\,\ell\bigr)^{n},\;\;
      \Im\bigl(\gamma-\sqrt{-d}\,\ell\bigr)^{n}
    \;\Bigr\rangle_{\QQ},
  \]
  and every Weil class on $A$ is a polynomial with rational coefficients in
  the three divisor classes $\beta,\widehat{\beta},\ell$, hence is algebraic.
\end{enumerate}
\end{theorem}

\begin{proof}
Write $U=H^{1}(X,\QQ)$ with basis $x_{1},\dots,x_{2n}$ chosen so that
$\beta=\sum_{j\le n}x_{j}\wedge x_{n+j}$, let $\xi_{1},\dots,\xi_{2n}$ be the
dual basis of $H^{1}(\widehat{X},\QQ)=U^{\vee}$, and let $B\colon U\to U^{\vee}$
be the isomorphism induced by $\beta$, so that $Bx_{j}=\xi_{n+j}$ and
$Bx_{n+j}=-\xi_{j}$. The endomorphism \eqref{eq:Mdef} acts on homology, and
its transpose acts on $H^{1}(A,\QQ)=U\oplus U^{\vee}$ by $u\mapsto -Bu$ and
$\phi\mapsto d\,B^{-1}\phi$;%
\footnote{In the basis $u_{1},\dots,u_{2n}$ of $H_{1}(X,\QQ)$ dual to the
$x_{i}$, the map $\beta\colon H_{1}(X)\to H^{1}(X)$ sends $u_{j}\mapsto
x_{n+j}$ and $u_{n+j}\mapsto-x_{j}$, so \eqref{eq:Mdef} sends
$u_{j}\mapsto dx_{n+j}$, $u_{n+j}\mapsto-dx_{j}$, $x_{n+j}\mapsto-u_{j}$ and
$x_{j}\mapsto u_{n+j}$; pairing with the dual basis gives the stated action
on $H^{1}(A,\QQ)$, where $\xi_{i}$ is the class dual to $x_{i}$ regarded as an
element of $H_{1}(\widehat X,\QQ)=H^{1}(X,\QQ)$.} its
square is $-d$; so, writing $\delta=\sqrt{-d}$,
\begin{equation}\label{eq:splitVpm}
  V_{\pm}=\Bigl\{\,u\pm\tfrac{\delta}{d}\,Bu\;:\;u\in U_{\CC}\Bigr\},
\end{equation}
since $M^{*}\bigl(u\pm\tfrac{\delta}{d}Bu\bigr)
=-Bu\pm\delta u=\pm\delta\bigl(u\pm\tfrac{\delta}{d}Bu\bigr)$.

(i) Take the sum over $j\le n$ of the wedges of the two eigenvectors attached
to $x_{j}$ and to $x_{n+j}$:
\[
  \sum_{j=1}^{n}
  \Bigl(x_{j}\pm\tfrac{\delta}{d}Bx_{j}\Bigr)\wedge
  \Bigl(x_{n+j}\pm\tfrac{\delta}{d}Bx_{n+j}\Bigr)
  =\beta\pm\tfrac{\delta}{d}\Sigma-\tfrac{1}{d}\widehat{\beta},
\]
using $\delta^{2}=-d$, where
$\Sigma=\sum_{j}\bigl(x_{j}\wedge Bx_{n+j}+Bx_{j}\wedge x_{n+j}\bigr)$ and
$\sum_{j}Bx_{j}\wedge Bx_{n+j}=\sum_{j}\xi_{n+j}\wedge(-\xi_{j})
=\widehat{\beta}$. Now
\[
  \Sigma=\sum_{j=1}^{n}\bigl(x_{j}\wedge(-\xi_{j})
        +\xi_{n+j}\wedge x_{n+j}\bigr)
   =-\sum_{i=1}^{2n}x_{i}\wedge\xi_{i}=-\ell .
\]
So the displayed sum is $\tfrac1d\bigl(d\beta-\widehat{\beta}\mp\delta\ell\bigr)
=\tfrac1d(\gamma\mp\delta\ell)$. It is a sum of wedges of pairs of vectors of
$V_{\pm}$, hence lies in $\bigwedge^{2}V_{\pm}$, which proves (i); the
statement about characters is \Cref{prop:grading} at $(r,s)=(2,0)$ and
$(0,2)$.

(ii) By (i) the $n$-th power $(\gamma\mp\delta\ell)^{n}$ lies in
$\bigwedge^{2n}V_{\pm}$, which is a line by \Cref{prop:weilline}. It is
nonzero: expanding $\gamma\mp\delta\ell=d\beta-\widehat{\beta}\mp\delta\ell$,
the only term of the $n$-th power lying in $\bigwedge^{2n}U$ is
$d^{\,n}\beta^{n}$, and $\beta^{n}\ne0$ because $\beta$ is nondegenerate;
the remaining terms all involve at least one $\xi$, so no cancellation can
occur. Hence $(\gamma\mp\delta\ell)^{n}$ spans $\bigwedge^{2n}V_{\pm}$ and is a
scalar multiple of $\omega_{\pm}=\bigwedge_{i=1}^{2n}\bigl(x_{i}\pm\tfrac{\delta}{d}Bx_{i}\bigr)$.
To find the scalar, compare the components in $\bigwedge^{2n}U$: that of
$\omega_{\pm}$ is $x_{1}\wedge\cdots\wedge x_{2n}$, while
\[
  d^{\,n}\beta^{n}=d^{\,n}\,n!\bigwedge_{j=1}^{n}(x_{j}\wedge x_{n+j})
  =d^{\,n}\,n!\,(-1)^{n(n-1)/2}\,x_{1}\wedge\cdots\wedge x_{2n},
\]
the sign being that of the permutation carrying
$(1,n+1,2,n+2,\dots,n,2n)$ to $(1,2,\dots,2n)$. Inverting gives
\eqref{eq:splitclosed}.

(iii) The two displayed classes are the real and imaginary parts of
$(\gamma-\delta\ell)^{n}$, which up to the nonzero rational factor in
\eqref{eq:splitclosed} are $\omega_{-}+\omega_{+}$ and
$(\omega_{-}-\omega_{+})/\delta$; by \Cref{prop:weilline}(i) these span
$\HW(A)$ over $\QQ$. Expanding $(\gamma-\delta\ell)^{n}$ by the binomial
theorem and separating even from odd powers of $\delta$ writes each of them as
a rational polynomial in $\beta,\widehat{\beta},\ell$, because
$\delta^{2m}=(-d)^{m}$ and $\delta^{2m+1}=(-d)^{m}\delta$. Products of divisor
classes are algebraic, so both are algebraic.
\end{proof}

\begin{example}\label{ex:splitsmall}
At $n=2$ the formula reads
\[
  \omega_{1}=d\,\ell^{2}-\gamma^{2},
  \qquad
  \omega_{2}=2\,\gamma\,\ell,
\]
and at $n=3$,
\[
  \omega_{1}=\gamma\,\ell^{2}-\tfrac13\gamma^{3},
  \qquad
  \omega_{2}=\gamma^{2}\ell-\tfrac13\,d\,\ell^{3},
\]
with $\gamma=d\beta-\widehat{\beta}$ throughout. Both cases, and the expansion
for every $n\le5$ and several $d$, are confirmed in
\texttt{split\_locus.py}; see item~(X) of \Cref{sec:verification}.
\end{example}

\begin{corollary}[Algebraicity on the split locus]\label{cor:splitalgebraic}
For every $n\ge1$, every squarefree $d>0$ and every principally polarised
abelian $n$-fold $X$, the Weil classes of $A=X\times\widehat{X}$ with the
$K$-action \eqref{eq:Mdef} are algebraic, and an algebraic representative is
given by \eqref{eq:splitclosed} as an explicit rational combination of
intersections of three divisors.
\end{corollary}

\begin{proof}
\Cref{thm:splitclosed}(iii), together with the Lefschetz theorem on $(1,1)$
classes, which makes $\beta,\widehat{\beta},\ell$ algebraic, and the fact that
the algebraic classes form a subring of $H^{*}(A,\QQ)$.
\end{proof}

\begin{remark}\label{rem:splitconsistent}
\Cref{cor:splitalgebraic} does not contradict \Cref{thm:character}. At the
split point $\NS(A)\otimes\QQ$ contains the three independent classes
$\beta,\widehat{\beta},\ell$, so the hypothesis
$\NS(A)\otimes\QQ=\QQ\eta$ of \Cref{prop:notpoly} fails, and this is the
situation of \Cref{rem:special}. What \Cref{thm:character} forbids is a
$\Kx$-eigenclass for the norm character meeting the Weil line, and
$\gamma\mp\sqrt{-d}\,\ell$ is an eigenclass for $\tau^{2}$ and for
$\bbar{\tau}^{2}$, not for the norm. Indeed \eqref{eq:splitclosed} is the
first place in this paper where a class in the right eigenspace is exhibited
at all, and it is exhibited only because the split member has enough divisors
to build one.
\end{remark}

\subsection{Products, and a base case in every dimension}\label{ssec:products}

The split member is one way to produce a member with enough divisors. There is
a second, and it behaves well under the one operation the Weil line respects.
Throughout this subsection $K$ is fixed and all Weil structures are for that
$K$.

\begin{lemma}[Products of Weil type]\label{lem:prodweil}
Let $(A_{1},\iota_{1},\eta_{1})$ be of $(K,-1,n_{1})$-Weil type and
$(A_{2},\iota_{2},\eta_{2})$ of $(K,-1,n_{2})$-Weil type. Then
$A=A_{1}\times A_{2}$, with the diagonal action
$\iota=\iota_{1}\times\iota_{2}$ and the product polarisation
$\eta=p_{1}^{*}\eta_{1}+p_{2}^{*}\eta_{2}$, is of
$(K,-1,n_{1}+n_{2})$-Weil type.
\end{lemma}

\begin{proof}
Write $V_{i}=H^{1}(A_{i},\QQ)$, so $V=H^{1}(A,\QQ)=V_{1}\oplus V_{2}$ as
$K$-vector spaces and $\dim_{K}V=2n_{1}+2n_{2}$. The Rosati involution of the
product polarisation is the product of the two Rosati involutions, so it
restricts on $K$ to conjugation, which is the sign $-1$ of
\Cref{def:weiltype}. For the multiplicities, the Hodge decomposition of $V$ is
the direct sum of those of $V_{1}$ and $V_{2}$, and the eigenspace
decomposition is likewise $V_{\pm}=V_{1,\pm}\oplus V_{2,\pm}$. Hence
\[
  \dim_{\CC}V_{+}^{1,0}
  =\dim_{\CC}V_{1,+}^{1,0}+\dim_{\CC}V_{2,+}^{1,0}
  =n_{1}+n_{2},
\]
and the same for the other three, which is \Cref{lem:balanced} for
$n=n_{1}+n_{2}$.
\end{proof}

The Weil line is the top exterior power over $K$, and a top exterior power of
a direct sum is the tensor product of the top exterior powers. That is the
whole content of the next theorem; what has to be checked is what it does to
the rational generators, since those are defined through the real structure
and not over $K$.

By \Cref{prop:weilline}(i) and (iii) the Weil line $\HW(A)$ is free of rank
one over $K$, and the basis $\omega_{1},\omega_{2}$ of
\Cref{thm:explicitgens} identifies $\HW(A)$ with $K$ itself by
$a\omega_{1}+b\omega_{2}\mapsto a+b\sqrt{-d}$. Write
$\omega(A)=\omega_{1}(A)+\sqrt{-d}\,\omega_{2}(A)$ for the resulting
$K$-valued invariant.

\begin{theorem}[The Weil class is multiplicative]\label{thm:weilmult}
With the notation of \Cref{lem:prodweil}, and cup product on
$H^{*}(A,\QQ)$ taken after pullback along the two projections,
\begin{equation}\label{eq:weilmult}
\begin{aligned}
  \omega_{1}(A)&=\tfrac12\bigl(\omega_{1}(A_{1})\,\omega_{1}(A_{2})
    -\omega_{2}(A_{1})\,\omega_{2}(A_{2})\bigr),\\
  \omega_{2}(A)&=\tfrac12\bigl(\omega_{1}(A_{1})\,\omega_{2}(A_{2})
    +\omega_{2}(A_{1})\,\omega_{1}(A_{2})\bigr),
\end{aligned}
\end{equation}
that is $\omega(A_{1})\,\omega(A_{2})=2\,\omega(A_{1}\times A_{2})$ in the
$K$-valued notation. For $k$ factors,
$\prod_{j}\omega(A_{j})=2^{k-1}\,\omega(A_{1}\times\cdots\times A_{k})$.
\end{theorem}

\begin{proof}
Let $\alpha_{\pm}^{i}$ generate $\bigwedge^{2n_{i}}V_{i,\pm}$. Since
$V_{\pm}=V_{1,\pm}\oplus V_{2,\pm}$,
\[
  \textstyle\bigwedge^{2n}V_{+}=\bigwedge^{2n_{1}}V_{1,+}\otimes
  \bigwedge^{2n_{2}}V_{2,+},
\]
so $\alpha_{+}=\alpha_{+}^{1}\wedge\alpha_{+}^{2}$ generates it, and likewise
$\alpha_{-}=\alpha_{-}^{1}\wedge\alpha_{-}^{2}$; the classes
$\alpha_{\pm}^{i}$ have even degree $2n_{i}$, so no sign enters when the two
blocks of generators are interleaved. By
\Cref{thm:explicitgens}(iii), $\omega_{1}=\alpha_{+}+\alpha_{-}$ and
$\omega_{2}=\sqrt{-d}\,(\alpha_{+}-\alpha_{-})/(-d)$ up to the normalisation
fixed there, so it suffices to expand. Writing $\omega_{1}^{i}$ and
$\omega_{2}^{i}$ for the generators of the factors,
\[
  \omega_{1}^{1}\omega_{1}^{2}-\omega_{2}^{1}\omega_{2}^{2}
  =(\alpha_{+}^{1}+\alpha_{-}^{1})(\alpha_{+}^{2}+\alpha_{-}^{2})
  +(\alpha_{+}^{1}-\alpha_{-}^{1})(\alpha_{+}^{2}-\alpha_{-}^{2})
  =2(\alpha_{+}^{1}\alpha_{+}^{2}+\alpha_{-}^{1}\alpha_{-}^{2})
  =2\,\omega_{1}(A),
\]
the cross terms cancelling in pairs, and the same computation with one factor
conjugated gives the second identity of \eqref{eq:weilmult}. In the $K$-valued
notation the two identities together say
$\omega(A_{1})\omega(A_{2})=2\omega(A)$, since the map
$a\omega_{1}+b\omega_{2}\mapsto a+b\sqrt{-d}$ carries the right-hand sides of
\eqref{eq:weilmult} to the real and imaginary parts of the product. The
statement for $k$ factors follows by induction, each application contributing
one factor of $2$. The identities are verified in exact arithmetic over
$\QQ(\sqrt{-d})$ for all $n_{1},n_{2}\le3$ and for the four-factor products
listed in \Cref{sec:verification}, item (XV).
\end{proof}

\begin{corollary}[Algebraicity propagates along products]\label{cor:prodalg}
If the Weil classes of $A_{1}$ and of $A_{2}$ are algebraic, then so are those
of $A_{1}\times A_{2}$, and an algebraic representative is the rational
combination of external products given by \eqref{eq:weilmult}.
\end{corollary}

\begin{proof}
Algebraic classes are closed under pullback along the projections and under
cup product, and \eqref{eq:weilmult} expresses $\omega_{1}(A)$ and
$\omega_{2}(A)$ in those operations applied to the $\omega_{j}(A_{i})$.
\end{proof}

The base case costs nothing.

\begin{lemma}[Dimension two]\label{lem:basecase}
Let $A_{0}$ be of $(K,-1,1)$-Weil type, so $\dim A_{0}=2$. Then
$\HW(A_{0})\subset H^{2}(A_{0},\QQ)$ and both generators are algebraic.
\end{lemma}

\begin{proof}
The Weil classes lie in $H^{2n}=H^{2}$ and are of type $(1,1)$ by
\Cref{prop:weilline}(ii), so they are algebraic by the Lefschetz theorem on
$(1,1)$ classes. No hypothesis on $A_{0}$ beyond the Weil structure is used.
\end{proof}

\begin{theorem}[An algebraic Weil class in every dimension]\label{thm:everyn}
Let $n\ge1$, let $d>0$ be squarefree, and let $A_{1},\dots,A_{n}$ be abelian
surfaces of $(K,-1,1)$-Weil type. Then $A=A_{1}\times\cdots\times A_{n}$ is of
$(K,-1,n)$-Weil type and its Weil classes are algebraic; explicitly, by
\Cref{thm:weilmult} and \Cref{lem:basecase} they are
$2^{-(n-1)}$ times the real and imaginary parts of
$\prod_{j=1}^{n}\bigl(\omega_{1}(A_{j})+\sqrt{-d}\,\omega_{2}(A_{j})\bigr)$,
a polynomial with integer coefficients in $2n$ divisor classes.
\end{theorem}

\begin{proof}
\Cref{lem:prodweil} gives the Weil structure, \Cref{lem:basecase} gives
algebraicity of each factor's classes, and \Cref{cor:prodalg} applied $n-1$
times gives the conclusion; the displayed form is \Cref{thm:weilmult} for $k=n$
factors. Each $\omega_{j}(A_{i})$ is a divisor class by
\Cref{lem:basecase}, so the product is a polynomial in $2n$ of them.
\end{proof}

\begin{remark}[What this locus is, and what it is not]\label{rem:prodscope}
\Cref{thm:everyn} is unconditional and holds at every $n$, and it is the only
statement in this paper that produces an algebraic Weil class in every
dimension without quoting anything. It should not be mistaken for more than it
is. The family of polarised abelian $2n$-folds of $(K,-1,n)$-Weil type has
dimension $n^{2}$, and the locus of products of $n$ surfaces has dimension
$n$, so for $n\ge2$ it is a proper closed subvariety, and a smaller one than
the quaternionic locus of \Cref{thm:quatweil}, whose dimension is
$n(n+1)/2$. A general member of the family is simple, so it is not isogenous
to a product at all, and neither \Cref{thm:everyn} nor
\Cref{cor:splitalgebraic} touches it. What the product construction adds is a
second family of base points, defined for every $n$ and every $d$, at which the
class is not merely known to be algebraic but is written down; and it adds the
observation that the algebraicity locus is closed under the product operation,
so that any future case $n_{0}$ propagates to every multiple of $n_{0}$ and to
every sum of solved cases. Combined with Markman's theorem at $n=2$
\cite{Mar25a,Mar25b} this gives, at every $n$, algebraicity on the locus of
products of abelian fourfolds and surfaces of Weil type, of dimension $2n$ when
$n$ is even. The general member remains outside every one of these loci, and
by \Cref{thm:residual} the gap between them and the general member is exactly
the propagation statement that this paper does not prove.
\end{remark}

\begin{figure}[htbp]
\centering
  \includegraphics[width=0.96\textwidth]{fig_product}
\caption{\Cref{thm:weilmult} and its reach. Left: the Weil line of each factor
is a $K$-line, the Weil line of the product is the tensor product of the two
over $K$, and in the $K$-valued notation the class is multiplicative up to the
factor $2$. Algebraicity therefore passes from the factors to the product, and
at $n=1$ the factors are divisor classes, so a product of $n$ abelian surfaces
of Weil type carries an explicit algebraic Weil class in every dimension.
Right: the dimensions of the loci on which the class is known to be algebraic,
against the dimension $n^{2}$ of the family. Both the quaternionic locus and
the product locus are proper subvarieties once $n\ge2$, and a general member
lies on neither.}
\label{fig:product}
\end{figure}

\subsection{The residual statement}\label{ssec:residual}

Everything above lets the remaining question be stated in one sentence with no
auxiliary hypotheses, and we state it, because a problem in that form is
worth more than a programme.

\begin{setupenv}[The family]\label{setup:family}
Fix $n\ge2$ and a squarefree $d>0$. Let $\cD_{n,n}$ be the period domain of
\Cref{prop:weildomain} and let
\[
  \pi\colon\cA\longrightarrow\cD_{n,n}
\]
be the family of polarised abelian $2n$-folds of $(K,-1,n)$-Weil type it
parametrises, with $\omega$ the flat section of $R^{2n}\pi_{*}\QQ$ given
fibrewise by the class $\omega_{1}$ of \eqref{eq:omegagens}. Let
$\cD^{\mathrm{spl}}\subset\cD_{n,n}$ be the locus of split members
$X\times\widehat{X}$ carrying the action \eqref{eq:Mdef}.
\end{setupenv}

\begin{question}[The residual statement]\label{q:residual}
Is $\omega_{s}$ algebraic for every $s\in\cD_{n,n}$?
\end{question}

\begin{theorem}[The reduction]\label{thm:residual}
Fix $n\ge2$ and a squarefree $d>0$, and keep the notation of
\Cref{setup:family}.
\begin{enumerate}[label=\textup{(\alph*)},nosep]
\item If \Cref{q:residual} has a positive answer, then the Hodge conjecture
  holds, in every degree, for every abelian $2n$-fold of $(K,-1,n)$-Weil type
  with $\Hg(A)=\SU(V,H)$.
\item Conversely, if the Hodge conjecture holds for every abelian $2n$-fold of
  $(K,-1,n)$-Weil type, then \Cref{q:residual} has a positive answer.
\item $\cD^{\mathrm{spl}}$ is nonempty, $\omega_{s}$ is algebraic at every
  $s\in\cD^{\mathrm{spl}}$ with the explicit representative
  \eqref{eq:splitclosed}, and $\cD_{n,n}$ is connected.
\end{enumerate}
So \Cref{q:residual} is the variational Hodge conjecture for one flat class on
one connected family, it is necessary as well as sufficient, and its base case
is supplied. \Cref{thm:everyfamily} supplies one in every family, for every
discriminant, so the hypothesis of nonemptiness is never the obstacle.
\end{theorem}

\begin{proof}
(a) Let $A$ be of $(K,-1,n)$-Weil type with $\Hg(A)=\SU(V,H)$. By
\Cref{prop:weildomain} its Hodge structure is a point $s$ of $\cD_{n,n}$, so
$A$ is a fibre of $\pi$, and the Weil line of $A$ is generated over $K$ by
$\omega_{s}$. A positive answer makes $\omega_{s}$ algebraic, hence makes the
whole Weil line algebraic by the argument of \Cref{cor:whatisleft}, and by
\Cref{thm:hdgcount} the Weil line together with the powers of $\eta$ exhausts
$\Hdg^{*}(A)$.

(b) For every $s$ the class $\omega_{s}$ is a Hodge class on the fibre $A_{s}$,
by \Cref{prop:weilline}(ii), and $A_{s}$ is of $(K,-1,n)$-Weil type by
construction. So the Hodge conjecture for $A_{s}$ gives its algebraicity.

(c) $\cD^{\mathrm{spl}}$ is nonempty because the construction \eqref{eq:Mdef}
applies to any principally polarised abelian $n$-fold and produces a point of
$\cD_{n,n}$ by \Cref{thm:explicitweil}; algebraicity there is
\Cref{cor:splitalgebraic}; and $\cD_{n,n}$ is connected because it is a
bounded symmetric domain, by \Cref{prop:weildomain}.
\end{proof}

\begin{remark}[Why this is the precise target]\label{rem:target}
Three things have been removed from the problem by the preceding sections, and
we list what is left after each.

\Cref{thm:hdgcount} removes the question of how many classes are at issue:
one, in one degree, on each fibre. \Cref{thm:divisor} and
\Cref{thm:character} remove two families of would-be constructions, so no
effort need be spent on them. \Cref{thm:splitclosed} removes the base case: the
class is algebraic somewhere in the family, with a representative one can
write on a single line.

What is left is propagation, and it is exactly Grothendieck's variational
Hodge conjecture for the pair $(\pi,\omega)$ of \Cref{setup:family}. That
statement is open in general; what \Cref{thm:residual} adds is that for the
Hodge conjecture on abelian varieties of Weil type it is not merely sufficient
but necessary, so no other route can avoid an equivalent difficulty.
\Cref{ssec:properness} takes the propagation apart. The passage from a formal
deformation to an algebraic one, which is usually the awkward half, is
supplied there by properness of the Hilbert scheme, and what survives is the
infinitesimal half at a single point, or else a bound on the size of the
cycles already constructed.
\end{remark}

\subsection{The divisorial criterion, and exactly what it reaches}
\label{ssec:divcrit}

\Cref{thm:splitclosed} is a computation on one member of the family. The
mechanism behind it is not special to that member, and isolating the mechanism
turns the base case into a criterion that can be tested at any point. The
criterion, and the precise determination of where it holds, is the content of
this subsection.

\begin{theorem}[The divisorial criterion]\label{thm:divcriterion}
Let $(A,\iota,\eta)$ be of $(K,-1,n)$-Weil type and suppose there exists
$\Theta\in\bigwedge^{2}V_{-}$ with
\begin{enumerate}[label=\textup{(\roman*)},nosep]
\item $\Theta$ nondegenerate as an alternating form on $V_{-}^{\vee}$,
  equivalently $\Theta^{n}\ne0$;
\item $\Theta$ of Hodge type $(1,1)$;
\item $\Theta+\bbar{\Theta}$ rational, that is in $H^{2}(A,\QQ)$.
\end{enumerate}
Then the Weil classes of $A$ are algebraic. Explicitly, writing
$g=\Theta+\bbar{\Theta}$ and $\lambda=(\Theta-\bbar{\Theta})/\sqrt{-d}$, both
$g$ and $\lambda$ are divisor classes and
\begin{equation}\label{eq:divrep}
  \HW(A)=\Bigl\langle\;
    \Re\bigl(g+\sqrt{-d}\,\lambda\bigr)^{n},\;\;
    \Im\bigl(g+\sqrt{-d}\,\lambda\bigr)^{n}\;\Bigr\rangle_{\QQ},
\end{equation}
so every Weil class is a rational polynomial in two divisor classes.
\end{theorem}

\begin{proof}
Conjugation interchanges $\bigwedge^{2}V_{-}$ and $\bigwedge^{2}V_{+}$, so
$\bbar{\Theta}\in\bigwedge^{2}V_{+}$, and both $g$ and $\lambda$ are fixed by
conjugation, hence real. Rationality of $g$ is (iii); rationality of
$\lambda$ needs an argument, and it is the following. Fix
$\tau=a+b\sqrt{-d}\in\Kx$ with $ab\ne0$. By \Cref{prop:grading} the operator
$\iota(\tau)^{*}$ acts on $\bigwedge^{2}V_{+}$ by $\tau^{2}=a^{2}-db^{2}+2ab\sqrt{-d}$
and on $\bigwedge^{2}V_{-}$ by $\bbar{\tau}^{\,2}$, so
\[
  J:=\frac{\iota(\tau)^{*}-(a^{2}-db^{2})}{2ab}
\]
is a $\QQ$-rational endomorphism of $H^{2}(A,\QQ)$ preserving $R$ and acting
on $\bigwedge^{2}V_{\pm}$ by $\pm\sqrt{-d}$. Applying it to
$g=\Theta+\bbar{\Theta}$ gives $J(g)=-\sqrt{-d}\,\Theta+\sqrt{-d}\,\bbar{\Theta}
=-\sqrt{-d}\cdot\sqrt{-d}\,\lambda=d\,\lambda$, so $\lambda=J(g)/d$ is
rational. Both $g$ and $\lambda$ are of type $(1,1)$ by (ii) and conjugation
symmetry, so both lie in $\NS(A)\otimes\QQ$ by the Lefschetz theorem on
$(1,1)$ classes, and both are algebraic.

By (i) the form $\Theta$ on the $2n$-dimensional space $V_{-}$ has nonzero
Pfaffian, so $\Theta^{n}$ is a nonzero element of the line
$\bigwedge^{2n}V_{-}$, and it therefore spans it; conjugating,
$\bbar{\Theta}^{\,n}$ spans $\bigwedge^{2n}V_{+}$. Hence
$\Theta^{n}+\bbar{\Theta}^{\,n}$ and
$(\Theta^{n}-\bbar{\Theta}^{\,n})/\sqrt{-d}$ span $\HW(A)$ over $\QQ$, by
\Cref{prop:weilline}(i). Finally
$2^{n}\Theta^{n}=(g+\sqrt{-d}\lambda)^{n}$, and expanding by the binomial
theorem with $(\sqrt{-d})^{2m}=(-d)^{m}$ writes the real and imaginary parts
as rational polynomials in $g$ and $\lambda$. Products of divisor classes are
algebraic.
\end{proof}

\begin{corollary}\label{cor:splitsatisfies}
The split member of \Cref{thm:splitclosed} satisfies the criterion, with
$\Theta=\gamma-\sqrt{-d}\,\ell$ and $\gamma=d\beta-\widehat{\beta}$; the two
divisor classes of \eqref{eq:divrep} are $2\gamma$ and $-2\ell$.
\end{corollary}

\begin{proof}
\Cref{thm:splitclosed}(i) gives $\Theta\in\bigwedge^{2}V_{+}$ for the sign
$\gamma-\sqrt{-d}\ell$, and the conjugate lies in $\bigwedge^{2}V_{-}$;
$\gamma$ and $\ell$ are divisor classes, hence of type $(1,1)$, so $\Theta$ is
of type $(1,1)$ and $\Theta+\bbar{\Theta}=2\gamma$ is rational.
Nondegeneracy is the nonvanishing of $\Theta^{n}$, proved in
\Cref{thm:splitclosed}(ii).
\end{proof}

So the base case is one instance of a general criterion. The next two results
say exactly how far the criterion reaches, and both are negative in a way that
is worth knowing.

\begin{proposition}[The criterion fails at a general point]\label{prop:divfails}
Let $(A,\iota,\eta)$ be of $(K,-1,n)$-Weil type with $n\ge2$ and
$\Hg(A)=\SU(V,H)$. Put
\[
  R=\Bigl(\textstyle\bigwedge^{2}V_{+}\oplus\bigwedge^{2}V_{-}\Bigr)
      \cap H^{2}(A,\QQ).
\]
Then $R\cap\NS(A)_{\QQ}=0$, and no $\Theta$ as in \Cref{thm:divcriterion}
exists.
\end{proposition}

\begin{proof}
By \Cref{prop:normpart} the space $R$ is a rational sub-Hodge structure of
$H^{2}(A,\QQ)$ admitting no proper nonzero rational sub-Hodge structure, and
by \Cref{prop:h2pieces} it has $h^{2,0}=2\binom{n}{2}=n(n-1)$, which is
positive for $n\ge2$. The intersection
$R\cap\NS(A)_{\QQ}$ is a rational sub-Hodge structure of $R$, so it is $0$ or
all of $R$; it is of pure type $(1,1)$, and $R$ is not, so it is $0$. A
$\Theta$ as in \Cref{thm:divcriterion} would give the nonzero element
$\Theta+\bbar{\Theta}$ of $R\cap\NS(A)_{\QQ}$.
\end{proof}

\begin{proposition}[What the split locus is]\label{prop:splitgeom}
Fix $n\ge1$ and a squarefree $d>0$, and let $A=X\times\widehat{X}$ carry the
$K$-action \eqref{eq:Mdef} determined by a nondegenerate
$\beta\in\NS(X)\otimes\QQ$ of elementary divisors $(d_{1},\dots,d_{n})$.
\begin{enumerate}[label=\textup{(\roman*)},nosep]
\item Let $B$ be the Gram matrix of $\beta$ in a basis of $U$, and let
  $\widehat\beta$ be the form on $U^{\vee}$ with Gram matrix $-B^{-1}$ in the
  dual basis, which for principal $\beta$ is the class of
  \Cref{not:threeclasses}. The polarisation of $A$ is
  $\eta=\beta+\widehat{\beta}/d$, a positive rational multiple of
  $d\beta+\widehat\beta$; its Gram matrix on $H_{1}(A,\QQ)=U\oplus U^{\vee}$
  is the block form $Q=B\oplus(-B^{-1}/d)$, and $M$ is anti-self-adjoint for
  it.
\item In the $K$-basis $u_{1},\dots,u_{2n}$ of $H_{1}(A,\QQ)$ the hermitian
  form of \Cref{lem:hermform} has Gram matrix $\sqrt{-d}\,B$, so
  \[
    \det H=(-d)^{n}\det B=(-d)^{n}(d_{1}\cdots d_{n})^{2},
  \]
  and $(-1)^{n}\det H=d^{\,n}(d_{1}\cdots d_{n})^{2}$ is a norm from $K$.
  Every split member therefore lies in the Weil family of trivial
  discriminant, whatever $\beta$ is chosen.
\item The split members form a locus of dimension $n(n+1)/2$ inside
  $\cD_{n,n}$, of codimension $n(n-1)/2$, which is zero only for $n=1$ and is
  a divisor for $n=2$.
\item For $X$ with $\Hg(X)=\mathrm{Sp}(H^{1}(X,\QQ))$ one has
  $\NS(A)_{\QQ}=\langle\beta,\widehat{\beta},\ell\rangle$, of rank three, and
  in degree $2n$ the monomials in those three classes span the whole space of
  Hodge classes of $A$.
\end{enumerate}
\end{proposition}

\begin{proof}
(i) In the chosen basis of $U$ and the dual basis of $U^{\vee}$ the map
$\beta\colon U\to U^{\vee}$ has matrix $B^{\top}=-B$ and its inverse has
matrix $-B^{-1}$, so \eqref{eq:Mdef} becomes
$M=\begin{psmallmatrix}0&B^{-1}\\-dB&0\end{psmallmatrix}$, with
$M^{2}=-d$ and
$M^{\top}=\begin{psmallmatrix}0&dB\\-B^{-1}&0\end{psmallmatrix}$, using
$B^{\top}=-B$ and $(B^{-1})^{\top}=-B^{-1}$. For a block form
$Q=\begin{psmallmatrix}P&\Lambda\\-\Lambda^{\top}&S\end{psmallmatrix}$ the
Rosati relation $M^{\top}Q+QM=0$ reads
$B\Lambda^{\top}+\Lambda B=0$, $dBS+PB^{-1}=0$, $B^{-1}P+dSB=0$ and
$B^{-1}\Lambda+\Lambda^{\top}B^{-1}=0$. Taking $\Lambda=0$ and $P=B$ forces
$S=-B^{-1}/d$ from the second relation, and the third then reads $-I+I=0$.
So $Q=B\oplus(-B^{-1}/d)$, which is the Gram matrix of
$\beta+\widehat\beta/d$. For principal $\beta$ one may take
$B=\begin{psmallmatrix}0&I\\-I&0\end{psmallmatrix}$, so that $-B^{-1}=B$
and $\widehat\beta=\sum_{j}\xi_{j}\wedge\xi_{n+j}$ as in
\Cref{not:threeclasses}. The computation is repeated entry by entry for
$n\le4$ in item (XI) of \Cref{sec:verification}, together with the check
that the induced action on $H^{2}(A,\QQ)$ multiplies $d\beta+\widehat\beta$
by $d$ and $\gamma$, $\ell$ by $-d$, and that $\beta+d\widehat\beta$ is not
an eigenclass when $d>1$.

(ii) The set $u_{1},\dots,u_{2n}$ is a $K$-basis because
$MU=d\beta(U)=U^{\vee}$, and $H(u_{i},u_{j})=Q(Mu_{i},u_{j})+\sqrt{-d}\,Q(u_{i},u_{j})$
by \eqref{eq:QfromH}; the first term vanishes since $Mu_{i}\in U^{\vee}$,
$u_{j}\in U$ and $Q$ has no cross block by (i), and the second is
$\sqrt{-d}\,\beta_{ij}$. Hence the Gram matrix is $\sqrt{-d}\,B$ and
$\det H=(\sqrt{-d})^{2n}\det B=(-d)^{n}\det B$. The determinant of an
alternating matrix is the square of its Pfaffian, here $(d_{1}\cdots
d_{n})^{2}$. Now $d=\Nm_{K/\QQ}(\sqrt{-d})$ is a norm and every square is a
norm, and the norms form a subgroup of $\QQ^{\times}$, so
$d^{\,n}(d_{1}\cdots d_{n})^{2}$ is a norm.

(iii) A split member is determined by $X$ together with $\beta$, and the
polarised moduli of $X$ has dimension $\dim\cA_{n}=n(n+1)/2$ while the type of
$\beta$ is discrete. By \Cref{prop:weildomain} the ambient domain has
dimension $n^{2}$, and $n^{2}-n(n+1)/2=n(n-1)/2$.

(iv) The Hodge classes of $A=X\times\widehat{X}$ are the invariants of
$\Hg(A)=\Hg(X)=\mathrm{Sp}(U)$ acting on
$\bigwedge^{*}(U\oplus U^{\vee})$. In degree two the invariants are the three
lines spanned by $\beta$, $\widehat{\beta}$ and $\ell$, one from each Kunneth
summand $\bigwedge^{2}U$, $U\otimes U^{\vee}$ and $\bigwedge^{2}U^{\vee}$.
In degree $2n$ the dimension of the invariant space is computed from the
primitive decomposition $\bigwedge^{a}U=\bigoplus_{j}\bigwedge^{a-2j}_{0}U$
together with Schur's lemma, and it agrees with the rank of the set of
monomials $\beta^{i}\widehat{\beta}^{j}\ell^{k}$ with $i+j+k=n$; both numbers
are computed for $n\le3$ in item (XI) of \Cref{sec:verification} and agree.
\end{proof}

\begin{remark}[Where this leaves the residual statement]\label{rem:residualmap}
Putting \Cref{prop:splitgeom}(ii) beside \Cref{thm:residual} changes the shape
of what is left. The Weil families for a fixed $K$ and fixed $n$ are
classified by the discriminant of the hermitian form, an element of
$\QQ^{\times}_{>0}/\Nm_{K/\QQ}(K^{\times})$, which is an infinite group. The
split construction of \Cref{ssec:split} produces a base point in exactly one of
them, the one with trivial discriminant; products of abelian surfaces of Weil
type (\Cref{thm:everyn}, \Cref{lem:products}) reach every discriminant, but
only on a locus of dimension $n$. So \Cref{q:residual} is really two
questions, and they are of different kinds.
\begin{enumerate}[label=\textup{(R\arabic*)},leftmargin=2.8em,itemsep=2pt]
\item \emph{Trivial discriminant.} A base point exists, with the explicit
  cycle \eqref{eq:splitclosed}, and what is missing is propagation off a locus
  of codimension $n(n-1)/2$. At $n=2$ that locus is a divisor and the
  propagation is known, by Markman \cite{Mar25a,Mar25b} and by Floccari and Fu
  \cite{FF26}; combined with \Cref{thm:residual}(a) this recovers the Hodge
  conjecture for abelian fourfolds of Weil type of trivial discriminant with
  full Hodge group. At $n=3$ it is settled as well, by Markman's theorem on
  the sixfolds of trivial discriminant \cite{Mar25a,Mar25b}, for every $K$.
  For $n\ge4$ it is open, and by \Cref{prop:descent} it is then the only
  question left in the imaginary quadratic case.
\item \emph{Nontrivial discriminant.} The construction of this subsection
  never lands there, by \Cref{prop:splitgeom}(ii), so it supplies no base
  point; and by \Cref{prop:divfails} no divisorial argument can supply one at
  a general member. A base point must come from a special member at which $R$
  acquires Hodge classes.
\end{enumerate}
Question (R2) is settled in \Cref{ssec:everyfamily}: the special members
required are those with quaternionic multiplication extending $K$, they exist
in every family, and on them \Cref{thm:divcriterion} applies. After that only
(R1) remains, now with a base point in every family.
\end{remark}

\begin{remark}[The split locus is larger than a Hodge locus has to be]
\label{rem:nlcodim}
A component of the Noether--Lefschetz locus of a weight two variation has
codimension at most $h^{2,0}$, here $n(n-1)$. The split locus has codimension
$n(n-1)/2$ by \Cref{prop:splitgeom}(iii), exactly half of that, so it is a
component of twice the expected dimension. That is a reflection of
\Cref{prop:splitgeom}(iv): at a split point $R$ acquires two independent Hodge
classes at once, namely $\gamma$ and $\ell$, and each of them costs
$h^{2,0}/2$ conditions rather than $h^{2,0}$. Whether there are components of
the expected codimension carrying only one class, and whether the criterion of
\Cref{thm:divcriterion} can be met on them, is the concrete form of question
(R2).
\end{remark}

\subsection{A base point in every family}\label{ssec:everyfamily}

\Cref{prop:splitgeom}(ii) says the split construction reaches one discriminant
class. \Cref{thm:divcriterion} does not care where its $\Theta$ comes from,
and the source used below is an endomorphism rather than a pair of divisors
chosen by hand. It reaches every class, and with it question (R2) of
\Cref{rem:residualmap} is settled.

\begin{definition}[Discriminant]\label{def:disc}
For $(A,\iota,\eta)$ of $(K,-1,n)$-Weil type with hermitian form $H$ as in
\Cref{lem:hermform}, the \emph{discriminant} is the class of
$(-1)^{n}\det H$ in $\QQ^{\times}_{>0}/\Nm_{K/\QQ}(K^{\times})$. The sign
makes the representative positive, because $H$ has signature $(n,n)$; and the
class is well defined because changing the $K$-basis multiplies $\det H$ by a
norm. Hermitian forms of rank $2n$ over $K$ with signature $(n,n)$ are
classified by this class, so it labels the Weil families.
\end{definition}

The mechanism behind \Cref{thm:divcriterion} can be named. A class in $R$ is
the same thing as an endomorphism that conjugates $K$.

\begin{lemma}[What a class in $R$ is]\label{lem:Rendomorphism}
Let $(A,\iota,\eta)$ be of $(K,-1,n)$-Weil type. The map $\mu\mapsto
g_{\mu}$, $g_{\mu}(x,y)=\eta(\mu x,y)$, is a bijection between
\begin{enumerate}[label=\textup{(\alph*)},nosep]
\item the elements $\mu\in\End^{0}(A)$ with $\mu^{\dagger}=\mu$ and
  $\mu\,\iota(\tau)=\iota(\bbar{\tau})\,\mu$ for all $\tau\in K$, and
\item the elements of $R\cap\NS(A)_{\QQ}$.
\end{enumerate}
Under it, $\mu$ invertible corresponds to $g_{\mu}$ nondegenerate, hence to a
nondegenerate $\Theta$. If such a $\mu\ne0$ exists and the centraliser of
$\iota(K)$ in $\End^{0}(A)$ is $\iota(K)$, then $\mu^{2}\in\QQ^{\times}$ and
$\iota(K)+\iota(K)\mu$ is a quaternion algebra $(-d,\mu^{2})_{\QQ}$ inside
$\End^{0}(A)$.
\end{lemma}

\begin{proof}
A class in $\NS(A)_{\QQ}$ is a symmetric homomorphism $A\to\Av$ up to isogeny,
that is $\eta\mu$ with $\mu^{\dagger}=\mu$. The condition
$g_{\mu}(\iota(\tau)x,y)=g_{\mu}(x,\iota(\tau)y)$, which says that $g_{\mu}$ is
$K$-bilinear and hence lies in $R=\bigwedge^{2}_{K}V$, reads
$\eta(\mu\iota(\tau)x,y)=\eta(\mu x,\iota(\tau)y)=\eta(\iota(\tau)^{\dagger}\mu
x,y)=\eta(\iota(\bbar{\tau})\mu x,y)$, using \Cref{def:weiltype}(iii); so it is
exactly $\mu\iota(\tau)=\iota(\bbar{\tau})\mu$. For the last sentence,
$\mu^{2}$ commutes with $\iota(K)$, hence lies in $\iota(K)$ by hypothesis, and
$\mu\mu^{2}\mu^{-1}=\mu^{2}$ forces it to be fixed by conjugation, so
$\mu^{2}\in\QQ^{\times}$. Then $\iota(\sqrt{-d})$ and $\mu$ generate the
quaternion algebra with the stated invariants.
\end{proof}

So the criterion asks for quaternionic multiplication extending $K$. Such
abelian varieties exist in abundance, and their invariants can be computed.

\begin{setupenv}[The quaternionic model]\label{setup:quat}
Let $b\in\QQ_{>0}$ and let $B=(-d,b)_{\QQ}$ be the quaternion algebra with
basis $1,i,j,ij$ and $i^{2}=-d$, $j^{2}=b$, $ij=-ji$. Let
$V=B^{n}$ as a left $B$-module, of dimension $4n$ over $\QQ$, let
$K=\QQ(i)\subset B$ act through $i$, and fix $a_{1},\dots,a_{n}\in\QQ^{\times}$
and
\begin{equation}\label{eq:quatE}
  E(x,y)=\sum_{k=1}^{n}a_{k}\,\mathrm{trd}\bigl(\bbar{x_{k}}\,i\,y_{k}\bigr),
\end{equation}
where $\mathrm{trd}$ is the reduced trace and $\bbar{\phantom{x}}$ the
canonical involution of $B$.
\end{setupenv}

\begin{theorem}[Weil type from quaternionic multiplication]\label{thm:quatweil}
With the data of \Cref{setup:quat}:
\begin{enumerate}[label=\textup{(\roman*)},nosep]
\item $E$ is alternating and nondegenerate, and its Rosati involution is
  $\beta\mapsto i^{-1}\bbar{\beta}i$, which restricts to complex conjugation on
  $K$ and fixes $j$; it is positive, the form $\mathrm{trd}(\beta\beta^{\dagger})$
  being diagonal with entries $2,2d,2b,2db$ on the basis;
\item in the $K$-basis $e_{1},e_{1}j,\dots,e_{n},e_{n}j$ of $V$ the hermitian
  form of \Cref{lem:hermform} is
  \[
    H=\mathrm{diag}\bigl(2a_{1}d,\,-2a_{1}bd,\;\dots,\;2a_{n}d,\,-2a_{n}bd\bigr),
  \]
  of signature $(n,n)$;
\item there is a complex structure on $V_{\RR}$ commuting with $B$ and
  polarised by $E$, and for every such the resulting abelian variety $A$ is of
  $(K,-1,n)$-Weil type, the balanced signature being forced and not imposed;
\item the discriminant of $A$ is the class of $b^{n}$;
\item $\mu=$ left multiplication by $j$ satisfies \Cref{lem:Rendomorphism}(a),
  so $g=g_{\mu}$ lies in $R\cap\NS(A)_{\QQ}$ and is nondegenerate, the
  divisorial criterion \Cref{thm:divcriterion} applies, and the Weil classes
  of $A$ are algebraic.
\end{enumerate}
\end{theorem}

\begin{proof}
(i) $\mathrm{trd}(z)=\mathrm{trd}(\bbar{z})$ and
$\overline{\bbar{x}iy}=\bbar{y}\,\bbar{i}\,x=-\bbar{y}ix$ give
$E(y,x)=-E(x,y)$. Nondegeneracy is the nondegeneracy of the reduced trace
pairing. For the Rosati, $E(\beta x,y)=E(x,\beta^{\dagger}y)$ for all $x,y$
reads $\bbar{\beta}i=i\beta^{\dagger}$, that is
$\beta^{\dagger}=i^{-1}\bbar{\beta}i$. Then $i^{\dagger}=-i$ and, since
$(-j)i=ij$, also $j^{\dagger}=i^{-1}(ij)=j$. Positivity is the displayed
diagonal, which is positive because $d>0$ and $b>0$.

(ii) $V=Ke_{k}\oplus Ke_{k}j$ over all $k$, since $B=K\oplus Kj$. By
\eqref{eq:QfromH}, $H(x,y)=E(ix,y)+\sqrt{-d}\,E(x,y)$. On one factor,
$E(ix,y)=a\,\mathrm{trd}(\overline{ix}\,i\,y)=-a\,\mathrm{trd}(\bbar{x}i^{2}y)
=ad\,\mathrm{trd}(\bbar{x}y)$, so
\[
  H(x,y)=a\bigl[d\,\mathrm{trd}(\bbar{x}y)+\sqrt{-d}\,\mathrm{trd}(\bbar{x}iy)\bigr].
\]
Now $\mathrm{trd}(1)=2$, $\mathrm{trd}(i)=\mathrm{trd}(j)=\mathrm{trd}(ij)=0$,
$\bbar{j}j=-b$ and $\bbar{j}ij=ib$, which give $H(1,1)=2ad$,
$H(j,j)=-2abd$ and $H(1,j)=0$. Since $b>0$ the two entries of each block have
opposite signs, so the signature is $(n,n)$.

(iii) A complex structure commuting with $B$ and compatible with $E$ is the
same as a subspace $V^{1,0}\subset V_{\CC}$, stable under $B$, with
$V_{\CC}=V^{1,0}\oplus\overline{V^{1,0}}$ and the two Riemann relations. Being
$K$-stable it splits as $P\oplus P'$ with $P\subset V_{+}$ and
$P'\subset V_{-}$; left multiplication by $j$ is $\sigma$-semilinear, so it
carries $V_{+}$ to $V_{-}$, and $B$-stability forces $P'=jP$. Conversely, for
any $n$-dimensional $H$-negative definite $P\subset V_{+}$ the subspace
$P\oplus jP$ is $B$-stable and satisfies the Riemann relations, by the
computation in the proof of \Cref{thm:model} applied blockwise; such $P$
exists because $H$ has signature $(n,n)$ by (ii). Since $j$ is invertible and
carries $P$ isomorphically to $P'$, one gets $\dim P=\dim P'=n$, which is
\Cref{def:weiltype}(ii). The general theory of abelian varieties with
prescribed endomorphism structure is in \cite[Ch.~9]{BL04}.

(iv) By (ii), $\det H=\prod_{k}(2a_{k}d)(-2a_{k}bd)=(-1)^{n}b^{n}\bigl(2^{n}d^{n}
\prod_{k}a_{k}\bigr)^{2}$, so $(-1)^{n}\det H$ is $b^{n}$ times a square, and
squares are norms.

(v) Left multiplication by $j$ satisfies $\mu^{\dagger}=\mu$ by (i) and
$\mu\,\iota(\tau)=\iota(\bbar{\tau})\,\mu$ because $ji=-ij$; it is invertible
because $\mu^{2}=b\ne0$. \Cref{lem:Rendomorphism} puts $g_{\mu}$ in
$R\cap\NS(A)_{\QQ}$ and makes it nondegenerate, and
\Cref{thm:divcriterion} applies to the component $\Theta$ of $g_{\mu}$ in
$\bigwedge^{2}V_{-}$.
\end{proof}

\begin{proposition}[Nothing numerical forbids the object]%
\label{prop:notforbidden}
Let $t$ be a principal polarisation on $X$ and let
$\ch=au+bv$ be a rational point of the secant plane of \Cref{thm:plane},
with $a,b\in\ZZ$. Then
\begin{enumerate}[label=\textup{(\roman*)},nosep]
\item every $\ch_{k}$ is an integral multiple of the integral class
  $t^{k}/k!$, the multiple being $\pm a d^{j}$ in even degree $k=2j$ and
  $\pm b d^{j}$ in odd degree $k=2j+1$;
\item the Chern classes $c_{1},\dots,c_{n}$ obtained from $\ch$ by Newton's
  identities are integral, for every $n\le8$, every squarefree $d\le11$ and
  every $a\le4$, $|b|\le4$;
\item the discriminant is
  $\Delta=2rc_{2}-(r-1)c_{1}^{2}=\bigl(a^{2}d+b^{2}\bigr)t^{2}
  =\Nm_{K/\QQ}(b+a\sqrt{-d})\,t^{2}$, which is positive, so Bogomolov's
  inequality never binds.
\end{enumerate}
\end{proposition}

\begin{proof}
(i) is the definition of $u$ and $v$: the coefficient of $t^{k}$ is
$a(-d)^{j}/(2j)!$ or $b(-d)^{j}/(2j+1)!$, and multiplying by $k!$ clears the
factorial and leaves $\pm ad^{j}$ or $\pm bd^{j}$. For (ii), Newton's
identities express $k!\,c_{k}$ as an integer polynomial in the $i!\,\ch_{i}$
apart from one division by $k$ at each step, and the computation in item (XX)
of \Cref{sec:verification} carries out those divisions exactly over the stated
range and finds them exact every time. (iii) is the computation recorded in
\Cref{thm:secantinvariants}, and the norm form of an imaginary quadratic field
is positive definite.
\end{proof}

\begin{remark}[What this settles, and what it does not]%
\label{rem:notforbidden}
The content is negative and that is the point. One might hope that the object
the construction needs at $n\ge4$ is forbidden by its numbers, in which case
the secant route would be dead on arithmetic alone and the residual question
would change shape. \Cref{prop:notforbidden} says it is not: the invariants
are integral, the discriminant is positive, and nothing in Riemann--Roch or in
Bogomolov's inequality stands in the way. Taken with \Cref{prop:objectsize},
which bounds the size of the object from below without forbidding it, and with
the parameter count of \Cref{rem:markmancheck}, which fixes the rank one shape
at $n=3$, the position is that every count we know how to
make is consistent with the object existing and none of them produces it. What
is missing at $n\ge4$ is the object, not a number.
\end{remark}

\subsection{The quaternionic locus is a Lagrangian Grassmannian}
\label{ssec:lagrangian}

\Cref{thm:quatweil} produces a point of the family for each admissible
$\psi$, and the footnote to \Cref{rem:ntwo} records the dimension of the
locus of all such points by a Morita argument. The locus is more than a
dimension, and saying what it is puts the residual difficulty in one line.

\begin{proposition}[The quaternionic locus]\label{prop:lagrangian}
Keep \Cref{setup:quat} and let $\psi$ be $K$-semilinear, invertible,
Rosati-symmetric, with $\psi^{2}=b$. For $x,y\in V_{+}$ put
\[
  \omega_{\psi}(x,y)\;=\;E(x,\psi y).
\]
\begin{enumerate}[label=\textup{(\roman*)},nosep]
\item $\omega_{\psi}$ is a nondegenerate alternating form on $V_{+}$, and it
  is alternating precisely because $\psi$ is Rosati-symmetric.
\item A point $P=V_{+}^{1,0}$ of $\cD_{n,n}$ lies on the locus where $\psi$
  is a morphism of Hodge structures if and only if $P$ is isotropic for
  $\omega_{\psi}$.
\item Since $\dim_{\CC}V_{+}=2n$ and $\dim_{\CC}P=n$, such a $P$ is
  Lagrangian, so the locus is the open subset of
  $\mathrm{LG}(n,V_{+},\omega_{\psi})$ cut out by negative definiteness of
  $H$. Its dimension is $n(n+1)/2$ and its codimension in $\cD_{n,n}$ is
  $n(n-1)/2$.
\end{enumerate}
\end{proposition}

\begin{proof}
(i) $\psi$ carries $V_{+}$ to $V_{-}$, since it is semilinear, and $E$ pairs
$V_{+}$ with $V_{-}$ perfectly, so $\omega_{\psi}$ is defined and
nondegenerate. Rosati symmetry is the identity $E(\psi x,y)=E(x,\psi y)$;
with $E$ alternating it gives
$\omega_{\psi}(y,x)=E(y,\psi x)=E(\psi x,y)\cdot(-1)=-\omega_{\psi}(x,y)$.

(ii) Suppose $\psi$ preserves the Hodge structure. Being semilinear it
carries $V_{+}^{1,0}$ into $V_{-}^{1,0}$, so $\psi P\subseteq V^{1,0}$, and
$E$ vanishes on $V^{1,0}$ by the first Riemann relation; hence
$E(P,\psi P)=0$, which is the isotropy of $P$. Conversely, if $P$ is
isotropic then $\psi P$ lies in the $E$-annihilator of $P$ inside $V_{-}$.
That annihilator has dimension $n$, because $E$ pairs $V_{+}$ with $V_{-}$
perfectly, and it contains $V_{-}^{1,0}$, which also has dimension $n$ and is
$E$-orthogonal to $P$ since both lie in $V^{1,0}$. The two therefore agree,
so $\psi P=V_{-}^{1,0}$ and $\psi$ preserves $V^{1,0}$.

(iii) A Lagrangian Grassmannian $\mathrm{LG}(n,2n)$ has dimension $n(n+1)/2$,
its tangent space at a Lagrangian $P$ being the symmetric forms on $P$;
negative definiteness is an open condition and does not change the dimension.
Subtracting from $\dim\cD_{n,n}=n^{2}$ gives $n(n-1)/2$. The chart by
symmetric forms, the isotropy of its members, and the two dimensions are
checked for $n\le5$ in item (XVII) of \Cref{sec:verification}.
\end{proof}

\begin{figure}[htbp]
\centering
  \includegraphics[width=0.80\textwidth]{fig_lagsweep}
\caption{\Cref{prop:lagrangian}, drawn as a pencil. Each admissible $\psi$ gives
one leaf, an open piece of $\mathrm{LG}(n,V_{+},\omega_{\psi})$ of dimension
$n(n+1)/2$, and the leaves through a common axis form a continuous family as
$\psi$ varies over the reals, so they sweep $\cD_{n,n}$. The two drawn heavy
are the ones whose $\psi$ is rational, and only those are quaternionic loci of
an actual abelian variety. The picture is exact for the pencil it draws; in
the family the leaves have dimension $n(n+1)/2$ inside $n^{2}$ and the
parameter is not one real number but a variety of them. \protect\footnotemark}
\label{fig:lagsweep}
\end{figure}
\footnotetext{Drawing the leaves as planes through one axis is faithful to the feature
that matters, that the family is continuous and its rational members are not.
It is not faithful to the codimension, which is $n(n-1)/2$ and therefore zero
only at $n=1$; at $n=2$ the leaves are divisors, which is \Cref{rem:ntwo}.}

\begin{remark}[What this says about the residual difficulty]%
\label{rem:onlyrational}
The description is worth having because of what it makes visible. Fix a
point $P$ of $\cD_{n,n}$. Complete a basis of $P$ to a basis of $V_{+}$ and
let $\omega$ be the alternating form that is standard in that basis; then $P$
is Lagrangian for $\omega$, and $\omega$ is $\omega_{\psi}$ for a real
$\psi$ satisfying every condition of \Cref{prop:lagrangian} except
rationality. So the quaternionic loci, taken over the \emph{real} $\psi$,
cover $\cD_{n,n}$ completely; they cover it with room to spare, since each
has dimension $n(n+1)/2$ and they are parametrised by a positive dimensional
family. What stops the covering is that $\psi$ must lie in
$\End^{0}(A)$, hence be rational, and the rational $\psi$ form a countable
set.

That is \Cref{cor:allornothing} again, with everything else stripped away. It
is not a deformation-theoretic difficulty and not a positivity difficulty;
it is a countable set of parameters inside an uncountable one. We state it in
this form because it is the sharpest version we can give of what a new idea
would have to supply: a mechanism that moves the Lagrangian Grassmannians
continuously, or equivalently that produces from one rational $\psi$ a family
of cycles over a neighbourhood rather than over its own locus, is the whole
of \Cref{q:residual}.
\end{remark}

\begin{lemma}[Products]\label{lem:products}
Let $(A_{1},\iota_{1},\eta_{1})$ and $(A_{2},\iota_{2},\eta_{2})$ be of
$(K,-1,n_{1})$- and $(K,-1,n_{2})$-Weil type. Then
$A_{1}\times A_{2}$ with the diagonal $K$-action and the product polarisation
is of $(K,-1,n_{1}+n_{2})$-Weil type, its hermitian form is the orthogonal
sum, its discriminant is the product of the two discriminants, and if both
factors satisfy the divisorial criterion then so does the product.
\end{lemma}

\begin{proof}
Signatures and hermitian forms add over an orthogonal sum, so the signature is
$(n_{1}+n_{2},n_{1}+n_{2})$ and $\det H=\det H_{1}\det H_{2}$, whence
$(-1)^{n_{1}+n_{2}}\det H$ is the product of the two normalised determinants.
For the criterion, $V_{-}=V_{1,-}\oplus V_{2,-}$, and
$\Theta=\Theta_{1}+\Theta_{2}$ lies in $\bigwedge^{2}V_{-}$, is of type
$(1,1)$, has $\Theta+\bbar{\Theta}$ rational, and is nondegenerate on the
direct sum because each summand is nondegenerate on its factor.
\end{proof}

\begin{theorem}[A base point in every Weil family]\label{thm:everyfamily}
Let $K=\QQ(\sqrt{-d})$, let $n\ge1$, and let $\delta$ be any class in
$\QQ^{\times}_{>0}/\Nm_{K/\QQ}(K^{\times})$. Then there is an abelian
$2n$-fold $A$ of $(K,-1,n)$-Weil type with discriminant $\delta$ whose Weil
classes are algebraic, with an explicit representative as a rational
polynomial in two divisor classes.
\end{theorem}

\begin{proof}
Choose $b\in\QQ_{>0}$ representing $\delta$.

If $n$ is odd, take the quaternionic model of \Cref{setup:quat} with this $b$.
Its discriminant is $b^{n}=b\cdot(b^{(n-1)/2})^{2}$, which is $\delta$, and
its Weil classes are algebraic by \Cref{thm:quatweil}(v).

If $n$ is even, take $A=A_{1}\times A_{2}$ with $A_{1}$ the quaternionic model
at $n_{1}=1$ and parameter $b$, and $A_{2}$ the quaternionic model at
$n_{2}=n-1$ and parameter $d$. Both factors satisfy the criterion, so the
product does by \Cref{lem:products}; and the discriminant is
$b^{1}\cdot d^{\,n-1}$, which is $\delta$ because $d=\Nm_{K/\QQ}(\sqrt{-d})$ is
a norm. The dimension is $2(1+(n-1))=2n$.

In both cases the representative is \eqref{eq:divrep} for the class
$g_{\mu}$ of \Cref{thm:quatweil}(v), or the sum of the two such classes.
\end{proof}

\begin{corollary}[Question (R2) is settled]\label{cor:R2}
Every Weil family contains a member at which the Weil classes are algebraic.
Consequently \Cref{q:residual} reduces to propagation alone: the Hodge
conjecture for abelian $2n$-folds of Weil type, for every $n$ and every
discriminant, holds if and only if the class $\omega$ of
\Cref{setup:family} propagates from the nonempty locus where it is known to be
algebraic to the rest of its connected family.
\end{corollary}

\begin{proof}
The first sentence is \Cref{thm:everyfamily} applied to each class, together
with the fact that the families are indexed by those classes
(\Cref{def:disc}). The second is \Cref{thm:residual} with the base locus now
known to be nonempty in every family.
\end{proof}

\begin{remark}[What is left, and what it is not]\label{rem:whatsleftnow}
The residual statement is now a single one. It is the variational Hodge
conjecture for one flat class on one connected family, with a base point
supplied in every family, and by \Cref{thm:residual} it is necessary as well
as sufficient. Two features of it are worth recording, because they say which
tools cannot settle it.

The base locus is thin. By \Cref{prop:divfails} the divisorial criterion fails
at every point with full Hodge group, so the locus where it holds is contained
in the Noether--Lefschetz locus of $R$, a countable union of proper closed
algebraic subvarieties of $\cD_{n,n}$ by \Cref{thm:cdk}. Enlarging the base
locus further by the same method is therefore impossible; any new point it
produces is again special.

The propagation is not formal. By \Cref{rem:variational} the infinitesimal
form of it would follow from semiregularity of the cycle at a base point, and
the cycle at a base point is now explicit in every family, being an
intersection of two divisors raised to the $n$-th power. Whether it is
semiregular is a computation in the normal bundle that we have not carried
out. The second half of that objection, that a positive answer would give
deformation only in the formal category, is answered in
\Cref{ssec:properness}: once the algebraic locus is known to be a countable
union of closed subvarieties, a formal deformation at one point is enough.
\end{remark}

\subsection{The exact form of what remains}\label{ssec:exactform}

The algebraicity of a Hodge class is usually approached through an
intermediate property. A rational class of type $(p,p)$ is \emph{absolutely
Hodge} if, after spreading the variety out over a finitely generated field and
applying any automorphism of $\CC$, the de Rham component of the class remains
of type $(p,p)$ in every comparison. The implications

\begin{equation}\label{eq:threechain}
  \text{algebraic}
  \;\Longrightarrow\;
  \text{absolutely Hodge}
  \;\Longrightarrow\;
  \text{Hodge}
\end{equation}
are both elementary; the Hodge conjecture is the reverse of the first, and the
reverse of the second is the weaker question that Deligne settled for abelian
varieties.

\begin{theorem}[Deligne \cite{Del82}]\label{thm:deligneAH}
Every Hodge class on a complex abelian variety is absolutely Hodge. Moreover,
\emph{Principle B}: if $\pi\colon\cX\to S$ is smooth projective over a
connected smooth base, and $\alpha$ is a flat section of $R^{2p}\pi_{*}\QQ(p)$
that is Hodge at every point of $S$ and absolutely Hodge at one point, then it
is absolutely Hodge at every point.
\end{theorem}

Put beside \Cref{thm:residual} and \Cref{thm:everyfamily}, this locates the
remaining problem precisely, and the location is worth stating because it is
not where one might guess.

\begin{corollary}[What is left, in one implication]\label{cor:exactform}
Let $\pi\colon\cA\to\cD_{n,n}$ and $\omega$ be as in \Cref{setup:family}. Then
$\omega_{s}$ is absolutely Hodge for every $s$, and the following are
equivalent:
\begin{enumerate}[label=\textup{(\roman*)},nosep]
\item the Hodge conjecture holds, in every degree, for every abelian
  $2n$-fold of $(K,-1,n)$-Weil type with $\Hg(A)=\SU(V,H)$;
\item the absolutely Hodge class $\omega_{s}$ is algebraic for every $s$.
\end{enumerate}
This holds for every $K$, every $n$ and every discriminant.
\end{corollary}

\begin{proof}
The first assertion is \Cref{thm:deligneAH}; alternatively, and this is the
point of the present section, it follows from Principle B together with
\Cref{thm:everyfamily}, which supplies a point of $\cD_{n,n}$ at which
$\omega$ is algebraic and hence absolutely Hodge, and with the connectedness of
$\cD_{n,n}$ from \Cref{prop:weildomain}. The equivalence is
\Cref{thm:residual}(a) and (b).
\end{proof}

\begin{remark}[Why the remaining implication resists the methods used here]
\label{rem:whyterminal}
The two halves of \eqref{eq:threechain} behave differently in families, and
the difference is the whole of the difficulty.

Principle B holds because being absolutely Hodge is a condition one can test
after spreading out over a finitely generated field, the set of Hodge classes
of bounded norm is finite by \Cref{lem:boundednorm}, and the locus where a
class stays Hodge is algebraic by \Cref{thm:cdk}; those three facts let one
propagate the property along $S$. Being algebraic has no counterpart of the
first: a cycle class is not determined by finitely many numerical conditions
that a field automorphism preserves, and the group of cycles is not locally
constant in a family. That is exactly the content of the variational Hodge
conjecture, and nothing in this paper supplies it.

An obstruction of this shape is removed in one of three ways: by a correspondence that carries the problem to a variety where
the class is a divisor or an endomorphism, which is what Kuga and Satake do
for weight two and what \Cref{cor:ksboundary} shows is unavailable here beyond
$n=1$; by a construction that produces the cycle directly at the general
point, which \Cref{thm:divisor} and \Cref{thm:character} close for two natural
families of constructions and \Cref{prop:divfails} closes for every divisorial
one; or by turning the statement over the whole family into a statement at one
point, which is what \Cref{ssec:properness} does.

The last of those is as far as we get. \Cref{thm:degreebound} and
\Cref{thm:semiregclosure} each reduce \Cref{cor:exactform}(ii), in its
entirety, to a property of the cycles already constructed in
\Cref{ssec:everyfamily}: a bound on their size along one orbit, or the
injectivity of one linear map attached to one of them. Neither reduction
appeals to a general point, to a limit, or to an unproved conjecture. We have
not established either property, so we do not claim the Hodge conjecture; what
we claim is that after \Cref{rem:tworemain} the outstanding question is of a
different kind from the one \Cref{q:residual} began with.
\end{remark}

\subsection{The algebraic locus is closed, and what that costs}
\label{ssec:properness}

Everything so far has been pointwise: a class at a point, a criterion at a
point, a cycle at a point, and in \Cref{cor:exactform} an implication to be
proved at every point. The propagation question of \Cref{q:residual} is about
the set
\[
  \Sigma \;=\; \bigl\{\,s\in\cD_{n,n} \;:\;
    \omega_{s}\ \text{is algebraic on}\ A_{s}\,\bigr\},
\]
and nothing has yet been said about what kind of set that is. It turns out to
be a very restricted kind, and the restriction changes the shape of the
remaining problem: instead of a class to construct at every point, one is left
with a bound to prove along one orbit, or a rank to compute at one point.
The mechanism is that the Hilbert scheme is proper, so each piece of $\Sigma$
cut out by bounded data is Zariski closed.\footnote{The mechanism is not new.
That the locus where a fixed Hodge class becomes algebraic is a countable
union of closed subvarieties is the standard companion to the statement that
the Hodge locus itself is a countable union of closed subvarieties, for which
see \cite[Chapter~5]{Voi03} and, in the form we use in \Cref{thm:cdk},
\cite{CDK95}. What is specific to the present situation is that the class
whose algebraic locus is being described has been pinned down completely by
\Cref{thm:hdgcount} and written out by \Cref{thm:explicitgens}, that the locus
is known to be nonempty in every family by \Cref{thm:everyfamily}, and that
the group acting has been identified. Those three inputs are what turn a soft
statement into a criterion one can try to verify.}

\begin{setupenv}[The quotient, and the universal family]\label{setup:shimura}
Keep \Cref{setup:family}. Fix an integer $N\ge3$ and let $\Gamma$ be the
subgroup of $\SU(V,H)(\QQ)$ preserving the lattice $L\subset V$ used to build
the family and acting trivially on $L/NL$. Write
\[
  p\colon\cD_{n,n}\longrightarrow\cS=\Gamma\backslash\cD_{n,n} .
\]
By Baily and Borel \cite{BB66} the quotient $\cS$ is a smooth quasi-projective
complex variety, and it is irreducible because $\cD_{n,n}$ is connected. Since
$N\ge3$ the level structure rigidifies, so $\cS$ carries a universal family
\[
  f\colon\cA\longrightarrow\cS,
\]
and the relative polarisation $3\eta$ is relatively very ample, which fixes a
relative projective embedding and hence a notion of degree in the fibres.
Write $G=\SU(V,H)$, an algebraic group over $\QQ$, so that
$G(\RR)$ acts transitively on $\cD_{n,n}$ and $\Gamma\subset G(\QQ)$.
\end{setupenv}

Two small facts have to be recorded before the main one, because both are used
without comment later and neither is quite automatic.

\begin{lemma}[The section descends]\label{lem:omegadescends}
The section $\omega$ of \Cref{setup:family} is $\Gamma$-invariant, hence
descends to a flat section of $R^{2n}f_{*}\QQ$ over $\cS$ that is of type
$(n,n)$ at every point.
\end{lemma}

\begin{proof}
The Weil line is $\HW=\bigwedge^{2n}_{K}V$, and an element of $\SU(V,H)$ is by
definition a $K$-linear automorphism of $V$ of $K$-determinant $1$. Its action
on the top $K$-exterior power is multiplication by that determinant, so
$\Gamma$ acts trivially on $\HW$ and in particular fixes $\omega$. Flatness is
then immediate, and the Hodge type at every point is
\Cref{prop:weilline}(ii).
\end{proof}

\begin{lemma}[One component is enough]\label{lem:onecomponent}
Let $(A,\iota,\eta)$ be of $(K,-1,n)$-Weil type with
$\iota(\cO_{K})\subseteq\End(A)$, and let $\HW\subset H^{2n}(A,\QQ)$ be its
Weil line. If one nonzero element of $\HW$ is algebraic, every element of
$\HW$ is algebraic.\footnote{This is the step compressed into the last
sentence of the proof of \Cref{cor:whatisleft}, written out because the
witness $\tau$ has to be chosen and because the choice is the same one the
Lean theorem \texttt{irrational\_witness\_exists} certifies. The determinant
that makes the two classes independent is $b(w_{1}^{2}+dw_{2}^{2})$ for
$\tau^{2n}=a+b\sqrt{-d}$ and $w=w_{1}+w_{2}\sqrt{-d}$, and that is the content
of \texttt{weil\_line\_spanned}.}
\end{lemma}

\begin{proof}
Among the elements $\tau_{m}=m+\sqrt{-d}$ of $\cO_{K}$ with $1\le m\le7$ the
ratios $\tau_{m}/\bbar{\tau}_{m}$ are pairwise distinct, since
$\tau_{m}\bbar{\tau}_{m'}$ is rational only when $m=m'$. The group of roots of
unity in an imaginary quadratic field has order $2$, $4$ or $6$, so some $m$
in that range has $\tau=\tau_{m}$ with $\tau/\bbar\tau$ not a root of unity,
and then $\tau^{2n}\ne\bbar{\tau}^{2n}$, that is $\tau^{2n}\notin\QQ$, for
every $n\ge1$.\footnote{The exceptional witnesses are real and are the reason
the argument names a range rather than taking $m=1$: for $d=1$ the ratio
attached to $1+\sqrt{-1}$ is a primitive fourth root of unity and for $d=3$
the ratio attached to $1+\sqrt{-3}$ is a primitive sixth root of unity, so at
those two values of $d$ the naive witness fails for the divisible $n$. The
Lean theorems \texttt{bad\_witness} and \texttt{bad\_witness\_exceptions} of
\Cref{sec:lean} record exactly which pairs fail and that no others do in the
range searched.}

Since $\tau\in\cO_{K}$, the element $\iota(\tau)$ is an endomorphism of $A$,
so $\iota(\tau)^{*}$ on $H^{2n}(A,\QQ)$ is induced by the graph
correspondence and therefore carries algebraic classes to algebraic classes.
On $H^{1}$ it is multiplication by $\tau$, so on
$\bigwedge^{2n}_{K}V$ it is multiplication by $\tau^{2n}$. As
$\tau^{2n}\notin\QQ$ and $\HW$ is a one dimensional $K$-vector space, the pair
$w,\;\tau^{2n}w$ is a $\QQ$-basis of $\HW$ for every $w\ne0$. If $w$ is
algebraic then so is $\iota(\tau)^{*}w=\tau^{2n}w$, and the two together span.
\end{proof}

\begin{theorem}[The algebraic locus is a countable union of closed
subvarieties]\label{thm:closedlocus}
Keep \Cref{setup:shimura}. For a finite tuple $\underline{P}=(P_{1},\dots
,P_{m})$ of Hilbert polynomials and a tuple $\underline{q}\in\QQ^{m}$ put
\begin{align*}
  \cH_{\underline{P}}
  &\;=\;
  \Hilb^{P_{1}}(\cA/\cS)\times_{\cS}\cdots\times_{\cS}\Hilb^{P_{m}}(\cA/\cS),
  \\[2pt]
  \Sigma_{\underline{P},\underline{q}}
  &\;=\;
  \mathrm{pr}\Bigl(\bigl\{\,\underline{Z}\in\cH_{\underline{P}}
    \;:\;\textstyle\sum_{i}q_{i}[Z_{i}]=\omega\,\bigr\}\Bigr)
  \;\subseteq\;\cS .
\end{align*}
Then each $\Sigma_{\underline{P},\underline{q}}$ is Zariski closed in
$\cS$,\footnote{The bound on the Hilbert polynomials is what the statement
lives on, and it cannot be dropped. Without it one is looking at the image of
an infinite disjoint union, and an infinite union of closed sets need be
neither closed nor constructible; indeed the whole difficulty of the Hodge
conjecture in families is that $\Sigma$ itself is not known to be either. The
relative Chow variety of effective $n$-cycles of bounded degree would serve
equally well in place of the Hilbert schemes, with the same proof, and the
choice between them is immaterial here because all that is used is properness
over the base together with the local constancy of the cycle class. What is
used later, and what forces the bookkeeping with tuples rather than a single
subscheme, is that a rational cycle is a rational combination of several
integral ones and the several may have different Hilbert polynomials.} and
\[
  p^{-1}\Bigl(\,\bigcup_{\underline{P},\underline{q}}
    \Sigma_{\underline{P},\underline{q}}\Bigr)\;=\;\Sigma ,
\]
the union running over a countable index set.
\end{theorem}

\begin{proof}
Each $\Hilb^{P_{i}}(\cA/\cS)$ is projective over $\cS$ by Grothendieck's
existence theorem \cite{Gro61}, so the fibre product $\cH_{\underline{P}}$ is
projective over $\cS$, and being of finite type it has finitely many connected
components.

Fix a connected component $T$, replace it by its reduction, which changes
neither the underlying space nor any cohomology group, and let
$\cZ_{i}\subset\cA\times_{\cS}T$ be the pullback of the $i$-th universal
subscheme. It is flat over $T$ and its fibres have Hilbert polynomial $P_{i}$
of degree $n$, so the associated $n$-cycles, the top-dimensional components
counted with their multiplicities, are the fibres of a cycle on
$\cA\times_{\cS}T$ whose class restricts on the fibre over $t$ to $[Z_{i,t}]$.
Restriction of a global class to the fibres of a topologically
locally trivial family is a flat section of the local system, so
$t\mapsto\sum_{i}q_{i}[Z_{i,t}]$ is flat on $T$; and $\omega$ is flat on $\cS$
by \Cref{lem:omegadescends}, hence flat after pullback. Two flat sections of a
local system over a connected base agree everywhere as soon as they agree at
one point, so the locus where they agree is a union of connected components of
$\cH_{\underline{P}}$. That locus is therefore itself projective over $\cS$,
and the image of a proper morphism is Zariski closed.

For the equality of sets, let $s\in\cS$ with $\omega_{s}$ algebraic. Then
$\omega_{s}=\sum_{i}q_{i}[Z_{i}]$ for finitely many closed integral subschemes
$Z_{i}\subset A_{s}$ of dimension $n$ and rationals $q_{i}$; taking
$P_{i}$ to be the Hilbert polynomial of $Z_{i}$ puts $s$ in
$\Sigma_{\underline{P},\underline{q}}$. The converse holds because the class
of a closed subscheme of pure dimension $n$ is algebraic. Countability is
clear: a Hilbert polynomial is determined by its finitely many integer
coordinates in the binomial basis, and $\QQ$ is countable.
\end{proof}

\begin{corollary}[All, or almost nothing]\label{cor:allornothing}
Exactly one of the following holds.
\begin{enumerate}[label=\textup{(\alph*)},nosep]
\item $\Sigma=\cD_{n,n}$, and the Hodge conjecture holds in every degree for
  every abelian $2n$-fold of $(K,-1,n)$-Weil type with $\Hg(A)=\SU(V,H)$ in
  the family.
\item $\Sigma$ has empty interior in the analytic topology. Its complement is
  then a dense $G_{\delta}$, so $\omega_{s}$ fails to be algebraic for every
  $s$ outside a meagre set, and the Hodge conjecture fails.
\end{enumerate}
In particular, if $\omega_{s}$ is algebraic for all $s$ in some nonempty open
subset of $\cD_{n,n}$, then it is algebraic for all $s$.
\end{corollary}

\begin{proof}
Suppose $\Sigma$ contains a nonempty open set; then so does its image in
$\cS$, and we may pick a nonempty connected open $U\subseteq\cS$ contained in
that image. By \Cref{thm:closedlocus} the countably many sets
$\Sigma_{\underline{P},\underline{q}}\cap U$ are closed in $U$ and cover it.
A connected complex manifold is a Baire space, so one of them has nonempty
interior in $U$. A Zariski closed subset of the irreducible variety $\cS$
either is all of $\cS$ or is a proper closed subvariety, and a proper closed
subvariety of a connected complex manifold has empty interior. Hence that
$\Sigma_{\underline{P},\underline{q}}$ is all of $\cS$, so $\Sigma=\cD_{n,n}$
and (a) holds by \Cref{thm:residual}(a). If $\Sigma$ has empty interior, its
complement is the intersection of the countably many dense open sets
$\cD_{n,n}\setminus p^{-1}(\Sigma_{\underline{P},\underline{q}})$, and Baire
again gives that this intersection is dense.
\end{proof}

\begin{figure}[htbp]
\centering
\includegraphics[width=0.88\textwidth]{fig_properness}
\caption{The dichotomy of \Cref{cor:allornothing}. Each stratum
$\Sigma_{\underline{P},\underline{q}}$ is Zariski closed by
\Cref{thm:closedlocus}, and there are countably many of them; the base point
$s_{0}$ of \Cref{thm:everyfamily} lies on one. By \Cref{prop:densealg} the
union of the strata is dense, so a point $s$ off all of them is squeezed
between strata however small a neighbourhood is taken. The conjecture for the
family is the assertion that no such $s$ exists, and by
\Cref{cor:allornothing} it is enough that one stratum has nonempty interior.}
\label{fig:properness}
\end{figure}

The next fact says that case (b) of \Cref{cor:allornothing}, if it occurs,
occurs in an uncomfortable way: the algebraic locus is always dense, so the
two sets in the dichotomy would both be dense and the distinction between them
would be entirely one of measure rather than of position.

\begin{lemma}[Algebraicity is a rational invariant]\label{lem:Ginvariance}
$\Sigma$ is stable under the action of $G(\QQ)$ on $\cD_{n,n}$.
\end{lemma}

\begin{proof}
A point of $\cD_{n,n}$ is a complex structure $J$ on $V_{\RR}$ compatible with
$H$, and $g\in G(\QQ)$ carries $J$ to $gJg^{-1}$. The map $g$ is then
$\QQ$-rational and $\CC$-linear from $(V_{\RR},J)$ to $(V_{\RR},gJg^{-1})$, so
after clearing denominators it induces an isogeny $A_{J}\to A_{gJ}$, whose
graph is an algebraic correspondence. Algebraicity of a rational cohomology
class is preserved by such a correspondence and by its transpose. Finally $g$
lies in $\SU(V,H)$, so by the computation in \Cref{lem:omegadescends} it fixes
the Weil line pointwise, and the class transported from $\omega_{J}$ is
$\omega_{gJ}$.
\end{proof}

\begin{proposition}[The algebraic locus is dense]\label{prop:densealg}
Let $n\ge2$. Then $\Sigma$ is dense in $\cD_{n,n}$ in the analytic topology,
and its image in $\cS$ is dense.\footnote{We should be clear about which
of the two standard density statements this is, because they behave in
opposite ways. For a class that is not flat, the locus where it stays Hodge is
a countable union of proper closed subvarieties, and the content of
\Cref{thm:cdk} is that each of them is algebraic; density there would be a
statement about how special points distribute. For the class $\omega$ of
\Cref{setup:family} the Hodge locus is the whole family by
\Cref{prop:weilline}(ii), so \Cref{thm:cdk} says nothing at all, and the set
that carries information is the algebraic locus instead. The proposition below
is a statement about that set, and its proof uses no Hodge theory: the
algebraic locus is a union of orbits of a group whose rational points are
dense in its real points, and the density is inherited from the group.}
\end{proposition}

\begin{proof}
The hermitian form $H$ has rank $2n\ge4$ over $K$ and signature $(n,n)$, so it
is isotropic over $\RR$; over every $\QQ_{p}$ a hermitian form of rank at
least three over a quadratic extension is isotropic; and hermitian forms over
an imaginary quadratic field satisfy the Hasse principle
\cite[Ch.~10, \S 6]{Sch85}. So $H$ is isotropic over $\QQ$ and $G$ is
$\QQ$-isotropic. The group $G=\SU(V,H)$ is simply connected and almost
$\QQ$-simple, so by Platonov's theorem $G(\QQ)$ is generated by the rational
points of the unipotent radicals of its proper parabolic $\QQ$-subgroups
\cite[Thm.~7.6 and \S 7.2]{PR94}. For a unipotent group $U$ over $\QQ$ the
exponential identifies $U(\QQ)$ with a $\QQ$-vector space inside the
corresponding $\RR$-vector space $U(\RR)$, so $U(\QQ)$ is dense in $U(\RR)$.
The closure of $G(\QQ)$ in $G(\RR)$ therefore contains $U(\RR)$ for every such
$U$, and those generate the identity component $G(\RR)^{+}$, which acts
transitively on $\cD_{n,n}$. Hence every $G(\QQ)$-orbit in $\cD_{n,n}$ is
dense.

By \Cref{thm:everyfamily} the set $\Sigma$ is nonempty, and by
\Cref{lem:Ginvariance} it is $G(\QQ)$-stable, so it contains a dense orbit.
The projection $p$ is open and surjective, so the image is dense.
\end{proof}

\begin{corollary}[The conjecture for Weil type is a non-meagreness
statement]\label{cor:meagre}
Let $n\ge2$. The Hodge conjecture holds for every abelian $2n$-fold of
$(K,-1,n)$-Weil type in the family of \Cref{setup:family} if and only if
$\Sigma$ is non-meagre in $\cD_{n,n}$, and if and only if $\Sigma$ has
nonempty interior.
\end{corollary}

\begin{proof}
If $\Sigma=\cD_{n,n}$ both conditions hold. If not, \Cref{cor:allornothing}(b)
says $\Sigma$ is a countable union of closed sets with empty interior, hence
meagre, and \Cref{thm:residual}(b) says the conjecture fails.
\end{proof}

\subsubsection*{The first criterion: a bound along one orbit}

The proof of \Cref{prop:densealg} produces the dense algebraic locus by moving
one cycle around by rational isogenies, and an isogeny multiplies degrees. So
the dense set it produces is stratified by unbounded data, which is exactly
what stops \Cref{cor:allornothing} from applying. Reversing this gives a
criterion, and it is the sharpest consequence of \Cref{thm:closedlocus} we
know how to state.

\begin{theorem}[A bound along the orbit closes the family]\label{thm:degreebound}
Let $n\ge2$ and let $s_{0}\in\Sigma$ be one of the base points supplied by
\Cref{thm:everyfamily}. Suppose that there is a finite set $\mathcal{T}$ of pairs
$(\underline{P},\underline{q})$ such that for every $g\in G(\QQ)$ the class
$\omega_{g\cdot s_{0}}$ is represented by some cycle whose Hilbert data lies
in $\mathcal{T}$. Then $\Sigma=\cD_{n,n}$, and the Hodge conjecture holds in every
degree for every abelian $2n$-fold of $(K,-1,n)$-Weil type with
$\Hg(A)=\SU(V,H)$ in that family.
\end{theorem}

\begin{remark}[The cycle is chosen, not transported]\label{rem:chosencycle}
The hypothesis quantifies over representing cycles at each point of the orbit
and asks only that one of them, at each point, be of bounded complexity. An
earlier form of this statement fixed the cycle, taking at $g\cdot s_{0}$ the
transport of the cycle at $s_{0}$ along the isogeny of
\Cref{lem:Ginvariance}. That form is stronger, and \Cref{thm:p1forces} below
shows that it is too strong to be usable: for a base cycle with a component of
finite stabiliser the transported data is unbounded, so its hypothesis is never
satisfied. The proof below uses only the weaker hypothesis, which is what the
covering argument needs.
\end{remark}

\begin{proof}
The orbit $G(\QQ)\cdot s_{0}$ is dense by the proof of \Cref{prop:densealg},
and by hypothesis it is the union of the finitely many sets
$G(\QQ)\cdot s_{0}\cap p^{-1}(\Sigma_{\underline{P},\underline{q}})$ with
$(\underline{P},\underline{q})\in\mathcal{T}$. A finite
union of closed sets is dense only if one of them is, so one
$p^{-1}(\Sigma_{\underline{P},\underline{q}})$ is dense; it is closed in the
analytic topology, hence equals $\cD_{n,n}$. Now apply
\Cref{thm:residual}(a).
\end{proof}

The hypothesis of \Cref{thm:degreebound} can be sharpened, and the sharpening
says which half of the Hilbert data is at issue. The leading coefficient is
not.

\begin{proposition}[The leading coefficient is constant on the orbit]
\label{prop:leadingconstant}
Let $s_{0}\in\Sigma$ and $g\in G(\QQ)$, and let $\omega_{s}$ denote the Weil
class at $s$ and $\eta$ the polarisation class. Then
\[
  \int_{A_{g\cdot s_{0}}}\omega_{g\cdot s_{0}}\cdot\frac{\eta^{n}}{n!}
  \;=\;\int_{A_{s_{0}}}\omega_{s_{0}}\cdot\frac{\eta^{n}}{n!},
\]
so the degree of the transported class with respect to the polarisation, which
is the leading coefficient of the Hilbert polynomial of any cycle
representing it, is the same at every point of the orbit.
\end{proposition}

\begin{proof}
By \Cref{lem:Ginvariance} the correspondence transporting the class is induced
by a $\QQ$-rational map which lies in $\SU(V,H)$, so it preserves the
hermitian form and hence the polarisation class, and by
\Cref{lem:omegadescends} it fixes the Weil line pointwise. An intersection
number of classes preserved by a correspondence is preserved by it. The
identification of the number with the leading coefficient is
Hirzebruch--Riemann--Roch on $A$, the higher coefficients involving the
structure sheaf of the cycle and not only its class.
\end{proof}

\begin{remark}[Where the unboundedness can only come from]
\label{rem:whereunbounded}
Write the transported cycle as a difference $\xi^{+}-\xi^{-}$ of effective
cycles. By \Cref{prop:leadingconstant} the difference of their degrees is
constant along the orbit, and by the boundedness of the Chow variety of
effective cycles of bounded degree in a fixed polarised variety, if
$\deg\xi^{+}$ were bounded along the orbit then only finitely many tuples
$(\underline{P},\underline{q})$ could occur and \Cref{thm:degreebound} would
apply. So the hypothesis of that theorem is equivalent to a bound on the
effective part alone.

That is not a formality, and the mechanism that threatens it is visible. The
transport at $g$ is by the isogeny obtained from $g$ by clearing denominators;
if $g$ has denominator $m$ on the reference lattice, which has rank $4n$, the
isogeny has degree a divisor of $m^{4n}$, and pushing a cycle forward along it
multiplies the multiplicities by factors of that order. The rational points of
$G$ are dense precisely because their denominators are unbounded, so along any
orbit dense enough to be useful the isogeny degrees are unbounded as well.
What \Cref{thm:degreebound} asks is that the cycles nevertheless simplify: that
after each transport the class $\omega$, whose degree has not changed, admits
an effective representation no more complicated than before. Nothing in this
paper decides that, and we do not know how to decide it. The content of the
present remark is only that the question has been reduced from the whole
Hilbert datum to the degree of the effective part, and that the part of the
datum which is determined by cohomology is already constant.
\end{remark}

\subsubsection*{What the transport does, exactly}

The mechanism of \Cref{rem:whereunbounded} can be computed rather than
described, and computing it decides one of the two readings of (P1). Write
$\Lambda$ for the reference lattice, of rank $4n$ over $\ZZ$, and $E$ for the
alternating form of the polarisation, so that a point $s$ of $\cD_{n,n}$ is a
complex structure on $\Lambda\otimes\RR$ and $A_{s}=(\Lambda\otimes\RR)/\Lambda$.
Throughout, $G=\dim A_{s}=2n$ and a component of a cycle representing
$\omega_{s}$ has dimension $n$.

\begin{lemma}[The transport, in coordinates]\label{lem:transport}
Let $g\in G(\QQ)$ and let $c=c(g)$ be the least positive integer with
$cg\Lambda\subseteq\Lambda$. Then $\varphi_{g}:=cg$ induces an isogeny
$A_{s_{0}}\to A_{g\cdot s_{0}}$ with
\begin{enumerate}[label=\textup{(\roman*)},nosep,leftmargin=2.6em]
\item $\deg\varphi_{g}=c^{4n}$;
\item $\varphi_{g}^{*}\eta_{g\cdot s_{0}}=c^{2}\eta_{s_{0}}$ and
  $\varphi_{g}^{*}\omega_{g\cdot s_{0}}=c^{2n}\omega_{s_{0}}$;
\item the transported cycle representing $\omega_{g\cdot s_{0}}$ is
  $c^{2n-4n}\,\varphi_{g*}(Z)$, where $Z$ represents $\omega_{s_{0}}$.
\end{enumerate}
Moreover $\sup_{g\in G(\QQ)}c(g)=\infty$.
\end{lemma}

\begin{proof}
(i) The index of $cg\Lambda$ in $\Lambda$ is $|\det_{\QQ}(cg)|
=c^{4n}|\det_{\QQ}g|$, and $\det_{\QQ}g=\Nm_{K/\QQ}(\det_{K}g)=1$ for
$g\in\SU(V,H)$.

(ii) $g$ preserves $E$, so $(cg)^{*}E=c^{2}E$, which is the first identity.
On $H^{2n}=\bigwedge^{2n}(\Lambda^{\vee}\otimes\QQ)$ the map
$\varphi_{g}^{*}$ is $c^{2n}\bigwedge^{2n}g^{\vee}$, and $g\in\SU(V,H)$ fixes
the Weil line pointwise by \Cref{lem:omegadescends}.

(iii) $\varphi_{g*}\varphi_{g}^{*}=\deg\varphi_{g}$, so
$\omega_{g\cdot s_{0}}=(\deg\varphi_{g})^{-1}\varphi_{g*}
\varphi_{g}^{*}\omega_{g\cdot s_{0}}
=c^{2n-4n}\varphi_{g*}\omega_{s_{0}}$ by (i) and (ii).

For the last sentence, $G$ is $\QQ$-isotropic by the proof of
\Cref{prop:densealg}, so it has a proper parabolic $\QQ$-subgroup whose
unipotent radical $U$ satisfies $U(\QQ)\cong\QQ^{k}$ with $k\ge1$ through the
exponential; the entries of $u\in U(\QQ)$ are $\QQ$-linear in that coordinate,
so $c(u)$ is unbounded on $U(\QQ)$. Item (XXXVII) of \Cref{sec:verification}
verifies (i) and (ii) on an explicit sample of such $u$ with $c$ up to $29$,
for $2n=4,6,8$.
\end{proof}

\begin{lemma}[The multiplicity is the stabiliser met by the kernel]
\label{lem:multiplicity}
Let $\varphi\colon A\to A'$ be an isogeny of abelian varieties and $Z\subseteq A$
an integral closed subvariety. Put
$G(Z)=\{x\in A:Z+x=Z\}$, a closed algebraic subgroup. Then
\[
  \varphi_{*}[Z]=m\,[\varphi(Z)],\qquad
  m=\bigl|\ker\varphi\cap G(Z)\bigr| .
\]
If moreover $\dim Z=n$ and $\varphi^{*}\eta'=c^{2}\eta$, then
\[
  \deg_{\eta'}\varphi(Z)=\frac{c^{2n}}{m}\,\deg_{\eta}Z .
\]
\end{lemma}

\begin{proof}
The map $\varphi|_{Z}\colon Z\to\varphi(Z)$ is finite and surjective, and its
degree $m$ is the number of points of a general fibre. For $z\in Z$ the fibre
through $z$ is $Z\cap(z+\ker\varphi)$, which as a set of translations is
$\{t\in\ker\varphi:z+t\in Z\}$. For a fixed $t$ the locus
$\{z\in Z:z+t\in Z\}=Z\cap(Z-t)$ is all of $Z$ when $t\in G(Z)$ and a proper
closed subset otherwise; since $\ker\varphi$ is finite, a general $z$ avoids
the finitely many proper subsets, and there
$\{t\in\ker\varphi:z+t\in Z\}=\ker\varphi\cap G(Z)$. This is the first
identity, and $\varphi_{*}[Z]=m[\varphi(Z)]$ by the definition of the
pushforward of cycles.

For the degree, the projection formula gives
\[
  m\deg_{\eta'}\varphi(Z)=\int\varphi_{*}[Z]\cdot\frac{\eta'^{n}}{n!}
  =\int[Z]\cdot\frac{(\varphi^{*}\eta')^{n}}{n!}
  =c^{2n}\int[Z]\cdot\frac{\eta^{n}}{n!}=c^{2n}\deg_{\eta}Z .
\]
Item (XXXVII) verifies both identities for subtori of an abelian fourfold,
with $m$ computed as a lattice index by Smith normal form and the degrees as
Pfaffians; there the two cases are visible side by side, a subtorus preserved
by $\varphi$ having $m=c^{2n}$ and image degree constant, and one not
preserved having $m=c^{2n-1}$ and image degree growing linearly in $c$.
\end{proof}

\begin{theorem}[What \textup{(P1)} forces on the base cycle]\label{thm:p1forces}
Let $s_{0}\in\Sigma$ and let $Z=\sum_{i}q_{i}Z_{i}$, with $q_{i}\ne0$ and the
$Z_{i}$ integral of dimension $n$, represent $\omega_{s_{0}}$.
\begin{enumerate}[label=\textup{(\roman*)},leftmargin=2.6em,itemsep=2pt]
\item For $g\in G(\QQ)$ the transported cycle is
  $\sum_{i}c(g)^{2n-4n}m_{i}(g)\,q_{i}\,[\varphi_{g}(Z_{i})]$ with
  $m_{i}(g)=|\ker\varphi_{g}\cap G(Z_{i})|$, and the $i$-th component has
  degree $c(g)^{2n}\deg_{\eta}(Z_{i})/m_{i}(g)$.
\item If some $Z_{i}$ has finite stabiliser $G(Z_{i})$, then
  $\deg\varphi_{g}(Z_{i})\ge c(g)^{2n}\deg_{\eta}(Z_{i})/|G(Z_{i})|$, which is
  unbounded on $G(\QQ)$. The Hilbert data of the transported cycle then takes
  infinitely many values, and the hypothesis of \Cref{thm:degreebound} in the
  form that fixes the transported cycle is not satisfied by $Z$.
\item Consequently, a base cycle for which the transported data is bounded
  must have every component of positive-dimensional stabiliser, that is, every
  $Z_{i}$ must be stable under translation by a positive-dimensional abelian
  subvariety of $A_{s_{0}}$.
\end{enumerate}
\end{theorem}

\begin{proof}
(i) is \Cref{lem:transport}(iii) and \Cref{lem:multiplicity}.

(ii) $\ker\varphi_{g}\cap G(Z_{i})\subseteq G(Z_{i})$, so
$m_{i}(g)\le|G(Z_{i})|$, a constant; the displayed bound is (i), and $c(g)$ is
unbounded by \Cref{lem:transport}. Distinct degrees of the components give
distinct Hilbert polynomials, so the data takes infinitely many values.

(iii) is the contrapositive of (ii), together with the fact that an algebraic
subgroup of positive dimension in an abelian variety has an abelian subvariety
of positive dimension as its identity component.
\end{proof}

\begin{remark}[What this settles and what it leaves]\label{rem:p1meaning}
Three things follow, and we separate them.

The statement (P1) has two readings, and one of them is now closed. Read with
the cycle fixed at $s_{0}$ and transported, it fails for every base cycle this
paper writes down: those cycles are intersections of divisors, and an
intersection of $n$ divisors on an abelian $2n$-fold has finite stabiliser
unless the divisors are pulled back from a quotient, which the ones here are
not. Read with the cycle chosen afresh at each point of the orbit, which is
the form \Cref{thm:degreebound} now takes, it is untouched by
\Cref{thm:p1forces} and remains open.

The two readings are different, and the difference is the content of
\Cref{lem:multiplicity}. The class $\omega_{g\cdot s_{0}}$ has constant degree
along the orbit, by \Cref{prop:leadingconstant}; what grows is the complexity
of the particular representative the transport produces, and it grows because
the isogeny spreads each component out by the factor $c^{2n}/m_{i}$. A
representative of bounded complexity may well exist at every point; the
transport does not find it.

The third is a demand on any future attempt at (P1). It must either produce a
base cycle supported on subvarieties with positive-dimensional stabiliser,
which for a Weil class is a strong and possibly impossible requirement, or
abandon the transport and bound instead the minimal complexity of a
representative, which is a statement about the Chow variety of the family and
not about any one cycle. We do not know how to do either.
\end{remark}

The bound on the minimal complexity is the form in
which \Cref{thm:degreebound} now stands. It is the weakest hypothesis under
which the covering argument runs, and it is therefore the form in which (P1)
would have to be attacked. The next theorem decides it, and the answer is that
there is nothing there to attack: the hypothesis is not a sufficient condition
for the conjecture on the family but another way of writing it.

\begin{theorem}[The bounded form is the conjecture itself]\label{thm:p1equivalent}
Let $n\ge2$, let $s_{0}\in\Sigma$ and keep \Cref{setup:family}. The following
are equivalent.
\begin{enumerate}[label=\textup{(\alph*)},leftmargin=2.6em,itemsep=2pt]
\item There is a finite set $\mathcal{T}$ of pairs
  $(\underline{P},\underline{q})$ such that at every point of the orbit
  $G(\QQ)\cdot s_{0}$ the class $\omega$ is represented by a cycle whose
  Hilbert data lies in $\mathcal{T}$; that is, \textup{(P1)} holds at
  $s_{0}$.
\item There is a single pair $(\underline{P},\underline{q})$ such that at
  every point of $\cD_{n,n}$, not merely of the orbit, the class $\omega$ is
  represented by a cycle with that datum.
\item $\Sigma=\cD_{n,n}$.
\item $\mathbf{W}(K,n,\delta)$ holds, that is, the Weil classes are algebraic
  on every member of the family.
\end{enumerate}
In particular \Cref{thm:degreebound} is an equivalence, and \textup{(P1)} is a
reformulation of the conjecture for the family and not a reduction of it. The
same holds for every CM field, with \Cref{thm:cmpropagation}(iv) in place of
\Cref{thm:degreebound}.
\end{theorem}

\begin{proof}
(a) implies (c) is \Cref{thm:degreebound}. (b) implies (a) is immediate, the
orbit being contained in $\cD_{n,n}$ and $\mathcal{T}$ being the one pair.
(c) is equivalent to (d) by \Cref{thm:residual}(a) and (b).

It remains to deduce (b) from (c). If $\Sigma=\cD_{n,n}$ then $\Sigma$ has
nonempty interior, and the first paragraph of the proof of
\Cref{cor:allornothing} shows that some
$\Sigma_{\underline{P},\underline{q}}$ is then all of $\cS$: the countably
many sets $\Sigma_{\underline{P},\underline{q}}\cap U$ are closed in a
nonempty connected open $U\subseteq\cS$ and cover it, a connected complex
manifold is a Baire space, so one of them has nonempty interior in $U$, and a
proper Zariski closed subvariety of the irreducible variety $\cS$ has empty
interior. By the definition of $\Sigma_{\underline{P},\underline{q}}$ in
\Cref{thm:closedlocus} as the image of a proper morphism, every $s\in\cS$
then carries a tuple $\underline{Z}$ of subschemes with Hilbert polynomials
$\underline{P}$ and $\sum_{i}q_{i}[Z_{i}]=\omega$, which is (b).

For a CM field the same argument applies verbatim, using
\Cref{thm:cmpropagation}(iii) in place of \Cref{thm:closedlocus} and
\Cref{cor:allornothing}.
\end{proof}

\begin{remark}[What is left of the first criterion]\label{rem:p1collapse}
\Cref{thm:p1equivalent} removes (P1) from the list of things that would close
the family, and we should be exact about what it does and does not say.

It does not say that \Cref{thm:degreebound} is false or that its proof is
defective. The implication it proves is correct, and the covering argument is
the only place in this paper where density of the rational orbit is converted
into a global statement. What the theorem says is that the implication has a
converse, so that the hypothesis is no weaker than the conclusion. A criterion
of that kind cannot be used: verifying it at one family is exactly as hard as
proving the conjecture for that family, no more and no less.

Nor does it say that the boundedness is uninteresting. It is a concrete
finiteness statement about the Chow variety of the universal family, and it is
the only reformulation of $\mathbf{W}(K,n,\delta)$ we know that mentions
neither Hodge theory nor deformation theory. What it is not is a step.

The two forms of the propagation statement are therefore not on the same
footing, and we had them on the same footing until now. (P2) remains a genuine
sufficient condition: injectivity of one semiregularity map at one point
implies the conjecture for the whole family, and nothing here shows the
converse, nor do we see how it could, since $\Sigma=\cD_{n,n}$ says nothing
about the deformation theory of any particular perfect complex. So the
propagation half of the frontier is the single statement (P2), and the
constraints \Cref{thm:p2forces} and \Cref{cor:p2shape} place on it are the
constraints on the whole of that half.
\end{remark}

At the quaternionic points of \Cref{setup:quat} the mechanism of
\Cref{rem:whereunbounded} can be seen explicitly. Take $n=2$ and the integral
lattice
$\Lambda_{b}=\cO_{b}^{2}$, $\cO_{b}=\ZZ\langle 1,i,j,ij\rangle$, for
$b\in\ZZ_{>0}$. Put
\[
  g_{b}(x,y)=E_{b}(jx,y),\qquad
  \lambda_{b}(x,y)=-\tfrac{1}{d}\,g_{b}(ix,y),\qquad
  Z_{b}=g_{b}^{2}-d\lambda_{b}^{2},\qquad Z'_{b}=2g_{b}\lambda_{b},
\]
so that $Z_{b}+\sqrt{-d}\,Z'_{b}=(g_{b}+\sqrt{-d}\,\lambda_{b})^{2}$ and
$Z_{b},Z'_{b}$ span the Weil plane. Write $x\in\cO_{b}$ as $x=u+vj$ with
$u,v\in\ZZ[i]$; the pair $(\Lambda_{b},i)$ is then the same $\ZZ[i]$-module
$\ZZ[i]^{4}$ for every $b$, and we use this identification throughout.

\begin{proposition}[The quaternionic Weil cycle is a square multiple]
\label{prop:quatdivisible}
With these data, for every $b\in\ZZ_{>0}$:
\begin{enumerate}[label=\textup{(\roman*)},nosep]
\item $E_{b}=E_{u}+b\,E_{v}$, $g_{b}=b\,g_{1}$ and $\lambda_{b}=b\,\lambda_{1}$,
  where $E_{u},E_{v},g_{1},\lambda_{1}$ are integral on $\ZZ[i]^{4}$ and do not
  depend on $b$; in particular $g_{1},\lambda_{1}\in\mathrm{NS}(A_{b})$;
\item $Z_{b}=b^{2}Z_{1}$ and $Z'_{b}=b^{2}Z'_{1}$ in $\mathrm{CH}^{2}(A_{b})$
  modulo algebraically trivial cycles, for every choice of line bundles
  representing $g_{b},\lambda_{b}$; in cohomology the equalities hold in
  $\wedge^{4}\Lambda^{*}$;
\item $\int\eta_{b}^{4}=b^{2}\int\eta_{1}^{4}$,
  $\int g_{b}^{2}\eta_{b}^{2}=-\tfrac{b}{3}\int\eta_{b}^{4}$ and
  $\int Z_{b}^{2}=\tfrac{4b^{2}}{3}\int\eta_{b}^{4}$, so that
  $(\int g_{b}^{2}\eta_{b}^{2})^{2}=\tfrac{1}{12}\int Z_{b}^{2}\int\eta_{b}^{4}$;
\item for $\alpha\in K^{\times}$ the element $\alpha j$ generates the same
  algebra at the same point, and replacing $j$ by $\alpha j$ multiplies the
  two ratios in \textup{(iii)} by $N(\alpha)$ and $N(\alpha)^{2}$.
\end{enumerate}
\end{proposition}

\begin{proof}
For $x=u+vj$ one has $\bbar{x}=\bbar{u}-vj$, and the relations $ju=\bbar{u}j$,
$j^{2}=b$ give
$\bbar{x}\,i\,y\equiv\bbar{u}\,i\,u'+b\,v\,i\,\bbar{v'}$ modulo $Kj$. Since
$\mathrm{trd}$ vanishes on $Kj$ and restricts to $\mathrm{tr}_{K/\QQ}$ on $K$,
\[
  E_{b}(x,y)=\textstyle\sum_{k}a_{k}\bigl(\mathrm{tr}(i\,\bbar{u_{k}}u'_{k})
  +b\,\mathrm{tr}(i\,v_{k}\bbar{v'_{k}})\bigr).
\]
From $jx=b\bbar{v}+\bbar{u}j$ one reads off
$g_{b}(x,y)=b\sum_{k}a_{k}\bigl(\mathrm{tr}(i\,v_{k}u'_{k})
+\mathrm{tr}(i\,\bbar{u_{k}}\bbar{v'_{k}})\bigr)$, which is $b\,g_{1}$, and
$\lambda_{b}=b\lambda_{1}$ follows. The traces take integral values on
$\ZZ[i]$, and $g_{1}=b^{-1}g_{b}$ is a rational multiple of a Hodge class,
which proves (i). Then $g_{b}+\sqrt{-d}\,\lambda_{b}=b(g_{1}+\sqrt{-d}\,\lambda_{1})$
and (ii) holds in cohomology. Choosing $L,M$ with $c_{1}(L)=g_{1}$,
$c_{1}(M)=\lambda_{1}$, the cycle $c_{1}(L^{b})^{2}-d\,c_{1}(M^{b})^{2}$ equals
$b^{2}(c_{1}(L)^{2}-d\,c_{1}(M)^{2})$ in $\mathrm{CH}^{2}(A_{b})$; another
choice changes both sides by algebraically trivial cycles.

For (iii), $E_{u}$ and $E_{v}$ live on complementary halves $U,V'$ of
$\Lambda\otimes\QQ$ of rank four, so $\eta_{u}^{3}=\eta_{v}^{3}=0$ and
$\int\eta_{b}^{4}=6b^{2}\int\eta_{u}^{2}\eta_{v}^{2}$. The form $g_{1}$ pairs
$U$ with $V'$, so $g_{1}^{2}\in\wedge^{2}U^{*}\otimes\wedge^{2}V'^{*}$, and
only the mixed term survives in
$\int g_{b}^{2}\eta_{b}^{2}=2b^{3}\int g_{1}^{2}\eta_{u}\eta_{v}$; likewise
$\int Z_{b}^{2}=b^{4}\int Z_{1}^{2}$. Hence the ratios
$\int g_{b}^{2}\eta_{b}^{2}/\int\eta_{b}^{4}$ and
$\int Z_{b}^{2}/\int\eta_{b}^{4}$ are $b$ and $b^{2}$ times their values at
$b=1$, which are $-1/3$ and $4/3$ by direct evaluation. For (iv), $(\alpha j)^{2}=N(\alpha)b$ and
$g_{\alpha j}+\sqrt{-d}\,\lambda_{\alpha j}$ is $\alpha$ or $\bbar{\alpha}$
times $g_{j}+\sqrt{-d}\,\lambda_{j}$.
\end{proof}

Item (XXIII) of \Cref{sec:verification} confirms (i) to (iv) in exact
arithmetic for $d\in\{1,3\}$, two weight vectors and fifteen values of $b$ up
to $101$. It also computes the integral Weil lattice
$\HW_{\ZZ}=(\QQ Z_{b}+\QQ Z'_{b})\cap\wedge^{4}\Lambda^{*}$ and the
sublattice spanned inside it by products of two integral divisor classes: for
$d\in\{1,2,3,7\}$, three weight vectors and seven values of $b$, the Gram
matrix of $\HW_{\ZZ}$ is $\mathrm{diag}(8d,8d^{2})$ and the index is
$2(a_{1}a_{2})^{2}$, both independent of $b$.

\begin{remark}[The multiple sits in the polarisation]\label{rem:quatdivisible}
At these points the Weil class carries no multiple that grows with $b$ once
it is written in integral terms: by \Cref{prop:quatdivisible}(ii) it is
$b^{2}Z_{1}$, with $Z_{1}$ one integral algebraic cycle built from the two
fixed divisor classes $g_{1},\lambda_{1}$. The factor that the orbit argument
would otherwise have to absorb into the coefficients $\underline{q}$ of
\Cref{thm:degreebound} is therefore absent. What grows is the polarisation:
$\int\eta_{b}^{4}=b^{2}\int\eta_{1}^{4}$, so the points $A_{b}$ carry
polarisations of unbounded degree, and an isogeny carrying $\eta_{b}$ to the
polarisation of one fixed family has degree $\int\eta_{b}^{4}/\int\eta^{4}$,
which grows like $b^{2}$. This is the unbounded isogeny degree of
\Cref{rem:whereunbounded}, located at an explicit family of points, and it is
untouched by the integrality of $Z_{1}$.
\end{remark}

The base points of \Cref{thm:everyfamily}, and every algebraicity statement
for Weil classes proved in this paper, represent $\omega$ through products of
divisor classes. The next result says that a representation of that kind can
hold with bounded degrees only on a nowhere dense set, so the bound asked for
in \Cref{thm:degreebound} has to come from cycles that are not products of
divisors. For a class $x\in H^{2}(A_{s},\RR)$ put
\[
  t(x)=\frac{\int x\,\eta^{2n-1}}{\int\eta^{2n}},\qquad
  I(x)=t(x)^{2}\!\int\eta^{2n}-\int\bigl(x-t(x)\eta\bigr)^{2}\eta^{2n-2}.
\]
Since $\cD_{n,n}$ is contractible, the local system $R^{2}\pi_{*}\ZZ$ is
trivial; write $H^{2}_{\ZZ}$ for its fibre, identified with $H^{2}(A_{s},\ZZ)$
at every $s$. Both $t$ and $I$ are polynomial functions on $H^{2}_{\ZZ}$ and
do not depend on $s$.

\begin{lemma}[The Hodge norm of a divisor]\label{lem:nefnorm}
Let $h_{s}$ be the Hodge norm on $H^{2}(A_{s},\RR)$ defined by the
polarisation, normalised by $h_{s}(\eta)=\int\eta^{2n}$. Then
\begin{enumerate}[label=\textup{(\roman*)},nosep]
\item $h_{s}(x)=I(x)$ for every real class $x$ of type $(1,1)$ at $s$;
\item if $D$ is an effective divisor on $A_{s}$ then
  $I([D])\le 2\bigl(\int[D]\eta^{2n-1}\bigr)^{2}\big/\!\int\eta^{2n}$.
\end{enumerate}
\end{lemma}

\begin{proof}
Write $x=t\eta+x_{0}$ with $t=t(x)$, so that $x_{0}$ is primitive. The
Lefschetz decomposition is orthogonal for $h_{s}$, and on primitive classes
$h_{s}(x_{0})=-\int x_{0}\wedge C_{s}x_{0}\wedge\eta^{2n-2}$, where $C_{s}$ is
the Weil operator, positive by the Hodge--Riemann relations. For $x_{0}$ of type
$(1,1)$, $C_{s}x_{0}=x_{0}$, which gives (i). An effective divisor on an
abelian variety is nef, since a translate of it meets any given curve
properly, so $\int[D]^{2}\eta^{2n-2}\ge0$. Expanding,
$\int[D]^{2}\eta^{2n-2}=t^{2}\int\eta^{2n}+\int x_{0}^{2}\eta^{2n-2}$, hence
$-\int x_{0}^{2}\eta^{2n-2}\le t^{2}\int\eta^{2n}$ and
$I([D])\le2t^{2}\int\eta^{2n}$, which is (ii).
\end{proof}

\begin{proposition}[Divisor products of bounded degree]\label{prop:divroute}
In \Cref{setup:family}, fix $C>0$ and let $\Sigma_{C}\subset\cD_{n,n}$ be the
set of points $s$ at which $\omega_{s}$ is a $\QQ$-linear combination of
classes $[D_{1}]\cdots[D_{n}]$ with every $D_{i}$ an effective divisor on
$A_{s}$ of degree $\int[D_{i}]\eta^{2n-1}\le C$. Then $\Sigma_{C}$ is
contained in a locally finite union of proper closed analytic subsets of
$\cD_{n,n}$. In particular $\Sigma_{C}$ is nowhere dense, and on any dense
subset of $\cD_{n,n}$ at which $\omega$ is a combination of products of
effective divisors the degrees of those divisors are unbounded on every
nonempty open set.
\end{proposition}

\begin{proof}
For $x\in H^{2}_{\ZZ}$ let $\mathrm{Hdg}(x)\subset\cD_{n,n}$ be the
closed analytic set where $x$ is of type $(1,1)$. If $x\notin\QQ\eta$ it is
proper, because $\NS(A_{s})\otimes\QQ=\QQ\eta$ at a very general point by
\Cref{lem:genericNS}. Let $s\in\Sigma_{C}$. Not every $[D_{i}]$ in the
representation lies in $\QQ\eta$, since $\QQ\eta^{n}$ and the Weil line carry
different characters by \Cref{lem:charsdiffer}. So $s\in\mathrm{Hdg}(x)$ for
some $x$ in
\[
  X_{C}=\bigl\{x\in H^{2}_{\ZZ}\setminus\QQ\eta\;:\;
  I(x)\le 2C^{2}/\textstyle\int\eta^{2n}\bigr\},
\]
by \Cref{lem:nefnorm}(ii). Now let $U\subset\cD_{n,n}$ be open with compact
closure and $s_{0}\in U$. The Hodge norm depends continuously on the point, so
$h_{s_{0}}\le c_{U}\,h_{s}$ for some $c_{U}$ and all $s\in U$. If
$\mathrm{Hdg}(x)$ meets $U$ at $s$ and $x\in X_{C}$, then by
\Cref{lem:nefnorm}(i)
$h_{s_{0}}(x)\le c_{U}h_{s}(x)=c_{U}I(x)\le 2c_{U}C^{2}/\int\eta^{2n}$, and
only finitely many integral classes satisfy this. So the sets
$\mathrm{Hdg}(x)$, $x\in X_{C}$, form a locally finite family of proper closed
analytic subsets. Their union is closed with empty interior in the connected
manifold $\cD_{n,n}$, and it contains $\Sigma_{C}$.
\end{proof}

\begin{remark}[What the proposition closes]\label{rem:divroute}
The loci on which this paper proves $\omega$ algebraic are the split,
quaternionic and product loci and their translates under $G(\QQ)$, and on
each of them $\omega$ is a rational polynomial in divisor classes, since an
isogeny carries $\NS\otimes\QQ$ isomorphically. \Cref{prop:divroute} shows that no
bounded family of such cycles covers a dense set, so the finiteness in
\Cref{thm:degreebound} cannot be obtained by bounding the divisors that
produce $\omega$ at the points where it is known to be algebraic. It has to
come from the transported cycles of \Cref{lem:Ginvariance} themselves, which
are push-forwards along isogenies and not products of divisors on the target.
That is the unbounded isogeny degree of \Cref{rem:whereunbounded} and
\Cref{rem:quatdivisible}, now shown to be unavoidable for the divisorial
constructions. The argument uses nothing about $\omega$ beyond the fact that it
is not a multiple of $\eta^{n}$, and nothing about the family beyond the
generic Picard number, so it applies equally to any variational argument that
feeds on divisor products.
\end{remark}

At $n=2$ the phenomenon can be exhibited on one fixed polarised lattice. Take
$\Lambda=\cO^{2}$ with $\cO=\ZZ\langle1,i,j,ij\rangle$, $i^{2}=-d$, $j^{2}=1$,
and $E$ as in \eqref{eq:quatE}; this is one member of the family, with its
lattice and polarisation fixed. Let
$R=\{x:\ x(iu,v)=x(u,iv)\}\subset H^{2}(\Lambda,\QQ)$ be the space of
$K$-bilinear classes, where the divisor classes of the quaternionic points
live, and for $x\in R$ let $\phi_{x}=E^{-1}x$, so that $x(u,v)=E(\phi_{x}u,v)$.

\begin{proposition}[A pencil of quaternionic loci]\label{prop:pencil}
In this setting, for $d\in\{1,3\}$:
\begin{enumerate}[label=\textup{(\roman*)},nosep]
\item $R\cap H^{2}(\Lambda,\ZZ)$ has rank $12$, every $x\in R$ has $t(x)=0$,
  and $I(x)=\tfrac{1}{24}\,\mathrm{tr}(\phi_{x}^{2})\int\eta^{4}$; the form $I$
  has signature $(8,4)$ on $R$;
\item there are $x,y\in R\cap H^{2}(\Lambda,\ZZ)$ such that
  $\phi_{x+ky}^{2}=\beta(k)$ is a positive rational scalar for every
  $k\in\ZZ$, with $\beta(k)=\tfrac12+\tfrac12k+\tfrac14k^{2}$ for $d=1$ and
  $\beta(k)=\tfrac19+\tfrac29k^{2}$ for $d=3$;
\item for every $k$ the set $\mathrm{Hdg}(x+ky)$ is a quaternionic locus
  with algebra $(-d,\beta(k))$, nonempty of dimension $3$, and loci with
  different values of $\beta$ are not related by any automorphism of
  $(\Lambda,E,i)$;
\item at a very general point of $\mathrm{Hdg}(x+ky)$ one has
  $\NS\otimes\QQ=\QQ\eta\oplus\QQ x_{k}\oplus\QQ\,x_{k}(i\cdot,\cdot)$ with
  $x_{k}=x+ky$, and for every $k\in\ZZ$ the minimum $\mu_{k}$ of $I$ on
  $\NS\setminus\QQ\eta$ satisfies
  $\mu_{k}\ge\beta(k)\big/\bigl(3N_{d}^{2}\int\eta^{4}\bigr)$, with $N_{1}=2$
  and $N_{3}=4$.
\end{enumerate}
Consequently, at a very general point of $\mathrm{Hdg}(x+ky)$ every effective
divisor whose class is not a multiple of $\eta$ has degree at least
$\sqrt{\beta(k)/6}\,/N_{d}$, which grows linearly in $|k|$.
\end{proposition}

\begin{proof}
(i) and (ii) are finite exact computations, item (XXIV) of
\Cref{sec:verification}: the trace identity is checked on all pairs of a basis
of $R\cap H^{2}(\Lambda,\ZZ)$, which proves it on $R$ by bilinearity, and in
(ii) the three matrices $\phi_{x}^{2}$, $\phi_{y}^{2}$ and
$\phi_{x}\phi_{y}+\phi_{y}\phi_{x}$ are scalar, so
$\phi_{x+ky}^{2}=\phi_{x}^{2}+k(\phi_{x}\phi_{y}+\phi_{y}\phi_{x})+k^{2}\phi_{y}^{2}$
is scalar for every $k$; $\beta$ has positive leading coefficient and negative
discriminant. For (iii) write $\phi=\phi_{x+ky}$ and $\beta=\beta(k)$. The map
$\phi$ is $E$-self-adjoint, anticommutes with $i$, and $\phi^{2}=\beta>0$, so
$V_{\RR}=V_{+}\oplus V_{-}$ with $V_{\pm}$ the $\pm\sqrt{\beta}$ eigenspaces,
$E$-orthogonal to each other, and $i$ exchanges them. A complex structure $J$
on $V_{+}$ compatible with $E|_{V_{+}}$, extended to $V_{-}=iV_{+}$ by
$J(iw)=iJw$, commutes with $i$ and $\phi$ and is compatible with $E$, and $i$
has eigenvalues $\pm\sqrt{-d}$ each with multiplicity $2$ on $V^{1,0}$; so
these $J$ are points of $\cD_{2,2}$ at which $x+ky$ is of type $(1,1)$, and
they form a copy of the Siegel space of $V_{+}$, of dimension $3$, in
agreement with \Cref{prop:lagrangian}. An automorphism $\gamma$
of $(\Lambda,E,i)$ sends $\phi_{x}$ to $\gamma\phi_{x}\gamma^{-1}$ and so
preserves $\mathrm{tr}(\phi_{x}^{2})$.

For the first half of (iv), let $S_{k}\subset\cD_{2,2}$ be the Siegel space
just described and $\mathcal{C}_{k}\subset\Sp(V_{\RR},E)$ the centraliser of
$i$ and $\phi$. An element of $\mathcal{C}_{k}$ preserves $V_{+}$ and is
determined by its restriction there, because it commutes with $i$, so
$\mathcal{C}_{k}\cong\Sp(V_{+})$, acting on $V_{\RR}=V_{+}\oplus iV_{+}$ as
$g\oplus igi^{-1}$. A real class is of type $(1,1)$ at $J$ exactly when it is
invariant under $J$. The points of $S_{k}$ form one orbit
$\{gJ_{0}g^{-1}:g\in\mathcal{C}_{k}\}$ of a fixed point $J_{0}$. Let $z$ be a
real $2$-form invariant under every point of $S_{k}$, and write
$\mathfrak{c}=\mathrm{Lie}(\mathcal{C}_{k})\cong\mathfrak{sp}(V_{+})$,
$\mathfrak{k}$ for the centraliser of $J_{0}$ in $\mathfrak{c}$ and
$\mathfrak{p}=[\mathfrak{c},J_{0}]$. For $X\in\mathfrak{c}$ the curve
$g_{t}=\exp(tX)$ gives $(g_{t}J_{0}g_{t}^{-1})^{*}z=z$ for all $t$, and the
derivative at $t=0$ says that $[X,J_{0}]$ annihilates $z$, so $\mathfrak{p}$
annihilates $z$. The annihilator of $z$ in $\mathfrak{c}$ is a Lie
subalgebra, so it contains $\mathfrak{p}+[\mathfrak{p},\mathfrak{p}]$. Now
$\mathfrak{c}=\mathfrak{k}\oplus\mathfrak{p}$ with
$[\mathfrak{k},\mathfrak{p}]\subseteq\mathfrak{p}$ and
$[\mathfrak{p},\mathfrak{p}]\subseteq\mathfrak{k}$, so
$\mathfrak{p}+[\mathfrak{p},\mathfrak{p}]$ is an ideal of $\mathfrak{c}$;
it is nonzero, and $\mathfrak{sp}_{4}(\RR)$ is simple, so it is all of
$\mathfrak{c}$.\footnote{Simplicity of $\mathfrak{sp}_{4}(\RR)$ is the one
fact from Lie theory used here, and it is elementary: the complexification
$\mathfrak{sp}_{4}(\CC)$ has the irreducible root system $C_{2}$, a complex
semisimple Lie algebra with irreducible root system is simple, and an ideal of a
real form is an ideal of the complexification after tensoring with $\CC$, so
a real form of a simple complex Lie algebra is simple. Nothing finer than this
is needed, in particular no statement about the abstract group
$\mathrm{PSp}_{4}(\RR)$; and the conclusion
$\mathfrak{p}+[\mathfrak{p},\mathfrak{p}]=\mathfrak{c}$ is also checked by
direct computation in item (XXIV), at the two members of the pencil that carry
a rational point.} Hence $z$ is annihilated by $\mathfrak{c}$ and, since
$\Sp(V_{+})$ is connected, invariant under all of $\mathcal{C}_{k}$.
So the classes of type $(1,1)$ at every point of $S_{k}$ are the invariants of
$\Sp(V_{+})$ in $\bigwedge^{2}(V_{+}\otimes\RR^{2})
=\bigl(\bigwedge^{2}V_{+}\otimes S^{2}\RR^{2}\bigr)
\oplus\bigl(S^{2}V_{+}\otimes\bigwedge^{2}\RR^{2}\bigr)$, a space of dimension
$3$ spanned by $\eta$, $x_{k}$ and $x_{k}(i\cdot,\cdot)$. Item (XXIV)
confirms both steps over $\QQ$: for $0\le k\le4$ the centraliser of $i$ and
$\phi$ in $\mathfrak{sp}(V,E)$ has dimension $10$ and exactly three invariants
in $\bigwedge^{2}V^{*}$; and at the values of $k$ where $\beta(k)$ is a
rational square, $k=-1$ for $d=1$ and $k=-2$ for $d=3$, a rational point
$J_{0}$ of $S_{k}$ is written down, $\mathfrak{p}=[\mathfrak{c},J_{0}]$ has
dimension $6$, and $\mathfrak{p}+[\mathfrak{p},\mathfrak{p}]$ has dimension
$10$. Every other rational class $z$ is of type $(1,1)$ only on
a proper closed analytic subset of the connected manifold $S_{k}$, and a very
general point avoids these countably many subsets.

For the bound in (iv), write $e=\int\eta^{4}$, $\iota x=x(i\cdot,\cdot)$ and
$M_{k}=(\QQ x_{k}+\QQ\iota x_{k})\cap H^{2}(\Lambda,\ZZ)$. Since
$\phi_{\iota x}=\phi_{x}\,i$ and $i\phi_{x}=-\phi_{x}i$, the identity of (i)
gives $I((a+b\iota)x_{k})=(a^{2}+db^{2})I(x_{k})$ and
$I(x_{k})=\beta(k)e/3$. The index of $\ZZ x_{k}+\ZZ\iota x_{k}$ in $M_{k}$
divides the gcd $G(k)$ of the $2\times2$ minors of the matrix with rows
$x_{k}$, $\iota x_{k}$. These minors are integer polynomials of degree at most
two in $k$, so $G(k)$ divides every element of the ideal they generate in
$\ZZ[k]$ that is a constant; item (XXIV) finds such constants, resultants of
pairs of minors, whose gcd is $N_{d}$. Hence $N_{d}M_{k}\subset\ZZ x_{k}+\ZZ\iota
x_{k}$, and every nonzero $w\in M_{k}$ has $I(w)\ge I(x_{k})/N_{d}^{2}$. Now let
$y\in\NS\setminus\QQ\eta$ and $z=y-t(y)\eta$. Then $0\ne ez\in M_{k}$, because
$e\,y$ and $e\,t(y)=\int y\eta^{3}$ are integral, and
$I(y)=t(y)^{2}e+I(z)\ge I(ez)/e^{2}\ge\beta(k)/(3N_{d}^{2}e)$. The final
assertion follows from \Cref{lem:nefnorm}(ii).
\end{proof}

The bound in (iv) is far from sharp. The exact values computed in item (XXIV)
for $0\le k\le10$ are $\mu_{k}=\beta(k)\int\eta^{4}/3$ or half of it for
$d=1$, and $\mu_{k}=\beta(k)\int\eta^{4}/12$ for $d=3$; the proof loses the
factor $(\int\eta^{4})^{2}$ in passing from $y$ to its primitive part, and
only the growth in $k$ matters here. Each
locus of the pencil carries algebraic Weil classes, by the divisorial
criterion, and on the $k$-th of them the divisors that produce those classes
have degree growing like $k$. This is \Cref{prop:divroute} seen on explicit
quaternionic loci in one fixed polarised lattice; \Cref{fig:pencil} draws
it.

\begin{figure}[t]
\centering
  \includegraphics[width=\textwidth]{fig_pencil}
\caption{\Cref{prop:pencil} at $d=3$. Left: the members $S_{k}$ of the
pencil, each a Siegel threefold inside the four-dimensional family, with the
member $k=-2$, where $\beta=1$ and a rational point $J_{0}$ can be written
down, drawn apart. Right: along the pencil the ratio $\mu_{k}/\beta(k)$ is
the constant $2592$, so the Weil class is the same size at every member, while
the least degree of a divisor not proportional to $\eta$ grows linearly in
$k$. The class does not change; the divisors that produce it must.}
\label{fig:pencil}
\end{figure}

\begin{remark}[At $n=2$ the hypothesis is weaker still]\label{rem:ntwo}
Let $n=2$, so $\dim\cS=n^{2}=4$. The quaternionic locus of
\Cref{setup:quat} is the symmetric space of $\Sp(4,\RR)$ inside
$\cD_{2,2}$,\footnote{Left multiplication by the quaternion algebra
$B=(-d,b)$ commutes with $K$, and $B\otimes\RR\cong M_{2}(\RR)$ because
$b>0$, so Morita equivalence writes $V_{\RR}$ as $W\otimes_{\RR}\RR^{2}$ with
$W$ symplectic of dimension $2n$. The complex structures commuting with $B$
are then the complex structures on $W$, that is the points of the Siegel
space $\mathfrak{H}_{n}$, of dimension $n(n+1)/2$. The codimension in
$\cD_{n,n}$ is $n^{2}-n(n+1)/2=n(n-1)/2$, which is the count verified in
item (XII) of \Cref{sec:verification} for $n\le10$ and in
\texttt{quat\_locus\_codimension} of \Cref{sec:lean}.} of dimension
$n(n+1)/2=3$, so it is a divisor. Its $G(\QQ)$-translates are infinitely many
pairwise distinct irreducible closed subvarieties of $\cS$ of dimension $3$,
since a finite union of proper closed analytic subsets cannot be dense. If
infinitely many of them carry cycles with one and the same data
$(\underline{P},\underline{q})$, then the closed set
$\Sigma_{\underline{P},\underline{q}}$ has an irreducible component containing
infinitely many distinct irreducible threefolds, so that component has
dimension at least $4$ and therefore equals $\cS$. No density is needed at
$n=2$, only infinitude. Combined with the theorem of Moonen and Zarhin
\cite{MZ99} that every Hodge class on an abelian fourfold lies in the
subalgebra generated by divisor classes and Weil classes, the conclusion would
be the Hodge conjecture for all abelian fourfolds of Weil type, hence for all
abelian fourfolds. That conclusion is now Markman's theorem
\cite[Corollary 1.6.1]{Mar25a}, reached by a different route; at $n=2$ the
argument above would be a second proof, and it is at $n\ge3$ that it has
something to add.
\end{remark}

\subsubsection*{The second criterion: a rank at one point}

The other way to get from a dense algebraic locus to an open one is to make
the cycle move, and the classical tool for that is semiregularity. What
follows is not a new theorem about semiregularity; it is the observation that
after \Cref{thm:hdgcount}, \Cref{thm:everyfamily} and \Cref{thm:closedlocus},
the semiregularity of a single object at a single point implies the Hodge
conjecture for the whole family, with no further propagation argument. The
first step is that the object exists, and that its Chern character is as rigid
as it can be.

\begin{proposition}[The $K$-theory class at a base point]\label{prop:ktheory}
Let $s_{0}\in\Sigma$ be a base point of \Cref{thm:everyfamily} and write
$A_{0}=A_{s_{0}}$. There are an integer $N\ge1$ and a perfect complex $E$ on
$A_{0}$ with
\[
  \ch_{k}(E)=0\ \ (k\ne n),\qquad \ch_{n}(E)=N\,\omega_{1} .
\]
The tuple $v=\ch(E)$ extends to a flat section of
$\bigoplus_{k}R^{2k}f_{*}\QQ$ over $\cS$ which is of type $(k,k)$ at every
point.
\end{proposition}

\begin{proof}
By \Cref{thm:everyfamily} there is $\xi\in\CH^{n}(A_{0})_{\QQ}$ with
$\cl(\xi)=\omega_{1}$. For a smooth quasi-projective variety over a field the
Chern character $\ch\colon K_{0}(A_{0})_{\QQ}\to\CH^{*}(A_{0})_{\QQ}$ is a
ring isomorphism \cite[Ex.~15.2.16(b) and Cor.~18.3.2]{Ful98}, so there is a
unique $u\in K_{0}(A_{0})_{\QQ}$ with $\ch(u)=\xi$, and its Chern character
has no component outside degree $n$. Choose $N$ clearing the denominators of
$u$; then $Nu\in K_{0}(A_{0})$ is a difference of classes of vector bundles
and is therefore the class of a perfect complex $E$, with
$\ch(E)=N\xi$. Passing to cohomology gives the displayed identities.

For the last sentence, $\ch(E)$ has a single nonzero component, equal to
$N\omega_{1}$, and $\omega$ is a flat section of type $(n,n)$ everywhere by
\Cref{lem:omegadescends}.
\end{proof}

\begin{theorem}[Buchweitz and Flenner \cite{BF03}; Pridham \cite{Pri24}]\label{thm:BF}
Let $X$ be a smooth projective complex variety, $E$ a perfect complex on $X$,
and $X\hookrightarrow\cX\to\Spec R$ a deformation of $X$ over an Artin local
$\CC$-algebra along which $\ch(E)$ stays of Hodge type. If the semiregularity
map
\[
  \sigma_{E}\colon\Ext^{2}_{X}(E,E)
  \longrightarrow\bigoplus_{q\ge0}H^{q+2}(X,\Omega^{q}_{X}),
  \qquad
  \sigma_{E}(\zeta)=\sum_{q}\tfrac{1}{q!}\,
    \mathrm{tr}\bigl(\zeta\circ\mathrm{At}(E)^{q}\bigr),
\]
is injective, then $E$ extends to a perfect complex on $\cX$.
\end{theorem}

In this form, for perfect complexes and with every obstruction annihilated by
$\sigma_{E}$ and not only the curvilinear ones, the theorem is proved by
Pridham \cite{Pri24}, who identifies $\sigma_{E}$ with the tangent map of a
generalised Abel--Jacobi map on the derived moduli of perfect complexes; for
coherent sheaves see also \cite{BLM23}. Perry \cite{Per26} extends it to
noncommutative varieties with group actions, which covers twisted derived
categories and answers a question of Markman.

\begin{theorem}[Semiregularity at one point closes the family]
\label{thm:semiregclosure}
Let $n\ge2$, let $s_{0}$ and $E$ be as in \Cref{prop:ktheory}, and suppose the
semiregularity map $\sigma_{E}$ of \Cref{thm:BF} is
injective.\footnote{The hypothesis is not vacuous and it is not free, and the
degree-zero component shows both at once. For a line bundle $L$ one has
$\Ext^{2}(L,L)=H^{2}(\cO_{X})$ and $\sigma_{0}=\mathrm{tr}$ is the identity
on it, so $\sigma_{L}$ is injective for every line bundle on every variety,
and \Cref{thm:BF} returns the classical fact that a divisor class deforms
wherever it stays of type $(1,1)$. The object of \Cref{prop:ktheory} has rank
zero, so its trace component is not an isomorphism and in general not
injective, the source having no reason to be as small as
$h^{2}(\cO_{A_{0}})=\binom{2n}{2}$. The entire hypothesis therefore rests on
the components with $q\ge1$, which is exactly the range in which
semiregularity stops being formal and starts being a computation.} Then
$\Sigma=\cD_{n,n}$, and the Hodge conjecture holds in every degree for every
abelian $2n$-fold of $(K,-1,n)$-Weil type with $\Hg(A)=\SU(V,H)$ in that
family.
\end{theorem}

\begin{remark}[The hypothesis is not met by the object above]
\label{rem:semiregunavailable}
By \Cref{prop:concentrated} the complex of \Cref{prop:ktheory} has
$\dim\ker\sigma_{E}\ge n^{2}$ at every base point constructed in this paper,
so the hypothesis of \Cref{thm:semiregclosure} is not satisfied by it. The
theorem is stated because the hypothesis concerns a choice of perfect complex
representing a class in $K_{0}$, and a different representative is not
excluded; but no representative that is presently available satisfies it. See
\Cref{rem:killsreduction}, and \Cref{thm:p2forces} below for what any
representative satisfying it would have to look like.
\end{remark}

\begin{proof}
Let $\mathfrak{M}$ be the stack of perfect complexes with vanishing negative
self-extensions on the fibres of $f\colon\cA\to\cS$. It is algebraic and
locally of finite presentation over $\cS$ \cite{Lie06}, and at a point $[E]$
over $s_{0}$ its deformation theory relative to $\cS$ has tangent space
$\Ext^{1}(E,E)$ and obstruction space $\Ext^{2}(E,E)$. If
$\Ext^{<0}(E,E)\ne0$, which that stack excludes,\footnote{Lieblich's
stack parametrises the complexes he calls universally gluable, those with
$\Ext^{<0}(E,E)=0$; without that condition a complex can have automorphisms
of negative degree, and the natural moduli object is the derived stack of
To\"en and Vaqui\'e.} use instead the analytic Kuranishi germ of pairs of the
proof of \Cref{thm:p2false}, with the same tangent and obstruction spaces; the argument below then gives a nonempty
analytic open set in place of a Zariski open one, which is all that
\Cref{cor:allornothing} needs.

By \Cref{prop:ktheory} the Chern character of $E$ stays of Hodge type along
every direction in $\cS$, because it is a flat section of Hodge classes on the
whole base. So \Cref{thm:BF} applies to every infinitesimal deformation of
$A_{0}$ inside the family, and injectivity of $\sigma_{E}$ kills every
obstruction. The morphism $\mathfrak{M}\to\cS$ is therefore formally smooth at
$[E]$, and being locally of finite presentation it is smooth at $[E]$; in
particular it is dominant near $[E]$, so by Chevalley its image contains a
nonempty Zariski open $U\subseteq\cS$.

For $s\in U$ pick a perfect complex $E_{s}$ on $A_{s}$ in the fibre. Its
Chern character is $N\omega_{1}$ in degree $n$ and zero elsewhere, and Chern
characters of perfect complexes are algebraic, so $N\omega_{1}$ and hence
$\omega_{1}$ is algebraic on $A_{s}$; the rest of the Weil line follows by
\Cref{lem:onecomponent}. So $p^{-1}(U)\subseteq\Sigma$, which therefore
contains a nonempty open set, and \Cref{cor:allornothing} gives
$\Sigma=\cD_{n,n}$. Apply \Cref{thm:residual}(a).
\end{proof}

\subsubsection*{What the rank at one point forces}

The hypothesis of \Cref{thm:semiregclosure} is a rank, and a rank is not
open to inspection until the object is written down. It nevertheless has
consequences that can be stated for every representative at once, because the
semiregularity map composed with the evaluation map of Hochschild cohomology
is contraction into the Chern character (\Cref{setup:criterion}), and the
Chern character here is a Weil class. The annihilator of a Weil class in
$HT^{2}$ is large, and injectivity forces the evaluation map to vanish on the
whole of it. We compute the annihilator and draw the consequences; they are the
exact analogue, for (P2), of what \Cref{thm:candidatefails} and
\Cref{thm:nosplit} do for the secant route, and they use the same local
products.

\begin{theorem}[What injectivity of the semiregularity map forces]
\label{thm:p2forces}
Let $A$ be an abelian $2n$-fold of $(K,-1,n)$-Weil type, $n\ge2$, and write
$H^{1}(A,\CC)=V_{+}\oplus V_{-}$ for the eigenspaces of $K$,
$V_{\pm}^{1,0}$, $V_{\pm}^{0,1}$ for their Hodge components and
$T_{\pm}\subset H^{0}(T_{A})=T_{0}A$ for the eigenspaces of $K$ on the tangent
space, each of dimension $n$. Let $\omega=a\alpha_{+}+b\alpha_{-}$ be a
class on the Weil line with $ab\ne0$, for instance any nonzero rational one,
and let $E$ be a nonzero perfect complex on $A$ with $\ch(E)=N\omega$,
$N\ne0$, and no other component.
\begin{enumerate}[label=\textup{(\roman*)},leftmargin=2.6em,itemsep=2pt]
\item The annihilator of $\ch(E)$ in $HT^{2}(A)$ is
  \[
    \Ann_{HT^{2}}(\ch E)
    =\bigl(V_{+}^{0,1}\otimes V_{-}^{0,1}\bigr)
    \;\oplus\;\Hom_{K}\bigl(H^{1,0},H^{0,1}\bigr)
    \;\oplus\;\bigl(T_{+}\otimes T_{-}\bigr),
  \]
  of dimension $n^{2}+2n^{2}+n^{2}=4n^{2}$: the mixed block of $H^{0,2}$,
  which is the annihilator of \Cref{lem:annihilator}; the $K$-linear first
  order deformations of $A$, which contain the tangent space of the polarised
  Weil family; and the mixed bivectors.
\item If $\sigma_{E}$ is injective, then $\operatorname{ev}_{E}$ vanishes on
  this space. Explicitly: $z\cdot\id_{E}=0$ for every $z\in V_{+}^{0,1}\otimes
  V_{-}^{0,1}$; $E$ extends to first order along every $K$-linear deformation
  of $A$; and the translation classes satisfy $A_{\theta}A_{\theta'}=0$ in
  $\Ext^{2}(E,E)$ for every $\theta\in T_{+}$ and $\theta'\in T_{-}$.
\item Under the same hypothesis the algebra map $H^{*}(\cO_{A})\to\Ext^{*}(E,E)$,
  $z\mapsto z\cdot\id_{E}$, kills the ideal generated by the mixed block,
  which is $\bigoplus_{a,b\ge1}\bigwedge^{a}V_{+}^{0,1}\otimes\bigwedge^{b}V_{-}^{0,1}$;
  so the image of $H^{0,k}(A)$ has dimension at most $2\binom{n}{k}$ for
  $1\le k\le n$ and is zero for $k>n$, and every endomorphism $f$ of $E$ has
  $\operatorname{tr}(f)=0$ in $H^{0}(\cO_{A})=\CC$.
\item Under the same hypothesis, if at some point $p$ of its support $E$ is
  locally isomorphic to $i_{*}F$ with $F$ locally free of positive rank on a
  smooth closed subvariety $Z\ni p$ of codimension at least two, then
  $T_{+}\subseteq T_{p}Z$ or $T_{-}\subseteq T_{p}Z$.
\end{enumerate}
\end{theorem}

\begin{proof}
(i) Write $H^{*}(A,\CC)=\bigwedge^{*}(H^{1,0}\oplus H^{0,1})$ and
$\alpha_{\pm}$ for a generator of $\bigwedge^{2n}V_{\pm}
=\bigwedge^{n}V_{\pm}^{1,0}\otimes\bigwedge^{n}V_{\pm}^{0,1}$. The three
summands of $HT^{2}$ act as follows: $z\in H^{0,2}=\bigwedge^{2}H^{0,1}$ by
wedge product; $v\in H^{1}(T_{A})=\Hom(H^{1,0},H^{0,1})$ as the derivation
that replaces one $(1,0)$ factor by its image; $\pi\in H^{0}(\bigwedge^{2}T_{A})
=\bigwedge^{2}(H^{1,0})^{\vee}$ by double contraction of $(1,0)$ factors. For
$z$, the computation is \Cref{lem:annihilator}. For $v$, decompose
$\Hom(H^{1,0},H^{0,1})$ into its four blocks by eigenspaces. A block
$\Hom(V_{+}^{1,0},V_{+}^{0,1})$ replaces a factor of $\bigwedge^{n}V_{+}^{1,0}$
in $\alpha_{+}$ by an element of $V_{+}^{0,1}$, which is already exhausted in
$\alpha_{+}$, so it kills $\alpha_{+}$; it acts trivially on $\alpha_{-}$,
which has no $V_{+}^{1,0}$ factor. Likewise $\Hom(V_{-}^{1,0},V_{-}^{0,1})$
kills $\omega$. The block $\Hom(V_{+}^{1,0},V_{-}^{0,1})$ carries $\alpha_{+}$
to a nonzero element of $\bigwedge^{n-1}V_{+}^{1,0}\otimes V_{-}^{0,1}
\otimes\bigwedge^{n}V_{+}^{0,1}$ and kills $\alpha_{-}$, and
$\Hom(V_{-}^{1,0},V_{+}^{0,1})$ carries $\alpha_{-}$ to a nonzero element of
a different summand of $H^{n-1,n+1}$; since $ab\ne0$ and the two summands are
independent, a general $v$ kills $\omega$ if and only if its two mixed blocks
vanish, that is, if and only if $v$ is $K$-linear. The space of $K$-linear
maps has dimension $2n^{2}$, and it contains the $n^{2}$-dimensional tangent
space of the polarised family, which is its symmetric part. For $\pi$, a
mixed bivector $\theta\wedge\theta'$ with $\theta\in T_{+}=(V_{+}^{1,0})^{\vee}$
and $\theta'\in T_{-}=(V_{-}^{1,0})^{\vee}$ kills $\alpha_{+}$ because
$\theta'$ kills every factor of $\bigwedge^{n}V_{+}^{1,0}$, and kills
$\alpha_{-}$ symmetrically; a bivector in $\bigwedge^{2}T_{+}$ carries
$\alpha_{+}$ to a nonzero element of $\bigwedge^{n-2}V_{+}^{1,0}
\otimes\bigwedge^{n}V_{+}^{0,1}$ and kills $\alpha_{-}$, and symmetrically
for $\bigwedge^{2}T_{-}$. Since the images of the three summands of $HT^{2}$
lie in the three different Hodge components $H^{n,n+2}$, $H^{n-1,n+1}$ and
$H^{n-2,n}$ of $H^{2n+2}(A,\CC)$, the annihilator is the direct sum of the
three annihilators computed, of dimensions $n^{2}$, $2n^{2}$ and $n^{2}$. Item
(XXXV) of \Cref{sec:verification} carries out the whole computation in the
exterior algebra for $n=2$ and $n=3$.

(ii) By the commutative square $\sigma_{E}\circ\operatorname{ev}_{E}
=(\,\cdot\,)\lrcorner\ch(E)$ of \Cref{setup:criterion}, $\sigma_{E}$
vanishes on $\operatorname{ev}_{E}(\Ann_{HT^{2}}(\ch E))$; if $\sigma_{E}$
is injective, $\operatorname{ev}_{E}$ vanishes there. The three explicit
statements are \eqref{eq:evgeneral} on the three summands of (i): on
$H^{0,2}$ the evaluation map is $z\mapsto z\cdot\id_{E}$; on $H^{1}(T_{A})$
it is the obstruction map $v\mapsto v\lrcorner\mathrm{At}(E)$, whose
vanishing on $v$ is the extension of $E$ to first order along $v$
\cite{Tod09}; and on a bivector $\theta\wedge\theta'$ it is
$\tfrac12(\theta\wedge\theta')\lrcorner\mathrm{At}(E)^{2}
=A_{\theta}A_{\theta'}$, as in the proof of \Cref{thm:candidatefails}.

(iii) The map $z\mapsto z\cdot\id_{E}$ is a homomorphism of graded algebras
from $H^{*}(\cO_{A})=\bigwedge^{*}H^{0,1}$ to the Yoneda algebra, because the
classes from $H^{*}(A,\CC)$ are graded central there, so its kernel is an
ideal, and by (ii) it contains the mixed block $V_{+}^{0,1}\otimes V_{-}^{0,1}$.
The ideal generated by the mixed block in $\bigwedge^{*}(V_{+}^{0,1}\oplus
V_{-}^{0,1})=\bigwedge^{*}V_{+}^{0,1}\otimes\bigwedge^{*}V_{-}^{0,1}$ is the
sum of the summands $\bigwedge^{a}V_{+}^{0,1}\otimes\bigwedge^{b}V_{-}^{0,1}$
with $a,b\ge1$, since each such summand is a multiple of the mixed block and
the remaining summands, $a=0$ or $b=0$, meet the ideal in zero; the quotient in
degree $k\ge1$ is $\bigwedge^{k}V_{+}^{0,1}\oplus\bigwedge^{k}V_{-}^{0,1}$,
of dimension $2\binom{n}{k}$, and zero for $k>n$. In particular the top class
of $H^{2n}(\cO_{A})$, which is mixed for $n\ge1$, acts by zero. Serre duality
on $A$, whose canonical bundle is trivial, is the perfect pairing
$\Ext^{2n}(E,E)\times\Hom(E,E)\to H^{2n}(\cO_{A})\cong\CC$,
$(\xi,f)\mapsto\operatorname{tr}(\xi\circ f)$; for $\xi=w\cdot\id_{E}$
with $w$ the top class this is $w\cdot\operatorname{tr}(f)$, where
$\operatorname{tr}(f)\in H^{0}(\cO_{A})=\CC$, so $w\cdot\id_{E}=0$ forces
$\operatorname{tr}(f)=0$ for every $f$.

(iv) Let $\cO$ be the local ring of $A$ at $p$ and choose coordinates with
$Z=\{x_{1}=\dots=x_{c}=0\}$, $c\ge2$. Locally $E\cong(\cO/J)^{\oplus r}$
with $J=(x_{1},\dots,x_{c})$, a regular sequence, so
$\Ext^{*}_{\cO}(\cO/J,\cO/J)=\bigwedge^{*}_{\cO/J}N$ as a graded algebra,
where $N=\Hom(J/J^{2},\cO/J)$ is the normal module with basis the classes
$\partial_{1},\dots,\partial_{c}$ dual to $x_{1},\dots,x_{c}$; the
$\Ext^{1}$ is computed on the Koszul resolution exactly as in the proof of
\Cref{lem:localproducts}, and the product is the wedge because the Yoneda
algebra of a Koszul complex is the exterior algebra on the dual of the
generators. For $E$ of rank $r$ the local Yoneda algebra is
$\bigwedge^{*}N\otimes M_{r}(\cO/J)$. By \Cref{lem:edge} the germ of
$A_{\theta}$ at $p$ is the jet class along $\theta$; on the Koszul complex the
jet class of a derivation $D$ is the chain map $-D(d)$, and for
$D=\theta=\sum_{i}c^{i}\partial_{i}$ its component in $\Ext^{1}=N\otimes M_{r}(\cO/J)$
is $\sum_{i\le c}(c^{i}\bmod J)\,\partial_{i}\otimes\id$, the class of
$\theta$ modulo $T_{Z}$. Multiplicativity of the germ map then gives the germ
of $A_{\theta}A_{\theta'}$ as
$(\theta\bmod T_{Z})\wedge(\theta'\bmod T_{Z})\otimes\id$ in
$\bigwedge^{2}N\otimes M_{r}(\cO/J)$, whose value at $p$ is
$\bar\theta(p)\wedge\bar\theta'(p)\in\bigwedge^{2}(T_{p}A/T_{p}Z)$. If
$A_{\theta}A_{\theta'}=0$ for all $\theta\in T_{+}$, $\theta'\in T_{-}$,
then $\bar a\wedge\bar b=0$ for all $\bar a$ in the image $I_{+}$ of $T_{+}$
and $\bar b$ in the image $I_{-}$ of $T_{-}$ in $N_{p}=T_{p}A/T_{p}Z$, and
$I_{+}+I_{-}=N_{p}$ since $T_{+}\oplus T_{-}=T_{p}A$. If both $I_{\pm}$ were
nonzero, every nonzero $\bar b\in I_{-}$ would be proportional to every
nonzero $\bar a\in I_{+}$, so $I_{+}=I_{-}$ would be a line and
$N_{p}=I_{+}+I_{-}$ of dimension one, against $c\ge2$. Hence $I_{+}=0$ or
$I_{-}=0$, which is the claim. The two pieces of linear algebra, the shape of
the ideal in (iii) and the dichotomy just used, are checked in item (XXXV).
\end{proof}

\begin{corollary}[Where (P2) can and cannot be met]\label{cor:p2shape}
Let $A$ be as in \Cref{thm:p2forces} and suppose that no proper abelian
subvariety $B\subset A$ has $T_{0}B\supseteq T_{+}$ or $T_{0}B\supseteq T_{-}$;
equivalently, that $A$ has no $K$-stable abelian subvariety on which $K$ acts
through a single embedding, which holds whenever $A$ is simple, at every
split member $X\times\widehat X$ with $X$ simple of endomorphism algebra
$\QQ$, and at a general point of the quaternionic locus of \Cref{setup:quat}.
Let $E$ be a perfect complex on $A$ with $\ch(E)=N\omega$ as above and
$\sigma_{E}$ injective. Then no irreducible component $Z_{0}$ of the support
of $E$ of codimension at least two has a dense open subset along which $E$ is
a locally free sheaf of positive rank on $Z_{0}$. In particular $E$ is not a
sheaf, and along every component of its support of codimension at least two
it has, generically, either at least two nonvanishing cohomology sheaves or a
cohomology sheaf that is not locally free on the reduced support.
\end{corollary}

\begin{proof}
Suppose $Z_{0}$ is such a component, with $U\subseteq Z_{0}^{\mathrm{sm}}$ dense
open and $E|_{U}$ a locally free sheaf of positive rank on $U$. By
\Cref{thm:p2forces}(iv), $U=U_{+}\cup U_{-}$ with
$U_{\pm}=\{p\in U:T_{\pm}\subseteq T_{p}Z_{0}\}$, both closed in $U$; as $U$
is irreducible, $U=U_{+}$ say. Then every constant vector field
$\theta\in T_{+}$ is tangent to $Z_{0}$ along $U$, hence along
$Z_{0}^{\mathrm{sm}}$. The set
$S=\{(p,t)\in Z_{0}\times\CC:p+\exp(t\theta)\in Z_{0}\}$ is a closed
analytic subset of the irreducible space $Z_{0}\times\CC$, and it contains
an open neighbourhood of $Z_{0}^{\mathrm{sm}}\times\{0\}$, because the local
flow of the restriction of $\theta$ to the submanifold $Z_{0}^{\mathrm{sm}}$
exists on such a neighbourhood; so $S=Z_{0}\times\CC$, and
$Z_{0}+\exp(t\theta)\subseteq Z_{0}$, hence $=Z_{0}$, for every $t$. Let
$G=\{x\in A:Z_{0}+x=Z_{0}\}$, a closed algebraic subgroup of $A$, and $B$
its identity component, an abelian subvariety. The connected set
$\exp(T_{+})$ contains $0$ and lies in $G$, hence in $B$, so
$T_{+}\subseteq T_{0}B$. By hypothesis $B=A$, so $Z_{0}=A$, against the
codimension. If $E$ were a sheaf, its support would have codimension exactly
$n\ge2$ and along a dense open set of every component it would be locally
free of positive rank on the reduced support, which is excluded; the last
sentence is the negation of the same condition.

For the equivalence in the hypothesis, let $\sqrt{-d}$ act on $A$ through
an order of $K$. If $B\ne A$ has $T_{0}B\supseteq T_{+}$, then
$H_{1}(B,\QQ)\cap\sqrt{-d}\,H_{1}(B,\QQ)$ is a $K$-stable sub-Hodge
structure, hence $H_{1}(C,\QQ)$ for a $K$-stable abelian subvariety
$C\subseteq B$, and its real points $T_{0}B\cap\sqrt{-d}\,T_{0}B$ still
contain $T_{+}$, on which $\sqrt{-d}$ acts by a scalar. The $K$-stable
complement $C'$ of $C$ for the polarisation is nonzero, since $C\subseteq B\ne A$,
and $T_{0}C'$ is a $K$-stable complement of $T_{0}C\supseteq T_{+}$, so
$T_{0}C'\subseteq T_{-}$ and $K$ acts on $C'$ through a single embedding.
Conversely, if a nonzero $K$-stable $C'\ne A$ has $T_{0}C'\subseteq T_{-}$,
its $K$-stable complement $B$ is proper with $T_{0}B\supseteq T_{+}$. For the
three cases: a simple $A$ has no proper abelian subvariety; for
$X\times\widehat X$ with $X$ simple and $\End^{0}(X)=\QQ$ the sub-Hodge
structures of $H_{1}(X)\otimes\QQ^{2}$ are $H_{1}(X)\otimes W$ with
$W\subseteq\QQ^{2}$, and $W$ is $K$-stable only for $W=0$ or $\QQ^{2}$, $K$
acting on $\QQ^{2}$ through a quadratic field; and at a general quaternionic
member $A$ is either simple or, when the quaternion algebra of
\Cref{setup:quat} splits, isogenous to $C\times C$ with $C$ a general
principally polarised abelian $n$-fold, whose sub-Hodge structures are again
$H_{1}(C)\otimes W$ with the same conclusion.
\end{proof}

\subsubsection*{The whole Hochschild action, and the two branches}

\Cref{thm:p2forces} computes the annihilator degree by degree and reads off
what the vanishing of the evaluation map on it means for $E$. The annihilator
has a structure that the degree-by-degree description hides, and bringing it
out costs nothing and gives more: it is the mixed part of a splitting of
$HH^{1}(A)$ into two halves, one for each of the two conjugate pieces of the
Weil class. Everything in this subsection is about the Hochschild cohomology
of $A$ alone, and the consequences for $E$ follow from it formally.

\begin{setupenv}[Hochschild cohomology in these terms]\label{setup:hh}
For an abelian variety $A$ the Hochschild--Kostant--Rosenberg isomorphism
identifies
\[
  HH^{*}(A)\;\cong\;HT^{*}(A)=\bigoplus_{p,q}H^{p}\bigl(\textstyle\bigwedge^{q}T_{A}\bigr)
  =\textstyle\bigwedge^{*}\bigl(H^{0,1}(A)\oplus H^{0}(T_{A})\bigr),
\]
as a graded ring, the grading being $p+q$; the Todd class of $A$ is trivial,
so the twist by $\td^{1/2}$ in \cite{Cal05}, \cite{CRV12} is the identity and
the identification is a ring isomorphism on the nose. Write
$HH^{1}(A)=H^{0,1}(A)\oplus H^{0}(T_{A})$, of dimension $4n$, so that
$HH^{*}(A)$ is the exterior algebra on it, of dimension $2^{4n}$. It acts on
$H^{*}(A,\CC)=\bigwedge^{*}\bigl(H^{1,0}\oplus H^{0,1}\bigr)$ by wedging with
the classes of $H^{0,1}$ and contracting the classes of $H^{1,0}$ against
$H^{0}(T_{A})$; the two families of operators involve disjoint generators, so
they pairwise anticommute and the action is an action of the exterior algebra.
In degree two this is the operator $x\lrcorner$ of \Cref{setup:criterion}, and
$\operatorname{ev}_{E}\colon HH^{*}(A)\to\Ext^{*}(E,E)$ is a homomorphism of
graded rings into the graded centre, with
$\sigma_{E}\circ\operatorname{ev}_{E}=(\,\cdot\,)\lrcorner\ch(E)$
\cite[Section 6]{BF08}, \cite{Cal05}, \cite{CRV12}.
\end{setupenv}

\begin{theorem}[The annihilator is the mixed part of a splitting]
\label{thm:hhsplitting}
Let $A$ be of $(K,-1,n)$-Weil type with $n\ge2$, and let
$\omega=a\alpha_{+}+b\alpha_{-}$ lie on the Weil line with $ab\ne0$. Put
\[
  P=\Ann_{HH^{1}}(\alpha_{+})=V_{+}^{0,1}\oplus T_{-},
  \qquad
  Q=\Ann_{HH^{1}}(\alpha_{-})=V_{-}^{0,1}\oplus T_{+} .
\]
\begin{enumerate}[label=\textup{(\roman*)},leftmargin=2.6em,itemsep=2pt]
\item $\dim P=\dim Q=2n$ and $HH^{1}(A)=P\oplus Q$; in particular
  $\Ann_{HH^{1}}(\omega)=0$.
\item $\Ann_{HH^{2}}(\omega)=P\wedge Q$, of dimension $4n^{2}$. This is the
  space computed in \Cref{thm:p2forces}(i), written as the mixed part of the
  splitting.
\item Contraction with $\alpha_{-}$ is injective on $\bigwedge^{*}P$, and
  contraction with $\alpha_{+}$ is injective on $\bigwedge^{*}Q$; the two
  images meet in exactly one line, in degree $2n$.
\item For every $k\ge1$,
  \[
    \Ann_{HH^{k}}(\omega)
    =\Bigl(\bigoplus_{\substack{i+j=k\\ i,j\ge1}}
      \textstyle\bigwedge^{i}P\wedge\bigwedge^{j}Q\Bigr)\;\oplus\;\epsilon_{k},
  \]
  where $\epsilon_{k}=0$ for $k\ne2n$ and $\epsilon_{2n}$ is the line of
  \textup{(iii)}. So
  $\dim\Ann_{HH^{k}}(\omega)=\binom{4n}{k}-2\binom{2n}{k}+\delta_{k,2n}$ for
  $k\ge1$, while $\Ann_{HH^{0}}(\omega)=0$, and the annihilator is the ideal
  generated by its degree-two part together with that one line.
\end{enumerate}
\end{theorem}

\begin{proof}
Choose bases $x_{1},\dots,x_{n}$ of $V_{+}^{1,0}$ and $y_{1},\dots,y_{n}$ of
$V_{-}^{1,0}$, and bases $\bbar x_{i}$ of $V_{+}^{0,1}$ and $\bbar y_{j}$ of
$V_{-}^{0,1}$ (the bar is notation only: complex conjugation carries
$V_{+}^{1,0}$ to $V_{-}^{0,1}$), so that $\alpha_{+}=x_{1}\cdots x_{n}\bbar x_{1}\cdots\bbar x_{n}$
and $\alpha_{-}=y_{1}\cdots y_{n}\bbar y_{1}\cdots\bbar y_{n}$, and write
$\partial_{x_{i}},\partial_{y_{j}}$ for the dual basis of
$H^{0}(T_{A})=T_{+}\oplus T_{-}$.

(i) Of the $4n$ basis operators, $\bbar x_{i}\wedge$ kills $\alpha_{+}$
because $\bbar x_{i}$ already occurs in it, and $\partial_{y_{j}}\lrcorner$
kills $\alpha_{+}$ because no $y$ occurs in it; so $P\supseteq
V_{+}^{0,1}\oplus T_{-}$. Conversely $\bbar y_{j}\wedge\alpha_{+}\ne0$ and
$\partial_{x_{i}}\lrcorner\alpha_{+}\ne0$, and these images lie in different
summands of $H^{*}$, so no nonzero combination of them kills $\alpha_{+}$:
$P=V_{+}^{0,1}\oplus T_{-}$, of dimension $2n$. The computation for $Q$ is the
same with the roles of the two eigenspaces exchanged, and
$P\cap Q=0$ with $\dim P+\dim Q=4n$, so $HH^{1}=P\oplus Q$. A class killing
$\omega$ in degree one kills $\alpha_{+}$ and $\alpha_{-}$ separately, the two
images lying in different summands and $ab\ne0$, so it lies in $P\cap Q=0$.

(iii) A monomial $\bbar x_{I}\partial_{y_{J}}\in\bigwedge^{*}P$ contracts
$\alpha_{-}$ to $\pm\,\bbar x_{I}\,y_{[n]\setminus J}\,\bbar y_{[n]}$, and
these are linearly independent for distinct pairs $(I,J)$, which is the
injectivity; the statement for $Q$ is symmetric. For the intersection, the
image of $\bigwedge^{*}P$ consists of classes with no $x$ and at most $n$
factors $\bbar x$, and the image of $\bigwedge^{*}Q$ of classes with $n$
factors $\bbar x$ and $n-\#\{\partial_{x}\}$ factors $x$. A common class
therefore has no $x$ and all $n$ factors $\bbar x$ and, by the same count in
the other eigenspace, no $y$ and all $n$ factors $\bbar y$: it is a multiple
of $\bbar x_{[n]}\bbar y_{[n]}$, and it comes from
$\pi_{P}=\bbar x_{[n]}\partial_{y_{[n]}}$ and
$\pi_{Q}=\bbar y_{[n]}\partial_{x_{[n]}}$, both of degree $2n$. So the
intersection is the single line spanned by that class.

(ii) and (iv) Decompose $\bigwedge^{k}HH^{1}=\bigoplus_{i+j=k}\bigwedge^{i}P
\wedge\bigwedge^{j}Q$. A summand with $i,j\ge1$ annihilates $\alpha_{+}$,
because it has a factor in $P$ and the action is by an algebra, and
annihilates $\alpha_{-}$ for the same reason with $Q$; so it lies in the
annihilator of $\omega$. On the remaining two summands, $\xi_{P}\in\bigwedge^{k}P$
has $\xi_{P}\lrcorner\omega=b\,\xi_{P}\lrcorner\alpha_{-}$ and
$\xi_{Q}\in\bigwedge^{k}Q$ has $\xi_{Q}\lrcorner\omega=a\,\xi_{Q}\lrcorner\alpha_{+}$,
so by (iii) a combination $\xi_{P}+\xi_{Q}$ annihilates $\omega$ if and only
if $b\,\xi_{P}\lrcorner\alpha_{-}=-a\,\xi_{Q}\lrcorner\alpha_{+}$, which by
(iii) happens only for $k=2n$ and then only on the line
$b^{-1}\pi_{P}-a^{-1}\pi_{Q}$ up to scale. Taking $k=2$ and $n\ge2$, so that
$2\ne2n$, gives (ii), the dimension being
$\dim P\cdot\dim Q=4n^{2}$; the general count is
$\binom{4n}{k}-\binom{2n}{k}-\binom{2n}{k}$ plus the exceptional line. Item
(XXXVIII) of \Cref{sec:verification} carries out the whole computation in the
exterior algebra over $\QQ$ for $n=1$ and $n=2$ in every degree, for two
choices of $(a,b)$, and for $n=3$ in degrees up to three.
\end{proof}

\begin{corollary}[What the criterion forces on the Hochschild action]
\label{cor:hhfactor}
Let $A$ and $\omega$ be as above and let $E$ be a nonzero perfect complex on
$A$ with $\ch(E)=N\omega$, $N\ne0$.
\begin{enumerate}[label=\textup{(\roman*)},leftmargin=2.6em,itemsep=2pt]
\item Unconditionally, $\dim\Ext^{2}(E,E)\ge2n(2n-1)$.
\item If $\sigma_{E}$ is injective then $\operatorname{ev}_{E}$ vanishes on
  $P\wedge Q$, hence on the ideal it generates, and the ring homomorphism
  $\operatorname{ev}_{E}\colon HH^{*}(A)\to\Ext^{*}(E,E)$ factors through
  \[
    R\;=\;\textstyle\bigwedge^{*}P\;\oplus_{\CC}\;\bigwedge^{*}Q,
    \qquad \dim_{\CC}R=2^{2n+1}-1,
  \]
  the two exterior algebras glued along their scalars, inside $HH^{*}(A)$ of
  dimension $2^{4n}$. At $n=2$ that is $31$ out of $256$.
\item Granting the compatibility
  $\sigma_{E}\circ\operatorname{ev}_{E}=(\,\cdot\,)\lrcorner\ch(E)$ in degree
  $k$ \cite{Cal05}, \cite{CRV12}, the induced map on $R$ is injective away
  from the line of \Cref{thm:hhsplitting}(iii), so
  \[
    \dim\Ext^{k}(E,E)\;\ge\;2\binom{2n}{k}\quad(1\le k<2n),
    \qquad
    \sum_{k\ge1}\dim\Ext^{k}(E,E)\;\ge\;2^{2n+1}-3 .
  \]
\end{enumerate}
\end{corollary}

\begin{proof}
(i) By \Cref{setup:hh}, $\ker\operatorname{ev}_{E}\subseteq\Ann_{HH^{2}}(\ch E)$
in degree two, since $\sigma_{E}(\operatorname{ev}_{E}(\xi))=\xi\lrcorner\ch(E)$
and the right side is nonzero off the annihilator. Hence
$\dim\Ext^{2}(E,E)\ge\dim HH^{2}(A)-\dim\Ann_{HH^{2}}(\omega)
=\binom{4n}{2}-4n^{2}=2n(2n-1)$ by \Cref{thm:hhsplitting}(ii). No hypothesis
on $\sigma_{E}$ is used.

(ii) $\operatorname{ev}_{E}$ vanishes on the annihilator by
\Cref{thm:p2forces}(ii), its kernel is an ideal because it is a ring map, and
the ideal generated by $P\wedge Q$ in $\bigwedge^{*}(P\oplus Q)
=\bigwedge^{*}P\otimes\bigwedge^{*}Q$ is the sum of the summands with
$i,j\ge1$, whose complement is $\bigwedge^{*}P\oplus_{\CC}\bigwedge^{*}Q$ of
dimension $2\cdot2^{2n}-1$.

(iii) With the compatibility in degree $k$, the argument of (i) gives
$\ker\operatorname{ev}_{E}\subseteq\Ann_{HH^{k}}(\omega)$, and by
\Cref{thm:hhsplitting}(iv) that annihilator meets $\bigwedge^{k}P\oplus
\bigwedge^{k}Q$ only in the exceptional line. So the map induced on $R$ has
kernel contained in that one line and its image has dimension at least
$\dim R-1=2^{2n+1}-2$. One of those dimensions sits in degree zero, where the
image is $\CC\cdot\id_{E}$, nonzero because $E$ is; the rest is distributed
over the positive degrees, which is the second display, and in degree $k$ the
image contains the $2\binom{2n}{k}$ dimensions of
$\bigwedge^{k}P\oplus\bigwedge^{k}Q$ for $k<2n$, which is the first.
\end{proof}

\begin{remark}[The two branches, and what they are not]\label{rem:twobranches}
The ring $R$ of \Cref{cor:hhfactor}(ii) is the ring of functions on two odd
affine spaces of dimension $2n$ crossing at the origin. Read that way, the
statement is that if $E$ satisfies the criterion then the deformations of
$D^{b}(A)$ along which its Hochschild action survives form two branches, $P$
and $Q$, meeting only at the trivial deformation; and the two branches are the
annihilators of the two conjugate pieces $\alpha_{+}$ and $\alpha_{-}$ of the
Weil class. The criterion says that $E$ sees the two halves of the class
separately and their interaction not at all.

Two things this is not. The splitting $HH^{1}=P\oplus Q$ is not the splitting
into eigenspaces of $K$: the eigenspaces are
$V_{+}^{0,1}\oplus T_{+}$ and $V_{-}^{0,1}\oplus T_{-}$, and $P$ and $Q$ cross
them. It is the splitting into the annihilators of $\alpha_{+}$ and of
$\alpha_{-}$, and it exists because the Weil class is a sum of two pieces each
of which is a top form on an eigenspace. And the conclusion is not a
contradiction: $R$ is a perfectly good ring, the bound of (iii) is consistent
with $\chi(E,E)$ computed in \Cref{prop:objectsize}, since
$\sum_{k}(-1)^{k}2\binom{2n}{k}=0$, and nothing here forbids an $E$ of the
required shape. What the corollary does is say how large such an $E$ must be
and how little of the Hochschild cohomology it may see. We have not
constructed one, and we do not claim that one exists.
\end{remark}

\subsubsection*{The criterion as a dimension, and where the object can live}

The same square gives two more things, and between them they change what one
would look for. The first converts the rank of $\sigma_{E}$ into a single
dimension, which is the form in which a candidate can actually be tested. The
second says where the support of $E$ may lie, and it excludes the place where
anyone would look first, namely a cycle representing the Weil class.

\begin{theorem}[The criterion is a dimension]\label{thm:p2numerical}
Let $A$, $\omega$ and $E$ be as in \Cref{cor:hhfactor}, so that
$\ch(E)=N\omega$ with $N\ne0$ and no other component.
\begin{enumerate}[label=\textup{(\roman*)},leftmargin=2.6em,itemsep=2pt]
\item If $\dim\Ext^{2}(E,E)=2n(2n-1)$, the least value that
  \Cref{cor:hhfactor}(i) allows, then $\operatorname{ev}_{E}\colon
  HH^{2}(A)\to\Ext^{2}(E,E)$ is surjective with kernel exactly $P\wedge Q$,
  the map $\sigma_{E}$ is injective, and every class of $\Ext^{2}(E,E)$ is a
  sum of products $\operatorname{ev}_{E}(p)\operatorname{ev}_{E}(p')
  +\operatorname{ev}_{E}(q)\operatorname{ev}_{E}(q')$ with $p,p'\in P$ and
  $q,q'\in Q$, hence lies in the graded centre of $\Ext^{*}(E,E)$.
\item Conversely, if $\sigma_{E}$ is injective, then
  $\operatorname{ev}_{E}(HH^{2}(A))$ has dimension exactly $2n(2n-1)$, and
  $\sigma_{E}$ carries it isomorphically onto $HH^{2}(A)\lrcorner\,\omega$.
\item Consequently, if at one base point $s_{0}$ of a family some perfect
  complex $E$ with $\ch(E)=N\omega_{1}$ has $\dim\Ext^{2}(E,E)=2n(2n-1)$,
  then $\Sigma=\cD_{n,n}$, and the Hodge conjecture holds in every degree for
  every member of that family with $\Hg(A)=\SU(V,H)$.
\end{enumerate}
\end{theorem}

\begin{proof}
(i) The square of \Cref{setup:hh} gives $\ker\operatorname{ev}_{E}
\subseteq\Ann_{HH^{2}}(\omega)=P\wedge Q$ in degree two, the equality being
\Cref{thm:hhsplitting}(ii). Hence $\operatorname{rank}\operatorname{ev}_{E}
\ge\binom{4n}{2}-4n^{2}=2n(2n-1)=\dim\Ext^{2}(E,E)$, so
$\operatorname{ev}_{E}$ is onto and its kernel, of dimension $4n^{2}$ and
contained in $P\wedge Q$, is all of $P\wedge Q$. Every class of
$\Ext^{2}(E,E)$ is then $\operatorname{ev}_{E}(\xi)$ for some $\xi$, and
$\sigma_{E}(\operatorname{ev}_{E}(\xi))=\xi\lrcorner\,\omega$ vanishes only for
$\xi\in P\wedge Q=\ker\operatorname{ev}_{E}$; so $\sigma_{E}$ is injective.
Since $HH^{2}(A)=\bigwedge^{2}P\oplus(P\wedge Q)\oplus\bigwedge^{2}Q$ and
$\operatorname{ev}_{E}$ is a ring map into the graded centre, the image is
$\operatorname{ev}_{E}(\bigwedge^{2}P)+\operatorname{ev}_{E}(\bigwedge^{2}Q)$,
which is the description of the classes.

(ii) By \Cref{thm:p2forces}(ii) the kernel of $\operatorname{ev}_{E}$ in
degree two contains $P\wedge Q$, and by the argument of (i) it is contained in
it, so it equals $P\wedge Q$ and the image has dimension $2n(2n-1)$. On that
image $\sigma_{E}\circ\operatorname{ev}_{E}$ is contraction into $\omega$,
whose kernel is exactly $P\wedge Q$, so $\sigma_{E}$ maps the image
isomorphically onto $HH^{2}(A)\lrcorner\,\omega$.

(iii) By (i) the map $\sigma_{E}$ is injective, and \Cref{thm:semiregclosure}
applies.
\end{proof}

\Cref{thm:p2numerical}(iii) is a sufficient condition for
$\mathbf{W}(K,n,\delta)$ that mentions no map at all: one object, at one
point, and one number.%
\footnote{The dimension of $\Ext^{2}(E,E)$ is computable in a way the rank of
$\sigma_{E}$ is not. For a complex built from sheaves whose cohomology is
known, the local to global spectral sequence
$E_{2}^{p,q}=H^{p}(A,\cE xt^{q}(E,E))\Rightarrow\Ext^{p+q}(E,E)$ reduces the
count to the cohomology of the local $\cE xt$ sheaves, and on an abelian
variety the cohomology of a sheaf supported on an abelian subvariety, or of a
homogeneous bundle, is a question of linear algebra. The rank of $\sigma_{E}$,
by contrast, requires the trace of the exponential of the Atiyah class on
every class of $\Ext^{2}$, which is a computation with the object and not only
with its invariants.}
At $n=2$ the number is $12$, at $n=3$ it is $30$, and the two computations
behind it are item (XXXIX) of \Cref{sec:verification}. Part (ii) says that
nothing is lost in reading the criterion this way on the central part of
$\Ext^{2}$; whether an injective $\sigma_{E}$ forces the whole of
$\Ext^{2}(E,E)$ to be central is a statement about $\Ext^{2n-2}$, by Serre
duality, which the square in degree two does not reach, and nothing below
uses it.

The second consequence rests on a local fact about complexes supported at a
point, which we state and prove first.

\begin{lemma}[Two jet classes on a complex of finite length]
\label{lem:finiteproduct}
Let $R=\CC[[x_{1},\dots,x_{c}]]$ with $c\ge2$, let $M$ be a bounded complex of
finitely generated free $R$-modules whose cohomology has finite length, and
put $\chi(M)=\sum_{i}(-1)^{i}\ell(H^{i}(M))$. For $a\in\CC^{c}$ let
$J(a)\in\Ext^{1}_{R}(M,M)$ be the jet class of $M$ along the constant vector
field $\sum_{i}a_{i}\partial_{i}$, represented by the chain map
$-\sum_{i}a_{i}\partial_{i}(d)$ as in the proof of \Cref{thm:p2forces}(iv).
If $\chi(M)\ne0$ and $a,b$ are linearly independent, then $J(a)J(b)\ne0$ in
$\Ext^{2}_{R}(M,M)$.
\end{lemma}

\begin{proof}
Let $T$ be an abelian variety of dimension $c$ and $t\in T$ a point. Linear
coordinates on the universal cover identify the completed local ring
$\widehat{\cO}_{T,t}$ with $R$ in such a way that the translation invariant
vector fields become the constant fields $\sum_{i}a_{i}\partial_{i}$.%
\footnote{The exponential map of $T$ is a local biholomorphism from
$T_{t}T=\CC^{c}$, and it carries the constant vector fields on $\CC^{c}$ to
the translation invariant ones. Its germ at the origin identifies
$\widehat{\cO}_{T,t}$ with $\CC[[x_{1},\dots,x_{c}]]$, and the identification
is compatible with derivations. Nothing about $T$ beyond its being a compact
torus carrying a projective structure is used.}
A complex of $R$-modules with finite length cohomology is the same thing as a
complex of $\cO_{T,t}$-modules with finite length cohomology, since both are
modules over $\cO_{T,t}/\mathfrak{m}^{N}$ for $N$ large, so $M$ defines a
perfect complex $M_{T}$ on $T$ supported at $t$. Its local $\cE xt$ sheaves are
skyscrapers at $t$, so the local to global spectral sequence gives
$\Ext^{*}_{T}(M_{T},M_{T})=\Ext^{*}_{R}(M,M)$, and by \Cref{lem:edge} the
translation classes $A_{a}$, $A_{b}$ of $M_{T}$ restrict to $J(a)$, $J(b)$.
Apply the square of \Cref{setup:criterion} on $T$ to the bivector $a\wedge b$:
\[
  \sigma_{M_{T}}\bigl(A_{a}A_{b}\bigr)=(a\wedge b)\lrcorner\,\ch(M_{T}).
\]
A sheaf of length $\ell$ supported at a point has Chern character
$\ell\cdot[\mathrm{pt}]$, and Chern characters are additive over the
cohomology of a complex, so $\ch(M_{T})=\chi(M)\,[\mathrm{pt}]$. The class of a
point is a nonzero multiple of $dz_{1}\cdots dz_{c}\,d\bbar z_{1}\cdots
d\bbar z_{c}$, and contracting it with $a\wedge b$ removes two holomorphic
factors and leaves a nonzero class of type $(c-2,c)$, because $a$ and $b$ are
independent. So the right side is nonzero when $\chi(M)\ne0$, hence so is
$A_{a}A_{b}$, and its restriction $J(a)J(b)$ is nonzero.
\end{proof}

The hypothesis on $\chi(M)$ cannot be dropped, and here is why.%
\footnote{The Euler characteristic is the only additive invariant here: the
Grothendieck group of complexes with finite length cohomology over a regular
local ring is $\ZZ$, generated by the residue field, since every module of
finite length has a composition series with that one factor. The lemma says
that the product of two independent jet classes detects the generator. What
it cannot detect is a class of Euler characteristic zero, and the example in
the text shows that it does not.}
Item (XL) of \Cref{sec:verification} takes the Koszul resolution $K$ of the
residue field $k$ of $\QQ[x_{1},x_{2}]$ and forms the cone $C$ of the class
$J(\partial_{2})J(\partial_{1})$, viewed as a chain map $K\to K[2]$. The
complex $C$ has two cohomology modules, each isomorphic to $k$, so
$\chi(C)=0$, and on $C$ the product $J(\partial_{2})J(\partial_{1})$ is null
homotopic: the gluing has killed it. On the split complex $k\oplus k[1]$,
whose Euler characteristic is also zero, the same product is not null
homotopic. The same item tests the lemma on eighteen finite length modules in
two and three variables, among them powers of the maximal ideal, monomial and
non-monomial complete intersections, a Gorenstein algebra of socle degree two
and three modules that are not cyclic, and on three complexes with several
cohomology modules, for every pair of coordinate directions and for a random
pair; in every case with $\chi\ne0$ the product is nonzero.

\begin{theorem}[Where an object meeting the criterion can live]
\label{thm:p2support}
Let $A$ be an abelian $2n$-fold of $(K,-1,n)$-Weil type with $n\ge2$, let
$\omega=a\alpha_{+}+b\alpha_{-}$ with $ab\ne0$, and let $E$ be a perfect
complex on $A$ with $\ch(E)=N\omega$, $N\ne0$, no other component, and
$\sigma_{E}$ injective. For an irreducible component $Z$ of the support of $E$
write
\[
  m_{Z}(E)=\sum_{i}(-1)^{i}\,\ell_{\cO_{A,Z}}\bigl(\cH^{i}(E)_{Z}\bigr)
\]
for the multiplicity of $Z$ in the cycle of $E$.
\begin{enumerate}[label=\textup{(\roman*)},leftmargin=2.6em,itemsep=2pt]
\item If $Z$ has codimension at least two and $m_{Z}(E)\ne0$, then $Z$ is a
  union of cosets of an abelian subvariety $B\subsetneq A$ with
  $T_{0}B\supseteq T_{+}$ or $T_{0}B\supseteq T_{-}$.
\item If $A$ has no proper abelian subvariety $B$ with $T_{0}B\supseteq T_{+}$
  or $T_{0}B\supseteq T_{-}$, which by \Cref{cor:p2shape} holds when $A$ is
  simple, at every split member $X\times\widehat X$ with $X$ simple and
  $\End^{0}(X)=\QQ$, and at a general point of the quaternionic locus, then
  $m_{Z}(E)=0$ for every component $Z$ of codimension at least two, and the
  cycle of $E$ is a divisor.
\item At every member of the family, whatever its endomorphisms, the support
  of $E$ has an irreducible component of codimension less than $n$.
\end{enumerate}
In particular no perfect complex whose support has codimension at least $n$,
at any member of any Weil family, satisfies the criterion of
\Cref{thm:semiregclosure}. This includes every complex assembled from the
structure sheaves of the components of a cycle representing $\omega$, the base
cycle of \Cref{thm:everyfamily} among them.
\end{theorem}

\begin{proof}
(i) Let $c\ge2$ be the codimension of $Z$ and let $U\subseteq Z$ be the dense
open set of smooth points of $Z_{\mathrm{red}}$ that lie on no other component
of the support. Fix $\theta\in T_{+}$ and $\theta'\in T_{-}$; by
\Cref{thm:p2forces}(ii), $A_{\theta}A_{\theta'}=0$ in $\Ext^{2}(E,E)$.
Suppose that at some $p\in U$ the images $\bar\theta,\bar\theta'$ of
$\theta,\theta'$ in $N_{p}=T_{p}A/T_{p}Z$ are linearly independent, and
complete them by images of constant fields $\nu_{3},\dots,\nu_{c}$ to a basis
of $N_{p}$. In linear coordinates on the universal cover the constant fields
are coordinate derivations, and the germ $S$ at $p$ of the translate of the
span of $\theta,\theta',\nu_{3},\dots,\nu_{c}$ is a smooth germ of dimension
$c$, transverse to $Z$ at $p$, to which $\theta$ and $\theta'$ are tangent.
Let $j\colon S\to A$ be the inclusion and $M=Lj^{*}E$ on
$\widehat{\cO}_{S,p}\cong\CC[[x_{1},\dots,x_{c}]]$. Derived restriction is a
ring map $\Ext^{*}(E,E)\to\Ext^{*}(M,M)$, and it carries the germ of
$A_{\theta}$ to the jet class of $M$ along $\theta|_{S}$, because the Atiyah
class is functorial for pullback and $\theta$ is tangent to $S$
\cite[Section 3]{BF03}. So $J(\theta)J(\theta')=0$ on $M$. Near $p$ the support
of $E$ is $Z$, and $Z\cap S=\{p\}$, so $M$ has finite length cohomology; its
Euler characteristic is the intersection multiplicity of the cycle of $E$
with $S$ at $p$, which, $S$ being transverse to $Z$ at a smooth point lying on
no other component, is $m_{Z}(E)$, by Serre's formula for intersection
multiplicities applied to each cohomology sheaf and summed with signs
\cite[Chapter~V]{Ser00}.%
\footnote{Serre's formula expresses $\sum_{i}(-1)^{i}\ell(\mathrm{Tor}_{i}(F,
\cO_{S}))$, for a module $F$ and a smooth germ $S$ meeting its support
properly at $p$, as a sum over the components $Z'$ of the support of
$\ell(F_{Z'})$ times the intersection multiplicity of $Z'$ and $S$ at $p$. For
$S$ transverse to a component smooth at $p$ the intersection multiplicity is
one. For a complex the Tor spectral sequence gives the same formula with the
lengths replaced by their alternating sum over the cohomology, which is the
definition of $m_{Z}(E)$.}
By \Cref{lem:finiteproduct}, $m_{Z}(E)=0$.

So if $m_{Z}(E)\ne0$, then at every $p\in U$ and for all $\theta\in T_{+}$,
$\theta'\in T_{-}$ the images satisfy $\bar\theta\wedge\bar\theta'=0$ in
$\bigwedge^{2}N_{p}$. As in the proof of \Cref{thm:p2forces}(iv), the images
$I_{\pm}$ of $T_{\pm}$ in $N_{p}$ span $N_{p}$, which has dimension $c\ge2$,
so one of them is zero: $T_{+}\subseteq T_{p}Z$ or $T_{-}\subseteq T_{p}Z$.
The two conditions cut out closed subsets of the irreducible set $U$ whose
union is $U$, so one of them, say the first, holds on all of $U$, and the
argument of \Cref{cor:p2shape} shows that $Z$ is stable under translation by
the identity component $B$ of its stabiliser, an abelian subvariety with
$T_{0}B\supseteq T_{+}$. Since $Z\ne A$, $B\ne A$.

(ii) is (i) when no such $B$ exists.

(iii) Suppose every component of the support has codimension at least $n$.
Then $\ch_{k}(E)=0$ for $k<n$ and $\ch_{n}(E)=\sum_{Z}m_{Z}(E)[Z]$, the sum
over the components of codimension exactly $n$; this is Riemann--Roch without
denominators in its simplest case \cite[Chapter~15]{Ful98}. By (i), each $Z$
with $m_{Z}(E)\ne0$ is $\pi^{-1}(\pi(Z))$ for the quotient $\pi\colon A\to A/B$
by an abelian subvariety $B$ with $T_{0}B\supseteq T_{+}$ or
$T_{0}B\supseteq T_{-}$, so $[Z]=\pi^{*}[\pi(Z)]$ lies in the subring of
$H^{*}(A,\CC)$ generated by $H^{1}(A/B,\CC)$, which consists of the classes
vanishing on $T_{0}B$ and its conjugate. When $T_{0}B\supseteq T_{+}$, a
$(1,0)$-form vanishing on $T_{+}$ lies in $V_{-}^{1,0}$, and a $(0,1)$-form
vanishing on the conjugate of $T_{+}$ lies in $V_{+}^{0,1}$, because complex
conjugation exchanges the two eigenvalues of $K$;%
\footnote{The field $K$ acts on $H_{1}(A,\QQ)$ by rational endomorphisms, so
on the complexification its action commutes with complex conjugation, which
therefore carries the eigenspace of $\sqrt{-d}$ for one embedding of $K$ to
the eigenspace for the other. On the tangent space the conjugate of $T_{+}$ is
the eigenspace in $T^{0,1}$ for the second embedding, and it pairs with
$V_{-}^{0,1}$ and not with $V_{+}^{0,1}$. Item (XXXIX)(D) checks this
bookkeeping in an explicit rational model.}
so the subring lies in $\bigwedge^{*}(V_{-}^{1,0}\oplus V_{+}^{0,1})$. When
$T_{0}B\supseteq T_{-}$ it lies in $\bigwedge^{*}(V_{+}^{1,0}\oplus
V_{-}^{0,1})$. Choose bases of the four pieces. Both exterior algebras are
spanned by monomials in these bases, so their sum is too, and $\alpha_{+}$,
a single monomial in $n$ generators of $V_{+}^{1,0}$ and $n$ of $V_{+}^{0,1}$,
lies in neither; nor does $\alpha_{-}$. As $ab\ne0$, the class $N\omega$ has
nonzero coordinates on both and is not in the sum, which contradicts
$\ch_{n}(E)=N\omega$. So some component has codimension less than $n$. The
last assertion follows, since a complex assembled from structure sheaves of
the components of a cycle of codimension $n$ is supported on that cycle.
\end{proof}

\begin{figure}[tbp]
\centering
  \includegraphics[width=\textwidth]{fig_support}
\caption{The proof of \Cref{thm:p2support}. The grey sheet stands for a
component of the support of codimension less than $n$, which the theorem
forces to exist; the clay curve $Z$ for a component of codimension at least
two. At a smooth point $p$ of $Z$ the directions $\theta\in T_{+}$ and
$\theta'\in T_{-}$ span the slice $S$, transverse to $Z$, on which $E$ restricts
to a complex of finite length with Euler characteristic $m_{Z}(E)$; the amber
parallelogram is the bivector $\theta\wedge\theta'$, independent modulo
$T_{p}Z$. Carried to a torus of the same dimension, the product of the two
jet classes is seen by the semiregularity map as $m_{Z}(E)$ times the class of
a point contracted with the bivector, which is not zero, so $m_{Z}(E)=0$ or
the bivector degenerates. The teal lines are cosets of an abelian subvariety
$B$ tangent to $T_{+}$, the only components of codimension at least two on
which a nonzero multiplicity survives. The scene is schematic in the
dimensions and exact in the incidences.}
\label{fig:support}
\end{figure}

\begin{remark}[What is left of the search]\label{rem:p2search}
Put together, \Cref{thm:p2forces}, \Cref{cor:hhfactor},
\Cref{thm:p2numerical} and \Cref{thm:p2support} describe an object meeting
the criterion at a general member as follows. It has rank zero and first Chern
class zero. Its support has a component of codimension less than $n$,
although its Chern character begins in codimension $n$; at $n=2$ that
component is a divisor or the whole of $A$. Along every component of its
support of codimension at least two its cohomology modules have alternating
length zero, so it has at least two of them there, glued as in the cone of
item (XL). It extends along all $2n^{2}$ $K$-linear deformations of $A$, its
endomorphisms have trace zero, and its Hochschild action factors through a
ring of dimension $2^{2n+1}-1$. The central part of $\Ext^{2}(E,E)$, the image
of $HH^{2}(A)$, has dimension exactly $2n(2n-1)$, and if $\Ext^{2}(E,E)$ is no
larger than its central part, the criterion holds.

The Weil class, which lives in codimension $n$, therefore has to come from the
part of the Chern character that the leading multiplicities do not see: from
twists that cancel in rank but not in higher degree. At $n=2$ the model is a
pair $\cO_{D}$ and $\cO_{D}(W)$ on a divisor $D$, with $W$ a divisor on $D$,
whose difference has Chern character $-i_{*}[W]$ plus terms of higher degree
and has multiplicity zero along $D$. That is a very specific shape, and none
of the objects in the literature or in this paper has it. We record it as the
form the search now takes, and not as a claim that such an object exists.
\end{remark}

\subsubsection*{What the dimension forces on the object}

The numerical criterion of \Cref{thm:p2numerical} fixes one number,
$\dim\Ext^{2}(E,E)$. Together with Serre duality, the lower bounds of
\Cref{cor:hhfactor} in the other degrees and the Hodge--Riemann relations for
the Weil class, it fixes a good deal more: the object that carries the class
can be taken indecomposable, and in the first open case it has more than one
endomorphism. This is where a search has to look.

\begin{theorem}[The shape of an object meeting the criterion]
\label{thm:p2endo}
Let $A$, $\omega$ and $E$ be as in \Cref{thm:p2numerical}, so that
$\ch(E)=N\omega$ with $N\ne0$ and no other component, and suppose that
$\dim\Ext^{2}(E,E)=2n(2n-1)$. Write $e_{k}=\dim\Ext^{k}(E,E)$.
\begin{enumerate}[label=\textup{(\roman*)},leftmargin=2.6em,itemsep=2pt]
\item $\chi(E,E)=\sum_{k}(-1)^{k}e_{k}=(-1)^{n}N^{2}\int_{A}\omega^{2}$, and
  this number is positive.
\item $E\cong E_{0}\oplus E'$ with $E_{0}$ indecomposable,
  $\ch(E_{0})=\ch(E)$, $\dim\Ext^{2}(E_{0},E_{0})=2n(2n-1)$, $\ch(E')=0$ and
  $\Ext^{2}(E_{0},E')=\Ext^{2}(E',E_{0})=\Ext^{2}(E',E')=0$. So $E_{0}$ meets
  the criterion as well.
\item If $E$ is indecomposable, then $\End(E)=\CC\cdot\id\oplus J$ with $J$
  the nilpotent radical, every $f\in J$ annihilates $\Ext^{2}(E,E)$ and
  $\Ext^{2n-2}(E,E)$ under the Yoneda product, and
  $\operatorname{ev}_{E}(\bigwedge^{2n}P)=\operatorname{ev}_{E}(\bigwedge^{2n}Q)$
  is a single line in $\Ext^{2n}(E,E)$. If moreover $e_{1}=4n$, then $J$
  annihilates $\Ext^{1}(E,E)$ as well.
\item If $n=2$ or $n=3$, then
  \[
    \dim\End(E)\;\ge\;2+\tfrac12\,\chi(E,E)\;>\;2,
  \]
  so $E$ has at least three linearly independent endomorphisms, and at
  $n=3$ also $e_{1}\le e_{0}+9$. In particular no simple object meets the
  criterion: no stable sheaf, no complex with $\End=\CC$, and no image of one
  under an autoequivalence of $D^{b}(A)$; and an indecomposable object that
  meets it has a nonzero nilpotent endomorphism. The smallest profiles
  $(e_{0},\dots,e_{2n})$ that the constraints allow are $(3,8,12,8,3)$ at
  $n=2$, with $\chi=2$, and $(3,12,30,40,30,12,3)$ and $(3,12,30,41,30,12,3)$
  at $n=3$, with $\chi=2$ and $\chi=1$.
\end{enumerate}
Parts (i) and (ii), and part (iv) at $n=2$, use the compatibility
$\sigma_{E}\circ\operatorname{ev}_{E}=(\,\cdot\,)\lrcorner\ch(E)$ of
\Cref{setup:hh} in degree two only, as \Cref{thm:p2numerical} does. Part
(iii), and part (iv) at $n=3$, use it in the further degrees in which
\Cref{cor:hhfactor}(iii) grants it: $2n-2$ and $2n$ for (iii), and three for
(iv).
\end{theorem}

\begin{proof}
(i) The Todd class of $A$ is trivial, so the Riemann--Roch theorem gives
$\chi(E,E)=\int_{A}\ch(E)^{\vee}\ch(E)$, and $\ch(E)^{\vee}=(-1)^{n}N\omega$
since $\omega$ has degree $2n$. The polarisation pairs $V_{+}$ only with
$V_{-}$ (\Cref{fig:tangent}), so its class $\theta$ lies in
$V_{+}\otimes V_{-}$ and $\theta\wedge\alpha_{\pm}=0$, because $\alpha_{\pm}$
spans the top exterior power of $V_{\pm}$. So $\omega$ is a nonzero rational
primitive class of type $(n,n)$ in the middle degree, and the Hodge--Riemann
relations give $(-1)^{n(2n-1)}\int_{A}\omega^{2}>0$, where
$(-1)^{n(2n-1)}=(-1)^{n}$. Item (L) of \Cref{sec:verification} checks the
primitivity and the sign of the intersection form on the whole rational Weil
plane in the coordinate model, for $n=2,3$ and several discriminants.

(ii) The category $D^{b}(A)$ has finite dimensional morphism spaces and split
idempotents, so $E=\bigoplus_{i}E_{i}$ with each $E_{i}$ indecomposable. The
evaluation map and the semiregularity map are compatible with direct sums,
and by \Cref{thm:p2numerical}(i) the kernel of $\operatorname{ev}_{E}$ in
degree two is $P\wedge Q$; so $\operatorname{ev}_{E_{i}}(P\wedge Q)=0$ and,
composing with $\sigma_{E_{i}}$, $(P\wedge Q)\lrcorner\ch(E_{i})=0$ for every
$i$. We claim that the only classes of Hodge type killed by $P\wedge Q$ are
those of $\CC\alpha_{+}\oplus\CC\alpha_{-}$. Write
$H^{*}(A,\CC)=\bigwedge^{*}(A_{+}\oplus B_{+}\oplus A_{-}\oplus B_{-})$ with
$A_{\pm}=V_{\pm}^{1,0}$ and $B_{\pm}=V_{\pm}^{0,1}$, and grade it by the four
degrees $(i,j,k,l)$. As in \Cref{setup:hh}, $P$ acts by wedging with $B_{+}$
and contracting $A_{-}$, and $Q$ by wedging with $B_{-}$ and contracting
$A_{+}$, so $P\wedge Q$ is spanned by four kinds of operators, each of which
shifts the four degrees by a fixed amount: wedging with $b_{+}\wedge b_{-}$
kills a component of degree $(i,j,k,l)$ for all choices exactly when $j=n$ or
$l=n$; wedging with $b_{+}$ after contracting $a_{+}$, exactly when $i=0$ or
$j=n$; contracting $a_{-}$ and wedging with $b_{-}$, exactly when $k=0$ or
$l=n$; and contracting $a_{-}$ and $a_{+}$, exactly when $i=0$ or $k=0$. Each
of these is a statement about a tensor product of maps that are injective
unless the degree in question is extreme. A class of type $(p,p)$ has
$i+k=j+l=p$ in every component. If $j=n$, then $l=n$ would force $i=k=n$,
against the fourth condition, so $k=0$, hence $i=n+l$, so $i=n$, $l=0$: the
component lies in $\bigwedge^{n}A_{+}\otimes\bigwedge^{n}B_{+}=\CC\alpha_{+}$.
If $j<n$, then $l=n$ and $i=0$, hence $k=n$ and $j=0$: the component lies in
$\CC\alpha_{-}$. This proves the claim, which item (L) also checks in every
degree for $n=2,3$. Since $\ch(E_{i})$ is rational, it lies on the rational
Weil plane, and when it is nonzero \Cref{cor:hhfactor}(i) gives
$\dim\Ext^{2}(E_{i},E_{i})\ge2n(2n-1)$, a nonzero rational Weil class having
both coefficients $a$ and $b$ nonzero. As
$\Ext^{2}(E,E)=\bigoplus_{i,j}\Ext^{2}(E_{i},E_{j})$ has dimension
$2n(2n-1)$ and $\sum_{i}\ch(E_{i})=N\omega\ne0$, exactly one summand $E_{0}$
has nonzero Chern character and every other block of $\Ext^{2}$ vanishes.

(iii) An indecomposable object of a category with finite dimensional
morphism spaces and split idempotents has a local endomorphism ring, so
$\End(E)=\CC\cdot\id\oplus J$ with $J$ nilpotent. By
\Cref{thm:p2numerical}(i), $\operatorname{ev}_{E}$ maps
$\bigwedge^{2}P\oplus\bigwedge^{2}Q$ isomorphically onto $\Ext^{2}(E,E)$.
By Serre duality $e_{2n-2}=e_{2}$, and contraction with $\omega$ is injective
on $\bigwedge^{2n-2}P\oplus\bigwedge^{2n-2}Q$ (item (L)), so, with the
compatibility in degree $2n-2$, $\operatorname{ev}_{E}$ maps that space
isomorphically onto $\Ext^{2n-2}(E,E)$. In degree $2n$ the semiregularity map is the trace, since
$H^{q+2n}(\Omega^{q}_{A})=0$ for $q>0$, so
$\operatorname{tr}\operatorname{ev}_{E}(\xi)$ is the $H^{0,2n}$ component of
$\xi\lrcorner\ch(E)$; for the generators $\pi_{P}$, $\pi_{Q}$ of
$\bigwedge^{2n}P$ and $\bigwedge^{2n}Q$ it is $Nc_{P}$ and $Nc_{Q}$ with
$c_{P}c_{Q}\ne0$ (item (L)). For $f\in\End(E)$ and
$x\in\bigwedge^{2}P\oplus\bigwedge^{2}Q$ write
$f\cdot\operatorname{ev}_{E}(x)=\operatorname{ev}_{E}(\varphi_{f}(x))$; then
$f\mapsto\varphi_{f}$ is a homomorphism of algebras. Let $x\in\bigwedge^{2}P$
and $y\in\bigwedge^{2n-2}P$, and write $x\wedge y=\beta(x,y)\pi_{P}$. Since
$\operatorname{ev}_{E}$ is a ring map that kills the ideal generated by
$P\wedge Q$,
\[
  \beta(x,y)\operatorname{tr}\bigl(f\operatorname{ev}_{E}(\pi_{P})\bigr)
  =\operatorname{tr}\bigl(f\operatorname{ev}_{E}(x)\operatorname{ev}_{E}(y)\bigr)
  =\operatorname{tr}\operatorname{ev}_{E}\bigl(\varphi_{f}(x)\wedge y\bigr)
  =Nc_{P}\,\beta\bigl(\varphi_{f}(x)_{P},y\bigr),
\]
and the pairing $\beta$ is perfect, so $\varphi_{f}(x)_{P}=\mu_{f}x$ with
$\mu_{f}=\operatorname{tr}(f\operatorname{ev}_{E}(\pi_{P}))/Nc_{P}$. The same
computation with $y\in\bigwedge^{2n-2}Q$, where $x\wedge y=0$ in the quotient,
gives $\varphi_{f}(x)_{Q}=0$. So $f$ acts on
$\operatorname{ev}_{E}(\bigwedge^{2}P)$ by a scalar $\mu_{f}$, and in the same
way on $\operatorname{ev}_{E}(\bigwedge^{2}Q)$ by a scalar $\nu_{f}$; $\mu$ and
$\nu$ are characters of $\End(E)$ and vanish on $J$. So $J$ kills
$\Ext^{2}(E,E)$, on both sides because $\operatorname{ev}_{E}$ lands in the
graded centre, and the same argument with the roles of the two degrees
exchanged shows that it kills $\Ext^{2n-2}(E,E)$. In particular
$\operatorname{tr}(f\operatorname{ev}_{E}(\pi_{P}))=\mu_{f}Nc_{P}=0$ for
$f\in J$, so $\operatorname{ev}_{E}(\pi_{P})$ and, likewise,
$\operatorname{ev}_{E}(\pi_{Q})$ lie in the annihilator of $J$ under the
perfect pairing of $\End(E)$ with $\Ext^{2n}(E,E)$. That annihilator is a
line, $J$ having codimension one, and both elements are nonzero, their traces
being $Nc_{P}$ and $Nc_{Q}$; so they span the same line. Finally, if
$e_{1}=4n$ then $\Ext^{1}(E,E)=\operatorname{ev}_{E}(HH^{1})$, because
$\operatorname{ev}_{E}$ is injective on $HH^{1}$, as shown in (iv). For
$f\in J$ and $h\in HH^{1}$ write $f\operatorname{ev}_{E}(h)=
\operatorname{ev}_{E}(h')$; then
$\operatorname{ev}_{E}(h'\wedge h'')=f\operatorname{ev}_{E}(h\wedge h'')=0$ for
every $h''\in HH^{1}$, so $h'\wedge HH^{1}\subseteq P\wedge Q$, which forces
$h'=0$ since $\dim P=\dim Q\ge2$.

(iv) By Serre duality $e_{k}=e_{2n-k}$. The map $\operatorname{ev}_{E}$ is
injective on $HH^{1}$ using degree two alone: if $\operatorname{ev}_{E}(h)=0$
then $\operatorname{ev}_{E}(h\wedge h')=0$ for every $h'$, so
$h\wedge HH^{1}\subseteq\Ann_{HH^{2}}(\omega)=P\wedge Q$, and $h=0$ as in
(iii). So $e_{1}\ge4n$, and \Cref{cor:hhfactor}(iii) gives
$e_{3}\ge2\binom{6}{3}=40$ at $n=3$. At $n=2$ the Euler
characteristic is $2e_{0}-2e_{1}+12$, and at $n=3$ it is
$2e_{0}-2e_{1}+60-e_{3}$. So $2e_{0}=\chi+2e_{1}-12\ge\chi+4$ at $n=2$, and
$2e_{0}=\chi+2e_{1}+e_{3}-60\ge\chi+4$ at $n=3$, where also
$e_{1}=e_{0}+30-\frac12(\chi+e_{3})\le e_{0}+9$. With (i) this gives the
bound. An autoequivalence preserves $\End$, and a stable sheaf is simple.
An indecomposable object with $\dim\End\ge3$ has $J\ne0$. The smallest
profiles are those with $e_{0}=3$, which item (L) lists by exhaustive search
over the admissible values, and Section 31 of the Lean certificate checks the
inequality $e_{0}\ge3$ for all natural numbers.
\end{proof}

\begin{remark}[What the shape does to the search]\label{rem:p2endo}
\Cref{thm:p2endo} is a necessary condition and nothing more; it does not say
that an object meeting the criterion exists. What it does is remove from the
search every object that anyone would construct first. Line bundles,
semi-homogeneous bundles, skyscraper sheaves, stable sheaves and their
Fourier--Mukai transforms are simple, and direct sums of them are excluded by
(ii) together with \Cref{thm:nosum}. At $n=3$, the first open case, what
remains is an indecomposable object with a local endomorphism algebra of
dimension at least three whose radical kills $\Ext^{2}$ and $\Ext^{4}$ and
pairs to zero with the line $\operatorname{ev}_{E}(\bigwedge^{6}P)$,
granting the compatibility of \Cref{cor:hhfactor}(iii); the smallest
conceivable one has $\Hom(E,E)$ of dimension three, $\Ext^{1}(E,E)$
equal to the image of $HH^{1}(A)$, of dimension twelve and killed by the
radical, and $\Ext^{3}(E,E)$ of dimension forty or forty-one. Iterated
self-extensions and thickenings are the natural objects with nilpotent
endomorphisms; an extension of a simple object $F$ by itself has
$\ch=2\ch(F)$, so $F$ would have to carry the class already, and the question
returns one level down. Nothing in this paper produces an object of the
required shape, and the remark is recorded as the form the search now takes.
For $n\ge4$ the same count only asks for an excess in some even degree other
than $2$ and $2n-2$ (item (L)), and the bound on $\End$ disappears.
\end{remark}

\subsubsection*{The pure form of the criterion is never met}

So far $\ch(E)=N\omega$ exactly. If $\sigma_{E}$ were injective, $E$ would
extend along every $K$-linear deformation of $A$ (\Cref{thm:p2forces}(ii)).
These reach non-algebraic tori, on which the Weil class is Hodge and the Chern
class of nothing. That settles the pure form of the criterion, negatively, in
every dimension.

\begin{setupenv}[Weil tori]\label{setup:weiltori}
Let $A$ be an abelian $2n$-fold of $(K,-1,n)$-Weil type, $V=H^{1}(A,\QQ)$ with
its action of $K$, and $V\otimes\CC=V_{+}\oplus V_{-}$ the two eigenspaces,
each of dimension $2n$. For complementary $n$-dimensional subspaces
$X,Y\subset V_{+}$ the subspace $X\oplus\bbar{Y}$ of $V\otimes\CC$ is the
space of holomorphic $1$-forms of a complex structure on the real torus
underlying $A$; call the resulting complex torus $T_{X,Y}$. On $T_{X,Y}$ the
field $K$ acts with signature $(n,n)$, so $T_{X,Y}$ is a Weil torus in the
sense of \cite{Voi02b}, and every complex structure on that real torus
commuting with $K$ and of signature $(n,n)$ arises in this way. The pairs
$(X,Y)$ form a connected complex manifold $S$ of dimension $2n^{2}$,
containing the point of $A$; the polarised Weil family is the submanifold of
dimension $n^{2}$ on which the polarisation of $A$ is of type $(1,1)$, and the
tangent space of $S$ is the space $\Hom_{K}(H^{1,0},H^{0,1})$ of $K$-linear
first order deformations in \Cref{thm:p2forces}(i). Since
$\alpha_{+}$ spans $\bigwedge^{2n}V_{+}=\bigwedge^{n}X\otimes\bigwedge^{n}Y$
and $X$ is holomorphic while $Y$ is antiholomorphic, the Weil classes are of
type $(n,n)$ on every $T_{X,Y}$. We say that a property holds for a very
general member of $S$ if it holds outside a countable union of proper
analytic subsets.
\end{setupenv}

\begin{lemma}[The Hodge classes of a very general Weil torus]
\label{lem:weiltorushodge}
For a very general member $T$ of $S$, $\Hdg^{k}(T)=0$ for
$1\le k\le2n-1$, $k\ne n$, and $\Hdg^{n}(T)$ is the rational Weil plane.
\end{lemma}

\begin{proof}
A rational class that is of Hodge type at a very general member is of Hodge
type on all of $S$, because its Hodge locus is analytic and $S$ is connected.
Let $G=\Res_{K/\QQ}\mathrm{GL}_{2n}$ be the group of $K$-linear automorphisms
of $V$. Its real points act transitively on $S$, since
$\mathrm{GL}(V_{+})$ acts transitively on pairs of complementary subspaces.
The Mumford--Tate group $M$ of a very general member lies in $G$, since $K$
acts by endomorphisms of every member, and it contains the Mumford--Tate group
of every member, whose Hodge classes include those of the very general one.
So the real Lie algebra of $M$ contains the complex structure $J_{s}$ of every
member, the derivative of its Hodge cocharacter, and therefore contains the $\mathrm{Ad}\,G(\RR)$-span of one $J_{s}$, which is an ideal of
$\mathfrak{g}_{\RR}=\mathfrak{gl}_{2n}(\CC)$ containing a non-central
element, $J_{s}$ acting on $V_{+}$ by $i$ on $X$ and $-i$ on $Y$; hence it
contains $\mathfrak{sl}_{2n}(\CC)$, and $M$ contains
$G^{\mathrm{der}}=\Res_{K/\QQ}\mathrm{SL}_{2n}$. Over $\CC$ this is
$\mathrm{SL}(V_{+})\times\mathrm{SL}(V_{-})$ acting on
$\bigwedge^{*}(V_{+}\oplus V_{-})=\bigwedge^{*}V_{+}\otimes\bigwedge^{*}V_{-}$,
whose invariants are spanned by $1$, $\alpha_{+}$, $\alpha_{-}$ and
$\alpha_{+}\alpha_{-}$. Hodge classes are invariant under $M$, which gives
the statement. Item (LI) of \Cref{sec:verification} confirms it for $n=2$
and $n=3$ by exhibiting explicit members at which no other rational class of
those degrees is of Hodge type.
\end{proof}

\begin{lemma}[A very general Weil torus has no subvarieties]
\label{lem:weiltorussubsets}
A very general member $T$ of $S$ has no analytic subsets other than finite
sets and $T$ itself.
\end{lemma}

\begin{proof}
Let $Z\subset T$ be an irreducible analytic subset of codimension $c$ with
$0<c<2n$. For every K\"ahler class $\kappa$ the class $[Z]$ satisfies
$\int_{T}[Z]\,\kappa^{2n-c}=\int_{Z}\kappa^{2n-c}>0$, so $[Z]$ is a nonzero
Hodge class of degree $2c$. By \Cref{lem:weiltorushodge} this forces $c=n$ and
$[Z]$ into the Weil plane. Take for $\kappa_{0}$ the class of a positive
definite hermitian form that is block diagonal for the splitting
$H^{1,0}=X\oplus\bbar{Y}$. Its blocks lie in $X\otimes\bbar{X}$ and
$\bbar{Y}\otimes Y$, both inside $V_{+}\otimes V_{-}$, so every power of
$\kappa_{0}$ has a factor from $V_{+}$ and a factor from $V_{-}$ in each term,
and $\alpha_{\pm}\wedge\kappa_{0}^{n}=0$ because $\alpha_{\pm}$ spans the top
exterior power of $V_{\pm}$ (item (LI)). So
$\int_{T}[Z]\,\kappa_{0}^{n}=0$, which is impossible.
\end{proof}

\begin{proposition}[Coherent sheaves on a very general Weil torus]
\label{prop:weiltorussheaves}
Let $n\ge2$ and let $T$ be a very general member of $S$. Every coherent
analytic sheaf $F$ on $T$ has $\ch_{k}(F)=0$ in $H^{2k}(T,\QQ)$ for
$1\le k\le2n-1$. For $n=2$ this is Voisin's theorem that every coherent sheaf
on a general Weil torus of dimension four has $c_{2}(F)=0$
\cite[Theorem~10]{Voi02b}, whose proof we follow.
\end{proposition}

\begin{proof}
Chern classes of coherent sheaves on $T$ exist in rational Deligne
cohomology, extend those of bundles and are additive \cite{Gri10}, so they
are Hodge classes. By \Cref{lem:weiltorussubsets} the torsion of $F$, the locus where the
torsion-free quotient is not locally free, and the cokernel of the map from
that quotient to its reflexive hull are all supported on finite sets, and a
sheaf supported on a finite set has Chern character in the top degree only. So
it suffices to treat a reflexive sheaf $G$. Every coherent sheaf on $T$ has
$c_{1}=0$, since $\Hdg^{1}(T)=0$ (\Cref{lem:weiltorushodge}); so every
torsion-free sheaf has slope zero for every K\"ahler class, and a torsion-free
sheaf is stable for one K\"ahler class exactly when it has no saturated
subsheaf of intermediate rank, which does not depend on the class. A
Jordan--H\"older filtration of $G$ therefore has stable torsion-free factors
$Q_{j}$, stable for every K\"ahler class, and $\ch(G)=\sum_{j}\ch(Q_{j})$.
By the theorem of Bando and Siu \cite{BS94} the stable reflexive sheaf
$Q_{j}^{**}$ carries an admissible Hermite--Einstein metric for every K\"ahler
class $\kappa$ and satisfies
$\int_{T}\bigl(2rc_{2}-(r-1)c_{1}^{2}\bigr)(Q_{j}^{**})\,\kappa^{2n-2}\ge0$,
with equality exactly when $Q_{j}^{**}$ is a projectively flat vector bundle.
Here $c_{1}=0$ and $c_{2}(Q_{j}^{**})\in\Hdg^{2}(T)$. For $n\ge3$ that group is
zero; for $n=2$ it is the Weil plane, and then
$\int_{T}c_{2}(Q_{j}^{**})\,\kappa_{0}^{2}=0$ for the class $\kappa_{0}$ of
the proof of \Cref{lem:weiltorussubsets}. In either case equality holds for
some $\kappa$, so $Q_{j}^{**}$ is a projectively flat bundle with $c_{1}=0$,
hence flat, and its rational Chern classes of positive degree vanish.
$Q_{j}$ and $Q_{j}^{**}$ differ on a finite set, so $\ch_{k}(Q_{j})=0$ for
$1\le k\le2n-1$, and the same holds for $G$ and for $F$.
\end{proof}

\begin{theorem}[The pure form of the criterion is never met]
\label{thm:p2false}
Let $n\ge2$, let $A$ be an abelian $2n$-fold of $(K,-1,n)$-Weil type, $\omega$
a nonzero rational Weil class, and $E$ a perfect complex on $A$ with
$\ch(E)=N\omega$, $N\ne0$, and no other component. Then $\sigma_{E}$ is not
injective. In particular $\dim\Ext^{2}(E,E)>2n(2n-1)$, the hypotheses of
\Cref{thm:semiregclosure} and \Cref{thm:p2numerical}(iii) are satisfied by no
object at any member of any Weil family of an imaginary quadratic field, and
the propagation statement \textup{(P2)} of \Cref{rem:tworemain}, in the form
stated there, is false.
\end{theorem}

\begin{proof}
Suppose $\sigma_{E}$ injective. Since $A$ is projective, $E$ is isomorphic to
a bounded complex $E^{\bullet}$ of holomorphic vector bundles. The class
$N\omega$ is of Hodge type on every member of $S$ (\Cref{setup:weiltori}), so
by the theorem of Buchweitz and Flenner in Pridham's form (\Cref{thm:BF},
\cite{BF03,Pri24}), every obstruction to extending $E$ along an
Artinian deformation of $A$ inside $S$ vanishes, exactly as in the proof of
\Cref{thm:semiregclosure}. The deformations of the complex structure inside
$S$ together with those of the holomorphic structures and differentials of
$E^{\bullet}$ are governed by a differential graded Lie algebra of Dolbeault
type with finite dimensional cohomology, and Kuranishi's method
\cite{Kur65}, in the form of \cite[Sections~2--3]{GM88}, gives an analytic
germ $(M,m_{0})$ with a morphism $\pi\colon(M,m_{0})\to(S,s_{0})$ and a
complex of holomorphic vector bundles on the pulled back family of tori
deforming $E^{\bullet}$, whose completion is a hull of the formal deformation
functor. The vanishing of the obstructions says that $\pi$ is formally
smooth, hence smooth, so its image contains a neighbourhood $U$ of $s_{0}$.
Thus every member $T_{s}$ with $s\in U$ carries a bounded complex of
holomorphic vector bundles $E_{s}^{\bullet}$ whose Chern character is the
parallel transport of $N\omega$. Choose $s\in U$ very general. By
\Cref{prop:weiltorussheaves},
$\ch_{n}(E_{s})=\sum_{i}(-1)^{i}\ch_{n}(E_{s}^{i})=0$, while it is the
transport of $N\omega$, which is not zero. This contradiction proves that
$\sigma_{E}$ is not injective. The argument uses only that $\ch(E)$ is of
Hodge type on every member of $S$, so it applies verbatim when $\ch(E)-N\omega$
has components in degrees $0$ and $4n$ as well. The remaining statements
follow, since $\dim\Ext^{2}(E,E)\ge2n(2n-1)$ by \Cref{cor:hhfactor}(i),
and equality forces $\sigma_{E}$ to be injective by
\Cref{thm:p2numerical}(i).
\end{proof}

\begin{remark}[What survives, and the corrected statement]\label{rem:p2prime}
\Cref{thm:p2false} turns \Cref{thm:p2forces}(ii) from a constraint into a
contradiction: an object that extends along every $K$-linear deformation of
$A$ lands on non-algebraic tori, where the Weil class is Hodge and is the
Chern class of nothing. The results of this part that describe the object,
\Cref{thm:p2forces}, \Cref{cor:hhfactor}(ii)--(iii), \Cref{thm:p2numerical},
\Cref{thm:p2support}, \Cref{thm:nosum} and \Cref{thm:p2endo}, remain true as
implications, and they are vacuous: no object satisfies their hypotheses. They
are kept as the record of what the pure form asked for.

The obstruction does not reach a Chern character with a component that is
Hodge on the polarised family and not beyond it, unless a twist removes that
component. Call
\textup{(P2$'$)} the statement that at one base point of each Weil family
some perfect complex $E$ with
\[
  \ch(E)=N\omega_{1}+\sum_{k}c_{k}\,\theta^{k},\qquad N\ne0,\ c_{k}\in\QQ,
\]
$\theta$ the polarisation, has injective semiregularity map. The proof of
\Cref{thm:semiregclosure} applies to it word for word, with the deformation
taken over the polarised family, along which $\ch(E)$ stays of Hodge type:
at a member $A_{s}$ it produces $\ch_{n}(E_{s})=N\omega_{1}+c_{n}\theta^{n}$
algebraic, and $\theta^{n}$ is algebraic, so $\omega_{1}$ is. The form above
is already the most general one allowed by the subring generated by $\theta$
and the Weil classes: a Weil class is primitive, so $\theta\omega_{1}=0$ and
that subring is spanned by the powers of $\theta$ and the Weil plane.
\Cref{thm:p2false} excludes directly a Chern character that remains of Hodge
type on all of $S$, which by \Cref{lem:weiltorushodge} is a multiple of the
Weil class plus terms in degrees $0$ and $4n$, and after a line bundle twist
the characters $c\,e^{t\theta}+N\omega_{1}+c'\theta^{2n}$ with $t\theta$
integral (\Cref{thm:p2primenumber}(v)). For every other corrected character
with some $c_{k}\ne0$, $1\le k\le2n-1$, the component of the Hodge locus of
$\ch(E)$ through $A$ is the polarised family itself, whose members are
abelian varieties, so nothing forces $E$ onto non-algebraic tori and the
argument cannot be run on $E$ (\Cref{prop:p2primelocus}); the one exception is
a ratio at $n=2$ at which the locus is five-dimensional
(\Cref{rem:p2primeexceptional}). The semiregular objects in Markman's
proofs for abelian fourfolds and for the sixfolds of trivial discriminant
\cite{Mar25a,Mar25b} are not of the pure form either: their Chern characters
are not multiples of a Weil class, and they are deformed only along families
on which those characters stay of Hodge type. We have not checked that they
have exactly the form above, and nothing below depends on it. From here on, in
\Cref{cor:fourremain}, \Cref{thm:closure} and the frontier, \textup{(P2)} is
read in this corrected form. The annihilator computations above depend on
$\ch(E)$ being a pure Weil class; they are redone for a Chern character with
components in every even degree at the end of this subsection
(\Cref{thm:p2primenumber}, \Cref{prop:p2primeprofile},
\Cref{prop:p2primelocus}).

The same obstruction does not reach the Mumford target of
\Cref{ssec:mumfordrigid}. The exceptional classes there are rigid: their Hodge
locus among all complex tori is the Mumford curve (\Cref{thm:mumfordrigid},
\Cref{prop:notwistor}), so a semiregular complex with Chern character
$N\omega$ on the square of a Mumford fourfold would only have to deform along
the curve, whose members are all algebraic.
\end{remark}

\begin{remark}[The two finite statements, and what each of them is]
\label{rem:tworemain}
After \Cref{thm:degreebound} and \Cref{thm:semiregclosure} the Hodge
conjecture for abelian varieties of Weil type follows from either of the
following, each of which is a statement about explicit finite data and neither
of which mentions a general point. The first is settled by
\Cref{thm:p1equivalent}, and we record it because it is the form in which the
question first presents itself. The second is settled by \Cref{thm:p2false}
in the form stated here, and remains open only in the corrected form of
\Cref{rem:p2prime}.

\begin{enumerate}[label=\textup{(P\arabic*)},leftmargin=*]
\item \emph{A bound.} At every point of the orbit $G(\QQ)\cdot s_{0}$ the
  class $\omega$ is represented by a cycle whose Hilbert data comes from one
  finite set. This one is settled, and negatively, in the sense that matters:
  \Cref{thm:p1forces} shows that the transport of the base cycle does not
  satisfy it, for any cycle whose components have finite stabiliser, and
  \Cref{thm:p1equivalent} shows that read on an arbitrary representative it is
  equivalent to $\mathbf{W}(K,n,\delta)$ itself. It is therefore a
  reformulation of the conjecture for the family, and we record it as such
  rather than as a criterion.
\item \emph{A rank.} The semiregularity map of a perfect complex with the
  Chern character of \Cref{prop:ktheory} at one base point is injective.
  In this form the statement is false (\Cref{thm:p2false}), and it survives
  only in the corrected form of \Cref{rem:p2prime}, where the Chern
  character may carry powers of the polarisation; what follows records what
  the first form asked for. Its
  source is $\Ext^{2}(E,E)$ for a single object, and its target has dimension
  $\sum_{q}h^{q,q+2}(A_{0})=\binom{4n}{2n-2}$, computed in
  \Cref{prop:target} and verified in item (V) of
  \Cref{sec:verification}. So the target is large, which is the direction that
  helps; what is not known to us is the rank of the map. What is known is
  what full rank would force (\Cref{thm:p2forces}): the evaluation map of
  Hochschild cohomology must vanish on the $4n^{2}$-dimensional annihilator of
  the Weil class, so $E$ must be killed by the mixed block of $H^{0,2}$, must
  extend along every $K$-linear deformation, must have all products of
  transverse translation classes zero, and must have no endomorphism of
  nonzero trace; and at a base point with no abelian subvariety tangent to an
  eigenspace it can be a locally free sheaf on no smooth germ of its support
  of codimension at least two (\Cref{cor:p2shape}). The representative that
  the construction supplies, a virtual sum of line bundles, fails at the first
  of these (\Cref{prop:concentrated}), and any representative that is a
  sheaf fails at the last. \Cref{thm:hhsplitting} adds the shape of that
  annihilator, the mixed part of a splitting $HH^{1}=P\oplus Q$ into the
  annihilators of the two conjugate pieces of the class, and
  \Cref{cor:hhfactor} the consequence: the whole Hochschild action of $E$
  factors through a ring of dimension $2^{2n+1}-1$ inside one of dimension
  $2^{4n}$, while $\dim\Ext^{2}(E,E)\ge2n(2n-1)$ for every representative,
  semiregular or not. \Cref{thm:p2numerical} turns the rank into that number:
  equality is enough. And \Cref{thm:p2support} places the support. At every
  member it has a component of codimension less than $n$, so no complex
  carried by a cycle representing $\omega$ can serve, and at a general member
  every component of codimension at least two has multiplicity zero.
  \Cref{thm:p2endo} fixes its shape: an indecomposable summand carries the
  class and meets the criterion alone, and at $n=2$, and at $n=3$ granting
  the compatibility of \Cref{cor:hhfactor}(iii), it has at least three
  independent endomorphisms, so no simple object serves.
\end{enumerate}

Two remarks on what these are not. Neither is the variational Hodge
conjecture: (P1) said nothing about deformations at all, and (P2) is the
infinitesimal statement at one point with the global propagation already
supplied by \Cref{thm:closedlocus}, which is what \Cref{rem:variational} said
was missing. And neither is implied by the results of
\Cref{ssec:everyfamily}: the base points are there, the cycles are there, and
both (P1) and (P2) are questions about those same cycles that the construction
does not answer. We have not settled either, and we do not claim the
conjecture.
\end{remark}

\subsubsection*{The corrected criterion, computed}

The annihilator computations of this subsection depend on $\ch(E)$ being a
pure Weil class, as \Cref{rem:p2prime} notes. We now redo them for a Chern character with components in every
even degree. In degree two the answer is complete and holds for every $n$;
in the other degrees it is a closed formula, proved away from the middle
degree and checked in it. The same computation gives the Hodge locus of the
character, and with it the list of obstructions of this paper that survive
the correction (\Cref{prop:p2primelocus}, \Cref{rem:p2primesurvive}). None
of this constructs an object or decides
whether one exists. It gives, for each shape of Chern character, the number
that $\dim\Ext^{2}(E,E)$ has to reach and the exact kernel that
$\operatorname{ev}_{E}$ must then have, which is the form in which a
candidate can be tested.

\begin{setupenv}[Corrected characters]\label{setup:p2prime}
Let $A$ be an abelian $2n$-fold of $(K,-1,n)$-Weil type, $n\ge2$, with
polarisation $\theta$, and let $\omega=u\alpha_{+}+\bbar u\,\alpha_{-}$ be a
nonzero real class on the Weil line, in the notation of \Cref{thm:p2forces}.
For complex numbers $c_{0},\dots,c_{2n}$ put
\[
  p(\theta)=\sum_{k=0}^{2n}c_{k}\theta^{k},\qquad
  \gamma=\omega+p(\theta),\qquad \mu_{m}=m!\,c_{m},
\]
and let $\rho\in\{0,1,2,3\}$ be the rank of the $(2n-1)\times3$ Hankel matrix
$H=[\mu_{j+i}]$, $0\le j\le2n-2$, $0\le i\le2$.\footnote{$H$ is a
catalecticant of the moment sequence $\mu_{m}=m!\,c_{m}$. A nonzero vector
$(x,\lambda,w)$ in its kernel is a linear recurrence
$x\mu_{j}+\lambda\mu_{j+1}+w\mu_{j+2}=0$ of order at most two, so $\rho\le2$
says that the moments satisfy one, and the shapes of
\Cref{thm:p2primenumber}(ii) are its solutions, read off from the roots of
$wT^{2}+\lambda T+x$.} Write $r(\gamma)$ for the
rank of contraction $HT^{2}(A)\to H^{*}(A,\CC)$, $\xi\mapsto\xi\lrcorner\gamma$.
For $a<n\le b$, in the basis of \Cref{thm:hhsplitting} adapted to an
eigenbasis $e_{j}$ of $V_{+}$ and $f_{j}$ of $V_{-}$ with $e_{j},f_{j}$ of
type $(1,0)$ for $j<n$ and $j\ge n$ respectively, write
$z_{ab}=w(e_{b})w(f_{a})\in V_{+}^{0,1}\otimes V_{-}^{0,1}$,
$v^{1}_{ab}=w(e_{b})c(e_{a})$ and $v^{2}_{ab}=w(f_{a})c(f_{b})$ in
$H^{1}(T_{A})$, and $\pi_{ab}=c(e_{a})c(f_{b})\in T_{+}\otimes T_{-}$, where
$w$ is wedge product with a $(0,1)$ class and $c$ contraction of a $(1,0)$
class.
\end{setupenv}

For a Chern character, $\gamma=N\omega_{1}+\sum_{k}c_{k}\theta^{k}$ with
rational $c_{k}$, and $u$ is determined by $N\omega_{1}$; intrinsically,
$|u|^{2}=(-1)^{n}\tfrac{(2n)!}{2}\int_{A}\omega^{2}\big/\!\int_{A}\theta^{2n}$
in the normalisation of item (LII). The rank $\rho$ does not change when
$\gamma$ is multiplied by $e^{t\theta}$, which acts on the moments by a
binomial transform, and $\omega e^{t\theta}=\omega$ because a Weil class is
primitive; so $\rho$ is an invariant of $\gamma$ up to twist.

\begin{theorem}[The corrected criterion is a number]\label{thm:p2primenumber}
In \Cref{setup:p2prime}:
\begin{enumerate}[label=\textup{(\roman*)},leftmargin=2.6em,itemsep=2pt]
\item For $n\ge3$,
  \[
    \Ann_{HT^{2}}(\gamma)=\bigoplus_{a<n\le b}
    \bigl\{x\,z_{ab}+y_{1}v^{1}_{ab}+y_{2}v^{2}_{ab}+w\,\pi_{ab}\;:\;
      (x,\;i(y_{1}-y_{2}),\;w)\in\ker H\bigr\},
  \]
  so that $\dim\Ann_{HT^{2}}(\gamma)=n^{2}(4-\rho)$ and
  $r(\gamma)=(4+\rho)n^{2}-2n$, whatever $u\ne0$ is.
\item $\rho=3$ for a general $p$. If $\rho=3$, and either $n\ge3$ or $n=2$ and
  $e=0$ in the notation of (iii), then $\Ann_{HT^{2}}(\gamma)$ is exactly
  the tangent space $T_{[A]}\cD_{n,n}$ of the polarised Weil family, spanned
  by the $n^{2}$ vectors $v^{1}_{ab}+v^{2}_{ab}$, inside $H^{1}(T_{A})$, with
  no component in $H^{0,2}$ or in $H^{0}(\bigwedge^{2}T_{A})$; and
  $r(\gamma)=\binom{4n}{2}-n^{2}=7n^{2}-2n$. Moreover $\rho\le2$ exactly when
  $p$ is one of $Ae^{s\theta}+Be^{t\theta}$, $(A+B\theta)e^{t\theta}$,
  $Ae^{t\theta}+B\theta^{2n}$, $A\theta^{2n-1}+B\theta^{2n}$, and $\rho\le1$
  exactly when $p\in\CC e^{t\theta}\cup\CC\theta^{2n}$; the pure case
  $\rho=0$ gives back $r=2n(2n-1)$ of \Cref{cor:hhfactor}(i).
\item For $n=2$, $\dim\Ann_{HT^{2}}(\gamma)=4(4-\rho)+e$ with
  $e=\dim\ker\bigl((HK)^{2}-|u|^{2}\bigr)\le2$ and $K$ the antidiagonal
  matrix with entries $1,-2,1$; so $r(\gamma)\in\{12,16,18,20,22,23,24\}$,
  and $e>0$ exactly when $X(X-I)^{2}=4J^{2}$, where $X=|u|^{2}$ and $I$, $J$
  are the two invariants of the binary quartic
  $\sum_{m}\binom{4}{m}\mu_{m}x^{4-m}y^{m}$.
\item Let $E$ be a perfect complex on $A$ with $\ch(E)=\gamma$ and $N\ne0$.
  Then $\dim\Ext^{2}(E,E)\ge r(\gamma)$ with no hypothesis on $\sigma_{E}$.
  If equality holds, then $\operatorname{ev}_{E}$ is onto with kernel
  $\Ann_{HT^{2}}(\gamma)$ and $\sigma_{E}$ is injective, so that
  \textup{(P2$'$)} holds at that point and, by \Cref{rem:p2prime} and
  \Cref{thm:semiregclosure}, the Hodge conjecture holds for every member of
  the family with full Hodge group. Conversely, if $\sigma_{E}$ is injective
  then $\ker\operatorname{ev}_{E}=\Ann_{HT^{2}}(\gamma)$.
\item If $p\in\CC e^{t\theta}+\CC\theta^{2n}$ with $t\theta$ the class of a line
  bundle $L$, a set of shapes with $\rho\le2$ that contains
  $\CC\theta^{2n}$, then no perfect complex $E$ with $\ch(E)=\gamma$ has
  injective $\sigma_{E}$.
\end{enumerate}
\end{theorem}

\begin{proof}
(i) Let the torus $(\CC^{\times})^{2n}$ act by $e_{j}\mapsto\tau_{j}e_{j}$,
$f_{j}\mapsto\tau_{j}^{-1}f_{j}$. It preserves the Hodge decomposition and
$\theta$, and scales $\alpha_{\pm}$ by $(\prod_{j}\tau_{j})^{\pm1}$;
contraction is equivariant. A weight of $HT^{2}$ has at most two nonzero
coordinates, so for $n\ge3$ two such weights never differ by a nonzero
multiple of $(1,\dots,1)$, and the images of one weight component in the
$\theta$-part and in the two $\alpha$-parts land in distinct weight spaces.
Hence $\Ann(\gamma)=\Ann(p(\theta))\cap\Ann(\alpha_{+})\cap\Ann(\alpha_{-})$,
and the last two meet in $P\wedge Q$ by \Cref{thm:hhsplitting}(ii). The space
$P\wedge Q$ is the sum over $a<n\le b$ of the four-dimensional weight spaces
spanned by $z_{ab},v^{1}_{ab},v^{2}_{ab},\pi_{ab}$, and $\Ann(p(\theta))$ is
torus-stable, so the computation splits pair by pair. Put
$m=e_{b}\wedge f_{a}$ and $\theta'=\theta-t_{a}-t_{b}$, where
$t_{j}=i s_{j}e_{j}\wedge f_{j}$ with $s_{j}=\pm1$ the sign that makes
$\theta=\sum_{j}t_{j}$ positive. Since the $t_{j}$ are even with
$t_{j}^{2}=0$, direct evaluation gives
\[
  z\lrcorner\theta^{k}=m\,\theta'^{k},\quad
  v^{1}\lrcorner\theta^{k}=ik\,m\,\theta'^{k-1},\quad
  v^{2}\lrcorner\theta^{k}=-ik\,m\,\theta'^{k-1},\quad
  \pi\lrcorner\theta^{k}=k(k-1)\,m\,\theta'^{k-2},
\]
so for $\xi=xz+y_{1}v^{1}+y_{2}v^{2}+w\pi$ and $\lambda=i(y_{1}-y_{2})$,
\[
  \xi\lrcorner p(\theta)=m\wedge\sum_{j=0}^{2n-2}\frac{\theta'^{j}}{j!}
  \bigl(x\mu_{j}+\lambda\mu_{j+1}+w\mu_{j+2}\bigr).
\]
The classes $m\wedge\theta'^{j}$, $0\le j\le2n-2$, are nonzero and of distinct
degrees, so $\xi$ kills $p(\theta)$ exactly when $(x,\lambda,w)\in\ker H$.
The map $(x,y_{1},y_{2},w)\mapsto(x,\lambda,w)$ is onto with kernel spanned by
$v^{1}+v^{2}$, which gives $1+(3-\rho)$ dimensions for each of the $n^{2}$
pairs. (ii) The vector $v^{1}_{ab}+v^{2}_{ab}$ is $K$-linear and kills
$\theta$, and these $n^{2}$ vectors span the tangent space of the polarised
family by \Cref{prop:firstorder}; when $\rho=3$ they are all of the
annihilator. If $\rho\le2$, a nonzero $(x,\lambda,w)\in\ker H$ turns the
equations $x\mu_{j}+\lambda\mu_{j+1}+w\mu_{j+2}=0$ into a linear recurrence of
full rank on $\mu_{0},\dots,\mu_{2n}$, whose two-dimensional space of
solutions is spanned by the listed shapes according to the roots of
$wT^{2}+\lambda T+x$; $\rho\le1$ is the intersection of two such spaces.
(iii) At $n=2$ the argument of (i) fails for exactly one pair of weights,
$\epsilon_{2}+\epsilon_{3}$ in $\bigwedge^{2}P$ and $-\epsilon_{0}-\epsilon_{1}$
in $\bigwedge^{2}Q$, whose difference is $(1,1,1,1)$; these give one coupled
block of size eight, whose kernel is isomorphic to $\ker(X-T)$ with $T$
conjugate on its image to $(HK)^{2}$, and the characteristic polynomial of
$HK$ is $y^{3}-Iy-2J$, which is traceless, so $e\le2$. (iv) is the argument of
\Cref{thm:p2numerical}(i)--(ii) with $\Ann_{HT^{2}}(\gamma)$ in place of
$P\wedge Q$: $\sigma_{E}\circ\operatorname{ev}_{E}$ is contraction into
$\gamma$, so $\ker\operatorname{ev}_{E}\subseteq\Ann_{HT^{2}}(\gamma)$. (v) If
$\gamma-\omega$ has components in degrees $0$ and $4n$ only, the proof of
\Cref{thm:p2false} applies verbatim. If $p=c\,e^{t\theta}+c'\theta^{2n}$ with
$t\theta=c_{1}(L)$, then
$\ch(E\otimes L^{-1})=\gamma e^{-t\theta}=c+\omega+c'\theta^{2n}$, because
$\omega\theta=0$ and $\theta^{2n+1}=0$, and
$\sigma_{E\otimes L^{-1}}=e^{-c_{1}(L)}\cdot\sigma_{E}$ is injective exactly
when $\sigma_{E}$ is, so the first case applies to $E\otimes L^{-1}$.
Item (LII) of \Cref{sec:verification} computes the annihilator structure-free,
exactly over $\QQ(i)$, for $n\le10$, and checks every statement above.
\end{proof}

So the number that the corrected criterion asks of one complex is $24$ at
$n=2$, $57$ at $n=3$ and $104$ at $n=4$ for a general shape, against $12$,
$30$ and $56$ for the pure form, and injectivity then forces nothing beyond
first-order extension along the polarised family, which is what the
criterion is for. The larger number is the price of the components
$c_{k}\theta^{k}$, and it is paid in the right place: every direction of
$HT^{2}$ other than the $n^{2}$ directions of the family moves $\gamma$ off
the Hodge classes, so an object meeting the bound has $\operatorname{ev}_{E}$
injective on a complement of $T_{[A]}\cD_{n,n}$.
\Cref{fig:hankel} plots the number against $n$ for the four Hankel ranks.

\begin{figure}[tbp]
\centering
  \includegraphics[width=0.86\textwidth]{fig_hankel}
\caption{The dimension $r(\gamma)=(4+\rho)n^{2}-2n$ that
$\Ext^{2}(E,E)$ must have for a complex with Chern character
$\gamma=N\omega_{1}+\sum_{k}c_{k}\theta^{k}$ to meet the corrected criterion,
for the four Hankel ranks $\rho$ of \Cref{thm:p2primenumber}; the dots are the
values computed exactly over $\QQ(i)$ by item (LII), through $n=10$, where
$HT^{2}$ has dimension $780$. The pure form, $\rho=0$, is the bottom curve
$2n(2n-1)$, which no complex attains (\Cref{thm:p2false}); a general shape
lies on the top curve $7n^{2}-2n$. At $n=2$ the weight argument fails for one
pair of weights and seven values occur.}
\label{fig:hankel}
\end{figure}

\begin{proposition}[The Hochschild profile of a corrected character]
\label{prop:p2primeprofile}
In \Cref{setup:p2prime}, let $\rho_{k}$ be the rank of contraction
$HT^{k}(A)\to H^{*}(A,\CC)$ into $\gamma$, put $C_{l}=l!\,c_{l}$, let
$r'$ be the rank of the $(n+1)\times(n+1)$ Hankel matrix $H'=[C_{p+q}]$, and
$d=\dim\ker\bigl((K'H')^{2}-|u|^{2}\bigr)$ with
$K'=\mathrm{antidiag}\bigl((-1)^{a}\binom{n}{a}\bigr)$.
\begin{enumerate}[label=\textup{(\roman*)},leftmargin=2.6em,itemsep=2pt]
\item $\rho_{0}=\rho_{2n}=1$, $\rho_{k}=0$ for $k>2n$, and for
  $1\le k\le2n-1$, $k\ne n$,
  \[
    \rho_{k}=2\binom{2n}{k}+M_{k}\min(k+1,\,2n+1-k,\,r'),\qquad
    M_{k}=\sum_{\substack{i+j=k\\1\le i,j\le n-1}}\binom{n}{i}\binom{n}{j};
  \]
  in the middle degree the same formula holds with $d$ subtracted, which is
  checked for $n\le5$. In particular $\rho_{k}=\rho_{2n-k}$, and the general
  profiles are $(1,8,24,8,1)$ at $n=2$ and $(1,12,57,112,57,12,1)$ at
  $n=3$.
\item $\chi(E,E)=V\,Q(c)+W$ for a perfect complex with $\ch(E)=\gamma$, where
  $V=\int_{A}\theta^{2n}>0$, $W=(-1)^{n}\int_{A}\omega^{2}>0$ and
  $Q(c)=\sum_{k}(-1)^{k}c_{k}c_{2n-k}$; and $\chi(E,E)$ is even.
\item Granting the compatibility of \Cref{cor:hhfactor}(iii) in every degree,
  $\dim\Ext^{k}(E,E)\ge\rho_{k}$. At $n=2$, at the points with $d=1$ one has
  $\rho_{2}=23$, which is odd, so $\dim\Ext^{2}(E,E)=r(\gamma)$ cannot hold
  there, since $\chi(E,E)$ is even and Serre duality gives
  $\chi(E,E)\equiv\dim\Ext^{2}(E,E)\bmod2$.
\end{enumerate}
\end{proposition}

\begin{proof}
(i) Split $HT^{1}=Q\oplus P$ as in \Cref{thm:hhsplitting}, so that
$HT^{k}=\bigoplus_{i+j=k}\bigwedge^{i}Q\otimes\bigwedge^{j}P$ and $Q$, $P$ act
on the two tensor factors of $H^{*}(A)=H_{+}\otimes H_{-}$ separately. A mixed
monomial kills $\omega$, and on $p(\theta)$ the block of mixed monomials with
fixed index sets $I$, $J$ has entries $C_{s+t+m+m'}$ depending only on
$s+t$ and $m+m'$, so its rank is that of a catalecticant of the binary form
$\sum_{k}c_{k}x^{k}y^{2n-k}/(2n-k)!$, which is $\min(k+1,2n+1-k,r')$ by the
classical Hilbert function of an apolar ideal. The pure sources
$\bigwedge^{k}Q$ and $\bigwedge^{k}P$ contract $\alpha_{+}$ and $\alpha_{-}$
injectively into targets that no other source reaches unless $k=n$. The
symmetry is \Cref{lem:rankduality}. (ii) $\ch(E)^{\vee}$ changes the sign of
the components of odd codegree, $\omega\theta=0$ kills the cross terms, and
Hodge--Riemann gives $W>0$; parity holds because the Chern character of a
perfect complex on a complex torus is integral, and the intersection form on
$\bigwedge^{2n}H^{1}(A,\ZZ)$ pairs each monomial only with its complement, so
$\int x^{2}$ is even. (iii) is the argument of \Cref{cor:hhfactor}(iii).
Item (LIII) of \Cref{sec:verification} checks (i) exactly over $\QQ(i)$ for
$n=2,3$ in every degree and for $n=4$ in every degree with \texttt{--extreme},
and (ii) and the parity at $n=2$.
\end{proof}

The profile says what the corrected criterion does not constrain. For the pure
form, \Cref{thm:p2endo}(iv) forced $\dim\End(E)\ge3$ at $n=2$ and $n=3$ out of
$\chi(E,E)>0$. For a corrected character $\chi(E,E)=VQ(c)+W$ can have either
sign, and when it is small enough the Euler characteristic forces nothing: at
$n=2$ in the model of \texttt{explicit\_weil.py} with $d=1$, the character with
$c=(2,1,0,1,-1)$ has $\chi(E,E)=-136$, and a simple object is not excluded.

\begin{proposition}[The Hodge locus of a corrected character]
\label{prop:p2primelocus}
\begin{enumerate}[label=\textup{(\roman*)},leftmargin=2.6em,itemsep=2pt]
\item Let $T$ be a complex torus, let $x$ be a real class of type $(k,k)$ on
  $T$, and for $v\in\Hom(H^{1,0},H^{0,1})=H^{1}(T_{T})$ let $J_{v}$ be the
  complex structure on the underlying real torus with $H^{1,0}_{v}=(1+v)H^{1,0}$.
  Then $x$ is of type $(k,k)$ for $J_{v}$ if and only if $v\lrcorner x=0$. So
  in this chart the Hodge locus of $x$ is the linear space
  $\Ann_{H^{1}(T_{T})}(x)$, and it coincides with its first-order locus.
\item Let $J$ be any complex structure on the real torus underlying $A$. If
  $\theta^{k}$ is of type $(k,k)$ for $J$ for some $1\le k\le2n-1$, then
  $\theta$ is of type $(1,1)$ for $J$.
\item In \Cref{setup:p2prime} with $n\ge3$, the germ at $A$ of the Hodge locus
  of $\gamma$ among all complex structures is the polarised Weil family
  $\cD_{n,n}$, of dimension $n^{2}$, if $c_{k}\ne0$ for some $1\le k\le2n-1$,
  and the family of all $K$-linear deformations, of dimension $2n^{2}$,
  otherwise. In the first case the connected component of the Hodge locus
  through $A$ is $\cD_{n,n}$, and all its members are abelian varieties
  polarised by $\theta$.
\item For $n=2$ the same holds, with one exception: if $c_{1}=c_{3}=0$ and
  $c_{2}^{2}\int_{A}\theta^{4}=\tfrac34\int_{A}\omega^{2}$, the germ has
  dimension $5$. In the rational model of \texttt{explicit\_weil.py} this is
  the character $N\omega_{1}\pm\tfrac{d}{2}N\eta^{2}+c_{0}+c_{4}\eta^{4}$;
  with rational data the ratio is attained exactly in the families of trivial
  discriminant.
\end{enumerate}
\end{proposition}

\begin{proof}
(i) Write $D_{v}$ for the derivation of $\bigwedge^{*}H^{1}$ extending $v$.
Since $v^{2}=0$ on $H^{1}$, $\bigwedge^{*}(1+v)=\exp(D_{v})$, so the Hodge
filtration of $J_{v}$ is $F^{p}_{v}=\exp(D_{v})F^{p}$. Now
$\exp(-D_{v})x=\sum_{m}(-1)^{m}D_{v}^{m}x/m!$ with $D_{v}^{m}x$ of type
$(k-m,k+m)$, so $x\in F^{k}_{v}$ exactly when $D_{v}x=v\lrcorner x=0$; and a
real class lies in $F^{k}_{v}\cap\overline{F^{k}_{v}}=H^{k,k}_{v}$ as soon as
it lies in $F^{k}_{v}$. (ii) Let $t\mapsto t_{J}$ be the circle acting on
$H^{p,q}$ by $t^{p-q}$. It acts on $\bigwedge^{*}H^{1}(A,\RR)$ by algebra
automorphisms and fixes $\theta^{k}$, so the path $t\mapsto t_{J}\theta$ lies in
$\{y\in\bigwedge^{2}H^{1}(A,\RR):y^{k}=\theta^{k}\}$. At a nondegenerate $y$
the differential $h\mapsto k\,y^{k-1}h$ of $y\mapsto y^{k}$ is injective on
$\bigwedge^{2}$ when $2+(k-1)\le2n$, by the linear hard Lefschetz theorem for
the symplectic form $y$ on a space of dimension $4n$,\footnote{For a
nondegenerate $y\in\bigwedge^{2}U^{\vee}$ on a real space $U$ of dimension
$2m$, multiplication by $y^{j}$ is injective from $\bigwedge^{p}$ to
$\bigwedge^{p+2j}$ whenever $p+j\le m$, because $y$ generates an
$\mathfrak{sl}_{2}$-action on $\bigwedge^{*}U^{\vee}$ in which $\bigwedge^{p}$
has weight $p-m$. Here $m=2n$, $p=2$ and $j=k-1$. The bound is sharp:
$\theta^{2n}$ is a multiple of the volume form for every complex structure.} so near $\theta$ that set
is discrete, the path is constant, and $\theta$ is of type $(1,1)$. (iii) By
(i), applied to each homogeneous component, the Hodge locus of $\gamma$ in the
chart is $\Ann_{H^{1}(T_{A})}(\gamma)$, the part of $\Ann_{HT^{2}}(\gamma)$ in
$H^{1}(T_{A})$. In the description of \Cref{thm:p2primenumber}(i) that part is
$x=w=0$ with $(0,\lambda,0)\in\ker H$; the vector $(0,1,0)$ lies in $\ker H$
exactly when $\mu_{1}=\dots=\mu_{2n-1}=0$, so the part is spanned by the $n^{2}$
vectors $v^{1}_{ab}+v^{2}_{ab}$, which span $T_{[A]}\cD_{n,n}$, when some
$c_{k}$ with $1\le k\le2n-1$ is nonzero, and by all $v^{1}_{ab},v^{2}_{ab}$,
the $K$-linear deformations, otherwise. A germ of submanifold of dimension
$n^{2}$ inside a linear space of the same dimension is that space. On
$\cD_{n,n}$ the class $\theta$ stays of type $(1,1)$ and nondegenerate, so the
signature of the hermitian form it defines is constant and every member is
polarised by $\theta$. The same argument at every member shows that
$\cD_{n,n}$ is open in the Hodge locus; it is closed in the space of complex
structures, since $K$-linearity is a closed condition and a limit of positive
hermitian forms that stays nondegenerate stays positive; and it is
connected. (iv) At $n=2$ the weight argument of
\Cref{thm:p2primenumber}(i) fails for one pair of weights, and the
first-order locus is computed directly: its dimension is $4$, $8$ when
$c_{1}=c_{2}=c_{3}=0$, and $5$ exactly at the stated ratio, which in the model
reads $|u|=4|c_{2}|$ with $\int\theta^{4}$ and $\int\omega^{2}$ equal to $24$
and $2|u|^{2}$ times the same volume. In the rational model $\int\omega_{1}^{2}=8d^{2}$
and $\int\eta^{4}=24$, which gives $c_{2}=\pm dN/2$. For a general rational
Weil class $\omega=x\omega_{1}+y\omega_{2}$ and hermitian form
$\mathrm{diag}(-a_{1},-a_{2},a_{3},a_{4})$ one has
$\int\omega^{2}=8d\,\Nm_{K/\QQ}(y+x\sqrt{-d})$ and
$\int\theta^{4}=24\,a_{1}a_{2}a_{3}a_{4}$, so the ratio asks
$(c_{2}/N)^{2}=d\,\Nm_{K/\QQ}(y+x\sqrt{-d})/(4a_{1}a_{2}a_{3}a_{4})$, a rational
square exactly when $a_{1}a_{2}a_{3}a_{4}$ is a norm from $K$. Items (LII)(E)
and (F) of \Cref{sec:verification} check the dimensions and the two integrals
exactly.
\end{proof}

So \Cref{thm:p2false} does not extend to the corrected form, and the reason is
visible. A semiregular complex deforms along the Hodge locus of its Chern
character; for the pure form that locus is the whole $2n^{2}$-dimensional
family of $K$-linear tori, whose very general member is not algebraic, while
for a corrected character with some $c_{k}\ne0$, $1\le k\le2n-1$, it is the
polarised family itself, all of whose members are abelian varieties, and
there is nothing for the argument through Weil tori to act on. That is a
statement about the argument run on $E$ itself; an autoequivalence can still
remove a component, and a line bundle twist is what reaches the corrected
characters of \Cref{thm:p2primenumber}(v). The one other place it can reach is
the exceptional ratio of (iv), where the Hodge locus leaves the polarised
family.

\begin{remark}[The exceptional ratio at $n=2$]\label{rem:p2primeexceptional}
At $n=2$, with $c_{1}=c_{3}=0$ and $c_{2}^{2}\int\theta^{4}=\frac34\int\omega^{2}$,
put $\gamma_{2}=\omega+c_{2}\theta^{2}$. Item (LII)(G) of
\Cref{sec:verification} computes its stabiliser $\mathfrak{g}$ in
$\mathfrak{gl}(H^{1}(A,\CC))$ exactly: it has dimension $21$, Hodge pieces of
dimensions $5$, $11$, $5$, commutant on $H^{1}$ equal to the scalars, and a
real form whose trace form has signature $(12,9)$; at an ordinary ratio the
stabiliser is $\mathfrak{su}(2,2)$, of dimension $15$. The scalars do not fix
$\gamma_{2}$, so $\mathfrak{g}$, acting irreducibly with a nondegenerate trace
form, is semisimple; a product of two factors would act on
$H^{1}=V_{1}\otimes V_{2}$ with $\dim V_{i}\in\{2,4\}$ and have dimension at
most $3+15$, so $\mathfrak{g}$ is simple of dimension $21$. The simple Lie
algebras of dimension $21$ are $\mathfrak{so}(7,\CC)$ and
$\mathfrak{sp}(6,\CC)$, and only the first has an irreducible representation
of dimension $8$, its spin representation; among its real forms only
$\mathfrak{so}(4,3)$ has a trace form of signature $(12,9)$. Its invariants in
$\bigwedge^{*}H^{1}$ are spanned by $1$, $\gamma_{2}$ and the class of a point,
so $\gamma_{2}$ is the invariant $4$-form of $\mathrm{Spin}(4,3)$, a split form
of the Cayley form.\footnote{In the compact case the stabiliser in
$\GL(8,\RR)$ of the Cayley $4$-form is $\mathrm{Spin}(7)$, acting on $\RR^{8}$
by its spin representation; here the same happens for the real form of
signature $(4,3)$. We know of no earlier appearance of this form among the
Hodge classes of abelian fourfolds, but we have not searched the literature on
Hodge loci for it systematically.} The Hodge locus of $\gamma$ near $A$ is the orbit of $A$
under the connected real group $G$ of $\mathfrak{g}$, since that orbit has
dimension $\dim\mathfrak{g}^{-1,1}=5$ and lies in the locus; the Hodge group
of a very general member $T$ has Lie algebra an ideal of $\mathfrak{g}$, hence
all of it, so the Hodge classes of $T$ are $1$, $\gamma_{2}$ and the class of a
point, and $\NS(T)=0$. The locus reaches tori that are not algebraic, and
\Cref{thm:p2false} can be started there. It does not finish: its last step,
\Cref{prop:weiltorussheaves}, needs a K\"ahler class on which the degree-four
Chern character integrates to zero, and here
\[
  \int_{A}\gamma_{2}\,\kappa_{h}^{2}
  =\frac{|c_{2}|}{6}\Bigl(\operatorname{sign}(c_{2})\,\sigma_{2}(h)
    +4\operatorname{Re}\bigl(\phi\det h_{01,23}\bigr)\Bigr)\int_{A}\theta^{4},
  \qquad|\phi|=1,
\]
for the K\"ahler class $\kappa_{h}=i\sum h_{jk}z_{j}\wedge\bbar z_{k}$ of a
positive hermitian $h$ in a unitary frame for $\theta$, which has the sign of
$c_{2}$. Indeed, writing $h=\sum_{j}\lambda_{j}w_{j}w_{j}^{*}$ with $(w_{j})$
orthonormal, Cauchy--Binet gives
$\sigma_{2}(h)\pm4\operatorname{Re}(\phi\det h_{01,23})
=\sum_{j<k}\lambda_{j}\lambda_{k}\bigl(1\pm4\operatorname{Re}(\phi\,\xi_{01}\bbar\xi_{23})\bigr)$
with $\xi=w_{j}\wedge w_{k}$ of norm one; the Pl\"ucker relation and the
inequality of the means give $|\xi_{01}\xi_{23}|\le\frac14$, so each term is
nonnegative, and Parseval gives $\sum_{j<k}$ of the brackets $=6$, so the sum is
at least $6\min_{j}\lambda_{j}^{2}>0$. Since $G$ fixes $\gamma_{2}$ and carries
K\"ahler cones to K\"ahler cones, the same holds on every member of the locus.
So the argument excludes at most a torsion-free sheaf with $c_{2}>0$, which
Bogomolov's inequality already forbids to be $\theta$-semistable, and not a
complex. Item (LII)(G) checks the displayed identity at $c_{2}=1$, $\phi=1$,
and the constancy of the sign at random K\"ahler classes, against a change of
sign at $R=3/16$. At $n=2$ Markman's theorem settles every Weil fourfold, so
this is a statement about the shape of (P2$'$), not about the conjecture.
\end{remark}

\begin{figure}[tbp]
\centering
  \includegraphics[width=0.86\textwidth]{fig_kahlersign}
\caption{The exceptional ratio of \Cref{rem:p2primeexceptional}. By the
identity displayed there, which is linear in $\gamma$, the infimum over the
K\"ahler cone of $6\int\gamma\kappa^{2}\big/\bigl(|c|\,\sigma_{2}(h)\int
\theta^{4}\bigr)$ is $m(R)=1-\sqrt{3/(4R)}$, approached at the boundary of the
cone by a positive hermitian form of rank two whose Pl\"ucker coordinates
satisfy $|\xi_{01}\xi_{23}|=\frac14$. It is negative exactly for $R<3/4$, and
$R=3/4$ is also where the Hodge locus of $\gamma$ jumps from dimension $4$ to
$5$ and its stabiliser becomes $\mathfrak{so}(4,3)$, with $r(\gamma)=23$; at
$R=3$ the annihilator has dimension $6$ and $r(\gamma)=22$. Items (LII)(D),
(E) and (G) check these values.}
\label{fig:kahlersign}
\end{figure}

\begin{remark}[Which obstructions survive the correction]
\label{rem:p2primesurvive}
The obstruction theorems of this section and of \Cref{sec:semireg} were
proved for a Chern character on the Weil line; for a corrected character they
divide as follows. Unchanged, with $r(\gamma)$ in place of $2n(2n-1)$ where a
number occurs: \Cref{thm:linebundle} and \Cref{cor:noline}, whose only
hypothesis is that $\ch(E)$ stay of Hodge type to first order along
$\cD_{n,n}$; \Cref{thm:positivity} and \Cref{cor:nofiltered}, since a
corrected character is a flat Hodge class in every degree;
\Cref{lem:finiteproduct}; and the first two parts of \Cref{thm:p2numerical},
which become \Cref{thm:p2primenumber}(iv). Unchanged only when
$c_{0}=\dots=c_{2n-2}=0$, so that $H^{0,2}$ still kills $\gamma$:
\Cref{thm:rankzero}, \Cref{prop:concentrated}, \Cref{cor:kernelfamily},
\Cref{thm:weakfails} and \Cref{thm:p2forces}(iii). Unchanged only when
$\ker H$ contains a vector with $w\ne0$, which needs $\rho\le2$, the vector
giving a compensated class of $P\wedge Q$ with a component in
$T_{+}\otimes T_{-}$: \Cref{thm:p2forces}(iv), \Cref{cor:p2shape} and
\Cref{thm:p2support}. Lost:
\Cref{thm:p2false}, outside \Cref{thm:p2primenumber}(v); the positivity of
$\chi(E,E)$ in \Cref{prop:objectsize} and with it \Cref{thm:p2endo}(i)--(iii),
by \Cref{prop:p2primeprofile}(ii); and \Cref{thm:nosum}. So for a general
shape, $\rho=3$, what this paper still excludes is a line bundle summand off
the polarisation line and a Chern character effective in the exponential
basis, and nothing else; \Cref{prop:scalaratiyah} below extends the first to
summands with a scalar Atiyah class. An object with no such summand, whose
Chern character is not effective, and with $\dim\Ext^{2}(E,E)=7n^{2}-2n$, is
not excluded by anything proved here, and none is known.
\end{remark}

\begin{proposition}[Summands with a scalar Atiyah class]\label{prop:scalaratiyah}
Let $A$ be of $(K,-1,n)$-Weil type with $n\ge2$, let $T=T_{[A]}\cD_{n,n}$, and
let $E$ be a perfect complex on $A$ with $v\lrcorner\ch(E)=0$ for every
$v\in T$. Suppose $E$ has a direct summand $F[p]$ with
$\mathrm{At}(F)=\delta\cdot\id_{F}$, where $\delta$ is a real class of type
$(1,1)$ that is not a real multiple of $\theta$, and $\rk F\ne0$. Then there is
$v\in T$ with $\operatorname{ev}_{E}(v)\ne0$ and
$\sigma_{E}(\operatorname{ev}_{E}(v))=0$. In particular no finite direct sum
of shifted bundles $\pi_{*}L$, $\pi$ an isogeny, among them the simple
semi-homogeneous bundles, has a corrected character $\gamma$ with $N\ne0$ and
injective semiregularity map.
\end{proposition}

\begin{proof}
As in \Cref{thm:linebundle}, $\operatorname{ev}_{E}(v)=v\lrcorner\mathrm{At}(E)$
on $H^{1}(T_{A})$ and the Atiyah class of a direct sum is block diagonal, so
the component of $\operatorname{ev}_{E}(v)$ in $\Ext^{2}(F,F)$ is
$(v\lrcorner\delta)\cdot\id_{F}$; the map $z\mapsto z\cdot\id_{F}$ on
$H^{0,2}$ is injective because $\mathrm{tr}(z\cdot\id_{F})=\rk(F)\,z$, and
some $v\in T$ has $v\lrcorner\delta\ne0$ by \Cref{prop:rigidity}(i). On the
other hand $\sigma_{E}(\operatorname{ev}_{E}(v))=v\lrcorner\ch(E)=0$. For the
last sentence, the proof of \Cref{lem:blocks}(i) gives
$\mathrm{At}(\pi_{*}L)=\delta\cdot\id$ with $\pi^{*}\delta=c_{1}(L)$ for any
line bundle $L$, and $\ch(\pi_{*}L)=\rk(\pi_{*}L)\,e^{\delta}$; if every
$\delta$ were a multiple of $\theta$ the Chern character of the sum would lie
in $\QQ[\theta]$, which does not contain $\omega$ because $\omega\theta=0$ and
$\theta^{n+1}\ne0$.
\end{proof}

\Cref{thm:semiregclosure} asks that $\sigma_{E}$ be injective. What its proof
uses is that the relative deformations of $E$ dominate the family, and
Pridham's form of \Cref{thm:BF}, in which $\sigma_{E}$ annihilates every
obstruction, gives that under a weaker hypothesis. It tolerates an object that
is obstructed to first order, which is exactly how \Cref{thm:linebundle},
\Cref{cor:noline} and \Cref{prop:scalaratiyah} fail an object.

\begin{proposition}[A flatness form of the criterion]\label{prop:p2flat}
Let $n\ge2$, $s_{0}\in\cD_{n,n}$, and let $E$ be a bounded complex of vector
bundles on $A_{s_{0}}$ with $\ch(E)=N\omega_{1}+\sum_{k}c_{k}\theta^{k}$ and
$N\ne0$. Let $K_{E}$ be the base of a miniversal deformation of $E$ on the
fixed variety $A_{s_{0}}$, put $e=\dim\Ext^{1}(E,E)$ and
$k=\dim\ker\sigma_{E}$, and grant \Cref{thm:BF} in Pridham's form, in which
$\sigma_{E}$ annihilates every obstruction, both on $A_{s_{0}}$ and along the
Artinian deformations of $A_{s_{0}}$ in $\cD_{n,n}$.
\begin{enumerate}[label=\textup{(\roman*)},leftmargin=2.6em,itemsep=2pt]
\item $\dim K_{E}\ge e-k$.
\item If $\dim K_{E}=e-k$, then $\Sigma=\cD_{n,n}$, and the Hodge conjecture
  holds in every degree for every member of the family with
  $\Hg(A)=\SU(V,H)$.
\item $\dim K_{E}=e-k$ as soon as the cone
  $\{\xi\in\Ext^{1}(E,E):\xi\circ\xi=0\}$ has codimension $k$.
\end{enumerate}
For $k=0$ this is \Cref{thm:semiregclosure}.
\end{proposition}

\begin{proof}
Let $R$ be the completed local ring of $\cD_{n,n}$ at $s_{0}$, regular of
dimension $n^{2}$. By hypothesis $\ker\sigma_{E}$ is a complete obstruction
space for the deformations of the pair $(A_{s},E)$ over Artinian $R$-algebras,
whose relative tangent space is $\Ext^{1}(E,E)$; so this functor has a hull
$P/J$ with $P=R[[x_{1},\dots,x_{e}]]$ and $J$ generated by at most $k$
elements \cite{FM98}, and its fibre over the closed point of $R$ is a hull of
the absolute functor, of dimension $\dim K_{E}$. Krull's height theorem gives
$\dim P/J\ge n^{2}+e-k$, and the same count on the fibre gives (i). If
$\dim K_{E}=e-k$, the dimension inequality for a local homomorphism
\cite[Theorem~15.1]{Mat86} gives $\dim P/J\le n^{2}+e-k$, so equality holds, $J$
has height $k$ and is generated by a regular sequence, $P/J$ is
Cohen--Macaulay, and it is flat over $R$ by miracle flatness
\cite[Theorem~23.1]{Mat86}.\footnote{A local homomorphism from a regular local
ring to a Cohen--Macaulay local ring is flat exactly when the dimension of the
target is the dimension of the source plus that of the fibre; here both sides
equal $n^{2}+e-k$.} The analytic Kuranishi germ of pairs, as in the
proof of \Cref{thm:p2false}, has $P/J$ as its completion, so it is flat over
the germ of $\cD_{n,n}$ at $s_{0}$ and therefore open onto a neighbourhood
$U$ of $s_{0}$. Every $A_{s}$ with $s\in U$ then carries a bounded complex of
vector bundles, algebraic by GAGA, whose Chern character is the parallel
transport of $\ch(E)$, so $N\omega_{1,s}$ differs from an algebraic class by a
polynomial in $\theta$ and is algebraic; \Cref{cor:allornothing} gives (ii).
For (iii), the quadratic part of the Kuranishi map of $E$ is
$\xi\mapsto\xi\circ\xi$, so the tangent cone of $K_{E}$ lies in the cone of the
hypothesis, and $\dim K_{E}\le e-k$. Every step uses only deformations along
$\cD_{n,n}$, which are projective, where \Cref{thm:p2false} needs
non-projective ones.
\end{proof}

No object is known that meets (iii) with $k>0$, and none is excluded by
anything in this paper; the proposition says only that injectivity of
$\sigma_{E}$ is more than propagation needs.

\subsection{What the object must be, and what it cannot be}
\label{ssec:whatobject}

\Cref{thm:semiregclosure} asks for one perfect complex at one base point.
The cheapest candidate is a sum of line bundles on a split member, because
there the Weil classes are already polynomials in $\beta$, $\widehat\beta$ and
$\ell$ by \Cref{thm:splitclosed}, so the Chern character one wants is
manifestly in reach. This subsection settles that candidate. It does not
exist, and the obstruction is not semiregularity: it is the Chern character
condition, and it fails because the norm form of $K$ is positive
definite.\footnote{The norm form has now closed three different doors in this
paper, by three different mechanisms, and they should be kept apart.
\Cref{thm:character} uses the norm \emph{character} of $\Kx$ and an
equivariance hypothesis. \Cref{prop:divfails} uses the fact that the part $R$
of $H^{2}$ on which $\Kx$ acts through $\tau^{2}$ and $\bbar\tau^{\,2}$ has no
Hodge classes at a general point, which is a statement about the Hodge group.
\Cref{thm:positivity} below uses neither: it uses that $\Nm_{K/\QQ}$ is a
positive definite quadratic form on $K$ and that the multiplicities of a
filtration are positive integers. The three are independent, and only the
third is available at a special point, which is exactly where the object of
\Cref{prop:ktheory} has to live.}

\begin{lemma}[The semiregularity map of a split object]\label{lem:splitsemireg}
Let $Y$ be a smooth projective complex variety and let
\[
  E=\bigoplus_{i=1}^{s}L_{i}[p_{i}]^{\oplus m_{i}},
  \qquad \lambda_{i}=c_{1}(L_{i}),
\]
with the $L_{i}$ line bundles, the $\lambda_{i}$ pairwise distinct and
$p_{i}\in\ZZ$. Then for $\zeta=(\zeta_{ij})\in\Ext^{2}(E,E)$,
\begin{equation}\label{eq:splitsigma}
  \sigma_{E}(\zeta)_{q}
  =\frac{1}{q!}\sum_{i=1}^{s}(-1)^{p_{i}}\,
    \mathrm{tr}(\zeta_{ii})\,\lambda_{i}^{\,q},
\end{equation}
so $\sigma_{E}$ annihilates $\Ext^{2}(L_{i}[p_{i}],L_{j}[p_{j}])$ for $i\ne j$
and the traceless part of each diagonal block. Consequently $\sigma_{E}$ is
injective if and only if all three of the following hold:
\begin{enumerate}[label=\textup{(\roman*)},nosep]
\item $m_{i}=1$ for every $i$;
\item $H^{\,2+p_{i}-p_{j}}\bigl(Y,L_{j}L_{i}^{-1}\bigr)=0$ for all $i\ne j$;
\item the map
  \[
    \Phi\colon H^{2}(\cO_{Y})^{\oplus s}\longrightarrow H^{*}(Y,\CC),
    \qquad
    (\zeta_{i})\longmapsto\sum_{i}e^{\lambda_{i}}\zeta_{i},
  \]
  is injective.
\end{enumerate}
\end{lemma}

\begin{proof}
The Atiyah class of a direct sum is the direct sum of the Atiyah classes, and
that of a line bundle is its first Chern class, so $\mathrm{At}(E)$ is
diagonal with entries $\lambda_{i}$ and $\mathrm{At}(E)^{q}$ is diagonal with
entries $\lambda_{i}^{\,q}$. The matrix product $\zeta\circ\mathrm{At}(E)^{q}$
therefore has $(i,j)$ entry $\zeta_{ij}\lambda_{i}^{\,q}$, and the trace keeps
only the diagonal; a shift by $p$ multiplies the trace by $(-1)^{p}$. That is
\eqref{eq:splitsigma}, and it exhibits the two spaces the map kills.

For the equivalence, the off-diagonal part of $\Ext^{2}(E,E)$ is
$\bigoplus_{i\ne j}\Ext^{2}(L_{i}[p_{i}],L_{j}[p_{j}])
=\bigoplus_{i\ne j}H^{2+p_{i}-p_{j}}(L_{j}L_{i}^{-1})$, which must vanish,
giving (ii); the diagonal block at $i$ is $H^{2}(\cO_{Y})\otimes
\mathrm{End}(\CC^{m_{i}})$ and only its trace survives, so $m_{i}=1$, giving
(i). What is left is $\bigoplus_{i}H^{2}(\cO_{Y})$, and by
\eqref{eq:splitsigma} an element $(\zeta_{i})$ of it dies exactly when
$\sum_{i}\lambda_{i}^{\,q}\zeta_{i}=0$ for every $q\ge0$, which is (iii) after
absorbing the signs $(-1)^{p_{i}}$ into the $\zeta_{i}$ and the factorials
into the grading.
\end{proof}

Condition (ii) is cheap, because the shifts $p_{i}$ are at our disposal: only
their parities are fixed by the Chern character, and a line bundle whose first
Chern class is nondegenerate of index $k$ has cohomology in degree $k$ alone,
so one may choose the $p_{i}$ to avoid every index that occurs. Condition
(iii) is a rank. Condition (i) is free. What is not free is the Chern
character itself, and that is what the rest of this subsection is about.

\begin{notation}[Coordinates on the divisor classes]\label{not:splitcoords}
Let $A_{0}=X\times\widehat{X}$ be a split member, $n\ge2$, and keep
$\gamma=d\beta-\widehat{\beta}$ from \Cref{not:threeclasses}. Write
\[
  \theta^{\pm}=\gamma\mp\sqrt{-d}\,\ell\in\textstyle\bigwedge^{2}V_{\pm},
  \qquad
  R_{0}=\QQ\gamma\oplus\QQ\ell\subset R .
\]
Every $\theta\in R_{0}$ is $\theta=u\theta^{+}+\bbar u\,\theta^{-}$ for a
unique $u\in K$, and in coordinates $u=g+h\sqrt{-d}$ corresponds to
$\theta=2g\gamma+2dh\,\ell$. By \Cref{prop:splitgeom}(iv) the group
$\NS(A_{0})_{\QQ}$ at a split member with $X$ general is
$\QQ\eta\oplus R_{0}$, so every divisor class is $e\eta+\theta$ with
$e\in\QQ$ and $\theta\in R_{0}$, and we write $u(\lambda)\in K$ for the
coordinate of its $R_{0}$-part.
\end{notation}

\begin{lemma}[A basis in degree four]\label{lem:sixmonomials}
For $n\ge2$ the six products
$\beta^{2},\widehat\beta^{2},\ell^{2},\beta\widehat\beta,\beta\ell,
\widehat\beta\ell$
are linearly independent in $H^{4}(A_{0},\QQ)$. Equivalently the six classes
$\eta^{2}$, $\eta\theta^{+}$, $\eta\theta^{-}$, $(\theta^{+})^{2}$,
$(\theta^{-})^{2}$, $\theta^{+}\theta^{-}$ are a basis of the degree four part
of the subring generated by the divisor classes.
\end{lemma}

\begin{proof}
Grade $H^{*}(A_{0},\QQ)=\bigwedge^{*}U\otimes\bigwedge^{*}U^{\vee}$ by
bidegree. Then $\beta$ has bidegree $(2,0)$, $\widehat\beta$ has $(0,2)$ and
$\ell$ has $(1,1)$, so the six products have bidegrees $(4,0)$, $(0,4)$,
$(2,2)$, $(2,2)$, $(3,1)$, $(1,3)$. All are distinct except for the pair
$\beta\widehat\beta$ and $\ell^{2}$, and those two are independent in
$\bigwedge^{2}U\otimes\bigwedge^{2}U^{\vee}$ because
$\beta\widehat\beta=\beta\otimes\widehat\beta$ is decomposable there while
$\ell^{2}=-2\sum_{i<j}(x_{i}x_{j})\otimes(\xi_{i}\xi_{j})$ has all
$\binom{2n}{2}$ terms and $\binom{2n}{2}>n$ for $n\ge2$. The second
description follows because the change of basis between the two sets of six
classes is invertible over $K$: $\eta=d\beta+\widehat\beta$ and
$\gamma=d\beta-\widehat\beta$ span the same plane as
$\beta$ and $\widehat\beta$ since $\det\begin{psmallmatrix}d&1\\
d&-1\end{psmallmatrix}=-2d\ne0$, and $\theta^{\pm}$ span the same plane
as $\gamma$ and $\ell$ over $\CC$.
\end{proof}

The next lemma says the same thing at every abelian variety of Weil type, not
only at a split member, and it says it without coordinates. It is the fact the
positivity argument really turns on.

\begin{lemma}[$\theta^{+}\theta^{-}$ is never proportional to $\eta^{2}$]
\label{lem:tensorrank}
Let $(A,\iota,\eta)$ be of $(K,-1,n)$-Weil type with $n\ge2$, and let
$0\ne\theta^{\pm}\in\bigwedge^{2}V_{\pm}$. Then $\theta^{+}\theta^{-}$ is not a
scalar multiple of $\eta^{2}$ in $H^{4}(A,\CC)$.
\end{lemma}

\begin{proof}
Both classes lie in the summand
$\bigwedge^{2}V_{+}\otimes\bigwedge^{2}V_{-}$ of $H^{4}(A,\CC)$, the isotypic
component for the character $\Nm^{2}$, so the question is one of tensor rank
there. The polarisation class $\eta$ lies in $V_{+}\otimes V_{-}$ and is the
hermitian form of \Cref{prop:weilline}, hence nondegenerate: in dual bases
$\eta=\sum_{a=1}^{2n}v_{a}\wedge w_{a}$ with $v_{a}$ a basis of $V_{+}$ and
$w_{a}$ a basis of $V_{-}$. Squaring,
\[
  \eta^{2}
  =\sum_{a,b}(v_{a}\wedge w_{a})(v_{b}\wedge w_{b})
  =-2\sum_{a<b}(v_{a}v_{b})\otimes(w_{a}w_{b}),
\]
which is diagonal in the induced bases of $\bigwedge^{2}V_{+}$ and
$\bigwedge^{2}V_{-}$ and therefore has tensor rank exactly $\binom{2n}{2}$.
The class $\theta^{+}\theta^{-}$ is, up to sign,
$\theta^{+}\otimes\theta^{-}$, of tensor rank one. For $n\ge2$ one has
$\binom{2n}{2}\ge6>1$, so the two cannot be proportional.
\end{proof}

\begin{theorem}[The positivity obstruction]\label{thm:positivity}
Let $(A,\iota,\eta)$ be of $(K,-1,n)$-Weil type with $n\ge2$, and suppose that
the only divisor classes on which $\Kx$ acts through the norm character are
the multiples of $\eta$, that is
\begin{equation}\label{eq:noextranorm}
  \NS(A)_{\QQ}\cap\bigl(V_{+}\otimes V_{-}\bigr)=\QQ\,\eta .
\end{equation}
Let $E$ be a perfect complex on $A$ whose Chern character is effective in the
exponential basis, meaning
\begin{equation}\label{eq:effectivech}
  \ch(E)=\sum_{i=1}^{s}m_{i}\,e^{\lambda_{i}},
  \qquad m_{i}\in\QQ_{>0},\quad\lambda_{i}\in\NS(A)_{\QQ},
\end{equation}
and write $\lambda_{i}=e_{i}\eta+\theta_{i}$ with $\theta_{i}$ in the part $R$
of $H^{2}(A,\QQ)$ carrying the characters $\tau^{2}$ and $\bbar\tau^{\,2}$, and
$u_{i}\in K$ for the $\theta^{+}$-coordinate of $\theta_{i}$. If $\ch_{2}(E)$
is a flat Hodge class, then
\begin{equation}\label{eq:normsum}
  \sum_{i=1}^{s}m_{i}\,\Nm_{K/\QQ}(u_{i})=0 ,
\end{equation}
and since $\Nm_{K/\QQ}$ is positive definite and every $m_{i}$ is positive,
$u_{i}=0$ for every $i$. Hence every $\lambda_{i}$ is a rational multiple of
$\eta$, and $\ch_{k}(E)\in\QQ\,\eta^{k}$ for every $k$. In particular
$\ch_{n}(E)$ has no component on the Weil line.
\end{theorem}

\begin{proof}
From \eqref{eq:effectivech},
$\ch_{2}(E)=\tfrac12\sum_{i}m_{i}\lambda_{i}^{2}$. Decompose $H^{4}(A,\CC)$
into $\Kx$-isotypic pieces and keep the one for $\Nm^{2}$, namely
$\bigwedge^{2}V_{+}\otimes\bigwedge^{2}V_{-}$. Writing
$\theta_{i}=u_{i}\theta^{+}+\bbar u_{i}\theta^{-}$ for a fixed choice of
generators $\theta^{\pm}$ of $\bigwedge^{2}V_{\pm}$, one has
\[
  \lambda_{i}^{2}
  =e_{i}^{2}\eta^{2}
   +2e_{i}\eta\theta_{i}
   +u_{i}^{2}(\theta^{+})^{2}+\bbar u_{i}^{2}(\theta^{-})^{2}
   +2\Nm_{K/\QQ}(u_{i})\,\theta^{+}\theta^{-},
\]
and of these five terms only the first and the last are in the $\Nm^{2}$
piece: $\eta\theta_{i}$ carries $\Nm\tau^{2}$ and $\Nm\bbar\tau^{\,2}$, and
$(\theta^{\pm})^{2}$ carry $\tau^{4}$ and $\bbar\tau^{\,4}$, all four distinct
from $\Nm^{2}$. Hypothesis \eqref{eq:noextranorm} is what makes the
decomposition $\lambda_{i}=e_{i}\eta+\theta_{i}$ available with $e_{i}$ a
scalar. So the $\Nm^{2}$-part of $\ch_{2}(E)$ is
\[
  \tfrac12\Bigl(\sum_{i}m_{i}e_{i}^{2}\Bigr)\eta^{2}
  +\Bigl(\sum_{i}m_{i}\Nm_{K/\QQ}(u_{i})\Bigr)\theta^{+}\theta^{-}.
\]

Now compare with the flat Hodge classes of degree four. For $n\ge3$ they are
$\QQ\eta^{2}$ by \Cref{thm:hdgcount}. For $n=2$ they are
$\QQ\eta^{2}\oplus\HW(A)$, and $\HW(A)$ spans the summands for $\tau^{4}$ and
$\bbar\tau^{\,4}$, which are not $\Nm^{2}$. In both cases the $\Nm^{2}$-part
of a flat Hodge class of degree four is a multiple of $\eta^{2}$. By
\Cref{lem:tensorrank} the class $\theta^{+}\theta^{-}$ is not such a multiple,
so its coefficient must vanish, which is \eqref{eq:normsum}.

Finally $\Nm_{K/\QQ}(g+h\sqrt{-d})=g^{2}+dh^{2}$ with $d>0$, so
$\Nm_{K/\QQ}(u)\ge0$ with equality only at $u=0$, and a sum of nonnegative
terms with positive weights vanishes only when every term does. Then
$\ch_{k}(E)=\frac{1}{k!}\bigl(\sum_{i}m_{i}e_{i}^{k}\bigr)\eta^{k}$, and
$\eta^{k}$ is not on the Weil line by \Cref{thm:hdgcount}.
\end{proof}

\begin{corollary}[The cheapest objects do not exist]\label{cor:nofiltered}
Let $A$ satisfy \eqref{eq:noextranorm}, which holds at every split member of
\Cref{thm:splitclosed} and at every quaternionic base point of
\Cref{thm:quatweil}.\footnote{At a split member $\NS(A)_{\QQ}$ has rank three
with basis $\beta,\widehat\beta,\ell$ by \Cref{prop:splitgeom}(iv), and of
these $\eta$ spans the norm-character part while $\gamma$ and $\ell$ span
$R$. At a quaternionic point $\NS(A)_{\QQ}$ is the space of Rosati-symmetric
elements of $B=(-d,b)_{\QQ}$, which is spanned by $1$, $j$ and $ij$; the first
gives $\eta$ and the other two are semilinear over $K$, hence lie in $R$. In
both cases the norm-character part of $\NS(A)_{\QQ}$ is exactly $\QQ\eta$.}
Then \Cref{thm:semiregclosure} cannot be applied at $A$ with any of the
following:
\begin{enumerate}[label=\textup{(\alph*)},nosep]
\item a coherent sheaf admitting a filtration by line bundles, in particular
  any direct sum of line bundles;
\item a simple semi-homogeneous vector bundle in the sense of Mukai
  \cite{Muk78}, since such a bundle of rank $r$ has $\ch(E)=r\,e^{c_{1}(E)/r}$;
\item any iterated extension of the objects in (a) and (b).
\end{enumerate}
In each case the failure is not of the semiregularity hypothesis but of the
Chern character hypothesis of \Cref{prop:ktheory}: no such $E$ has a Chern
character which is a flat section of Hodge classes with a nonzero Weil
component.
\end{corollary}

\begin{proof}
Each of (a), (b) and (c) has a Chern character of the shape
\eqref{eq:effectivech} with positive multiplicities, by additivity of $\ch$
along a filtration and by Mukai's computation of the Chern character of a
simple semi-homogeneous bundle. Apply \Cref{thm:positivity}.
\end{proof}

\begin{remark}[Where the theorem has force, and why it does not contradict
Markman]\label{rem:markmanpositivity}
At a very general member of the family $\NS(A)_{\QQ}=\QQ\eta$, so
\Cref{thm:positivity} is true for the empty reason that there are no
$\theta_{i}$ to begin with. Its content is at the special members, which is
exactly where \Cref{thm:semiregclosure} needs its object to live and where the
base points of \Cref{thm:everyfamily} are. That is the opposite of the usual
situation in this paper, where the special members are the easy ones.

The theorem also does not touch the objects of Markman \cite{Mar25a,Mar25b}.
Those are sheaves whose class in $K_{0}$ is not a positive combination of
exponentials of divisor classes, so \eqref{eq:effectivech} fails for them at
the first step; and they live at very general points, where as just noted the
statement is vacuous anyway. What the theorem does is remove the class of
objects one would reach for first, and it removes it for a reason which is
visible before any cohomology is computed.
\end{remark}

\begin{proof}
Immediate from \Cref{thm:positivity}, since the conclusion there contradicts
the requirement that $\ch_{n}(E)$ have a nonzero component on the Weil line.
\end{proof}

\begin{remark}[What is left of the candidate, and what it must look like]
\label{rem:whatitmustbe}
\Cref{thm:positivity} uses positivity of the multiplicities and nothing else,
so it says exactly where to go next: the object must be a complex whose
class in $K_{0}$ is not effective, that is one in which the line bundles
appear with multiplicities of both signs. Such objects exist. Writing
$E=F\oplus G[1]$ with $F$ and $G$ sums of line bundles, the class in $K_{0}$
is $[F]-[G]$ and \eqref{eq:normsum} becomes
$\sum_{L_{i}\subset F}\Nm(u_{i})=\sum_{L_{j}\subset G}\Nm(u_{j})$, which is
satisfiable; and \Cref{lem:splitsemireg} still applies verbatim, with the
shifts absorbing condition (ii).

Two further facts, both computational, say how little room is left at $n=2$,
and how much more there is beyond it. First, the map $\Phi$ of
\Cref{lem:splitsemireg}(iii) is injective at $n=2$ for $s\le3$ and, for a
general choice of $\lambda_{1},\dots,\lambda_{4}$, has rank $22$ rather than
$24$; so at $n=2$ no more than three line bundles may occur at all. This is
special to $n=2$: at $n=3$ the same map is injective for a general choice of
up to ten classes, and the dimension count permits up to $33$. Second, an
exhaustive search over the classes $a\beta+b\widehat\beta+c\ell$ with
$|a|,|b|,|c|\le3$, for $d=1$ and $d=2$, finds no signed tuple of at most
three line bundles whose Chern character is a flat section of Hodge classes
with a nonzero Weil component; the tuples of at most four that do satisfy the
Chern character condition with a nonzero Weil component, $888$ of them at
$d=1$ and $30$ at $d=2$, all have exactly four terms, and all of them fail
\Cref{lem:splitsemireg}(iii): the kernel has dimension $2$, $3$, $4$ or $6$
at $d=1$ and $3$ or $6$ at $d=2$, never $0$. Item (XIII) of
\Cref{sec:verification} carries both computations.

So the object of \Cref{thm:semiregclosure}, if it exists at a split member, is
not a sum of line bundles in any arrangement of signs and shifts: it must have
$\Ext^{2}$ that the trace does not see, that is, it must be indecomposable in
the derived category. That is a sharper description of the remaining
construction than \Cref{rem:tworemain} gives, and it is as far as we have been
able to push it.
\end{remark}

\subsection{The object, written down, and what it costs}
\label{ssec:pte}

\Cref{prop:ktheory} produced the object of \Cref{thm:semiregclosure} by an
existence argument: the Chern character is an isomorphism onto the Chow ring
with rational coefficients, so a $K$-theory class with the right character
exists. That is enough for the theorem and tells one nothing about the object.
This subsection writes it down. The answer turns the question into a
Diophantine one, and the Diophantine problem has a classical name.

\begin{setupenv}[Objects supported on $R$]\label{setup:pte}
Let $A$ be of $(K,-1,n)$-Weil type, $n\ge2$, satisfying
\eqref{eq:noextranorm}, and keep $\theta^{\pm}$ from
\Cref{not:splitcoords}. For $u\in K$ with $\theta(u)=u\theta^{+}+\bbar
u\,\theta^{-}\in\NS(A)_{\QQ}$, and for multiplicities $m_{i}\in\ZZ$, consider
\[
  v=\sum_{i=1}^{s}m_{i}\,e^{\theta(u_{i})},
  \qquad
  T_{q,r}=\sum_{i=1}^{s}m_{i}\,u_{i}^{q}\,\bbar u_{i}^{\,r}
  \quad (q,r\ge0).
\]
\end{setupenv}

\begin{proposition}[The Chern character in terms of the $u_{i}$]
\label{prop:chintermsofu}
In \Cref{setup:pte},
\begin{equation}\label{eq:chexpansion}
  v_{k}
  =\sum_{\substack{q+r=k\\ 0\le q,r\le n}}
    \frac{T_{q,r}}{q!\,r!}\,(\theta^{+})^{q}(\theta^{-})^{r},
\end{equation}
because $(\theta^{\pm})^{q}=0$ for $q>n$. Consequently $v$ is a flat section
of Hodge classes with a nonzero Weil component if
\begin{equation}\label{eq:pteconditions}
  T_{q,r}=0\ \ \text{for all } 0\le q,r\le n \text{ with }
  (q,r)\notin\{(0,0),(n,0),(0,n)\},
  \qquad T_{n,0}\ne0,
\end{equation}
and then $v_{0}=T_{0,0}=\sum_{i}m_{i}$, $v_{k}=0$ for $0<k<2n$ with $k\ne n$,
and
\[
  v_{n}=\frac{1}{n!}\Bigl(T_{n,0}(\theta^{+})^{n}
        +\bbar{T}_{n,0}(\theta^{-})^{n}\Bigr),
\]
a nonzero rational multiple of a generator of the Weil line, by
\Cref{thm:splitclosed}.
\end{proposition}

\begin{proof}
Expand $\theta(u_{i})^{k}$ by the binomial theorem, divide by $k!$ and collect;
the vanishing of $(\theta^{\pm})^{q}$ for $q>n$ is because
$\theta^{\pm}\in\bigwedge^{2}V_{\pm}$ and $\dim_{\CC}V_{\pm}=2n$. Under
\eqref{eq:pteconditions} the only surviving terms of \eqref{eq:chexpansion}
are the ones listed, and $(\theta^{+})^{n}$ spans $\bigwedge^{2n}V_{+}$, which
is the Weil line over $\CC$.
\end{proof}

So the search for the object of \Cref{thm:semiregclosure} is the search for a
finitely supported integer-valued function on $\cO_{K}$ that annihilates every
monomial $u^{q}\bbar u^{\,r}$ with $q,r\le n$ except three of them. Two things
follow at once, and they pull in opposite directions.

\begin{lemma}[Equal norms force the Prouhet--Tarry--Escott condition]
\label{lem:pte}
Suppose all the $u_{i}$ in \Cref{setup:pte} have the same norm
$N=\Nm_{K/\QQ}(u_{i})>0$. Then \eqref{eq:pteconditions} holds if and only if
\begin{equation}\label{eq:ptepure}
  \sum_{i}m_{i}u_{i}^{\,k}=0\ \ \text{for } k=0,1,\dots,n-1,
  \qquad
  \sum_{i}m_{i}u_{i}^{\,n}\ne0 .
\end{equation}
With $m_{i}=\pm1$ this says that the $u_{i}$ split into two multisets on the
circle $\Nm=N$ whose power sums agree through degree $n-1$ and differ at
degree $n$: a Prouhet--Tarry--Escott pair of degree $n-1$, constrained to a
norm circle of $\cO_{K}$.
\end{lemma}

\begin{proof}
For $q\ge r$ one has $u^{q}\bbar u^{\,r}=N^{r}u^{q-r}$, so
$T_{q,r}=N^{r}\sum_{i}m_{i}u_{i}^{\,q-r}$, and the pairs $(q,r)$ with
$q,r\le n$ realise every difference $k=q-r$ from $0$ to $n$. The excluded
pairs $(0,0)$ and $(n,0)$ are the ones with $r=0$ and $k=0$ or $k=n$; but
$k=0$ also occurs at $(1,1),\dots,(n,n)$, which are not excluded, so
$\sum_{i}m_{i}=0$ is required after all. The remaining differences
$k=1,\dots,n-1$ each occur at some non-excluded pair, and $k=n$ occurs only at
$(n,0)$.
\end{proof}

\begin{theorem}[Equal norms cannot be semiregular]\label{thm:equalnorms}
Let $E=\bigoplus_{i}L_{i}[p_{i}]$ realise \eqref{eq:ptepure} with all norms
equal. Then $\ch_{0}(E)=0$, and the semiregularity map $\sigma_{E}$ has a
kernel of dimension at least $n^{2}$.
\end{theorem}

\begin{proof}
By \Cref{lem:pte} the case $k=0$ of \eqref{eq:ptepure} gives
$\sum_{i}m_{i}=0$, which is $\ch_{0}(E)$. For the kernel, take
$\zeta_{i}=m_{i}\zeta$ for a single $\zeta\in H^{2}(\cO_{A})$; then by
\Cref{lem:splitsemireg}(iii)
\[
  \Phi\bigl((m_{i}\zeta)_{i}\bigr)
  =\Bigl(\sum_{i}m_{i}e^{\lambda_{i}}\Bigr)\zeta=\ch(E)\,\zeta ,
\]
which by \Cref{prop:chintermsofu} is a multiple of $\omega\,\zeta$ with
$\omega$ a generator of the Weil line. Now $\HW(A)_{\CC}$ is spanned by
$\bigwedge^{2n}V_{+}$ and $\bigwedge^{2n}V_{-}$, and the Weil type condition
says $\dim V_{\pm}^{0,1}=n$; writing
$H^{2}(\cO_{A})=\bigwedge^{2}(V_{+}^{0,1}\oplus V_{-}^{0,1})$, the product of
a generator of $\bigwedge^{2n}V_{+}$ with $\zeta$ kills the components of
$\zeta$ in $\bigwedge^{2}V_{+}^{0,1}$ and in
$V_{+}^{0,1}\otimes V_{-}^{0,1}$, and likewise with the roles exchanged. So
$\omega\,\zeta=0$ for every $\zeta$ in
$V_{+}^{0,1}\otimes V_{-}^{0,1}$, a space of dimension $n^{2}$, and all of
those lie in the kernel of $\Phi$.
\end{proof}

\begin{corollary}[At least $n+1$ norms]\label{cor:nplusone}
If $E$ as in \Cref{setup:pte} satisfies \eqref{eq:pteconditions} and
$\ch_{0}(E)\ne0$, then the $u_{i}$ take at least $n+1$ distinct norms.
\end{corollary}

\begin{proof}
Let $N_{1},\dots,N_{\nu}$ be the distinct norms and $M_{j}$ the sum of the
$m_{i}$ over the level $\Nm(u_{i})=N_{j}$. The conditions $T_{r,r}=0$ for
$r=1,\dots,n$ read $\sum_{j}M_{j}N_{j}^{\,r}=0$ for $r=1,\dots,n$. The matrix
$(N_{j}^{\,r})_{1\le r\le n,\,1\le j\le\nu}$ has rank $\min(n,\nu)$ because the
$N_{j}$ are distinct and nonzero, so $\nu\le n$ forces every $M_{j}=0$ and
hence $\ch_{0}(E)=\sum_{j}M_{j}=0$.
\end{proof}

\subsubsection*{The smallest solution, and what it does}

At $n=3$ and $K=\QQ(\sqrt{-3})$ the Prouhet--Tarry--Escott condition of
\Cref{lem:pte} has a solution of striking simplicity: the norm circle
$\Nm=1$ carries the six roots of unity of $K$, and the two cosets of the cube
roots of $-1$ and of $+1$ are a Prouhet--Tarry--Escott pair of degree
two.\footnote{Explicitly, with $\zeta$ a primitive sixth root of unity, the
cube roots of $-1$ are $\zeta,\zeta^{3},\zeta^{5}$ and those of $+1$ are
$1,\zeta^{2},\zeta^{4}$; each triple sums to zero and has zero sum of squares,
while the sums of cubes are $-3$ and $+3$. A search over the norm circles of
$\cO_{K}$ with $\Nm\le200$ and $|g|,|h|\le25$, for $K=\QQ(\sqrt{-d})$ with
$d=1,2,3,7$, finds this solution and its conjugate at $n=3$, one further pair
at $d=1$ and $n=4$, and nothing else; item (XIII) of \Cref{sec:verification}
records the search.}

\begin{proposition}[The object at $n=3$, written out]\label{prop:explicitobject}
Let $K=\QQ(\sqrt{-3})$, $n=3$, and let $A_{0}=X\times\widehat X$ be a split
member. In the basis $\beta,\widehat\beta,\ell$ of $\NS(A_{0})_{\QQ}$ put
\[
  \lambda_{1}=(3,-1,3),\quad
  \lambda_{2}=(-6,2,0),\quad
  \lambda_{3}=(3,-1,-3),
\]
and $\lambda_{3+i}=-\lambda_{i}$ for $i=1,2,3$, with $m_{i}=+1$ for
$i\le3$ and $m_{i}=-1$ for $i\ge4$. Then the six classes are distinct and
integral, the perfect complex
\[
  E=\bigl(L_{1}\oplus L_{2}\oplus L_{3}\bigr)\ \oplus\
    \bigl(L_{4}\oplus L_{5}\oplus L_{6}\bigr)[3]
\]
has
\[
  \ch_{k}(E)=0\ \ (k\ne3),
  \qquad
  \ch_{3}(E)=12\,\omega_{1},
\]
and every difference $\lambda_{j}-\lambda_{i}$ with $i\ne j$ is nondegenerate
of index exactly $3$, so that conditions (i) and (ii) of
\Cref{lem:splitsemireg} hold for this shift. Condition (iii) fails: the map
$\Phi$ has rank $75$ instead of $90$, with kernel of dimension
$15=\dim H^{2}(\cO_{A_{0}})$.
\end{proposition}

\begin{proof}
Everything asserted is a finite exact computation, carried out in item (XIII)
of \Cref{sec:verification}. The classes are the images under
$\theta(u)=2g\gamma+2dh\,\ell$, $u=g+h\sqrt{-3}$, of the six roots of unity
of $K$ scaled by $2$; the Chern character is \Cref{prop:chintermsofu} with the
Prouhet--Tarry--Escott data above; the index computation and the rank of
$\Phi$ are the two remaining items. For the shift, condition (ii) asks that
$2+p_{i}-p_{j}\ne3$ whenever $i\ne j$, that is $p_{i}-p_{j}\ne1$; within each
of the two blocks the difference is $0$, and between them it is $\pm3$.
\end{proof}

\begin{remark}[What this settles and what it does not]\label{rem:pteoutcome}
Three things are settled. The object of \Cref{prop:ktheory} is not an
abstraction: at $n=3$ over $\QQ(\sqrt{-3})$ it is six line bundles and a
shift, and its Chern character is exactly $12$ times a Weil class with every
other graded piece zero. The construction is completely general, and
\Cref{lem:pte} identifies the constraint as a Prouhet--Tarry--Escott problem
on a norm circle, which is a question of a kind that can be attacked. And
\Cref{cor:nplusone} says what any semiregular solution must look like: the
$u_{i}$ must take at least $n+1$ distinct norms, because otherwise the rank is
zero and \Cref{thm:equalnorms} produces a kernel of dimension $n^{2}$.

The remaining question was the existence of a solution with $n+1$ norms. The
level sums are then forced up to scale, and they compete with the supply of
elements of $\cO_{K}$ at each level; \Cref{thm:pterange} carries out that
competition in a range, and \Cref{cor:noline} then shows that no solution in
any range could have served, because a line bundle summand off the
polarisation line is obstructed along the family whatever the norms.
\end{remark}

\subsubsection*{The obstruction classes of the explicit object}

\Cref{cor:kernelfamily} and \Cref{rem:whatmeasures} say that the kernel of
$\sigma_{E}$ contains a copy of the tangent space of the Weil family, and
that this is a statement about where the semiregularity map is blind, not a
statement that the object is obstructed. For the explicit object the
obstruction classes can be written down, and they are not zero. The
computation is elementary and we record it in general.

Recall the conventions of \Cref{ssec:tangent}: a first-order deformation of
$A$ is $v\in\Hom(H^{1,0},H^{0,1})=H^{1}(A,T_{A})$, extended by zero on
$H^{0,1}$ and acting on $H^{*}(A,\CC)$ as a derivation $v\lrcorner$; a class
$\xi$ of type $(p,p)$ stays of type $(p,p)$ to first order along $v$ exactly
when $v\lrcorner\,\xi=0$. For a line bundle $L$ with $\lambda=c_{1}(L)$ the
Atiyah class is $\lambda\in H^{1}(\Omega^{1}_{A})$, its contraction with $v$
is $v\lrcorner\,\lambda\in H^{2}(\cO_{A})=H^{0,2}(A)$, and this is the
obstruction to deforming $L$ along $v$; a shift does not change it. For
$E=\bigoplus_{i}L_{i}[p_{i}]$ the Atiyah class is diagonal with entries
$\lambda_{i}$, as in the proof of \Cref{lem:splitsemireg}, so the
obstruction to deforming $E$ along $v$ is the diagonal element
\begin{equation}\label{eq:obdef}
  \ob_{E}(v)=\bigl(v\lrcorner\,\lambda_{i}\bigr)_{i}
  \in\bigoplus_{i}H^{0,2}(A)\subset\Ext^{2}(E,E).
\end{equation}

\begin{proposition}[The obstruction map of a split object]
\label{prop:splitobstruction}
Let $A_{0}=X\times\widehat X$ be a split member with $\beta$ principal and
$X$ general, so that $\NS(A_{0})_{\QQ}=\langle\beta,\widehat\beta,\ell\rangle$
by \Cref{prop:splitgeom}(iv). Let $T=T_{[A_{0}]}\cD_{n,n}$ be the space of
polarised first-order deformations commuting with $K$, of dimension $n^{2}$
by \Cref{prop:firstorder}, and put
\[
  T_{\mathrm{spl}}
  =\bigl\{\,v\in\Hom(H^{1,0},H^{0,1})\;:\;
    v\lrcorner\,\beta=v\lrcorner\,\widehat\beta=v\lrcorner\,\ell=0\,\bigr\}.
\]
Let $E=\bigoplus_{i=1}^{s}L_{i}[p_{i}]$ with $\lambda_{i}=c_{1}(L_{i})$ and
let $\ob_{E}\colon T\to\bigoplus_{i}H^{0,2}(A_{0})$ be the map
\eqref{eq:obdef} restricted to $T$.
\begin{enumerate}[label=\textup{(\roman*)},nosep]
\item $T_{\mathrm{spl}}$ is the tangent space to the split locus: it has
  dimension $n(n+1)/2$, it is contained in $T$, and it consists of the
  deformations of $X\times\widehat X$ induced by the polarised deformations
  of $(X,\beta)$.
\item If the classes $\lambda_{i}$ together with $\eta$ span
  $\NS(A_{0})_{\QQ}$, then $\ker\ob_{E}=T_{\mathrm{spl}}$ and
  $\operatorname{rank}\ob_{E}=n(n-1)/2$, the codimension of the split locus
  in $\cD_{n,n}$.
\item If $\ch_{k}(E)\in\QQ\eta^{k}$ for $k\ne n$ and
  $\ch_{n}(E)\in\QQ\eta^{n}\oplus\HW(A_{0})$, then
  $\sigma_{E}\circ\ob_{E}=0$ on $T$: for every $v\in T$,
  \begin{equation}\label{eq:sigmaob}
    \sigma_{E}\bigl(\ob_{E}(v)\bigr)=v\lrcorner\,\ch(E)=0 .
  \end{equation}
\item For the object of \Cref{prop:explicitobject}, at $n=3$ and
  $K=\QQ(\sqrt{-3})$: $\ker\ob_{E}=T_{\mathrm{spl}}$ has dimension $6$ and
  $\ob_{E}$ has rank $3$; the image of $\ob_{E}$ lies in $\ker\sigma_{E}$;
  the subspace $\Delta=\{z\cdot\id_{E}:z\in\Ann_{H^{0,2}}(\omega_{1})\}$ of
  \Cref{cor:kernelfamily} has dimension $9$ and meets the image of $\ob_{E}$
  only in zero, so that $\Delta\oplus\image\ob_{E}$ is a twelve-dimensional
  subspace of the fifteen-dimensional kernel of $\sigma_{E}$; and
  $\ob_{E}(v)$ has components $v\lrcorner\,\lambda_{i}$ on $L_{i}$ and
  $-v\lrcorner\,\lambda_{i}$ on $L_{i}^{-1}[3]$, so it lies on the line of
  vectors of the shape $(z,\dots,z)$ only when it vanishes. In particular
  $E$ has a nonzero first-order obstruction along every $v\in T$ not tangent
  to the split locus, and no element $z\cdot\id_{E}$ of $\Ext^{2}(E,E)$
  cancels it.
\end{enumerate}
\end{proposition}

\begin{proof}
(i) Let $P=H^{1,0}(X)\subset H^{1}(X,\CC)=U^{\vee}_{\CC}$ with basis
$p_{1},\dots,p_{n}$, so that $p_{1},\dots,p_{n},\bbar p_{1},\dots,\bbar p_{n}$
is a basis of $U^{\vee}_{\CC}$, and let $q_{a},\bbar q_{a}$ be the dual basis
of $U_{\CC}=H^{1}(\widehat X,\CC)$; the dual of a conjugate basis is the
conjugate of the dual basis, so the notation is consistent. The Poincar\'e
class is the identity element of $U^{\vee}\otimes U$,
\[
  \ell=\sum_{a=1}^{n}\bigl(p_{a}\wedge q_{a}+\bbar p_{a}\wedge\bbar q_{a}\bigr),
\]
and it is of type $(1,1)$, which forces $H^{1,0}(\widehat X)=\langle\bbar
q_{a}\rangle$ and $H^{0,1}(\widehat X)=\langle q_{a}\rangle$: the Hodge
structure of $\widehat X$ is the dual one. Since $\beta$ is of type $(1,1)$
and nondegenerate, $\beta=\sum_{a,b}h_{ab}\,p_{a}\wedge\bbar p_{b}$ with
$h\in\GL_{n}(\CC)$; and since $\beta$ is principal with Gram matrix $B$ in a
symplectic basis of $U$ and $\widehat\beta$ has Gram matrix $B=-B^{-1}$ in
the dual basis of $U^{\vee}$, the Gram matrix of $\widehat\beta$ in the
basis $p_{a},\bbar p_{a}$ is minus the inverse of the Gram matrix
$\begin{psmallmatrix}0&h\\-h^{\top}&0\end{psmallmatrix}$ of $\beta$ in the
basis $q_{a},\bbar q_{a}$, which gives
$\widehat\beta=-\sum_{a,b}(h^{-1})_{ab}\,\bbar q_{a}\wedge q_{b}$.

Now $H^{1,0}(A_{0})=\langle p_{a}\rangle\oplus\langle\bbar q_{a}\rangle$ and
$H^{0,1}(A_{0})=\langle\bbar p_{a}\rangle\oplus\langle q_{a}\rangle$, so a
deformation $v$ is four $n\times n$ matrices,
\[
  v(p_{a})=\sum_{c}\bigl(A_{ca}\,\bbar p_{c}+C_{ca}\,q_{c}\bigr),
  \qquad
  v(\bbar q_{a})=\sum_{c}\bigl(B_{ca}\,\bbar p_{c}+D_{ca}\,q_{c}\bigr),
\]
and $v$ kills $\bbar p_{a}$ and $q_{a}$. Contracting,
\[
  v\lrcorner\,\ell
  =\sum_{a,c}\Bigl(A_{ca}\,\bbar p_{c}\wedge q_{a}+C_{ca}\,q_{c}\wedge q_{a}
   +B_{ca}\,\bbar p_{a}\wedge\bbar p_{c}+D_{ca}\,\bbar p_{a}\wedge q_{c}\Bigr),
\]
and reading off the three kinds of monomials, $v\lrcorner\,\ell=0$ if and only
if $D=-A^{\top}$ and $B$, $C$ are symmetric. Next
\[
  v\lrcorner\,\beta
  =\sum_{a,b,c}h_{ab}\Bigl(A_{ca}\,\bbar p_{c}\wedge\bbar p_{b}
   +C_{ca}\,q_{c}\wedge\bbar p_{b}\Bigr),
\]
so $v\lrcorner\,\beta=0$ if and only if $Ch=0$, that is $C=0$, and $Ah$ is
symmetric. Likewise $v\lrcorner\,\widehat\beta=0$ if and only if $B=0$ and
$Dh^{-1}$ is symmetric. Given $D=-A^{\top}$, the last condition reads
$A^{\top}h^{-1}=h^{-\top}A$, which after multiplication by $h^{\top}$ on the
left and by $h$ on the right is $h^{\top}A^{\top}=Ah$, the symmetry of $Ah$
again. Hence
\[
  T_{\mathrm{spl}}=\bigl\{(A,B,C,D)=(A,0,0,-A^{\top})\;:\;Ah\ \text{symmetric}\bigr\},
\]
of dimension $n(n+1)/2$. It is the tangent space to the split locus: the
same computation on $X$ alone shows that the polarised first-order
deformations of $(X,\beta)$ are the matrices $A$ with $Ah$ symmetric, the
space $T_{[X]}\cA_{n}$ of dimension $n(n+1)/2$; the deformation of
$\widehat X$ induced by $A$ is the dual one, $-A^{\top}$; and the map
$A\mapsto(A,0,0,-A^{\top})$ is injective. Finally $T_{\mathrm{spl}}\subseteq
T$: for $v\in T_{\mathrm{spl}}$, $v\lrcorner\,\eta=0$ because
$\eta\in\langle\beta,\widehat\beta\rangle$, so $v$ is polarised, and
$v\lrcorner\,\omega_{1}=v\lrcorner\,\omega_{2}=0$ because $\omega_{1}$ and
$\omega_{2}$ are polynomials in $\beta,\widehat\beta,\ell$ by
\Cref{thm:splitclosed} and $v\lrcorner$ is a derivation; by
\Cref{prop:firstorder} such a $v$ commutes with $K$.

(ii) Every $v\in T$ has $v\lrcorner\,\eta=0$. If the $\lambda_{i}$ and
$\eta$ span $\NS(A_{0})_{\QQ}=\langle\beta,\widehat\beta,\ell\rangle$, then
$v\lrcorner\,\lambda_{i}=0$ for all $i$ is equivalent, for $v\in T$, to
$v\lrcorner\,\beta=v\lrcorner\,\widehat\beta=v\lrcorner\,\ell=0$, so
$\ker\ob_{E}=T\cap T_{\mathrm{spl}}=T_{\mathrm{spl}}$ by (i), and the rank
is $n^{2}-n(n+1)/2=n(n-1)/2$, which is the codimension of
\Cref{prop:splitgeom}(iii).

(iii) By \eqref{eq:splitsigma}, $\sigma_{E}(\zeta)=\sum_{i}(-1)^{p_{i}}
\zeta_{ii}\,e^{\lambda_{i}}$ for a diagonal $\zeta$. Since $v\lrcorner$ is a
derivation and $v\lrcorner\,\lambda_{i}$ has even degree,
$v\lrcorner\,\lambda_{i}^{\,q}=q\,\lambda_{i}^{\,q-1}(v\lrcorner\,\lambda_{i})$
and so $v\lrcorner\,e^{\lambda_{i}}=(v\lrcorner\,\lambda_{i})\,e^{\lambda_{i}}$.
Therefore
\[
  \sigma_{E}\bigl(\ob_{E}(v)\bigr)
  =\sum_{i}(-1)^{p_{i}}(v\lrcorner\,\lambda_{i})\,e^{\lambda_{i}}
  =v\lrcorner\sum_{i}(-1)^{p_{i}}e^{\lambda_{i}}
  =v\lrcorner\,\ch(E).
\]
For $v\in T$ one has $v\lrcorner\,\eta^{k}=k\eta^{k-1}(v\lrcorner\,\eta)=0$,
and $v\lrcorner\,\omega=0$ for every $\omega\in\HW(A_{0})$ by
\Cref{prop:firstorder}; under the hypothesis on $\ch(E)$ this gives
\eqref{eq:sigmaob}.

(iv) The three classes $\lambda_{1}=\gamma+3\ell$, $\lambda_{2}=-2\gamma$,
$\lambda_{3}=\gamma-3\ell$ span $\langle\gamma,\ell\rangle$, and
$\langle\eta,\gamma\rangle=\langle\beta,\widehat\beta\rangle$ by the
determinant in the proof of \Cref{lem:sixmonomials}, so (ii) applies and
gives the kernel and the rank; \Cref{prop:explicitobject} gives the
hypothesis of (iii), so the image lies in $\ker\sigma_{E}$. The subspace
$\Delta$ has dimension $\dim\Ann_{H^{0,2}}(\omega_{1})=n^{2}=9$ by
\Cref{lem:annihilator}. The components of $\ob_{E}(v)$ on $L_{i}$ and on
$L_{i}^{-1}[3]$ are $v\lrcorner\,\lambda_{i}$ and
$v\lrcorner\,(-\lambda_{i})=-v\lrcorner\,\lambda_{i}$, so an element of the
image of the shape $(z,\dots,z)$ has $z=-z$, hence is zero; this is
$\Delta\cap\image\ob_{E}=0$, and with $\dim\ker\sigma_{E}=15$ from
\Cref{prop:explicitobject} the dimension count follows. Every number in (iv)
is also computed directly in item (XXV) of \Cref{sec:verification}, in the
model of item (XIII).
\end{proof}

\begin{remark}[What the obstruction classes say]\label{rem:obstructionmeaning}
\Cref{prop:splitobstruction} sharpens \Cref{rem:whatmeasures} in one
direction and closes a door in another. The direction: for the explicit
object the obstruction classes along the family are nonzero, and they are
nonzero exactly in the $n(n-1)/2$ directions normal to the split locus,
which are the directions in which the Weil class has to be transported.
The door: those classes lie in the kernel of $\sigma_{E}$ by
\eqref{eq:sigmaob}, and \eqref{eq:sigmaob} holds for every object whose
Chern character is a flat section of Hodge classes along the family, so the
semiregularity map cannot detect the obstruction of any such object. The
obstruction is therefore genuine, and it is invisible to the criterion. For
the explicit object the situation is the worst possible: the twelve
dimensions $\Delta\oplus\image\ob_{E}$ of the kernel are accounted for by
the two mechanisms of \Cref{cor:kernelfamily} and \eqref{eq:obdef}, the
remaining three lie on the even side of \Cref{cor:pteobstructed} as well,
and none of the fifteen is an artefact. The classes of the shape
$z\cdot\id_{E}$ are the ones a change of the twisting class would
contribute, one and the same $z$ to every rank one summand; the last
assertion of (iv) says that no such change removes the obstruction. So the
explicit object does not move off the split locus to first order, as a
complex and as a twisted complex, and the question of \Cref{q:residual} for
it is not one of verifying a hypothesis but of replacing the object.

The positive content is the recipe. By (ii), a split object whose classes
span $\NS(A_{0})_{\QQ}$ with $\eta$ has obstruction map of rank
$n(n-1)/2$ against the family, and the spanning hypothesis is more than is
needed: in the model of item (XXV) a single class $\theta\in\NS(A_{0})_{\QQ}$
not proportional to $\eta$ already has $\{v\in T:v\lrcorner\,\theta=0\}=
T_{\mathrm{spl}}$, for each of the seven classes tested at
$(n,d)=(2,1),(2,3),(3,3)$. So every split object built from line bundles
other than powers of $\eta$ is obstructed in exactly the normal directions
to the split locus, and \Cref{thm:positivity} says that an object built
from powers of $\eta$ carries no Weil class. What survives is the statement
of \Cref{rem:familyuniform}: the object must have positive rank.
\end{remark}

\subsubsection*{The evaluation map of the explicit object}

\Cref{thm:weakfails} places the copy $\Delta$ of the tangent space inside the
image of the evaluation map. For a split object the whole of
$\operatorname{ev}_{E}$ can be written down, and for the explicit object the
entire kernel of $\sigma_{E}$ turns out to lie in its image. We identify
$HH^{2}(A_{0})$ with
\begin{equation}\label{eq:HT2}
  HT^{2}(A_{0})
  =H^{0,2}(A_{0})\;\oplus\;\Hom\bigl(H^{1,0},H^{0,1}\bigr)\;\oplus\;
   \textstyle\bigwedge^{2}\bigl(H^{1,0}\bigr)^{\vee}
\end{equation}
by the Hochschild--Kostant--Rosenberg isomorphism, the three summands being
$H^{2}(\cO)$, $H^{1}(T)$ and $H^{0}(\bigwedge^{2}T)$ for the trivial tangent
bundle. An element $x=(z,v,\pi)$ acts on $H^{*}(A_{0},\CC)$ by $z\wedge$, by
the derivation $v\lrcorner$ of \Cref{ssec:tangent}, and by the double
contraction $\pi\lrcorner$, which removes two $(1,0)$ slots; we write
$x\lrcorner\,\xi$ for the sum of the three.%
\footnote{The double contraction is a second order operator. For a class
$\lambda$ of type $(1,1)$, which has a single $(1,0)$ slot,
$\pi\lrcorner\,\lambda=0$ and $\pi\lrcorner\,\lambda^{k}=\binom{k}{2}
\lambda^{k-2}\,(\pi\lrcorner\,\lambda^{2})$, because $\pi\lrcorner$ is a
composite of two interior products, each of which is an antiderivation, and
the only nonzero terms are those in which the two interior products hit two
different factors. Summing over $k$ with the factorials gives
$\pi\lrcorner\,e^{\lambda}=\tfrac12(\pi\lrcorner\,\lambda^{2})\,e^{\lambda}$,
exactly parallel to $v\lrcorner\,e^{\lambda}=(v\lrcorner\,\lambda)\,e^{\lambda}$
for the first order operator $v\lrcorner$. This is the whole content of the
compatibility \eqref{eq:sigmaev} below: the semiregularity map of a line
bundle is multiplication by $e^{\lambda}$, and every summand of $HT^{2}$ acts
on $e^{\lambda}$ by a scalar-valued factor which is precisely the value of
$\operatorname{ev}_{L}$. The same identity, for arbitrary $E$, is the
commutative square of Buchweitz and Flenner \cite[Corollary 4.3]{BF03} and
its Hochschild form in \cite[Lemma 11.3]{Mar25b}; on an abelian variety the
Todd class is trivial and no correction to the HKR map is needed
\cite{Cal05,BF08}.}

\begin{proposition}[The evaluation map of a split object]\label{prop:evaluation}
Let $A_{0}$ be a split member with $X$ general and let
$E=\bigoplus_{i=1}^{s}L_{i}[p_{i}]$ with $\lambda_{i}=c_{1}(L_{i})$ and all
$\Ext^{2}(L_{i}[p_{i}],L_{j}[p_{j}])$, $i\ne j$, zero, so that
$\Ext^{2}(E,E)=\bigoplus_{i}H^{0,2}(A_{0})$.
\begin{enumerate}[label=\textup{(\roman*)},nosep]
\item Under \eqref{eq:HT2}, $\operatorname{ev}_{E}$ is diagonal:
  \begin{equation}\label{eq:evsplit}
    \operatorname{ev}_{E}(z,v,\pi)_{i}
    =z+v\lrcorner\,\lambda_{i}+\tfrac12\,\pi\lrcorner\,\lambda_{i}^{2}
    \;\in\;H^{0,2}(A_{0})=\Ext^{2}\bigl(L_{i}[p_{i}],L_{i}[p_{i}]\bigr).
  \end{equation}
\item For every $x\in HT^{2}(A_{0})$,
  \begin{equation}\label{eq:sigmaev}
    \sigma_{E}\bigl(\operatorname{ev}_{E}(x)\bigr)=x\lrcorner\,\ch(E),
  \end{equation}
  and consequently
  $\ker\sigma_{E}\cap\image\operatorname{ev}_{E}
   =\operatorname{ev}_{E}\bigl(\{x:x\lrcorner\,\ch(E)=0\}\bigr)$.
  The weakened hypothesis of \cite[Question 11.4]{Mar25b} holds for $E$ if
  and only if $\operatorname{ev}_{E}$ kills every $x$ that preserves
  $\ch(E)$ to first order.
\item $\dim HT^{2}(A_{0})=2\binom{2n}{2}+4n^{2}$ and
  $\dim\Ext^{2}(E,E)=s\binom{2n}{2}$; for $s\ge5$ and every $n\ge2$ the
  second exceeds the first, so $\operatorname{ev}_{E}$ is not surjective and
  \cite[Lemma 11.3]{Mar25b} cannot apply to a sum of five or more line
  bundles.
\item For the object of \Cref{prop:explicitobject}, with
  $\lambda_{i}=u_{i}\theta^{+}+\bbar u_{i}\theta^{-}$ as in
  \Cref{not:splitcoords}: the space $\{x:x\lrcorner\,\ch(E)=0\}$ is
  $\Ann_{H^{0,2}}(\omega_{1})\oplus\{v:v\lrcorner\,\omega_{1}=0\}\oplus
  \{\pi:\pi\lrcorner\,\omega_{1}=0\}$, of dimensions $9$, $18$ and $9$; the
  first and the third summands are both carried by $\operatorname{ev}_{E}$
  onto $\Delta$; the second is carried onto
  \[
    \Sigma=\bigl\{\,(u_{i}a+\bbar u_{i}b)_{i=1}^{6}\;:\;
      a\in\textstyle\bigwedge^{2}V_{+}^{0,1},\ b\in\bigwedge^{2}V_{-}^{0,1}\,\bigr\},
    \qquad \dim\Sigma=n(n-1)=6,
  \]
  which contains the image of the obstruction map of
  \Cref{prop:splitobstruction}; $\Delta\cap\Sigma=0$; and
  \[
    \ker\sigma_{E}=\Delta\oplus\Sigma\subseteq\image\operatorname{ev}_{E}.
  \]
  So every direction in which $\sigma_{E}$ is blind is the value of
  $\operatorname{ev}_{E}$ on a Hochschild class that preserves $\ch(E)$ to
  first order, and the weakened hypothesis fails for $E$ by all fifteen
  dimensions. The rank of $\operatorname{ev}_{E}$ is $45$. The same
  containment $\ker\sigma_{E}\subseteq\image\operatorname{ev}_{E}$, with
  $\dim\ker\sigma_{E}=15$, holds for each of the eighteen configurations of
  \Cref{rem:familyuniform}.
\end{enumerate}
\end{proposition}

\begin{proof}
(i) The evaluation map is the action of $HH^{*}(A_{0})$ on
$\Ext^{*}(E,E)$, which under HKR is contraction of a polyvector in
$H^{q}(\bigwedge^{p}T)$ with $\tfrac{1}{p!}\mathrm{At}(E)^{p}$
\cite{BF08}; for $p=0$ this is $z\cdot\id_{E}$ as in the proof of
\Cref{thm:weakfails}, for $p=1$ it is $v\lrcorner\,\mathrm{At}(E)$, the map
\eqref{eq:obdef}, and for $p=2$ it is $\tfrac12\pi\lrcorner\,\mathrm{At}(E)^{2}$.
The Atiyah class of $E$ is diagonal with entries $\lambda_{i}$, so all three
are diagonal with the entries displayed.

(ii) By \eqref{eq:splitsigma}, $\sigma_{E}(\zeta)=\sum_{i}(-1)^{p_{i}}\zeta_{i}
e^{\lambda_{i}}$ for diagonal $\zeta$. The three factors in \eqref{eq:evsplit}
are exactly the scalars by which $z\wedge$, $v\lrcorner$ and $\pi\lrcorner$
act on $e^{\lambda_{i}}$: the first trivially, the second because
$v\lrcorner$ is a derivation and $v\lrcorner\,\lambda_{i}$ has even degree,
the third by the footnote above. Hence
$\sigma_{E}(\operatorname{ev}_{E}(x))=\sum_{i}(-1)^{p_{i}}\,x\lrcorner\,
e^{\lambda_{i}}=x\lrcorner\,\ch(E)$. The description of
$\ker\sigma_{E}\cap\image\operatorname{ev}_{E}$ follows: $\operatorname{ev}_{E}(x)$
lies in $\ker\sigma_{E}$ if and only if $x\lrcorner\,\ch(E)=0$.

(iii) The three summands of \eqref{eq:HT2} have dimensions $\binom{2n}{2}$,
$4n^{2}$ and $\binom{2n}{2}$. The inequality
$s\binom{2n}{2}>2\binom{2n}{2}+4n^{2}$ is $(s-2)(2n-1)>4n$, and
$4n/(2n-1)\le8/3$ for $n\ge2$, so it holds for $s\ge5$; it is checked for
$2\le n\le30$ and $5\le s\le40$ in \Cref{sec:lean}.

(iv) Since $\ch(E)=12\omega_{1}$ is of type $(3,3)$, the three summands of
$x\lrcorner\,\ch(E)$ have types $(3,5)$, $(2,4)$ and $(1,3)$, so the kernel
is the direct sum of the three kernels. The first is $\Ann_{H^{0,2}}(\omega_{1})$,
of dimension $n^{2}=9$ by \Cref{lem:annihilator}. The second is
$\Hom(V_{+}^{1,0},V_{+}^{0,1})\oplus\Hom(V_{-}^{1,0},V_{-}^{0,1})$, of
dimension $2n^{2}=18$, by the proof of \Cref{prop:firstorder}, whose
argument for (i)$\Leftrightarrow$(iii) does not use the polarisation. For the
third, decompose $\bigwedge^{2}(H^{1,0})^{\vee}$ as
$\bigwedge^{2}(V_{+}^{1,0})^{\vee}\oplus(V_{+}^{1,0})^{\vee}\otimes(V_{-}^{1,0})^{\vee}
\oplus\bigwedge^{2}(V_{-}^{1,0})^{\vee}$; here $\omega_{1}$ is a combination
of $(\theta^{+})^{3}$ and $(\theta^{-})^{3}$ with both coefficients nonzero,
$(\theta^{\pm})^{3}$ spans $\bigwedge^{3}V_{\pm}^{1,0}\otimes\bigwedge^{3}V_{\pm}^{0,1}$,
the outer blocks of $\pi$ act injectively on the respective factor
$\bigwedge^{3}V_{\pm}^{1,0}$ and land in different summands of $H^{1,3}$,
and the middle block kills both; so the third kernel is the middle block,
of dimension $n^{2}=9$.

At a split point $\theta^{\pm}$ are of type $(1,1)$, so $\theta^{\pm}\in
V_{\pm}^{1,0}\otimes V_{\pm}^{0,1}$, and they are nondegenerate pairings
because $(\theta^{\pm})^{3}\ne0$. Write $\lambda_{i}^{2}=u_{i}^{2}(\theta^{+})^{2}
+\bbar u_{i}^{2}(\theta^{-})^{2}+2\Nm(u_{i})\,\theta^{+}\theta^{-}$. A middle
block $\pi$ kills $(\theta^{\pm})^{2}$, which have no $V_{\mp}^{1,0}$ slot, so
$\tfrac12\pi\lrcorner\,\lambda_{i}^{2}=\Nm(u_{i})\,\pi\lrcorner\,(\theta^{+}\theta^{-})$;
the six $u_{i}$ of \Cref{prop:explicitobject} are $\pm\zeta_{6},\pm1,
\pm\zeta_{6}^{-1}$, all of norm one, so this value is independent of $i$ and
$\operatorname{ev}_{E}(\pi)=z_{\pi}\cdot\id_{E}$ with
$z_{\pi}=\pi\lrcorner\,(\theta^{+}\theta^{-})$. The map $\pi\mapsto z_{\pi}$
from the middle block to $V_{+}^{0,1}\otimes V_{-}^{0,1}$ is the tensor
product of the two nondegenerate pairings, hence an isomorphism onto the
middle block of $H^{0,2}$, which is $\Ann_{H^{0,2}}(\omega_{1})$ by
\Cref{lem:annihilator}. So $\operatorname{ev}_{E}$ carries the third kernel
onto $\Delta$, as it carries the first by \Cref{thm:weakfails}.

For $v=(v_{+},v_{-})$ in the second kernel,
$v\lrcorner\,\lambda_{i}=u_{i}\,v_{+}\lrcorner\,\theta^{+}+\bbar u_{i}\,
v_{-}\lrcorner\,\theta^{-}$ with $a=v_{+}\lrcorner\,\theta^{+}\in
\bigwedge^{2}V_{+}^{0,1}$ and $b=v_{-}\lrcorner\,\theta^{-}\in\bigwedge^{2}V_{-}^{0,1}$;
as $v_{\pm}$ ranges over $\Hom(V_{\pm}^{1,0},V_{\pm}^{0,1})$, $a$ and $b$
range over the whole of $\bigwedge^{2}V_{\pm}^{0,1}$ because the pairings
$\theta^{\pm}$ are nondegenerate, which is the surjectivity of
$\operatorname{alt}(Vc)$ in \eqref{eq:phidelta}. So the image is $\Sigma$.
The map $(a,b)\mapsto(u_{i}a+\bbar u_{i}b)_{i}$ is injective, because
$(u_{1},\bbar u_{1})$ and $(u_{2},\bbar u_{2})$ are linearly independent in
$\CC^{2}$, $u_{1}/\bbar u_{1}=\zeta_{6}^{2}$ differing from
$u_{2}/\bbar u_{2}=1$; hence $\dim\Sigma=2\binom{n}{2}=n(n-1)=6$. The
obstruction map of \Cref{prop:splitobstruction} is the restriction to the
polarised $K$-commuting $v$, so its image lies in $\Sigma$. An element of
$\Sigma$ of the shape $(z,\dots,z)$ has $u_{i}a+\bbar u_{i}b=z=-z$ on
comparing $i$ with $3+i$, so $z=0$ and then $a=b=0$: $\Delta\cap\Sigma=0$.
By (ii), $\Delta\oplus\Sigma\subseteq\ker\sigma_{E}\cap\image\operatorname{ev}_{E}$,
and its dimension $9+6=15$ equals $\dim\ker\sigma_{E}$ by
\Cref{prop:explicitobject}, so the three spaces coincide. The rank of
$\operatorname{ev}_{E}$ and the statement for the eighteen configurations
are the computations of item (XXVI) of \Cref{sec:verification}, which also
recomputes every dimension in this proof in the model of item (XIII).
\end{proof}

\begin{figure}[tbp]
\centering
  \includegraphics[width=\textwidth]{fig_evaluation}
\caption{\Cref{prop:evaluation} for the explicit object at $n=3$. Left: the
three summands of $HT^{2}(A_{0})$, drawn with heights proportional to their
dimensions, and inside each the part that preserves $\ch(E)$ to first order:
the annihilator of $\omega_{1}$, the deformations killing $\omega_{1}$ with
the polarised $K$-commuting tangent space $T$ in front, and the middle block
of bivectors. Right: $\Ext^{2}(E,E)$ as six stacked copies of $H^{0,2}$, one
per summand of $E$, and inside it the kernel of $\sigma_{E}$: the column
$\Delta$, one and the same class in every layer, and the column $\Sigma$,
whose entries depend on the summand through $u_{i}$. The three arrows are the
evaluation map on the three kernels; the two teal ones both land on
$\Delta$, the ochre one on $\Sigma$, and together they cover the kernel.
Below, the commutative square that makes this a statement about the weakened
criterion of \cite[Question 11.4]{Mar25b}.}
\label{fig:evaluation}
\end{figure}

\begin{remark}[Question 11.4 for split objects, and what is left of it]
\label{rem:question114}
For a split object the two criteria stand or fall together, and for the
explicit object both fall completely. \Cref{thm:weakfails} makes the first
$n^{2}$ dimensions of the failure uniform over every concentrated object in
every dimension; \Cref{prop:evaluation}(iv) accounts for the remaining
$n(n-1)$ at $n=3$ by the deformations of the two halves $V_{\pm}$ of the
Hodge structure separately, and the equality $\ker\sigma_{E}=\Delta\oplus\Sigma$
says that the two mechanisms exhaust the kernel: nothing in it is
accidental. Nothing here touches the case the question was asked for. A
sheaf of positive rank has $z\cdot\ch(E)\ne0$ for every $z\ne0$, so its
$\Delta$ is zero, and whether $\operatorname{ev}_{E}$ kills the deformations
that preserve $\kappa(E)$ is a statement about the sheaf and not about its
Chern character. Markman's candidate in his Section 12, a rank one complex on
a principally polarised abelian fourfold with Chern character on the secant
$\langle e^{\pm\sqrt{-d}\Theta}\rangle$, is exactly such a sheaf, and he
records that the verification remains to be done. \Cref{thm:factor} reduces
that verification to one identity on the fourfold, and
\Cref{thm:candidatefails} shows that the identity fails: the candidate does
not satisfy the weakened criterion either.
\end{remark}

\begin{theorem}[Where a semiregular split object can live]\label{thm:pterange}
Keep \Cref{setup:pte} and suppose $E$ satisfies \eqref{eq:pteconditions} with
$\ch_{0}(E)\ne0$, so that the $u_{i}$ take $\nu\ge n+1$ distinct norms by
\Cref{cor:nplusone}. Write $N_{1}<\dots<N_{\nu}$ for those norms, $M_{j}$ for
the sum of the $m_{i}$ at level $N_{j}$, and
$r_{K}(N)=\#\{u\in\cO_{K}:\Nm_{K/\QQ}(u)=N\}$.
\begin{enumerate}[label=\textup{(\roman*)},leftmargin=2.6em]
\item If $\nu=n+1$ then, with $P=\prod_{k}N_{k}$, $S=\ch_{0}(E)$ and
  $D_{j}=N_{j}\prod_{k\ne j}(N_{j}-N_{k})$,
  \begin{equation}\label{eq:levelsums}
    M_{j}=(-1)^{n}\,\frac{S\,P}{D_{j}},
  \end{equation}
  and the least $S>0$ making every $M_{j}$ an integer is
  $S_{\min}=\operatorname{lcm}_{j}\bigl(|D_{j}|/\gcd(|D_{j}|,P)\bigr)$.
  Multiplicity one, \Cref{lem:splitsemireg}(i), forces
  $|M_{j}|\le r_{K}(N_{j})$ for every $j$.
\item For $4\le n\le8$, every $d\in\{1,2,3,7,11,19\}$ and every choice of
  norms below the bounds recorded in item (XXII) of \Cref{sec:verification},
  no set of $n+1$ norms and no set of $n+2$ norms satisfies those constraints.
  In that range no split object of the shape of \Cref{setup:pte} is
  semiregular, and the reason is arithmetic rather than geometric: the least
  integral scale outruns the supply of elements at the levels it demands.
\item At $n=3$ the constraints of (i) can be met. With norms at most $200$
  there are exactly eleven sets of four norms that meet them, all of them at
  $d=1$ or $d=2$, and after dividing out the common factor each has one of the
  three shapes $(1,2,3,4)$, $(1,2,4,5)$ and $(1,4,9,10)$. The smallest is
  $d=2$ at $(9,18,27,36)$, where $S=1$, $|M|=(4,6,4,1)$ and $\ch_{0}(E)=1$.
\item For all eleven, the full system \eqref{eq:pteconditions} has no
  solution: no assignment of signs to the elements of $\cO_{K}$ at those
  norms, with the level sums of \eqref{eq:levelsums}, satisfies the
  off-diagonal conditions. In particular no split object of the shape of
  \Cref{setup:pte} with norms at most $200$ is semiregular at $n=3$.
\end{enumerate}
\end{theorem}

\begin{proof}
(i) The conditions $T_{r,r}=0$ for $r=1,\dots,n$ say that the linear
functional $f\mapsto\sum_{j}M_{j}f(N_{j})$ kills $x,x^{2},\dots,x^{n}$, hence
sends every polynomial $f$ of degree at most $n$ to $S\,f(0)$. Applying it to
$f(x)=\prod_{k\ne j}(x-N_{k})$, which has degree $n$, gives
\[
  M_{j}\prod_{k\ne j}(N_{j}-N_{k})=S\,(-1)^{n}\prod_{k\ne j}N_{k},
\]
and multiplying by $N_{j}$ turns this into \eqref{eq:levelsums}. The value of
$S_{\min}$ is the definition of a least common multiple, and the bound on
$|M_{j}|$ holds because each $u$ of norm $N_{j}$ contributes at most one term
$m_{i}=\pm1$ to $M_{j}$.

(ii), (iii) and (iv) are the computations of item (XXII) of
\Cref{sec:verification}, carried out in exact integer arithmetic. For (ii) and
(iii) the script enumerates every set of $n+1$ norms represented by $\cO_{K}$
below the bound, forms $D_{j}$ and $S_{\min}$, and tests
$|M_{j}|\le r_{K}(N_{j})$; at $\nu=n+2$ the solution space of the diagonal
conditions is two dimensional and its integer points inside the same box are
swept by running two coordinates over their own ranges. For (iv) the
off-diagonal conditions are, after grouping by norm,
\[
  \sum_{j}N_{j}^{r}A_{1,j}=0\ \ (r=0,1,2),
  \qquad
  \sum_{j}N_{j}^{r}A_{2,j}=0\ \ (r=0,1),
  \qquad
  A_{k,j}=\sum_{\Nm(u)=N_{j}}m_{u}u^{k},
\]
which are additive over the levels; the four levels are therefore split into
one half that is tabulated and one that is streamed against the table, so the
enumeration is complete and not a sample. The four level sets attacked
directly are $d=2$ at $(9,18,27,36)$, $(27,54,81,108)$ and $(33,66,99,132)$
and $d=1$ at $(5,20,45,50)$; the other four of the eight are their images
under multiplication by an element of norm $2$ or $4$, which is a bijection on
the relevant norm levels because the element counts agree, so a solution for
one would give a solution for the other. The remaining three, $d=1$ at
$(13,52,117,130)$, $(17,68,153,170)$ and $(25,50,100,125)$, have levels of
$8,8,8,16$ and $12,12,12,16$ elements. For them the scale is $S=\pm S_{\min}$,
since $|S|\ge2$ violates the bound $|M_{j}|\le r_{K}(N_{j})$ and $S\mapsto-S$
is the symmetry $m\mapsto-m$. The first three levels are streamed against the
fourth, matched through a linear hash of the ten integer coordinates of the
off-diagonal conditions, and every match is checked on all ten coordinates
exactly. A solution has all ten coordinates zero and cannot be missed, so the
search is complete; it covers $4.29\times10^{14}$ sign patterns in the
largest case and finds none.
\end{proof}

\begin{remark}[What the range means]\label{rem:pterangemeaning}
The dimension that matters is $n\ge4$, since $n=2$ and $n=3$ are settled by
\cite{Mar25a,Mar25b} and \cite{FF26}. There
\Cref{thm:pterange}(ii) says that the split route closes before any geometry
is used: the level sums cannot exist, so there is nothing to test for
semiregularity. At $n=3$, where the geometry does produce an object, the level
sums do exist and the off-diagonal conditions are what remove them, which is a
different mechanism and a reminder that (ii) is a statement about a range and
not a theorem for all $n$. What both cases share is that the obstruction is
arithmetic: it lives in the norm form of $K$ and in the supply of integers of
a given norm, not in the cohomology of $A$.
\end{remark}

The range is now beside the point, because a second mechanism, this one in
the cohomology of $A$, closes the split route for every $n$ and every set of
norms at once.

\begin{corollary}[No sum of line bundles runs either argument]\label{cor:noline}
Let $A$ be of $(K,-1,n)$-Weil type with $n\ge2$, and let
$E=\bigoplus_{i}L_{i}[p_{i}]^{\oplus m_{i}}$ be a sum of shifted line bundles,
twisted by a common class or not, whose (twisted) Chern character remains of
Hodge type to first order along $\cD_{n,n}$ and has a nonzero Weil component.
Then $\sigma_{E}$ is not injective on the image of $\operatorname{ev}_{E}$:
$E$ is neither semiregular nor semiregular in the weakened sense of
\cite[Question 11.4]{Mar25b}, and it is obstructed to first order along some
direction of the Weil family. In particular every object of \Cref{setup:pte}
satisfying \eqref{eq:pteconditions}, for every $n\ge2$, every field $K$,
every rank and every set of norms, fails both criteria: the conclusion of
\Cref{thm:pterange}, that no split object of that shape is semiregular, holds
without any bound on the norms, in every dimension, and for the weakened
criterion as well.
\end{corollary}

\begin{proof}
Write $\lambda_{i}$ for the twisted first Chern classes of the summands. If
every $\lambda_{i}$ were a real multiple of $\eta$, it would be a rational
multiple, both classes being rational, and the twisted Chern character
$\sum_{i}m_{i}e^{\lambda_{i}}$ would be a polynomial in $\eta$; its component
in degree $2n$ would then lie in $\QQ\eta^{n}$, which meets the Weil line in
zero by \Cref{prop:grading}, against the hypothesis. So some $\lambda_{i}$ is
not a real multiple of $\eta$, and \Cref{thm:linebundle} applies to the
summand $L_{i}[p_{i}]$. For an object of \Cref{setup:pte} the classes are
$\lambda_{i}=\theta(u_{i})\in R_{0}$, and $R_{0}\cap\QQ\eta=0$; no twist is
needed because $c_{1}(E)=\theta(T_{1,0})=0$ by \eqref{eq:pteconditions}, as
$(1,0)$ is not an excluded pair for $n\ge2$; and \eqref{eq:pteconditions}
gives a flat Chern character with a nonzero Weil component by
\Cref{prop:chintermsofu}.
\end{proof}

\begin{remark}[What is closed and what is not]\label{rem:nolinemeaning}
\Cref{cor:noline} ends the search that \Cref{thm:pterange} began: there is no
sum of line bundles, at any member of any Weil family in any dimension, that
could run the semiregularity argument in either form, and the arithmetic of
\Cref{thm:pterange} is now a description of how the level sums fail rather
than the reason the route fails. The reason is \Cref{prop:rigidity}(i): the
only divisor class that stays of type $(1,1)$ across the whole family is the
polarisation, so every line bundle summand off that line is obstructed in
some direction of the family, and the semiregularity map, whose composite
with the obstruction is contraction with the flat class $\ch(E)$, cannot see
it. What is not closed is the object the argument needs: one with no rank one
summand off the polarisation line, of positive rank, indecomposable enough
that its Atiyah class contracted with every direction of $T$ vanishes.
Markman's reflexive sheaves at $n=3$ are of that kind, and nothing here
says that their analogues at $n\ge4$ do not exist.
\end{remark}

\subsubsection*{Markman's candidate in dimension eight}

There is one concrete candidate of that kind above $n=3$, proposed by Markman
in \cite[Section 12]{Mar25b} with the remark that it is yet to be checked
against the weakened criterion. Its Chern character is a finite computation
in the ring $\QQ[\Theta]/(\Theta^{5})$, and we carry it out, because it
places the candidate exactly at a rational point of the secant plane of
\Cref{def:secantplane} and reduces what has to be checked to a single
condition; that condition is then decided, in the negative, in
\Cref{thm:candidatefails}.

\begin{setupenv}[Markman's object]\label{setup:markman}
Let $(X,\Theta)$ be a principally polarised abelian fourfold, $d\ge3$ odd and
squarefree, $K=\QQ(\sqrt{-d})$ and $N=(d+9)/2$. Let $D_{i}$, $i\in\ZZ/N\ZZ$,
be $N$ translates of $\Theta$ in general position: any two, three or four of
them meet transversally, in a surface, a curve and $\Theta^{4}=24$ points,
and any five have empty intersection. Put $Z_{i}=D_{i}\cap D_{i+1}$,
$Z=\bigcup_{i}Z_{i}$, let $C_{i}=Z_{i}\cap Z_{i+1}=D_{i}\cap D_{i+1}\cap D_{i+2}$
be the curves along which consecutive components meet, and let
$\nu\colon\widetilde{Z}\to Z$ be the partial normalisation of $Z$ at
$m=(d+9)(2d-1)$ of its isolated points of self-intersection, the points of
$Z_{i}\cap Z_{j}$ for $i,j$ at cyclic distance at least three. Put
\[
  \cF=\bigl[\cO_{X}\longrightarrow\nu_{*}\cO_{\widetilde{Z}}\bigr]\otimes\cO_{X}(3\Theta),
\]
the complex placed in degrees $0$ and $1$.
\end{setupenv}

\begin{proposition}[The candidate sits at the point $(1,3)$ of the secant plane]
\label{prop:markman8}
In \Cref{setup:markman}, with $t=\Theta$ and $u_{t},v_{t}$ as in
\eqref{eq:uv} for $n=4$:
\begin{enumerate}[label=\textup{(\roman*)},nosep]
\item $\chi(\cO_{Z_{i}})=14$, $\chi(\cO_{C_{i}})=-36$, the isolated
  self-intersection points of $Z$ number $12N(N-5)=(d+9)(3d-3)$, and
  $\chi(\cO_{\widetilde Z})=110N-12N^{2}+2N(2d-1)=(d+9)(27-d)$;
\item $\ch(\nu_{*}\cO_{\widetilde Z})=N\Theta^{2}-2N\Theta^{3}
  +\chi(\cO_{\widetilde Z})\,[\mathrm{pt}]$;
\item $\ch(\cF)=u_{t}+3v_{t}$, that is
  \[
    \ch(\cF)=1+3\Theta-\tfrac{d}{2}\Theta^{2}-\tfrac{d}{2}\Theta^{3}
      +\tfrac{d^{2}}{24}\Theta^{4},
  \]
  the rational point $(a,b)=(1,3)$ of the secant plane $P_{\Theta}$; with
  the twist $\cO_{X}(\Theta)$ in place of $\cO_{X}(3\Theta)$ the degree two
  component would be $-\tfrac{d+8}{2}\Theta^{2}$ instead, so the twist in
  \cite[Section 12]{Mar25b} is to be read as $3\Theta$;
\item $\Nm_{K/\QQ}(3+\sqrt{-d})=d+9=2N$ is the number of divisors, and by
  \Cref{thm:secantinvariants} the Bogomolov discriminant of $\cF$ is
  $\Delta(\cF)=(d+9)\Theta^{2}$;
\item the objects $\Phi(\cF\boxtimes\cF)$ and $\Phi(\cF\boxtimes\cF^{\vee})$
  on $X\times\widehat X$ have ranks
  $\int_{X}\ch(\cF)^{2}=8d(d-9)$ and
  $\int_{X}\ch(\cF)\ch(\cF^{\vee})=8d(d+9)$, up to sign, both nonzero for
  squarefree $d$; the same formula gives $-8q$ for the pair of secant sheaves
  of \cite[Section 11]{Mar25b} on a threefold, which is his rank $8q$ up to
  the shift.
\end{enumerate}
Consequently $\cF\boxtimes\cF$ satisfies conditions (1), (2b) and (2c) of
Markman's strategy \cite[Section 4]{Mar25b}: the transform has nonzero rank,
its normalised Chern character is invariant under the special Mumford--Tate
group by \cite[Corollary 9.2]{Mar25b}, and its Weil component is nonzero by
\cite[Proposition 10.2]{Mar25b}, the relevant coefficient being
$\Nm(c)$ with $c=\tfrac12(1+3/\sqrt{-d})\ne0$. What remains is condition
(2a) in its weakened form: that $\operatorname{ev}_{E}$ kills every class in
$H^{1}(T)\oplus H^{0}(\bigwedge^{2}T)$ that preserves $\kappa(E)$ to first
order, for $E=\Phi(\cF\boxtimes\cF)$.
\end{proposition}

\begin{proof}
On a principally polarised abelian fourfold $\chi(\cO_{X}(k\Theta))=k^{4}$
and $\Theta^{4}=24$. The Koszul complex of $r$ transversal translates of
$\Theta$ gives $\ch$ of the structure sheaf of their intersection as
$(1-e^{-\Theta})^{r}$, so $\ch(\cO_{Z_{i}})=\Theta^{2}-\Theta^{3}+\tfrac{7}{12}\Theta^{4}$,
$\ch(\cO_{C_{i}})=\Theta^{3}-\tfrac32\Theta^{4}$ and $\ch$ of four divisors
is $\Theta^{4}$; integrating, $\chi(\cO_{Z_{i}})=14$, $\chi(\cO_{C_{i}})=-36$
and $24$ points, as also follows from $0-2\cdot1+16$, $0-3+48-81$ and
$0-4+96-324+256$.

(i), (ii). At every point $Z$ is locally a union of coordinate subspaces of
$\CC^{4}$: along $C_{i}$ two surfaces $\{x_{1}=x_{2}=0\}\cup\{x_{2}=x_{3}=0\}$,
at a point of $C_{i}\cap C_{i+1}$ three of them, at an isolated point
$\{x_{1}=x_{2}=0\}\cup\{x_{3}=x_{4}=0\}$. The ideals of such an arrangement
are monomial and distribute over sums, so the Mayer--Vietoris complex
\[
  0\to\cO_{Z}\to\bigoplus_{i}\cO_{Z_{i}}\to\bigoplus_{i<j}\cO_{Z_{i}\cap Z_{j}}
  \to\bigoplus_{i<j<k}\cO_{Z_{i}\cap Z_{j}\cap Z_{k}}\to\cdots
\]
with scheme-theoretic intersections is exact. Pairs at cyclic distance one
meet in $C_{i}$, $N$ of them; pairs at distance at least two meet in the $24$
points of four divisors, $N(N-3)/2$ of them; the only triples with nonempty
intersection are the consecutive ones $Z_{i}\cap Z_{i+1}\cap Z_{i+2}=C_{i}\cap C_{i+1}$,
$24$ points each, $N$ of them; and no four components meet, since that would
involve at least five divisors. Hence
\[
  \ch(\cO_{Z})=N\ch(\cO_{Z_{i}})-N\ch(\cO_{C_{i}})-\tfrac{N(N-3)}{2}\Theta^{4}+N\Theta^{4},
\]
whose degree two part is $N\Theta^{2}$, degree three part
$-N\Theta^{3}-N\Theta^{3}=-2N\Theta^{3}$, and degree four part
$\chi(\cO_{Z})[\mathrm{pt}]$ with $\chi(\cO_{Z})=14N+36N-12N(N-3)+24N=110N-12N^{2}$.
The isolated singular points are those of pairs at distance at least three,
of which there are $N(N-5)/2$, giving $12N(N-5)=3(d+9)(d-1)$ points; the
pairs at distance two lie on the double curves and are not isolated. At an
isolated point the two branches meet transversally, and separating them
gives $0\to\cO_{Z}\to\nu_{*}\cO_{\widetilde Z}\to k(p)\to0$, so each of the
$m$ normalisations raises $\chi$ by one and leaves the lower degrees of
$\ch$ unchanged. With $N=(d+9)/2$,
$\chi(\cO_{\widetilde Z})=110N-12N^{2}+2N(2d-1)=(d+9)(55-3(d+9)+2d-1)=(d+9)(27-d)$.

(iii) $\ch(\cF)=e^{3\Theta}\bigl(1-\ch(\nu_{*}\cO_{\widetilde Z})\bigr)$. The
degree two part is $(\tfrac92-N)\Theta^{2}=-\tfrac d2\Theta^{2}$, which is
where $N=(d+9)/2$ enters; with the twist $k\Theta$ it would be
$(\tfrac{k^{2}}{2}-N)\Theta^{2}$, equal to $-\tfrac d2\Theta^{2}$ exactly for
$k^{2}=9$, and $k=-3$ is excluded by the degree three part. The degree three
part is $(\tfrac{27}{6}-3N+2N)\Theta^{3}=(\tfrac92-N)\Theta^{3}=-\tfrac d2\Theta^{3}$,
and the degree four part is
$(\tfrac{81}{24}-\tfrac{9N}{2}+6N-\tfrac{\chi(\cO_{\widetilde Z})}{24})\Theta^{4}
=\tfrac{81+36N-\chi(\cO_{\widetilde Z})}{24}\Theta^{4}=\tfrac{d^{2}}{24}\Theta^{4}$
by (i). Comparing with \eqref{eq:uv} at $n=4$,
$u_{t}=1-\tfrac d2t^{2}+\tfrac{d^{2}}{24}t^{4}$ and $v_{t}=t-\tfrac d6t^{3}$,
gives $\ch(\cF)=u_{t}+3v_{t}$.

(iv) $\Nm(3+\sqrt{-d})=9+d$, and \eqref{eq:discnorm} with $(a,b)=(1,3)$.

(v) Orlov's equivalence is $\Phi=(\id\times\Phi_{\cP})^{-1}\circ\mu^{*}$ with
$\mu(x,y)=(x+y,y)$, so the restriction of $\Phi(\cF_{1}\boxtimes\cF_{2})$ to
$\{x\}\times\widehat X$ is the Fourier--Mukai transform of
$\tau_{x}^{*}\cF_{1}\otimes\cF_{2}$, whose derived fibre at $L\in\widehat X$
is $R\Gamma(X,\tau_{x}^{*}\cF_{1}\otimes\cF_{2}\otimes L^{\pm1})$; the rank
of a perfect complex is the Euler characteristic of its derived fibre, and
$\chi(\tau_{x}^{*}\cF_{1}\otimes\cF_{2}\otimes L)=\int_{X}\ch(\cF_{1})\ch(\cF_{2})$
by Riemann--Roch, translations and $L\in\Pic^{0}$ acting trivially on
cohomology. With $\ch(\cF^{\vee})=u_{t}-3v_{t}$, the involution $\tau$
changing the sign of the odd degrees, $\int u_{t}^{2}=8d^{2}$,
$\int u_{t}v_{t}=0$ and $\int v_{t}^{2}=-8d$ give the two values. On a
threefold $\int(u_{t}+v_{t})^{2}=2\int u_{t}v_{t}=-8q$.

The last paragraph: (2b) is the $\mathrm{Spin}(V_{\QQ})_{B}$-invariance of
$\kappa(E)$, which \cite[Corollary 9.2]{Mar25b} derives from
$\ch(\cF)\in B$ and the nonvanishing of the rank; (2c) is
\cite[Proposition 10.2]{Mar25b}, which for $K$ imaginary quadratic reduces
to the nonvanishing of the component of $\ch(\cF)\otimes\ch(\cF)$ in
$\ell_{W}\otimes\ell_{\bbar W}$, and writing $\ch(\cF)=c\,\ell_{W}+\bbar c\,\ell_{\bbar W}$
with $\ell_{W}=e^{\sqrt{-d}\Theta}$ gives that coefficient as $c\bbar c$.
Every number in (i) to (v) is recomputed for every odd squarefree $d\le101$
in item (XXVII) of \Cref{sec:verification}.
\end{proof}

\begin{remark}[The one condition left]\label{rem:onecondition}
\Cref{prop:markman8} does not verify Markman's candidate; it verifies
everything about it that is a statement about classes, and it shows that the
one remaining statement is about the object. That statement is, for
$E=\Phi(\cF\boxtimes\cF)$ on the eightfold $X\times\widehat X$ of Weil type
with $\kappa(E)=\ch(E)e^{-c_{1}(E)/\operatorname{rk}E}$: every
$x\in H^{1}(T)\oplus H^{0}(\bigwedge^{2}T)$ with $x\lrcorner\,\kappa(E)=0$
satisfies $\operatorname{ev}_{E}(x)=0$ in $\Ext^{2}(E,E)$. The $H^{2}(\cO)$
summand imposes nothing, because $\operatorname{rk}E\ne0$ makes
$z\cdot\kappa(E)\ne0$ for $z\ne0$. By \Cref{thm:linebundle} the condition
fails as soon as $E$ has a rank one direct summand whose class is not a
multiple of the polarisation; so the first thing to know about $E$ is
whether it is indecomposable. Markman's object at $n=3$ is reflexive of
rank $8q$ away from a divisor and he proves the condition by descending to
a quotient on which $\sigma_{E}$ is injective on an invariant subspace
\cite[Proposition 11.5]{Mar25b}. Whether the same can be done for
$\Phi(\cF\boxtimes\cF)$ is a question about Ext groups of one sheaf on one
eightfold, not about any Hodge class. \Cref{thm:factor} below moves it from
the eightfold to the fourfold and from the object $E$ to $\cF$ itself: it is
equivalent to the single identity \eqref{eq:heisenberg} in
$\Ext^{2}(\cF,\cF\otimes\Omega^{2}_{X})$, the other half of the condition
holding by construction; and \Cref{thm:candidatefails} shows that the
identity fails, so the answer is no.
\end{remark}

\subsubsection*{The one condition, reduced to the factor}

The condition left in \Cref{rem:onecondition} is about an object on an
eightfold. It is equivalent to a condition on the factor $\cF$ on the
fourfold, and on the fourfold it splits into two, one of which holds for
Markman's object because the object is constructed uniformly over the moduli
of principally polarised fourfolds. The reduction is not special to that
object: it works for every secant object on every abelian $n$-fold with
$n\ge3$, and its ingredients are three computations in the Hodge algebra of
one abelian variety, which is why they can be checked by machine. We fix the
conventions first.

\begin{setupenv}[The weakened criterion for one object]\label{setup:criterion}
Let $Y$ be an abelian variety and $F$ a nonzero perfect complex on $Y$. Write
$HT^{2}(Y)=H^{0,2}(Y)\oplus H^{1}(T_{Y})\oplus H^{0}(\bigwedge^{2}T_{Y})$ as
in \eqref{eq:HT2}, with $x=(z,v,\pi)$ acting on $H^{*}(Y,\CC)$ by
$x\lrcorner$, and likewise $HT^{1}(Y)=H^{0,1}(Y)\oplus H^{0}(T_{Y})$, with
$(z,\theta)$ acting by $z\wedge$ and by the interior product $\theta\lrcorner$.
Normalise the Atiyah class $\mathrm{At}(F)\in\Ext^{1}(F,F\otimes\Omega^{1}_{Y})$
so that $\ch(F)=\operatorname{tr}\exp\mathrm{At}(F)$; for a line bundle it is
then $c_{1}(L)$, as in \eqref{eq:obdef}.%
\footnote{This is the convention of \cite{Tod09} and of
\cite[footnote 14]{Mar25a}; Buchweitz and Flenner use the opposite sign, which
changes $\mathrm{At}(F)$ to $-\mathrm{At}(F)$ and leaves $\mathrm{At}(F)^{2}$
and every statement below unchanged except for the sign of the middle term of
$\operatorname{ev}_{F}$, which is irrelevant to the vanishing statements.}
The evaluation map of \Cref{thm:weakfails} is, under HKR,
\begin{equation}\label{eq:evgeneral}
  \operatorname{ev}_{F}(z,v,\pi)=z\cdot\id_{F}+v\lrcorner\,\mathrm{At}(F)
  +\tfrac12\,\pi\lrcorner\,\mathrm{At}(F)^{2}\;\in\;\Ext^{2}(F,F),
\end{equation}
the contraction of the polyvector with the exponential Atiyah class
\cite[Theorem A]{Hua21}, \cite{BF08}; its restriction to $H^{1}(T_{Y})$ is the
obstruction map, whose kernel is the space of first order deformations of $Y$
to which $F$ extends \cite{Tod09}, \cite[Section 8.3]{Mar25a}. Put
\[
  \Ann_{HT^{2}}(\ch F)=\bigl\{x\in HT^{2}(Y):x\lrcorner\,\ch(F)=0\bigr\},
\]
the Hochschild classes preserving $\ch(F)$ to first order. We say that $F$
\emph{satisfies the weakened criterion} if $\operatorname{ev}_{F}$ vanishes on
$\Ann_{HT^{2}}(\ch F)$. By the commutative square
$\sigma_{F}\circ\operatorname{ev}_{F}=(\,\cdot\,)\lrcorner\,\ch(F)$ of
\cite[Proposition 6.2.1 and Corollary 6.3.2]{BF08}, in the form quoted in
\cite[Lemma 11.3]{Mar25b},
\[
  \operatorname{ev}_{F}\bigl(\Ann_{HT^{2}}(\ch F)\bigr)
  =\ker\sigma_{F}\cap\image\operatorname{ev}_{F},
\]
so the weakened criterion is exactly the hypothesis of
\cite[Question 11.4]{Mar25b}, that $\sigma_{F}$ restricts to the image of
$\operatorname{ev}_{F}$ as an injective map; \Cref{prop:evaluation}(ii) is
the case of a split object.
\end{setupenv}

\begin{lemma}[The $B$-field transform]\label{lem:bfield}
Let $B\in H^{1,1}(Y)$. For $x=(z,v,\pi)\in HT^{2}(Y)$ put
\begin{equation}\label{eq:bfield}
  e^{B}x=\bigl(z+v\lrcorner\,B+\tfrac12\,\pi\lrcorner\,B^{2},\;
  v+\pi\lrcorner\,B,\;\pi\bigr),
\end{equation}
where $\pi\lrcorner\,B\in\Hom(H^{1,0},H^{0,1})=H^{1}(T_{Y})$ is the contraction
of one slot of $\pi$ with $B$: for $\pi=\partial_{a}\wedge\partial_{b}$ it sends
$p_{b}$ to $\partial_{a}\lrcorner\,B$, $p_{a}$ to $-\partial_{b}\lrcorner\,B$
and the other basis vectors to zero. Then $e^{B}$ is a linear automorphism of
$HT^{2}(Y)$ with inverse $e^{-B}$, and:
\begin{enumerate}[label=\textup{(\roman*)},nosep]
\item $x\lrcorner\,(w\,e^{B})=\bigl((e^{B}x)\lrcorner\,w\bigr)\,e^{B}$ for
  every even class $w\in H^{\mathrm{ev}}(Y,\CC)$;
\item for a line bundle $L$ with $c_{1}(L)=B$,
  $\operatorname{ev}_{F\otimes L}(x)=\operatorname{ev}_{F}(e^{B}x)$ and
  $\Ann_{HT^{2}}(\ch(F\otimes L))=e^{-B}\Ann_{HT^{2}}(\ch F)$; hence $F$
  satisfies the weakened criterion if and only if $F\otimes L$ does.
\end{enumerate}
The computation in (ii) is formal: it uses only that the Atiyah class of the
twisted object is $\mathrm{At}(F)+B\cdot\id_{F}$ and its Chern character
$\ch(F)e^{B}$. It therefore applies verbatim to the twisted sheaves of
\cite[Section 2]{Mar25b}, with $B$ rational: for $F$ of nonzero rank $r$, the
criterion for the twisted object $F\otimes L^{-1}$ with class
$\kappa(F)=\ch(F)e^{-c_{1}(F)/r}$ is the same condition as the criterion for
$F$ with class $\ch(F)$.
\end{lemma}

\begin{proof}
Everything takes place in a graded algebra on which the interior products
$\iota_{a}=\partial_{a}\lrcorner$ act as odd derivations, and the only
properties used are that $w$ is even, that $B$ is even and central with
$\iota_{a}\iota_{b}B=0$, which holds for $B$ of type $(1,1)$ and
$\partial_{a},\partial_{b}$ of type $(1,0)$, and that the odd scalar classes
$\iota_{a}B$ anticommute with every odd element. In $H^{*}(Y,\CC)$ this is
clear; the same argument applies in
$\bigoplus_{p,q}\Ext^{p}(F,F\otimes\Omega^{q}_{Y})=\Ext^{*}(F,F)\otimes
\bigwedge^{*}H^{1,0}$, on which the polyvectors act through the second factor
by odd derivations, because the classes coming from $H^{*}(Y,\CC)$ are graded
central in the Yoneda algebra. First $\iota_{a}(B^{k})=k(\iota_{a}B)B^{k-1}$,
since $\iota_{a}B$ is odd and $B$ is even, so
$\iota_{a}e^{B}=(\iota_{a}B)e^{B}$. For $w$ even this gives
$v\lrcorner\,(we^{B})=(v\lrcorner\,w)e^{B}+w\sum_{a}v(p_{a})(\iota_{a}B)e^{B}
=\bigl(v\lrcorner\,w+(v\lrcorner\,B)w\bigr)e^{B}$, and
$z(we^{B})=(zw)e^{B}$. For the bivector, with $\pi\lrcorner=\iota_{a}\iota_{b}$,
\[
  \iota_{b}(we^{B})=(\iota_{b}w)e^{B}+w(\iota_{b}B)e^{B},
\]
and applying $\iota_{a}$, using that $\iota_{b}w$ and $\iota_{b}B$ are odd,
\[
  \iota_{a}\iota_{b}(we^{B})
  =\bigl[\iota_{a}\iota_{b}w
   +(\iota_{a}B)\iota_{b}w-(\iota_{b}B)\iota_{a}w
   +w\bigl((\iota_{a}\iota_{b}B)-(\iota_{b}B)(\iota_{a}B)\bigr)\bigr]e^{B}.
\]
The middle two terms are $(\pi\lrcorner\,B)\lrcorner\,w$ for the element
$\pi\lrcorner\,B$ of $\Hom(H^{1,0},H^{0,1})$ displayed in the statement, acting
as a derivation, and the last bracket is
$\tfrac12\iota_{a}\iota_{b}(B^{2})=\iota_{a}\bigl((\iota_{b}B)B\bigr)$.
Summing the three actions gives (i). Applying (i) with $B$ replaced by $-B$
and then by $B$, and using
$(\pi\lrcorner\,B)\lrcorner\,B=2(\iota_{a}B)(\iota_{b}B)=\pi\lrcorner\,B^{2}$
when $\iota_{a}\iota_{b}B=0$, shows $e^{-B}e^{B}=\id$.

(ii) The Atiyah class of $F\otimes L$ is $\mathrm{At}(F)+B\cdot\id_{F}$,
the Atiyah class being additive on tensor products \cite[Section 3]{BF03},
and the two summands commute, so
$\exp\mathrm{At}(F\otimes L)=\exp\mathrm{At}(F)\cdot e^{B}$. By (i) in the
algebra $\Ext^{*}(F,F)\otimes\bigwedge^{*}H^{1,0}$, with $w=\exp\mathrm{At}(F)$,
\[
  x\lrcorner\,\exp\mathrm{At}(F\otimes L)
  =\bigl((e^{B}x)\lrcorner\,\exp\mathrm{At}(F)\bigr)e^{B},
\]
and the component in $\Ext^{2}(F,F)\otimes\bigwedge^{0}$ of the right hand
side is that of $(e^{B}x)\lrcorner\,\exp\mathrm{At}(F)$, because $e^{B}=1+B+\cdots$
raises the exterior degree. This is
$\operatorname{ev}_{F\otimes L}(x)=\operatorname{ev}_{F}(e^{B}x)$. By (i) in
$H^{*}(Y,\CC)$ with $w=\ch(F)$, $x\lrcorner\,\ch(F\otimes L)=0$ if and only
if $(e^{B}x)\lrcorner\,\ch(F)=0$, which is the statement on the annihilators.
Hence $\operatorname{ev}_{F\otimes L}(\Ann_{HT^{2}}(\ch(F\otimes L)))
=\operatorname{ev}_{F}(\Ann_{HT^{2}}(\ch F))$. The identity (i) is checked
on random data in item (XXVIII) of \Cref{sec:verification}.
\end{proof}

\begin{theorem}[Reduction to the factor]\label{thm:factor}
Let $(X,t)$ be a polarised abelian variety of dimension $n\ge3$, let $d\ge1$,
and let $F$ be a perfect complex on $X$ with
$\ch(F)=a\,u_{t}+b\,v_{t}$ for some $(a,b)\in\QQ^{2}\setminus\{0\}$, in the
notation of \eqref{eq:uv}. Then:
\begin{enumerate}[label=\textup{(\roman*)},nosep]
\item No nonzero class in $HT^{1}(X)$ preserves $\ch(F)$: the map
  $HT^{1}(X)\to H^{*}(X,\CC)$, $y\mapsto y\lrcorner\,\ch(F)$, is injective. Its
  two lowest components form the matrix
  $\bigl(\begin{smallmatrix}a&b\\ b&-ad\end{smallmatrix}\bigr)$ of
  determinant $-(a^{2}d+b^{2})=-\Nm_{K/\QQ}(b+a\sqrt{-d})\ne0$.
\item The classes in $HT^{2}(X)$ preserving $\ch(F)$ to first order are
  \begin{equation}\label{eq:secantkernel}
    \Ann_{HT^{2}}(\ch F)
    =\bigl\{(0,v,0):v\lrcorner\,t=0\bigr\}\;\oplus\;
     \bigl\{(\tfrac d2\,\pi\lrcorner\,t^{2},\,0,\,\pi):
       \pi\in H^{0}(\textstyle\bigwedge^{2}T_{X})\bigr\},
  \end{equation}
  the polarised first order deformations of $(X,t)$, of dimension $n(n+1)/2$,
  and the Poisson directions each compensated by a class in $H^{0,2}$, of
  dimension $n(n-1)/2$; the total is $n^{2}=\dim\cD_{n,n}$, and the space
  does not depend on $(a,b)$.
\item $F$ satisfies the weakened criterion if and only if
  \begin{enumerate}[label=\textup{(\alph*)},nosep]
  \item $v\lrcorner\,\mathrm{At}(F)=0$ in $\Ext^{2}(F,F)$ for every
    $v\in H^{1}(T_{X})$ with $v\lrcorner\,t=0$, that is, $F$ extends to first
    order along every polarised deformation of $(X,t)$; and
  \item $\mathrm{At}(F)^{2}=-d\,t^{2}\cdot\id_{F}$ in
    $\Ext^{2}(F,F\otimes\Omega^{2}_{X})$.
  \end{enumerate}
  The trace of (b) is automatic: $\operatorname{tr}\mathrm{At}(F)^{2}
  =2\ch_{2}(F)=-ad\,t^{2}=\operatorname{tr}(-d\,t^{2}\cdot\id_{F})$. So (b)
  says that the trace free part of the square of the Atiyah class vanishes.
\item Let $F_{1},F_{2}$ be two such complexes, $G=F_{1}\boxtimes F_{2}$ on
  $X\times X$ and $E=\Phi(G)$ on $X\times\widehat X$, with $\Phi$ Orlov's
  equivalence as in \Cref{prop:markman8}(v). Then
  $\Ann_{HT^{2}}(\ch G)=\Ann_{HT^{2}}(\ch F_{1})\otimes1\oplus
  1\otimes\Ann_{HT^{2}}(\ch F_{2})$, of dimension $2n^{2}$, and
  \[
    E\ \text{satisfies the weakened criterion}
    \iff G\ \text{does}
    \iff F_{1}\ \text{and}\ F_{2}\ \text{do}.
  \]
\end{enumerate}
\end{theorem}

\begin{proof}
Choose a basis $p_{1},\dots,p_{n}$ of $H^{1,0}(X)$ and a basis
$q_{1},\dots,q_{n}$ of $H^{0,1}(X)$ with $t=\sum_{a}p_{a}\wedge q_{a}$, which
is possible because $t$ is a positive class of type $(1,1)$; let
$\partial_{a}$ be the dual basis of $T$, so that $\iota_{a}p_{b}=\delta_{ab}$
and $\iota_{a}q_{b}=0$. Write $\ch(F)=\sum_{k}c_{k}t^{k}$ with
$c_{k}\,k!=a(-d)^{k/2}$ for $k$ even and $c_{k}\,k!=b(-d)^{(k-1)/2}$ for $k$
odd, so that
\begin{equation}\label{eq:crecursion}
  c_{k+1}\,(k+1)!=-d\,c_{k-1}\,(k-1)!\qquad(k\ge1).
\end{equation}
Two facts about the Hodge algebra are used throughout: for $v\in H^{1}(T)$
acting as a derivation, $v\lrcorner\,t^{k}=k\,(v\lrcorner\,t)\,t^{k-1}$; and
for a bivector $\pi$ of type $(2,0)$ in the tangent directions,
$\pi\lrcorner\,t=0$ because $t$ has a single $(1,0)$ slot in each term, so
$\pi\lrcorner\,t^{k}=\binom{k}{2}(\pi\lrcorner\,t^{2})\,t^{k-2}$ by the footnote
to \eqref{eq:HT2}. Finally, multiplication by $t$ is injective on $H^{0,1}$
when $n\ge2$ and on $H^{0,2}$ when $n\ge3$: $q_{a}\wedge t$ contains the
monomial $p_{c}q_{c}q_{a}$ for any $c\ne a$, and $q_{a}q_{b}\wedge t$ contains
$p_{c}q_{c}q_{a}q_{b}$ for any $c\notin\{a,b\}$, and distinct monomials are
independent.

(i) For $y=(z,\theta)$ the class $y\lrcorner\,\ch(F)$ has components of type
$(k-1,k)$ equal to $c_{k-1}z\,t^{k-1}+c_{k}\,\theta\lrcorner\,t^{k}
=\bigl(c_{k-1}z+k\,c_{k}\,(\theta\lrcorner\,t)\bigr)t^{k-1}$. For $k=1$ this
is $az+b\,(\theta\lrcorner\,t)$ and for $k=2$ it is
$\bigl(bz-ad\,(\theta\lrcorner\,t)\bigr)t$, by \eqref{eq:crecursion}. Both
$z$ and $\theta\lrcorner\,t=\sum_{a}\theta^{a}q_{a}$ lie in $H^{0,1}$, on which
$t$ acts injectively, so $y\lrcorner\,\ch(F)=0$ forces
$(z,\theta\lrcorner\,t)$ into the kernel of the displayed matrix, whose
determinant is $-a^{2}d-b^{2}<0$. Hence $z=0$ and $\theta\lrcorner\,t=0$, and
the latter gives $\theta=0$ because $t$ is nondegenerate.

(ii) For $x=(z,v,\pi)$ the component of $x\lrcorner\,\ch(F)$ of type
$(k-1,k+1)$ is
\[
  c_{k-1}z\,t^{k-1}+c_{k}\,v\lrcorner\,t^{k}+c_{k+1}\,\pi\lrcorner\,t^{k+1}
  =\Bigl(c_{k-1}z+k\,c_{k}\,W+\tbinom{k+1}{2}c_{k+1}\,P\Bigr)t^{k-1},
\]
with $W=v\lrcorner\,t$ and $P=\pi\lrcorner\,t^{2}$, both in $H^{0,2}$. By
\eqref{eq:crecursion}, $\binom{k+1}{2}c_{k+1}=-\tfrac d2c_{k-1}$, so the bracket
is $c_{k-1}Z+k\,c_{k}W$ with $Z=z-\tfrac d2P$. For $k=1$ and $k=2$ this gives
$aZ+bW$ and $(bZ-adW)\,t$, and $t$ is injective on $H^{0,2}$ because $n\ge3$;
the same determinant forces $Z=W=0$. Conversely $Z=W=0$ kills every
component. So $\Ann_{HT^{2}}(\ch F)$ is the set of $(z,v,\pi)$ with
$v\lrcorner\,t=0$ and $z=\tfrac d2\,\pi\lrcorner\,t^{2}$, which is
\eqref{eq:secantkernel}; the two summands are independent because the first
has $\pi=0$ and the second is the graph of a linear map from the $\pi$
component. A $v\in\Hom(H^{1,0},H^{0,1})$ has $v\lrcorner\,t=\sum_{a}v(p_{a})\wedge q_{a}$,
which vanishes exactly when the matrix of $v$ is symmetric; this is the
tangent space of the polarised deformations, of dimension $n(n+1)/2$, as in
\Cref{prop:splitobstruction}(i). The bivectors contribute
$\dim\bigwedge^{2}T=n(n-1)/2$. Nothing in the description involves $(a,b)$.

(iii) By (ii) and \eqref{eq:evgeneral}, $\operatorname{ev}_{F}$ vanishes on
$\Ann_{HT^{2}}(\ch F)$ if and only if it vanishes on both summands. On the
first it is $v\mapsto v\lrcorner\,\mathrm{At}(F)$, which is (a). On the
second it is
\[
  \tfrac d2(\pi\lrcorner\,t^{2})\cdot\id_{F}+\tfrac12\,\pi\lrcorner\,\mathrm{At}(F)^{2}
  =\tfrac12\,\pi\lrcorner\bigl(\mathrm{At}(F)^{2}+d\,t^{2}\cdot\id_{F}\bigr),
\]
since $\pi\lrcorner\,(t^{2}\cdot\id_{F})=(\pi\lrcorner\,t^{2})\cdot\id_{F}$.
Now $\Omega^{2}_{X}=\cO_{X}\otimes\bigwedge^{2}H^{1,0}$, so
$\Ext^{2}(F,F\otimes\Omega^{2}_{X})=\Ext^{2}(F,F)\otimes\bigwedge^{2}H^{1,0}$,
and the contractions with all $\pi\in\bigwedge^{2}(H^{1,0})^{\vee}$ vanish on
an element of this space if and only if the element is zero, the pairing
being perfect. This is (b). For the trace, $\ch(F)=\operatorname{tr}\exp\mathrm{At}(F)$
gives $\ch_{2}(F)=\tfrac12\operatorname{tr}\mathrm{At}(F)^{2}$, and
$\ch_{2}(F)=c_{2}t^{2}=-\tfrac{ad}{2}t^{2}$, while
$\operatorname{tr}(\id_{F})=\rk F=c_{0}=a$.

(iv) The Hochschild cohomology of a product is the tensor product,
$HT^{2}(X\times X)=HT^{2}\otimes HT^{0}\oplus HT^{1}\otimes HT^{1}\oplus
HT^{0}\otimes HT^{2}$, the Atiyah class of $G$ is
$\mathrm{At}(F_{1})\boxtimes1+1\boxtimes\mathrm{At}(F_{2})$ \cite[Proposition 3.14]{BF03},
\cite[(8.4.1)]{Mar25a}, so $\exp\mathrm{At}(G)=\exp\mathrm{At}(F_{1})\boxtimes
\exp\mathrm{At}(F_{2})$, and $\Ext^{2}(G,G)$ is the corresponding K\"unneth
sum of $\Ext^{i}(F_{1},F_{1})\otimes\Ext^{j}(F_{2},F_{2})$. A polyvector on a
factor contracts only the forms of that factor, so for $x_{1}\in HT^{i}(X)$
and $x_{2}\in HT^{j}(X)$ with $i+j=2$,
\[
  (x_{1}\otimes x_{2})\lrcorner\,\ch(G)=\pm(x_{1}\lrcorner\,\ch F_{1})\otimes
  (x_{2}\lrcorner\,\ch F_{2}),\qquad
  \operatorname{ev}_{G}(x_{1}\otimes x_{2})=\pm\operatorname{ev}_{F_{1}}(x_{1})
  \otimes\operatorname{ev}_{F_{2}}(x_{2}),
\]
with $\operatorname{ev}(1)=\id$ on $HT^{0}$ and Koszul signs that play no
role; this is the decomposition \cite[(8.4.2), (8.4.3)]{Mar25a}. Let
$x=x_{20}+x_{11}+x_{02}$ be a class in $\Ann_{HT^{2}}(\ch G)$ according to the
three summands. Since $\ch(F_{i})$ is even, $x_{11}\lrcorner\,\ch(G)$ lies in
$H^{\mathrm{odd}}\otimes H^{\mathrm{odd}}$ while the other two terms lie in
$H^{\mathrm{ev}}\otimes H^{\mathrm{ev}}$, so $x_{11}\lrcorner\,\ch(G)=0$
separately; and $x_{11}\lrcorner\,\ch(G)$ is the image of $x_{11}$ under the
tensor product of the two maps of (i), which is injective, so $x_{11}=0$. For
the even part write $x_{20}=x_{1}\otimes1$ and $x_{02}=1\otimes x_{2}$; then
$(x_{1}\lrcorner\,\ch F_{1})\otimes\ch(F_{2})+\ch(F_{1})\otimes(x_{2}\lrcorner\,\ch F_{2})=0$.
Applying a linear form on the second factor taking the value $1$ on
$\ch(F_{2})$ shows $x_{1}\lrcorner\,\ch(F_{1})=\lambda\,\ch(F_{1})$ for a
scalar $\lambda$, and then $x_{2}\lrcorner\,\ch(F_{2})=-\lambda\,\ch(F_{2})$.
But every component of $x_{1}\lrcorner\,\ch(F_{1})$ has Hodge type $(p,p+2)$
while $\ch(F_{1})\ne0$ has type $(p,p)$; so $\lambda=0$ and
$x_{i}\in\Ann_{HT^{2}}(\ch F_{i})$. This proves the description of
$\Ann_{HT^{2}}(\ch G)$ and its dimension $2n^{2}$ by (ii). On it
$\operatorname{ev}_{G}(x_{1}\otimes1+1\otimes x_{2})=\pm\operatorname{ev}_{F_{1}}(x_{1})
\otimes\id_{F_{2}}\pm\id_{F_{1}}\otimes\operatorname{ev}_{F_{2}}(x_{2})$, whose
two terms lie in different K\"unneth summands and vanish exactly when
$\operatorname{ev}_{F_{i}}(x_{i})=0$, as $\id_{F_{i}}\ne0$. So $G$ satisfies
the criterion if and only if $F_{1}$ and $F_{2}$ do.

For the transport under $\Phi$ we use the commutative diagram of
\cite[Lemma 8.3.10]{Mar25a}: an equivalence $\Phi\colon D^{b}(A)\to D^{b}(B)$
of derived categories of abelian varieties induces isomorphisms
$\Phi^{HT}\colon HT^{2}(A)\to HT^{2}(B)$ and $\Phi^{H}\colon H^{*}(A,\CC)\to
H^{*}(B,\CC)$ with $\Phi^{H}(x\lrcorner\,\ch G)=\Phi^{HT}(x)\lrcorner\,\ch\Phi(G)$
\cite{CRV12}, and $\Phi\circ\operatorname{ev}_{G}=\operatorname{ev}_{\Phi(G)}\circ\Phi^{HT}$
on $\Hom(G,G[2])\to\Hom(\Phi G,\Phi G[2])$ \cite[Theorem A]{Hua21}; on abelian
varieties the Todd class is trivial, so there is no Duflo correction
\cite[Remark 8.3.11]{Mar25a}. The first square gives
$\Phi^{HT}\bigl(\Ann_{HT^{2}}(\ch G)\bigr)=\Ann_{HT^{2}}(\ch E)$ and the second
$\operatorname{ev}_{E}\bigl(\Ann_{HT^{2}}(\ch E)\bigr)
=\Phi\bigl(\operatorname{ev}_{G}(\Ann_{HT^{2}}(\ch G))\bigr)$, with $\Phi$ an
isomorphism on $\Ext^{2}$. So $E$ satisfies the criterion if and only if $G$
does. Items (i) and (ii) are recomputed in the model of \Cref{sec:verification},
item (XXVIII), for $n=3,4,5$ and several $(d,a,b)$, together with the controls
that the bare Poisson class $(0,0,\pi)$ and the class compensated with the
opposite sign do not preserve $\ch(F)$.
\end{proof}

\begin{figure}[tbp]
\centering
  \includegraphics[width=\textwidth]{fig_factor}
\caption{\Cref{thm:factor} for $n=4$. Top: the K\"unneth stack of
$HT^{2}(X\times X)$, whose middle summand $HT^{1}\otimes HT^{1}$ meets the
kernel in zero by (i), so that the classes preserving $\ch(F_{1}\boxtimes F_{2})$
are two copies of the kernel on one factor; across Orlov's isomorphism,
$HT^{2}$ of the eightfold with the $2n^{2}$ classes preserving $\ch(E)$.
Bottom: $HT^{2}(X)$ as three slabs with heights proportional to their
dimensions, and inside them the $n^{2}$ classes of \eqref{eq:secantkernel}:
the polarised deformations in front of $H^{1}(T_{X})$, and the whole of
$H^{0}(\bigwedge^{2}T_{X})$, each bivector tilted into $H^{0,2}$ by its
compensation $\tfrac d2\,\pi\lrcorner t^{2}$; nothing in $H^{0,2}$ alone
preserves $\ch(F)$. On the right, the two conditions that the evaluation map
imposes on the two parts, (a) closed by the construction of $\cF$ and (b)
the identity \eqref{eq:heisenberg}.}
\label{fig:factor}
\end{figure}

\begin{corollary}[The one condition, on the fourfold]\label{cor:markmancondition}
In \Cref{setup:markman} let $E=\Phi(\cF\boxtimes\cF)$. Then $E$ satisfies the
weakened criterion of \cite[Question 11.4]{Mar25b} if and only if
\begin{equation}\label{eq:heisenberg}
  \mathrm{At}(\cF)^{2}=-d\,\Theta^{2}\cdot\id_{\cF}
  \qquad\text{in }\Ext^{2}(\cF,\cF\otimes\Omega^{2}_{X}).
\end{equation}
Equivalently, writing $\mathrm{At}(\cF)=\sum_{a}A_{a}\otimes p_{a}$ with
$A_{a}\in\Ext^{1}(\cF,\cF)$ in a basis $p_{a}$ of $H^{1,0}(X)$, all the
commutators $A_{a}A_{b}-A_{b}A_{a}\in\Ext^{2}(\cF,\cF)$ are the scalars
$\pm2d\,(q_{a}\wedge q_{b})\cdot\id_{\cF}$ prescribed by $\Theta^{2}$, the
sign being fixed by the Koszul rule.
\end{corollary}

\begin{proof}
By \Cref{prop:markman8}(iii), $\ch(\cF)=u_{\Theta}+3v_{\Theta}$, so
\Cref{thm:factor} applies with $n=4$, $(a,b)=(1,3)$ and $F_{1}=F_{2}=\cF$, and
$E$ satisfies the criterion if and only if $\cF$ satisfies (a) and (b). It
remains to see that (a) holds. The data of \Cref{setup:markman} are a
principally polarised fourfold $(X,\Theta)$ and $N$ points $x_{i}\in X$ with
$D_{i}=\tau_{x_{i}}^{*}\Theta$; the conditions imposed on the $D_{i}$, that
the intersections of two, three and four of them be smooth of the expected
codimension and that five of them be disjoint, are open in the parameter
space of such data over the local universal deformation $S$ of $(X,\Theta)$,
by the openness of smoothness for a proper flat family and the properness of
the projection; and the isolated points of $Z$, being transversal
intersections of four divisors, form a finite \'etale cover of the open set
where the conditions hold, so the $m$ points to be normalised can be followed
locally. The partial normalisation separates two smooth branches at each of
them and commutes with base change. So the construction produces a perfect
complex $\cF_{S}$ on the universal family over a neighbourhood of the base
point in $S$, whose fibre at the base point is $\cF$ and whose base has
Kodaira--Spencer image the whole space of polarised deformations, of
dimension $10=n(n+1)/2$. The obstruction $v\lrcorner\,\mathrm{At}(\cF)$ to
extending $\cF$ along $v$ vanishes for every $v$ along which $\cF$ does extend
\cite{Tod09}, which is (a).

For the reformulation, $\Omega^{1}_{X}=\cO_{X}\otimes H^{1,0}$ gives
$\mathrm{At}(\cF)=\sum_{a}A_{a}\otimes p_{a}$, and
$\mathrm{At}(\cF)^{2}=-\sum_{a<b}(A_{a}A_{b}-A_{b}A_{a})\otimes p_{a}\wedge p_{b}$
by the Koszul rule, while $\Theta^{2}=-2\sum_{a<b}(p_{a}\wedge p_{b})(q_{a}\wedge q_{b})$
for $\Theta=\sum_{a}p_{a}\wedge q_{a}$. Comparing coefficients of
$p_{a}\wedge p_{b}$ gives the commutator relations.
\end{proof}

\begin{remark}[What the reduction closes, and what it does not]\label{rem:factorclosed}
Three things are closed. First, the eightfold has disappeared: the condition
of \Cref{rem:onecondition}, on an object of rank $8d(d-9)$ on $X\times\widehat X$
of Weil type, is equivalent to one identity \eqref{eq:heisenberg} for a rank
one complex on the fourfold, in the Yoneda algebra $\Ext^{*}(\cF,\cF)$; the
dimension of the space of classes that has to be killed drops from $2n^{2}$
to $n^{2}$ and, by (a), to $n(n-1)/2=6$; at $n=3$ the count $n^{2}=9$ is the
nine-dimensional kernel of \cite[Corollaries 8.3.12 and 8.5.2]{Mar25a}, now
written out for every secant object in every dimension. Second, the half of that space made
of honest deformations imposes nothing, because Markman's construction is
uniform over the moduli of principally polarised fourfolds. Third, the two
forms of the criterion, for $E$ with $\ch(E)$ and for the twisted object with
$\kappa(E)$, are the same by \Cref{lem:bfield}, so nothing is lost by working
with the untwisted object.

The identity itself is a Heisenberg relation: the components of the Atiyah
class in the $n$ tangent directions must commute up to the scalars
prescribed by $\Theta^{2}$, and the only part of it that the Chern character
forces is the trace. It is exactly the statement that $\cF$ deforms, as an
object of the deformed category, along the Poisson directions of $X$
compensated by the gerbe classes $\tfrac d2\,\pi\lrcorner\,\Theta^{2}$, the
noncommutative first order deformations of $X$ with a gerby correction in
the language of \cite[Section 8.5]{Mar25a}. It holds
in the one case where the method is known to work: for Markman's sheaves
$\cF_{1}=I_{\bigcup C_{i}}(\Theta)$ on the Jacobian threefold of
\cite[Section 11.1]{Mar25b}, with $\ch=u_{\Theta}+v_{\Theta}$ and $d=q$, the
kernel of the obstruction map equals the annihilator of the Chern character
\cite[Corollary 8.3.12]{Mar25a}, which is stronger than the weakened
criterion, so by \Cref{thm:factor}(iii) these sheaves satisfy
$\mathrm{At}(\cF_{1})^{2}=-q\,\Theta^{2}\cdot\id$ on the threefold. For
the rank one complex of \Cref{setup:markman} on the fourfold the identity
is a computation of Yoneda products in $\Ext^{*}(\cF,\cF)$, and it can be
made local: the next paragraphs carry it out at one point of the support and
find that it fails.
\end{remark}

\subsubsection*{The identity fails at an isolated double point}

The identity \eqref{eq:heisenberg} can be tested locally. The evaluation map
is an algebra homomorphism \cite[(8.3.3)]{Mar25a}, so on a bivector
$\theta\wedge\theta'$ it is the Yoneda product of the classes
\[
  A_{\theta}=\operatorname{ev}_{\cF}(\theta)=\theta\lrcorner\,\mathrm{At}(\cF)
  \in\Ext^{1}(\cF,\cF),\qquad\theta\in H^{0}(T_{X}),
\]
which are the first order deformations of $\cF$ along the translations of
$X$; and the product of two global classes restricts to the product of their
germs. At a point where the support of $\cF$ has two smooth branches of
codimension two meeting transversally, the germs multiply nontrivially, and
since the compensating classes $z\cdot\id_{\cF}$ have zero germs, the
identity fails. Markman's construction has such points, and it cannot avoid
them.

\begin{lemma}[Localisation]\label{lem:edge}
Let $F$ be a perfect complex on a smooth projective variety $Y$, let
$\cE xt^{m}(F,F)$ be the sheaf associated with $U\mapsto\Hom_{D(U)}(F|_{U},F|_{U}[m])$,
and let $e\colon\Ext^{m}(F,F)\to H^{0}(Y,\cE xt^{m}(F,F))$ send a morphism
$F\to F[m]$ to its germs. Then $e$ is compatible with Yoneda products;
$e(z\cdot\id_{F})=0$ for $z\in H^{m}(\cO_{Y})$, $m\ge1$; and for
$\theta\in H^{0}(T_{Y})$ the germ of $A_{\theta}$ at $p$ is the class of the
jet extension $0\to F_{p}\to J_{\theta}(F_{p})\to F_{p}\to0$ of the germ
$F_{p}$ along the derivation $\theta$ of $\cO_{Y,p}$, which is
$\cO_{Y,p}$-linear in $\theta$.
\end{lemma}

\begin{proof}
Composition of morphisms commutes with restriction to open sets, which is the
first assertion. For the second, $z\cdot\id_{F}$ restricted to an affine open
$U$ is $z|_{U}\cdot\id_{F|_{U}}$, and $z|_{U}\in H^{m}(U,\cO_{U})=0$. For the
third, the Atiyah class of $F$ is the class of the first jet sequence
$0\to F\otimes\Omega^{1}_{Y}\to J^{1}(F)\to F\to0$ \cite[Section 3]{BF03},
its restriction to $U$ is the jet sequence of $F|_{U}$, and contracting with
$\theta$ is pushing out along $\id\otimes\theta\colon F\otimes\Omega^{1}\to F$,
which produces the extension $J_{\theta}(F)$ whose underlying sheaf is
$F\oplus F$ with $f\cdot(s,s')=(fs+\theta(f)s',fs')$. Linearity in $\theta$
is linearity of the contraction.
\end{proof}

\begin{lemma}[Two branches through a point]\label{lem:localproducts}
Let $\cO$ be the local ring of $\CC^{4}$ at the origin or its completion
$\CC[[x_{1},\dots,x_{4}]]$; let $P_{1}=\{x_{1}=x_{2}=0\}$ and
$P_{2}=\{x_{3}=x_{4}=0\}$, two smooth surfaces meeting transversally at the
origin; put $N_{1}=\langle\partial_{1},\partial_{2}\rangle$ and
$N_{2}=\langle\partial_{3},\partial_{4}\rangle$, the normal spaces of $P_{1}$
and of $P_{2}$ at the origin, so that
$\bigwedge^{2}T_{0}=\bigwedge^{2}N_{1}\oplus N_{1}\otimes N_{2}\oplus\bigwedge^{2}N_{2}$.
For a germ of vector field $\theta$ and a perfect complex $F$ of
$\cO$-modules write $a_{\theta}\in\Ext^{1}_{\cO}(F,F)$ for the class of the
jet extension along $\theta$.
\begin{enumerate}[label=\textup{(\roman*)},nosep]
\item For $F=I=I_{P_{1}}\cap I_{P_{2}}$, the ideal of the union,
  $\Ext^{2}_{\cO}(I,I)\cong N_{1}\otimes N_{2}$, a vector space of dimension
  four on which the maximal ideal acts trivially, and
  $a_{\theta}a_{\theta'}$ is, up to a fixed nonzero scalar, the component of
  $\theta(0)\wedge\theta'(0)$ in $N_{1}\otimes N_{2}$. In particular
  $a_{\theta}a_{\theta'}\ne0$ unless $\theta(0)\wedge\theta'(0)\in
  \bigwedge^{2}N_{1}\oplus\bigwedge^{2}N_{2}$.
\item For $F=K=[\cO\to\cO_{P_{1}}\oplus\cO_{P_{2}}]$, the complex in degrees
  $0$ and $1$ whose cohomology is $I$ in degree $0$ and the residue field $k$
  in degree $1$, $a_{\theta}a_{\theta'}\ne0$ unless
  $\theta(0)\wedge\theta'(0)\in N_{1}\otimes N_{2}$.
\end{enumerate}
In both cases the bilinear map $(\theta,\theta')\mapsto a_{\theta}a_{\theta'}$
is nonzero, and it depends only on $\theta(0)\wedge\theta'(0)$.
\end{lemma}

\begin{proof}
Write $A=\CC[x_{1},x_{2}]$, $B=\CC[x_{3},x_{4}]$, $R=A\otimes_{\CC}B$, with
maximal ideals $\mathfrak m_{A}$, $\mathfrak m_{B}$ at the origin, and
$k=\CC$. The ideal of the union is $I=\mathfrak m_{A}R\cap\mathfrak m_{B}R
=\mathfrak m_{A}\otimes_{\CC}\mathfrak m_{B}$, because the two ideals are
generated by disjoint sets of variables. Everything below is homogeneous, the
$\Ext$ groups in degree two are of finite length, and finite length modules
and jet classes are unchanged by the flat base change from $R$ to $\cO$; so we
compute over $R$. For finitely generated modules $M,M'$ over $A$ and $N,N'$
over $B$, the tensor product of free resolutions resolves $M\otimes_{\CC}N$
and the K\"unneth formula gives
\begin{equation}\label{eq:kunnethext}
  \Ext^{n}_{R}(M\otimes N,M'\otimes N')
  =\bigoplus_{i+j=n}\Ext^{i}_{A}(M,M')\otimes_{\CC}\Ext^{j}_{B}(N,N'),
\end{equation}
compatibly with Yoneda products up to the Koszul sign, since a product of
chain maps on the tensor product of resolutions is the tensor product of the
products.

The building block is one plane. The Koszul complex
$0\to A\xrightarrow{(-x_{2},x_{1})^{t}}A^{2}\xrightarrow{(x_{1},x_{2})}\mathfrak m_{A}\to0$
is a free resolution, so $\mathfrak m_{A}$ has projective dimension one and
$\Ext^{i}_{A}(\mathfrak m_{A},-)=0$ for $i\ge2$; and
$\Ext^{1}_{A}(\mathfrak m_{A},\mathfrak m_{A})=\coker\bigl(\Hom(A^{2},\mathfrak m_{A})
\to\Hom(A,\mathfrak m_{A})\bigr)=\mathfrak m_{A}/\mathfrak m_{A}^{2}$, a
two-dimensional space killed by $\mathfrak m_{A}$. On a complex of free
modules with differential $d$ the jet extension along a derivation $D$ is the
complex with terms $P_{i}\oplus P_{i}$ and differential
$\bigl(\begin{smallmatrix}d&-D(d)\\0&d\end{smallmatrix}\bigr)$, $D(d)$ the
entrywise derivative, because the twisted module $J_{D}(P_{i})$ is free on the
elements $(e,0)$ and $(0,e)$ and $d(0,e)=\sum_{l}d_{le}(0,e_{l})-\sum_{l}D(d_{le})(e_{l},0)$;
so the jet class is represented by the degree $-1$ chain map $-D(d)$. For
$D=\partial_{i}$ on the Koszul complex this is $e\mapsto\delta_{i2}f_{1}-\delta_{i1}f_{2}$
on the generator $e$ of the last term, whose image in
$\mathfrak m_{A}/\mathfrak m_{A}^{2}$ is $\delta_{i2}x_{1}-\delta_{i1}x_{2}$:
\[
  \alpha_{1}:=a_{\partial_{1}}=-[x_{2}],\qquad
  \alpha_{2}:=a_{\partial_{2}}=[x_{1}],
\]
a basis of $\Ext^{1}_{A}(\mathfrak m_{A},\mathfrak m_{A})$, which we identify
with $N_{1}$ by $\partial_{i}\mapsto\alpha_{i}$. Similarly $\beta_{3},\beta_{4}$
form a basis of $\Ext^{1}_{B}(\mathfrak m_{B},\mathfrak m_{B})$, identified
with $N_{2}$.

(i) By \eqref{eq:kunnethext} with $M=M'=\mathfrak m_{A}$, $N=N'=\mathfrak m_{B}$,
\[
  \Ext^{2}_{R}(I,I)=\Ext^{1}_{A}\otimes\Ext^{1}_{B}=N_{1}\otimes N_{2},
\]
the other two summands vanishing by the projective dimension, and it is killed
by $\mathfrak m_{A}$ and $\mathfrak m_{B}$. The derivation $\partial_{1}$ acts on
the factor $A$, so the jet extension of $I=\mathfrak m_{A}\otimes\mathfrak m_{B}$
along it is $J_{\partial_{1}}(\mathfrak m_{A})\otimes\mathfrak m_{B}$ and
$a_{\partial_{1}}=\alpha_{1}\otimes1$; likewise $a_{\partial_{2}}=\alpha_{2}\otimes1$,
$a_{\partial_{3}}=1\otimes\beta_{3}$, $a_{\partial_{4}}=1\otimes\beta_{4}$.
Hence $a_{\partial_{i}}a_{\partial_{j}}=\pm\alpha_{i}\otimes\beta_{j}$ for
$i\le2<j$, the four of them a basis of $N_{1}\otimes N_{2}$, while
$a_{\partial_{1}}a_{\partial_{2}}=\alpha_{1}\alpha_{2}\otimes1=0$ and
$a_{\partial_{3}}a_{\partial_{4}}=0$ because $\Ext^{2}_{A}(\mathfrak m_{A},\mathfrak m_{A})=0$.
For germs $\theta=\sum_{i}c^{i}\partial_{i}$ and $\theta'=\sum_{j}c'^{j}\partial_{j}$
the linearity of \Cref{lem:edge} and the triviality of the module structure
give $a_{\theta}a_{\theta'}=\sum_{i,j}c^{i}(0)c'^{j}(0)\,a_{\partial_{i}}a_{\partial_{j}}$,
the image of $\theta(0)\wedge\theta'(0)$ under the projection of
$\bigwedge^{2}T_{0}$ onto $N_{1}\otimes N_{2}$ followed by the isomorphism
$\partial_{i}\otimes\partial_{j}\mapsto\pm\alpha_{i}\otimes\beta_{j}$.

(ii) Let $Q=\cO_{P_{1}}\oplus\cO_{P_{2}}$. The map $\cO\to Q$ has kernel $I$
and cokernel $\cO/(I_{P_{1}}+I_{P_{2}})=k$, so $K$ sits in an exact triangle
$I\xrightarrow{\iota}K\xrightarrow{\pi}k[-1]\xrightarrow{\varepsilon}I[1]$.
The jet extension along $\theta$ is an exact functor, natural in its
argument, so the jet classes commute with $\iota$ and $\pi$; and for the
residue field $a^{k}_{\theta}$ depends only on $\theta(0)$, because
$\Ext^{1}(k,k)$ is killed by $\mathfrak m$ and the class is linear in
$\theta$. Therefore
\[
  \pi[2]\circ a^{K}_{\theta}a^{K}_{\theta'}
  =a^{k}_{\theta}a^{k}_{\theta'}\circ\pi
  =\pi^{*}\bigl(a^{k}_{\theta(0)}a^{k}_{\theta'(0)}\bigr),
  \qquad\pi^{*}\colon\Ext^{2}(k,k)\to\Ext^{1}(K,k),
\]
and it suffices to show that the right hand side is nonzero. The long exact
sequence of the triangle for $\Hom(-,k)$ reads
\[
  \Hom(I,k)\xrightarrow{\;\varepsilon^{*}\;}\Ext^{2}(k,k)\xrightarrow{\;\pi^{*}\;}
  \Ext^{1}(K,k),
\]
so $\ker\pi^{*}=\image\varepsilon^{*}$. Now $k=k_{A}\otimes k_{B}$ and, by
\eqref{eq:kunnethext}, $\Ext^{*}(k,k)$ is the tensor product of
$\Ext^{*}_{A}(k,k)$ and $\Ext^{*}_{B}(k,k)$. On the Koszul resolution
$A\to A^{2}\to A$ of $k_{A}$ the jet chain maps of $\partial_{1},\partial_{2}$
are $-D(d)$ as above, and their composite $A\to A$ is the scalar $-1$, whose
class in $\Ext^{2}_{A}(k,k)=\Hom(A,k)=k$ is nonzero; the degree one classes
$\gamma_{1},\gamma_{2}$ of $\partial_{1},\partial_{2}$ are the basis
$-e_{1}^{*},-e_{2}^{*}$ of $\Ext^{1}_{A}(k,k)=\Hom(A^{2},k)$, and their
product spans $\Ext^{2}_{A}(k,k)$. With the analogous classes
$\gamma_{3},\gamma_{4}$ over $B$, the products $\gamma_{1}\gamma_{2}\otimes1$,
$\gamma_{i}\otimes\gamma_{j}$ ($i\le2<j$) and $1\otimes\gamma_{3}\gamma_{4}$
form a basis of $\Ext^{2}(k,k)=k\oplus(k^{2}\otimes k^{2})\oplus k$, so
$\partial_{i}\wedge\partial_{j}\mapsto a^{k}_{\partial_{i}}a^{k}_{\partial_{j}}$
identifies $\bigwedge^{2}(N_{1}\oplus N_{2})$ with $\Ext^{2}(k,k)$, and
$a^{k}_{\theta(0)}a^{k}_{\theta'(0)}$ is the image of $\theta(0)\wedge\theta'(0)$. The
class $\varepsilon$ is the K\"unneth product $\varepsilon_{A}\otimes\varepsilon_{B}$
of the classes of $0\to\mathfrak m_{A}\to A\to k\to0$ and its $B$-analogue,
because $K$ is the two-step truncation of the tensor product of the complexes
$[A\to k_{A}]$ and $[B\to k_{B}]$ and $\varepsilon$ is the inclusion of the
omitted term $k_{A}\otimes k_{B}$. Hence $\varepsilon^{*}$ maps
$\Hom(I,k)=\Hom(\mathfrak m_{A},k)\otimes\Hom(\mathfrak m_{B},k)$ into
$\Ext^{1}_{A}(k,k)\otimes\Ext^{1}_{B}(k,k)=N_{1}\otimes N_{2}$, and onto it,
since $\Hom(\mathfrak m_{A},k)\to\Ext^{1}_{A}(k,k)$ is an isomorphism: it is
the connecting homomorphism of the sequence, injective because
$\Hom(A,k)\to\Hom(\mathfrak m_{A},k)$ is zero and surjective because
$\Ext^{1}_{A}(A,k)=0$. So $\ker\pi^{*}=N_{1}\otimes N_{2}$
and $\pi^{*}(\theta(0)\wedge\theta'(0))\ne0$ exactly when
$\theta(0)\wedge\theta'(0)\notin N_{1}\otimes N_{2}$.

Both computations are repeated by machine in item (XXIX) of
\Cref{sec:verification}, which also finds $\Ext^{2}_{\cO}(K,K)\cong
\bigwedge^{2}N_{1}\oplus\bigwedge^{2}N_{2}$ with $a_{\theta}a_{\theta'}$ the
projection of $\theta(0)\wedge\theta'(0)$ to it, the exact complement of (i).
\end{proof}

\begin{theorem}[Markman's candidate does not satisfy the weakened criterion]
\label{thm:candidatefails}
In \Cref{setup:markman}, for every odd $d\ge3$, every choice of the
translates $D_{i}$ in general position and every choice of the $m$ points to
be normalised:
\begin{enumerate}[label=\textup{(\roman*)},nosep]
\item the identity \eqref{eq:heisenberg} fails: there is a constant bivector
  $\pi\in H^{0}(\bigwedge^{2}T_{X})$ with
  $\tfrac12\pi\lrcorner\bigl(\mathrm{At}(\cF)^{2}+d\,\Theta^{2}\cdot\id_{\cF}\bigr)\ne0$
  in $\Ext^{2}(\cF,\cF)$;
\item $\cF$ does not satisfy the weakened criterion: the class
  $x_{\cF}=(\tfrac d2\,\pi\lrcorner\,\Theta^{2},0,\pi)$ has
  $x_{\cF}\lrcorner\,\ch(\cF)=0$ and $\operatorname{ev}_{\cF}(x_{\cF})\ne0$;
\item for every perfect complex $F_{2}$ on $X$ with Chern character on the
  secant plane, in particular for $F_{2}=\cF$ and $F_{2}=\cF^{\vee}$, the
  object $E=\Phi(\cF\boxtimes F_{2})$ on the eightfold $X\times\widehat X$ of
  Weil type does not satisfy the weakened criterion of
  \cite[Question 11.4]{Mar25b}: the class $x=\Phi^{HT}(x_{\cF}\otimes1)$
  satisfies $x\lrcorner\,\ch(E)=0$ and $\operatorname{ev}_{E}(x)\ne0$, so
  $\sigma_{E}$ is not injective on the image of $\operatorname{ev}_{E}$.
\end{enumerate}
\end{theorem}

\begin{proof}
(i) The surfaces $Z_{i}$ and $Z_{j}$ at cyclic distance at least three meet
in the $24$ points of $D_{i}\cap D_{i+1}\cap D_{j}\cap D_{j+1}$, and there are
$12N(N-5)=(d+9)(3d-3)$ such points, of which $m=(d+9)(2d-1)$ are normalised;
the remaining $(d+9)(d-2)\ge12$ are not. Let $p$ be one of them. The four
divisors through $p$ meet transversally, so their local equations
$x_{1},\dots,x_{4}$ form a regular system of parameters of $\cO_{X,p}$ and
identify its completion with $\CC[[x_{1},\dots,x_{4}]]$, with
$D_{i}=\{x_{1}=0\}$, $D_{i+1}=\{x_{2}=0\}$, $D_{j}=\{x_{3}=0\}$,
$D_{j+1}=\{x_{4}=0\}$, hence $Z_{i}=P_{1}$ and $Z_{j}=P_{2}$ in the
notation of \Cref{lem:localproducts}; the completion is flat, so a class
whose completed germ is nonzero is nonzero. No other
$D_{k}$ passes through $p$, five divisors having empty intersection, so no
other component of $Z$ does. Near $p$ the map $\nu$ is an isomorphism,
$\nu_{*}\cO_{\widetilde Z}=\cO_{Z}$, and $\cF$ is the ideal sheaf
$I_{Z}=I_{P_{1}}\cap I_{P_{2}}$ twisted by the locally trivial $\cO_{X}(3\Theta)$:
the germ $\cF_{p}$ is the module $I$ of \Cref{lem:localproducts}(i).

Let $\theta,\theta'\in H^{0}(T_{X})$ be constant vector fields whose values
at $p$ are $\partial_{1}$ and $\partial_{3}$; they exist because the constant
fields span every tangent space. Put $\pi=\theta\wedge\theta'$. Since
$\operatorname{ev}_{\cF}$ is an algebra homomorphism,
$\operatorname{ev}_{\cF}(\pi)=\tfrac12\pi\lrcorner\,\mathrm{At}(\cF)^{2}
=A_{\theta}A_{\theta'}$, and by \Cref{lem:edge} its germ at $p$ is
$a_{\theta}a_{\theta'}$, which is nonzero by \Cref{lem:localproducts}(i),
$\partial_{1}\wedge\partial_{3}$ having nonzero component in $N_{1}\otimes N_{2}$.
The germ of $\tfrac d2(\pi\lrcorner\,\Theta^{2})\cdot\id_{\cF}$ at $p$ is zero
by the same lemma. So the germ at $p$ of
$\tfrac12\pi\lrcorner(\mathrm{At}(\cF)^{2}+d\Theta^{2}\cdot\id_{\cF})$ is
$a_{\theta}a_{\theta'}\ne0$, and the class itself is nonzero.

(ii) By \Cref{thm:factor}(ii) the class $x_{\cF}$ preserves $\ch(\cF)$, and by
\eqref{eq:evgeneral} $\operatorname{ev}_{\cF}(x_{\cF})$ is the class of (i).

(iii) $\ch(\cF^{\vee})=u_{\Theta}-3v_{\Theta}$ lies on the secant plane. By
\Cref{thm:factor}(iv), $E$ satisfies the criterion only if both factors do,
and $\cF$ does not; the displayed class is the transport of
$x_{\cF}\otimes1\in\Ann_{HT^{2}}(\ch(\cF\boxtimes F_{2}))$, on which
$\operatorname{ev}$ is $\pm\operatorname{ev}_{\cF}(x_{\cF})\otimes\id_{F_{2}}\ne0$.
The last assertion is \Cref{setup:criterion}:
$\sigma_{E}(\operatorname{ev}_{E}(x))=x\lrcorner\,\ch(E)=0$ with
$\operatorname{ev}_{E}(x)\ne0$.
\end{proof}

\begin{remark}[The normalised points obstruct as well]\label{rem:normalisedpoints}
At a normalised point the germ of $\cF$ is the complex $K$ of
\Cref{lem:localproducts}(ii), and the same argument with $\theta(p)=\partial_{1}$,
$\theta'(p)=\partial_{2}$ produces a nonzero germ. So the failure does not
depend on which points are separated, and separating all of them would not
help even if it were allowed: normalising every isolated point changes
$\chi(\cO_{\widetilde Z})$ to $110N-12N^{2}+12N(N-5)=50N$, while the secant
plane requires $(d+9)(27-d)=72N-4N^{2}$, and the two agree only for
$N=11/2$. The two local computations are complementary: at a point of type
(i) the obstruction lives in $N_{1}\otimes N_{2}$, at a point of type (ii) in
$\bigwedge^{2}N_{1}\oplus\bigwedge^{2}N_{2}$, and together they exhaust
$\bigwedge^{2}T_{p}$.
\end{remark}

\begin{figure}[tbp]
\centering
  \includegraphics[width=\textwidth]{fig_doublepoint}
\caption{\Cref{lem:localproducts} and \Cref{thm:candidatefails}. Left: two
smooth surfaces of a fourfold meeting transversally at $p$, with the normal
spaces $N_{1}$ and $N_{2}$, each tangent to the other branch. Centre: the
bivectors at $p$, split into $\bigwedge^{2}N_{1}$, $N_{1}\otimes N_{2}$ and
$\bigwedge^{2}N_{2}$, with heights proportional to the dimensions. Right: the
local $\Ext^{2}$ of the germ of $\cF$ at $p$ in the two cases, and the part
of the bivectors that the products of the translation classes carry onto it:
the mixed part when the point is not normalised, the two same-plane parts
when it is. The compensating scalar has no germ, so in either case the
identity \eqref{eq:heisenberg} fails at $p$.}
\label{fig:doublepoint}
\end{figure}

\begin{corollary}[A local obstruction in every dimension]\label{cor:localobstruction}
Let $(X,t)$ be a polarised abelian variety of dimension $n\ge4$ and $F$ a
perfect complex on $X$ with $\ch(F)=au_{t}+bv_{t}$, $(a,b)\ne(0,0)$. Suppose
that at some point $p$ the germ $F_{p}$ is isomorphic, in the derived
category of $\cO_{X,p}$-modules, to the ideal of the union of two smooth germs
of codimension two meeting transversally along a germ of codimension four,
or, when $n=4$, to the complex $[\cO\to\cO_{P_{1}}\oplus\cO_{P_{2}}]$ of two
such germs. Then $F$ does not satisfy the weakened criterion, and neither does
$\Phi(F\boxtimes F_{2})$ on $X\times\widehat X$ for any secant $F_{2}$. In
particular, in every dimension $n\ge4$, the union $Z$ of the intersections
$D_{i}\cap D_{i+1}$ of consecutive translates of a polarisation in general
position contains, for every pair $(i,j)$ of nonconsecutive indices, a
subvariety $Z_{i}\cap Z_{j}$ of codimension four along which $Z$ has two
transversal branches; any secant object of the form $[\cO_{X}\to\nu_{*}\cO_{\widetilde Z}]\otimes L$
in which $\nu$ is an isomorphism at a general point of one such
$Z_{i}\cap Z_{j}$ fails the criterion.
\end{corollary}

\begin{proof}
In formal coordinates $x_{1},\dots,x_{n}$ at $p$ with the two branches
$\{x_{1}=x_{2}=0\}$ and $\{x_{3}=x_{4}=0\}$, which exist because the four
local equations of the branches form part of a regular system of
parameters, the ideal is
$I\otimes_{\CC}\CC[x_{5},\dots,x_{n}]$ with $I$ as in \Cref{lem:localproducts},
and \eqref{eq:kunnethext} with the third factor gives
$\Ext^{2}(I_{p},I_{p})=N_{1}\otimes N_{2}\otimes\cO_{W,p}$, $W$ the
intersection of the branches. The jet classes along $\partial_{5},\dots,\partial_{n}$
vanish, these derivations acting on a free factor, and the products of the
classes along $\partial_{1},\dots,\partial_{4}$ are as in the lemma; reducing
modulo the maximal ideal of $\cO_{W,p}$ shows that $a_{\theta}a_{\theta'}\ne0$
whenever $\theta(p)\wedge\theta'(p)$ has a nonzero component in
$N_{1}\otimes N_{2}$. The rest of the argument is that of
\Cref{thm:candidatefails}, using \Cref{thm:factor} for $n\ge3$. For the last
assertion, $Z_{i}\cap Z_{j}=D_{i}\cap D_{i+1}\cap D_{j}\cap D_{j+1}$ is a
transversal intersection of four divisors, of codimension four, and at a
general point of it no fifth divisor passes, so $Z$ has exactly the two
branches $Z_{i}$ and $Z_{j}$ there.
\end{proof}

\begin{remark}[What this settles]\label{rem:candidatemeaning}
\Cref{thm:candidatefails} answers the question left open in
\cite[Section 12]{Mar25b}: the secant objects proposed there do not satisfy
the weaker criterion of Question 11.4, for any discriminant. The reason is
local and geometric. Question 11.4 asks whether the semiregularity map may be
required to be injective only on the image of the evaluation map; for a
secant object that image contains, by \Cref{thm:factor}, the products
$A_{\theta}A_{\theta'}$ of the translation classes shifted by the scalars
$\tfrac d2(\theta\wedge\theta')\lrcorner\,\Theta^{2}$, and the semiregularity
map kills them all. For the criterion to hold these shifted products must
vanish, and they cannot: at an isolated double point of the support, the
translation classes of the two branches multiply to a nonzero local class,
while a scalar has no local class at all. What survives of the fourth
ingredient of \Cref{rem:deformation} is therefore sharper than before. A
secant object of positive rank on a fourfold that satisfies the weakened
criterion must have, at every point, germs whose translation classes multiply
to zero in the local $\Ext^{2}$; by \Cref{lem:localproducts} this excludes
every support with two transversal branches of codimension two through a
point, whether or not the branches are separated, and by
\Cref{cor:localobstruction} it excludes the chain of consecutive
intersections of translates in every dimension. An object supported on a
smooth surface, or on surfaces meeting only along curves, is not excluded by
this argument, but no such object with Chern character on the secant plane
is known.
\end{remark}

\subsection{Smooth support, and the four discriminants}
\label{ssec:smoothsupport}

\Cref{thm:candidatefails} and \Cref{cor:localobstruction} remove every secant
object whose support carries two transversal branches of codimension two
through a point. What they leave is a support that is smooth, or that meets
itself only along curves, and \Cref{rem:candidatemeaning} closes with the
remark that no such object is known. This subsection says why the local
argument cannot reach that case, and then determines what the case can be.
The answer is narrow: on the fourfold the invariants of the support are
forced outright, and the discriminant is confined to a finite list.

The first point is that at a smooth point there is nothing local left to
compute.

\begin{proposition}[No local obstruction at a smooth point]\label{prop:smoothlocal}
Let $Y$ be a smooth variety with $\omega_{Y}\cong\cO_{Y}$ and let $S\subset Y$
be smooth of codimension two. Then
\[
  \cE xt^{0}(\cI_{S},\cI_{S})=\cO_{Y},\qquad
  \cE xt^{1}(\cI_{S},\cI_{S})=N_{S/Y},\qquad
  \cE xt^{q}(\cI_{S},\cI_{S})=0\quad(q\ge2),
\]
and the same holds for $\cI_{S}\otimes L$ for any line bundle $L$.
Consequently, for $F=\cI_{S}\otimes L$ both sides of \eqref{eq:heisenberg}
have vanishing germs in $\cE xt^{2}$, and \Cref{lem:edge} decides nothing.
\end{proposition}

\begin{proof}
Twisting by $L$ changes neither side, so take $F=\cI_{S}$. The sheaf
$\cI_{S}$ is torsion free of rank one and reflexive, whence
$\cE xt^{0}(\cI_{S},\cI_{S})=\cO_{Y}$. Apply $\cH om(-,\cI_{S})$ to
\begin{equation}\label{eq:idealseq}
  0\to\cI_{S}\to\cO_{Y}\to\cO_{S}\to0 .
\end{equation}
Since $\cE xt^{q}(\cO_{Y},\cI_{S})=0$ for $q\ge1$, the long exact sequence
gives
\begin{equation}\label{eq:shift}
  \cE xt^{q}(\cI_{S},\cI_{S})\cong\cE xt^{q+1}(\cO_{S},\cI_{S})
  \qquad(q\ge1).
\end{equation}
Now apply $\cH om(\cO_{S},-)$ to \eqref{eq:idealseq}. Because $S$ is smooth of
codimension two, hence a local complete intersection, the Koszul resolution
gives $\cE xt^{k}(\cO_{S},\cO_{Y})=0$ for $k\ne2$ and
$\cE xt^{2}(\cO_{S},\cO_{Y})=\omega_{S}\otimes\omega_{Y}^{-1}=\omega_{S}$,
while $\cE xt^{k}(\cO_{S},\cO_{S})=\bigwedge^{k}N_{S/Y}$ for $0\le k\le2$ and
zero otherwise. Computing both from the same Koszul complex shows that the
map
\[
  \rho\colon\cE xt^{2}(\cO_{S},\cO_{Y})=\omega_{S}
  \longrightarrow
  \cE xt^{2}(\cO_{S},\cO_{S})=\textstyle\bigwedge^{2}N_{S/Y}=\omega_{S}
\]
induced by $\cO_{Y}\to\cO_{S}$ is the identity: in local coordinates in which
$S=\{x_{1}=x_{2}=0\}$, both groups are computed as $M/(x_{1},x_{2})M$ for
$M=\cO_{Y}$ and $M=\cO_{S}$, and the map is reduction modulo $(x_{1},x_{2})$,
which is the identity on $\cO_{S}$.

Take $k=2$ in the long exact sequence of $\cH om(\cO_{S},-)$. Since $\rho$ is
surjective and $\cE xt^{3}(\cO_{S},\cO_{Y})=0$, the segment
\[
  \cE xt^{2}(\cO_{S},\cO_{Y})\xrightarrow{\ \rho\ }
  \cE xt^{2}(\cO_{S},\cO_{S})\to
  \cE xt^{3}(\cO_{S},\cI_{S})\to
  \cE xt^{3}(\cO_{S},\cO_{Y})=0
\]
gives $\cE xt^{3}(\cO_{S},\cI_{S})=0$, so $\cE xt^{2}(\cI_{S},\cI_{S})=0$ by
\eqref{eq:shift}. For $k\ge3$ both $\cE xt^{k}(\cO_{S},\cO_{Y})$ and
$\cE xt^{k}(\cO_{S},\cO_{S})$ vanish, so
$\cE xt^{k+1}(\cO_{S},\cI_{S})=0$ and $\cE xt^{k}(\cI_{S},\cI_{S})=0$. Taking
$k=1$ and using that $\rho$ is injective, the segment
\[
  0=\cE xt^{1}(\cO_{S},\cO_{Y})\to
  \cE xt^{1}(\cO_{S},\cO_{S})=N_{S/Y}\to
  \cE xt^{2}(\cO_{S},\cI_{S})\to
  \cE xt^{2}(\cO_{S},\cO_{Y})\xrightarrow{\ \rho\ }
\]
exhibits $\cE xt^{2}(\cO_{S},\cI_{S})\cong N_{S/Y}$, which is
$\cE xt^{1}(\cI_{S},\cI_{S})$ by \eqref{eq:shift}.

For the last assertion, the right side of \eqref{eq:heisenberg} has vanishing
germs by \Cref{lem:edge}, and the left side lies in
$\Ext^{2}(F,F\otimes\Omega^{2}_{Y})$ whose sheaf of germs in the
$\cE xt^{2}$ direction is $\cE xt^{2}(F,F)\otimes\Omega^{2}_{Y}=0$.
\end{proof}

\begin{remark}[Where the obstruction goes instead]\label{rem:obstructionmoves}
\Cref{prop:smoothlocal} is the precise reason the argument of
\Cref{thm:candidatefails} stops at a smooth support, and it also says what
replaces it. With $\cE xt^{2}=0$ the local to global spectral sequence for
$\Ext^{2}(\cI_{S},\cI_{S})$ has only two surviving rows, $H^{1}(S,N_{S/Y})$
and $H^{2}(Y,\cO_{Y})$, and they sit in different steps of the filtration: the
first is the quotient that records the geometry of $S$, the second the
subgroup of scalars. A product $A_{\theta}A_{\theta'}$ of translation classes
therefore has a well defined image in $H^{1}(S,N_{S/Y})$, while the scalar on
the right of \eqref{eq:heisenberg} has image zero there. So the criterion
splits into two conditions of different kinds: a vanishing in
$H^{1}(S,N_{S/Y})$, which is a statement about the normal bundle and not about
any point, and an identity in $H^{2}(Y,\cO_{Y})$. At a double point the first
of these was already inconsistent locally. At a smooth support it is a global
question, and we do not decide it. What we can decide is which $S$ are even
available.
\end{remark}

The invariants of a smooth support are not free. On an abelian variety the
Todd class is trivial, so Grothendieck--Riemann--Roch turns the Chern
character of the object into the complete numerical package of the surface,
and three classical inequalities then close in on it.

\begin{theorem}[The invariants of a smooth support]\label{thm:smoothinvariants}
Let $(X,\Theta)$ be a principally polarised abelian fourfold, let $d>0$, let
$b\in\ZZ_{>0}$, and let $\cF=\cI_{S}(b\Theta)$ be a coherent sheaf with
\[
  \ch(\cF)=u_{\Theta}+b\,v_{\Theta}
\]
and $S\subset X$ smooth. Write $N=\tfrac12\Nm_{K/\QQ}(b+\sqrt{-d})
=\tfrac12(b^{2}+d)$, $\iota\colon S\hookrightarrow X$ and
$\lambda=\iota^{*}\Theta$. Then $S$ is a surface and
\begin{equation}\label{eq:smoothinvariants}
\begin{aligned}
  [S]&=N\Theta^{2}, &\qquad \iota_{*}K_{S}&=\tfrac{4bN}{3}\,\Theta^{3},
  &\qquad \chi(\cO_{S})&=4N(2b^{2}-N),\\
  K_{S}^{2}&=12N(4b^{2}-N), &\qquad e(S)&=12N(4b^{2}-3N),
  &\qquad K_{S}\cdot\lambda&=32bN,
\end{aligned}
\end{equation}
with $\lambda^{2}=24N$ and $[S]^{2}=24N^{2}$. Moreover $K_{S}$ is globally
generated, hence nef; $S$ is minimal of general type; and
\begin{equation}\label{eq:threebounds}
  \tfrac{4}{9}b^{2}\ \le\ N\ \le\ b^{2},
  \qquad\text{equivalently}\qquad
  -\tfrac{1}{9}b^{2}\ \le\ d\ \le\ b^{2} .
\end{equation}
The lower bound is the Hodge index theorem and is implied by $d>0$; the upper
bound is the Bogomolov--Miyaoka--Yau inequality. In terms of $b$ and $d$ the
five invariants read
\begin{equation}\label{eq:bdform}
  \chi(\cO_{S})=(b^{2}+d)(3b^{2}-d),\quad
  K_{S}^{2}=3(b^{2}+d)(7b^{2}-d),\quad
  e(S)=3(b^{2}+d)(5b^{2}-3d),
\end{equation}
and the defect in the Bogomolov--Miyaoka--Yau inequality is
\begin{equation}\label{eq:bmydefect}
  3e(S)-K_{S}^{2}=24\,(b^{4}-d^{2}),
\end{equation}
so that the inequality is exactly $d\le b^{2}$ and nothing else.
Finally $K_{S}$ is never numerically proportional to $\lambda$.
\end{theorem}

\begin{proof}
Since $\ch_{0}(\cF)=1$ and $c_{1}(\cF)=b\Theta$, the sheaf
$\cF(-b\Theta)=\cI_{S}$ has $\ch_{0}=1$, $\ch_{1}=0$, so $S$ has codimension
at least two, and
\[
  \ch(\cO_{S})=1-\ch(\cI_{S})=1-\ch(\cF)\,e^{-b\Theta}.
\]
Expanding the right side in $\QQ[\Theta]/(\Theta^{5})$ and collecting gives
\[
  \ch_{2}(\cO_{S})=\tfrac{b^{2}+d}{2}\Theta^{2}=N\Theta^{2},\quad
  \ch_{3}(\cO_{S})=-\tfrac{2bN}{3}\Theta^{3},\quad
  \ch_{4}(\cO_{S})=\tfrac{3b^{4}+2db^{2}-d^{2}}{24}\Theta^{4},
\]
which is the computation of item (XXX) of \Cref{sec:verification}. The first
identity says $S$ has codimension exactly two, so it is a surface, with class
$N\Theta^{2}$. Because $T_{X}$ is trivial, $\td(T_{X})=1$, and
Grothendieck--Riemann--Roch for $\iota$ reads
$\ch(\iota_{*}\cO_{S})=\iota_{*}\td(T_{S})$. Comparing degrees three and four,
\[
  \ch_{3}(\cO_{S})=\iota_{*}\tfrac{c_{1}(T_{S})}{2}=-\tfrac12\iota_{*}K_{S},
  \qquad
  \ch_{4}(\cO_{S})=\iota_{*}\tfrac{c_{1}^{2}+c_{2}}{12}
  =\chi(\cO_{S})\,[\mathrm{pt}],
\]
the second by Noether's formula. With $\Theta^{4}=24\,[\mathrm{pt}]$ these give
$\iota_{*}K_{S}=\tfrac{4bN}{3}\Theta^{3}$ and, after substituting
$d=2N-b^{2}$, $\chi(\cO_{S})=3b^{4}+2db^{2}-d^{2}=4N(2b^{2}-N)$.
Then $K_{S}\cdot\lambda=\int_{X}\iota_{*}K_{S}\cdot\Theta=\tfrac{4bN}{3}\cdot24=32bN$
and $\lambda^{2}=\int_{X}[S]\cdot\Theta^{2}=24N$.

For the last two invariants, the normal bundle sequence and the triviality of
$T_{X}$ give $c(T_{S})c(N_{S/X})=1$, so $c_{1}(N_{S/X})=K_{S}$ and
$c_{2}(N_{S/X})=K_{S}^{2}-e(S)$; the self-intersection formula
$[S]^{2}=\int_{S}c_{2}(N_{S/X})$ then reads
\[
  24N^{2}=K_{S}^{2}-e(S),
\]
while Noether's formula reads $12\chi(\cO_{S})=K_{S}^{2}+e(S)$. Solving the
two gives the stated $K_{S}^{2}$ and $e(S)$.

The conormal sequence presents $\Omega^{1}_{S}$ as a quotient of
$\Omega^{1}_{X}|_{S}=\cO_{S}^{\oplus4}$, so $\Omega^{1}_{S}$ is globally
generated, and therefore so is $K_{S}=\det\Omega^{1}_{S}$, a quotient of
$\bigwedge^{2}\cO_{S}^{\oplus4}$. A globally generated line bundle is nef. A
globally generated rank two bundle on a surface has $c_{2}\ge0$, the class
$c_{2}$ being represented by the vanishing scheme of a general section, which
is empty or finite; hence
\[
  e(S)=c_{2}(\Omega^{1}_{S})=12N(4b^{2}-3N)\ \ge\ 0,
\]
so $N\le\tfrac43b^{2}<4b^{2}$ and therefore $K_{S}^{2}=12N(4b^{2}-N)>0$. A nef
divisor of positive self-intersection is big, so $\kappa(S)=2$ and, $K_{S}$
being nef, $S$ is minimal of general type. The Bogomolov--Miyaoka--Yau
inequality $c_{1}^{2}\le3c_{2}$ then applies and reads
$K_{S}^{2}\le3e(S)$. Substituting $N=\tfrac12(b^{2}+d)$ in the two
expressions gives \eqref{eq:bdform}, hence
$3e(S)-K_{S}^{2}=24(b^{4}-d^{2})$, and the inequality is $d\le b^{2}$, that
is $N\le b^{2}$.

For the lower bound, $\lambda$ is ample on $S$, so the Hodge index theorem
gives $(K_{S}\cdot\lambda)^{2}\ge K_{S}^{2}\,\lambda^{2}$, that is
$(32bN)^{2}\ge12N(4b^{2}-N)\cdot24N$, which simplifies to $9N\ge4b^{2}$. Since
$N=\tfrac12(b^{2}+d)$ with $d>0$, this holds automatically. Equality in the
Hodge index theorem occurs exactly when $K_{S}$ is numerically proportional to
$\lambda$, hence exactly when $9N=4b^{2}$; as $9N=\tfrac92(b^{2}+d)>4b^{2}$ for
$d>0$, that never happens.
\end{proof}

\begin{corollary}[Four discriminants on the fourfold]\label{cor:fourdiscriminants}
In \Cref{setup:markman}, where $b=3$ by \Cref{prop:markman8}(iii), a secant
object $\cI_{S}(3\Theta)$ with smooth support exists only if $d\le9$; since
$d$ is squarefree this leaves $d\in\{1,2,3,5,6,7\}$, and if in addition
$N=(9+d)/2$ is an integer, which is what a support assembled from translates
of $\Theta$ requires, only
\[
  d\in\{1,\,3,\,5,\,7\},\qquad N\in\{5,\,6,\,7,\,8\}.
\]
For these the invariants are forced to the values in
\Cref{tab:smoothsupport}. In particular, for every imaginary quadratic field
of discriminant larger than $9$ there is no secant object of rank one with
smooth support on an abelian fourfold at all.
\end{corollary}

\begin{proof}
Immediate from \eqref{eq:threebounds} at $b=3$, together with the arithmetic
of \eqref{eq:smoothinvariants}; the table is item (XXX) of
\Cref{sec:verification}. The value $d=9$ is the equality case of
Bogomolov--Miyaoka--Yau, where $K_{S}^{2}=3e(S)=2916$ and $S$ would have to be
a quotient of the ball; it is excluded here because $9$ is not squarefree.
\end{proof}

\begin{table}[htbp]
\centering
\small
\renewcommand{\arraystretch}{1.12}
\begin{tabular}{@{}ccrrrrr@{}}
\toprule
$d$ & $N$ & $\chi(\cO_{S})$ & $K_{S}^{2}$ & $e(S)$ &
$K_{S}\cdot\lambda$ & $\lambda^{2}$ \\
\midrule
$1$ & $5$ & $260$ & $1860$ & $1260$ & $480$ & $120$ \\
$3$ & $6$ & $288$ & $2160$ & $1296$ & $576$ & $144$ \\
$5$ & $7$ & $308$ & $2436$ & $1260$ & $672$ & $168$ \\
$7$ & $8$ & $320$ & $2688$ & $1152$ & $768$ & $192$ \\
\bottomrule
\end{tabular}
\caption{The only invariants a smooth support can have on an abelian
fourfold, by \Cref{thm:smoothinvariants} and
\Cref{cor:fourdiscriminants}, at $b=3$ and with $N$ an integer. Every entry is
forced by $\ch(\cF)$ alone; none is a choice. The slope
$K_{S}^{2}/\chi(\cO_{S})$ runs from $93/13$ at $d=1$ to $42/5$ at $d=7$, and
reaches the Bogomolov--Miyaoka--Yau ceiling of $9$ only at the excluded value
$d=9$.}
\label{tab:smoothsupport}
\end{table}

\begin{figure}[htbp]
\centering
  \includegraphics[width=\textwidth]{fig_bmywindow}
\caption{\Cref{thm:smoothinvariants} over the quarter plane of $(b,d)$. The
height of the surface is the defect $3e(S)-K_{S}^{2}=24(b^{4}-d^{2})$ in the
Bogomolov--Miyaoka--Yau inequality for the surface the theorem attaches to the
pair, compressed by a signed cube root so that four orders of magnitude fit
one frame; the compression is odd and monotone, so the zero locus and the sign
are the true ones. The surface is drawn indigo where the defect is positive and
clay where it is negative, and it crosses zero along $d=b^{2}$, which is
therefore the whole content of the inequality. The weaker wall $d=\tfrac53
b^{2}$, where $e(S)$ itself changes sign, is drawn dashed and lies above it.
On the line $b=3$, which is the value \Cref{prop:markman8}(iii) forces on the
fourfold, the four squarefree discriminants $d=1,3,5,7$ with $N$ an integer sit
above the wall and are marked with their drop lines; $d=9$ sits exactly on it,
the equality case, and $d=11,13,15$ sit below and are excluded.}
\label{fig:bmywindow}
\end{figure}

\begin{remark}[What is left of the smooth case]\label{rem:smoothremains}
Three things have changed and one has not.

The local argument is out of reach here, and \Cref{prop:smoothlocal} says so
for a structural reason rather than as a failure of effort: the sheaf
$\cE xt^{2}$ of a smooth ideal vanishes, so both sides of
\eqref{eq:heisenberg} are locally zero and the germ comparison that decides
\Cref{thm:candidatefails} has nothing to compare.

The surface is pinned. Its class, its canonical class, its Euler
characteristic, its Euler number and its two intersection numbers against the
polarisation are all determined by $d$ alone, with no freedom, and they are
compatible with existence only for $d\le9$. That is the content of
\Cref{cor:fourdiscriminants}, and it is a genuine restriction: the secant
construction with smooth support is not available for all but finitely many
imaginary quadratic fields, whatever geometry one is willing to invent.

The residual question is now a single existence statement. For
$d\in\{1,3,5,7\}$, does a principally polarised abelian fourfold carry a smooth
surface $S$ with the invariants of \Cref{tab:smoothsupport}, and if so does
the image of $A_{\theta}A_{\theta'}$ in $H^{1}(S,N_{S/X})$ vanish for all
$\theta,\theta'\in H^{0}(T_{X})$? The first half is a question about surfaces
in abelian fourfolds with prescribed Chern numbers, of the kind that is
answered by construction or by a classification argument; the second is a
question about one cohomology group of one normal bundle. Neither is settled
here.

What has not changed is the standing of the conjecture. Even a positive answer
to both would produce Weil classes on abelian eightfolds for at most four
discriminants, and \Cref{cor:exactform} would still be open for every other
discriminant, for every $n\ge5$, and for every field of degree larger than
two.
\end{remark}

\subsection{The obstruction is a condition on depth}
\label{ssec:depth}

\Cref{prop:smoothlocal} says that at a smooth point of the support there is
nothing local to compare. That is a special case of something much more
general, and identifying the general case tells us exactly which supports
\Cref{thm:candidatefails} and \Cref{cor:localobstruction} remove and which
they cannot touch. The dividing line is not smoothness, and it is not being a
local complete intersection. It is depth.

\begin{theorem}[Cohen--Macaulay supports carry no local obstruction]%
\label{thm:cmvanishing}
Let $Y$ be a smooth variety with $\omega_{Y}\cong\cO_{Y}$, let $Z\subset Y$ be
a closed subscheme of codimension two, and let $p\in Z$. If the local ring
$\cO_{Z,p}$ is Cohen--Macaulay, then
\[
  \cE xt^{q}(\cI_{Z},\cI_{Z})_{p}=0\qquad\text{for every }q\ge2 ,
\]
and the same holds for $\cI_{Z}\otimes L$ for every line bundle $L$.
Consequently both sides of \eqref{eq:heisenberg} have vanishing germ at $p$,
and \Cref{lem:edge} decides nothing there.
\end{theorem}

\begin{proof}
Twisting by $L$ changes no $\cE xt$ sheaf, so take $F=\cI_{Z}$. Write
$A=\cO_{Y,p}$, a regular local ring of dimension $m$, and $I=\cI_{Z,p}$. By
hypothesis $A/I$ is Cohen--Macaulay of dimension $m-2$, so
$\operatorname{depth}A/I=m-2$, and the Auslander--Buchsbaum formula gives
\[
  \operatorname{pd}_{A}(A/I)=m-\operatorname{depth}_{A}(A/I)=2 .
\]
From the exact sequence $0\to I\to A\to A/I\to0$ with $A$ free it follows that
$\operatorname{pd}_{A}I=1$. A module of projective dimension one has a free
resolution of length one, so $\Ext^{q}_{A}(I,M)=0$ for every $A$-module $M$
and every $q\ge2$; taking $M=I$ gives the vanishing. The last assertion is
\Cref{lem:edge}: the right side of \eqref{eq:heisenberg} is
$z\cdot\id_{F}$ with $z\in H^{2}(\cO_{Y})$ and has vanishing germs, and the
left side lies in a sheaf whose $\cE xt^{2}$ factor is zero.
\end{proof}

\begin{remark}[The hypothesis is sharp, and it is the one that fails]%
\label{rem:cmsharp}
Two observations make \Cref{thm:cmvanishing} the right statement.

The first is that Cohen--Macaulayness is strictly weaker than the conditions
one might reach for instead. Item (XXXI) of \Cref{sec:verification} computes
$\cE xt^{2}(\cI_{Z},\cI_{Z})$ for ten codimension two germs in $\CC^{4}$. It
vanishes for the smooth germ $(x_{1},x_{2})$; for two planes meeting along a
line, $(x_{1},x_{2}x_{3})$, and for three planes meeting along a line,
$(x_{1},x_{2}x_{3}x_{4})$; for the double plane $(x_{1}^{2},x_{2})$ and the
quadric cone $(x_{1},x_{2}^{2}-x_{3}x_{4})$; and also for the fat plane
$(x_{1},x_{2})^{2}$ and for the cone over the twisted cubic, the ideal of
$2\times2$ minors of $\left(\begin{smallmatrix}x_{1}&x_{2}&x_{3}\\
x_{2}&x_{3}&x_{4}\end{smallmatrix}\right)$, neither of which is a local
complete intersection, both needing three generators in codimension two. So
reducedness, irreducibility and the complete intersection property are all
irrelevant; the projective dimension is what matters, and by
Auslander--Buchsbaum that is depth.

The second is that the germs where it does not vanish are exactly the ones
the construction of \Cref{setup:markman} produces. The same computation finds
$\cE xt^{2}\ne0$, with nonzero products of translation classes, for two planes
meeting at a point, for three planes meeting pairwise at a point, and for a
plane with an embedded point; the projective dimensions there are $3$, $3$ and
$4$. A scheme whose local ring is Cohen--Macaulay is connected in codimension
one at that point, by Hartshorne's connectedness theorem, so a support with
two components through $p$ meeting only at $p$ can never be Cohen--Macaulay
there. That is precisely the configuration \Cref{lem:localproducts} analyses.
\end{remark}

\begin{corollary}[What the local argument removes, in final form]%
\label{cor:cmdichotomy}
Let $(X,t)$ be a polarised abelian variety of dimension $n\ge4$ and $F$ a
perfect complex with $\ch(F)=au_{t}+bv_{t}$, $(a,b)\ne(0,0)$, whose germ at
each point is the germ of $\cI_{Z}\otimes L$ for a codimension two subscheme
$Z$. Then exactly one of the following holds.
\begin{enumerate}[label=\textup{(\roman*)},leftmargin=2.6em]
\item $Z$ fails to be Cohen--Macaulay at some point $p$ of one of the types of
  \Cref{rem:cmsharp}, and then $F$ does not satisfy the weakened criterion,
  nor does $\Phi(F\boxtimes F_{2})$ for any secant $F_{2}$, by
  \Cref{thm:candidatefails} and \Cref{cor:localobstruction}.
\item $Z$ is Cohen--Macaulay, and then there is no local obstruction at any
  point whatever. The criterion is then a single global condition: the image
  of $A_{\theta}A_{\theta'}$ in $H^{1}\bigl(Z,\cE xt^{1}(\cI_{Z},\cI_{Z})\bigr)$
  must vanish for all $\theta,\theta'\in H^{0}(T_{X})$, the identity in
  $H^{2}(\cO_{X})$ being all that the scalar on the right of
  \eqref{eq:heisenberg} can meet.
\end{enumerate}
In case (ii), $\cI_{Z}$ has projective dimension one at every point, so it is
locally the ideal of maximal minors of a Hilbert--Burch matrix, and when $Z$
is in addition smooth $\cE xt^{1}(\cI_{Z},\cI_{Z})=N_{Z/X}$ and the invariants
of $Z$ are those of \Cref{thm:smoothinvariants}.
\end{corollary}

\begin{proof}
The two cases are exclusive and exhaustive by definition. In case (i) the
local ring is not Cohen--Macaulay and the cited results apply, the germs being
those of \Cref{lem:localproducts} and \Cref{cor:localobstruction}. In case
(ii) \Cref{thm:cmvanishing} gives $\cE xt^{q}=0$ for $q\ge2$ at every point, so
the local to global spectral sequence for $\Ext^{2}(\cI_{Z},\cI_{Z})$ has only
the two rows $H^{1}(\cE xt^{1})$ and $H^{2}(\cO_{X})$, and these lie in
different steps of the filtration, the second being the subgroup of scalars.
The right side of \eqref{eq:heisenberg} lies in that subgroup, so it has zero
image in the quotient $H^{1}(\cE xt^{1})$, and the criterion forces the image
of the left side there to vanish. The projective dimension statement is the
proof of \Cref{thm:cmvanishing}, and Hilbert--Burch is the structure theorem
for modules of projective dimension one over a regular local ring; the last
clause is \Cref{prop:smoothlocal} and \Cref{thm:smoothinvariants}.
\end{proof}

\begin{remark}[Why the normalisation cannot repair the candidate]%
\label{rem:normalisationcm}
\Cref{cor:cmdichotomy} explains the tension in \Cref{rem:normalisedpoints} in
one word. The partial normalisation performed in \Cref{setup:markman}
separates the two branches at an isolated double point, and separating them is
exactly what would make the local ring Cohen--Macaulay there. So the
construction is already trying to reach case (ii). What
\Cref{rem:normalisedpoints} shows is that it cannot arrive: normalising every
isolated point moves $\chi(\cO_{\widetilde Z})$ to $50N$, while the secant
plane demands $(d+9)(27-d)=72N-4N^{2}$, and the two agree only at $N=11/2$.
Repairing the depth destroys the Chern character, and keeping the Chern
character keeps a point where the depth fails. That is the obstruction in its
final form, and it is why a construction assembled from intersections of
translates of a polarisation cannot be made to work: by
\Cref{cor:localobstruction} such a chain always has a point of the bad type,
and by \Cref{rem:normalisedpoints} it cannot be separated without leaving the
plane.
\end{remark}

\subsection{Cohen--Macaulay supports, and the two routes that remain}
\label{ssec:hilbertburch}

\Cref{cor:cmdichotomy} leaves one class of supports and one condition on them.
The class is the Cohen--Macaulay subschemes of codimension two, and for those
the proof of \Cref{thm:cmvanishing} gives more than the vanishing it was used
for: the ideal sheaf has projective dimension one at every point. That is the
Hilbert--Burch shape, and it can be searched.

\begin{setupenv}[A resolution of the support]\label{setup:burch}
Let $(X,\Theta)$ be a principally polarised abelian fourfold with
$\NS(X)_{\QQ}=\QQ\Theta$, let $Z\subset X$ be Cohen--Macaulay of codimension
two, and suppose $\cI_{Z}$ admits a global resolution
\begin{equation}\label{eq:burchres}
  0\longrightarrow E_{1}\stackrel{\varphi}{\longrightarrow}E_{0}
  \longrightarrow\cI_{Z}\longrightarrow0,
  \qquad \rk E_{0}=r+1,\quad\rk E_{1}=r,
\end{equation}
with $E_{0}=\bigoplus_{i=1}^{r+1}\cO(e_{i}\Theta)$ and
$E_{1}=\bigoplus_{j=1}^{r}\cO(f_{j}\Theta)$ sums of line bundles. Put
$p_{i}=e_{i}+b$ and $q_{j}=f_{j}+b$.
\end{setupenv}

\begin{lemma}[The resolution in moments]\label{lem:burchmoments}
In \Cref{setup:burch}, $\ch(\cI_{Z}(b\Theta))=u_{\Theta}+b\,v_{\Theta}$ holds
if and only if
\begin{equation}\label{eq:burchsystem}
  \sum_{i=1}^{r+1}p_{i}^{\,k}-\sum_{j=1}^{r}q_{j}^{\,k}=c_{k},
  \qquad k=0,1,2,3,4,
  \qquad (c_{k})=(1,\;b,\;-d,\;-db,\;d^{2}).
\end{equation}
The equation at $k=0$ is the rank count and is automatic. Moreover
$\varphi$ can be injective only if $\max_{i}e_{i}\ge\max_{j}f_{j}$, and
$\cO_{X}$ embeds in $E_{0}^{\vee}$ only if $\min_{i}e_{i}\le0$.
\end{lemma}

\begin{proof}
From \eqref{eq:burchres}, $\ch(\cI_{Z})=\ch(E_{0})-\ch(E_{1})$, so
$\ch(\cI_{Z}(b\Theta))=\sum_{i}e^{p_{i}\Theta}-\sum_{j}e^{q_{j}\Theta}$.
Comparing coefficients of $\Theta^{k}/k!$ against $u_{\Theta}+bv_{\Theta}$,
whose coefficients are the $c_{k}$ displayed, gives
\eqref{eq:burchsystem}. For the last two statements, on an abelian variety
$\Hom(\cO(f\Theta),\cO(e\Theta))=H^{0}(\cO((e-f)\Theta))$ vanishes unless
$e\ge f$; if $\varphi$ is injective then each summand of $E_{1}$ must map
nonzero to $E_{0}$, which needs $\max_{i}e_{i}\ge f_{j}$ for every $j$.
Dualising \eqref{eq:burchres} and using $\cI_{Z}^{\vee}=\cO_{X}$ gives an
injection $\cO_{X}\hookrightarrow E_{0}^{\vee}$, which needs
$-e_{i}\ge0$ for some $i$.
\end{proof}

The case $r=1$ is the complete intersections, and it closes outright, for
every $b$ and every $d$, by an argument that uses no search.

\begin{theorem}[No complete intersection supports a secant object]%
\label{thm:noci}
Let $(X,\Theta)$ be a principally polarised abelian fourfold with
$\NS(X)_{\QQ}=\QQ\Theta$, let $d>0$ and $b\ge1$. No complete intersection
$Z=D_{1}\cap D_{2}$ of two divisors satisfies
$\ch(\cI_{Z}(b\Theta))=u_{\Theta}+b\,v_{\Theta}$.
\end{theorem}

\begin{proof}
Such a $Z$ is \Cref{setup:burch} with $r=1$, the resolution being the Koszul
complex of the two equations. Put $w(x)=x^{2}+d$, which is strictly positive
on $\RR$ because $d>0$, and read \eqref{eq:burchsystem} in the combinations
$c_{k+2}+d\,c_{k}$ for $k=0,1,2$. Those three numbers are
\[
  c_{2}+dc_{0}=-d+d=0,\qquad
  c_{3}+dc_{1}=-db+db=0,\qquad
  c_{4}+dc_{2}=d^{2}-d^{2}=0,
\]
so that
\[
  \sum_{i}w(p_{i})\,p_{i}^{\,k}=\sum_{j}w(q_{j})\,q_{j}^{\,k},
  \qquad k=0,1,2 .
\]
The two sides are the moments of order $0$, $1$ and $2$ of the positive
measures $\nu_{P}=\sum_{i}w(p_{i})\delta_{p_{i}}$ and
$\nu_{Q}=\sum_{j}w(q_{j})\delta_{q_{j}}$, so those measures have the same
mass, the same mean and the same variance. At $r=1$ the measure $\nu_{Q}$ has
one atom, hence variance zero, so $\nu_{P}$ has variance zero as well, and a
positive measure with two atoms has variance zero only if the atoms coincide.
So $p_{1}=p_{2}$, and then equality of mass and mean forces $q_{1}=p_{1}$. The
resolution \eqref{eq:burchres} then cancels and $\cI_{Z}\cong\cO(e_{1}\Theta)$
is a line bundle, which the ideal sheaf of a subscheme of codimension two is
not.

Equivalently and by the same computation, the two divisor degrees $a_{1}$ and
$a_{2}$ would have to satisfy $a_{1}+a_{2}=\tfrac43b$ and
$(b-a_{1})(b-a_{2})=\tfrac12\bigl(\tfrac13b^{2}+d\bigr)$, and the quadratic
with that sum and product has discriminant $-\tfrac29b^{2}-2d<0$: the degrees
are not even real.
\end{proof}

For $r\ge2$ the moment conditions no longer decide by themselves, and
\eqref{eq:burchsystem} becomes a finite search over the twists. Item (XXXII)
of \Cref{sec:verification} carries it out exactly, by meet in the middle on
the four moments, and the outcome is \Cref{tab:burch}.

\begin{theorem}[The split resolutions, and where they can live]%
\label{thm:burchsearch}
In \Cref{setup:burch} with $b=3$, over every squarefree $d<24$ and every
choice of twists in the box recorded in \Cref{tab:burch}:
\begin{enumerate}[label=\textup{(\roman*)},leftmargin=2.6em]
\item at $r=2$ and $r=3$ the system \eqref{eq:burchsystem} has no solution at
  all, and the same holds for every $b\le6$ over seven fields;
\item at $r=4,5,6$ it has solutions, but only for $d=15$ and $d=23$;
\item every $d$ that survives at any rank exceeds $9$.
\end{enumerate}
\end{theorem}

\begin{proof}
The search is item (XXXII) of \Cref{sec:verification}, carried out in exact
integer arithmetic; the two conditions of \Cref{lem:burchmoments} on
$\max_{i}e_{i}$ and $\min_{i}e_{i}$ are imposed, and without the first the
system acquires solutions that are not resolutions, the smallest being
$p=(-8,-5,0,1,2)$, $q=(-7,-6,-4,4)$ at $r=4$ and $d=23$, where the summand of
$E_{1}$ with the largest twist admits no nonzero map to $E_{0}$.
\end{proof}

\begin{table}[htbp]
\centering
\small
\renewcommand{\arraystretch}{1.12}
\begin{tabular}{@{}clll@{}}
\toprule
$r$ & twists searched & $d$ surviving & first candidate, as twists of $\Theta$ \\
\midrule
$1$ & every $b$, every $d$ & none & \Cref{thm:noci} \\
$2$ & $[-18,18]$ & none & \\
$3$ & $[-14,14]$ & none & \\
$4$ & $[-11,11]$ & $23$ & $E_{0}:(-11,-6,0,2,7)$,\ \ $E_{1}:(-10,-9,5,6)$ \\
$5$ & $[-9,9]$ & $23$ & $E_{0}:(-8,-4,-3,-2,1,6)$,\ \ $E_{1}:(-7,-6,-6,4,5)$ \\
$6$ & $[-8,8]$ & $15,23$ & $E_{0}:(-8,-4,-4,-1,-1,1,4)$,\ \ $E_{1}:(-7,-6,-6,0,3,3)$ \\
\bottomrule
\end{tabular}
\caption{The split Hilbert--Burch resolutions of \Cref{setup:burch} at $b=3$,
over every squarefree $d<24$. Nothing survives below rank four, and what
survives above it lives only at $d=15$ and $d=23$. Both exceed the bound $d\le9$
of \Cref{cor:fourdiscriminants}, so none of these supports can be smooth.}
\label{tab:burch}
\end{table}


The search of \Cref{thm:burchsearch} records where a split resolution can
have the right Chern character. It is superseded by a statement that needs no
search and no bound on the twists: a split resolution never runs the
criterion, at any rank and any discriminant. The reason is a necessary
condition that is worth isolating, because it applies to every object and can
be tested on any candidate before any global computation is attempted.

\begin{lemma}[Morphisms to a line bundle are killed by $H^{2}(\cO)$]
\label{lem:annihilation}
Let $(X,t)$ be a polarised abelian variety of dimension $n\ge3$, let $F$ be
a perfect complex with $\ch(F)=au_{t}+bv_{t}$, $(a,b)\ne(0,0)$, which
satisfies the weakened criterion of \Cref{setup:criterion}, and let $L$ be a
line bundle with $c_{1}(L)\in\QQ\,t$. Then for every $m\in\ZZ$, every
morphism $p\colon F\to L[m]$, every morphism $q\colon L[m]\to F$, and every
$z\in H^{2}(\cO_{X})$,
\[
  z\cdot p=0\ \ \text{in }\Hom(F,L[m+2]),\qquad
  z\cdot q=0\ \ \text{in }\Hom(L[m],F[2]),
\]
where $z\cdot(\,\cdot\,)$ is the $H^{*}(\cO_{X})$-module structure on
morphisms in $D^{b}(X)$. In particular, if $F=I_{Z}\otimes L'$ for a closed
subscheme $Z$ and a line bundle $L'$ with $c_{1}(L')\in\QQ t$, then the
restriction $H^{n-2}(\cO_{X})\to H^{n-2}(\cO_{Z})$ is injective.
\end{lemma}

\begin{proof}
For a constant bivector $\pi$ the class
$x_{\pi}=(\tfrac d2\,\pi\lrcorner\,t^{2},0,\pi)$ preserves $\ch(F)$ by
\Cref{thm:factor}(ii), so $\operatorname{ev}_{F}(x_{\pi})=0$ by hypothesis.
The Hochschild action is by natural transformations $\id\to\id[2]$ of
$D^{b}(X)$, so for $p\colon F\to L[m]$,
$\operatorname{ev}_{L[m]}(x_{\pi})\circ p=p\circ\operatorname{ev}_{F}(x_{\pi})=0$,
and $\operatorname{ev}_{L[m]}=\operatorname{ev}_{L}[m]$. By \eqref{eq:evsplit},
$\operatorname{ev}_{L}(x_{\pi})=c_{\pi}\cdot\id_{L}$ with
$c_{\pi}=\tfrac12\,\pi\lrcorner\,(\lambda^{2}+d\,t^{2})\in H^{0,2}(X)$,
$\lambda=c_{1}(L)=e\,t$, that is $c_{\pi}=\tfrac{e^{2}+d}{2}\,\pi\lrcorner\,t^{2}$,
and $\pi\mapsto\pi\lrcorner\,t^{2}$ is an isomorphism from
$H^{0}(\bigwedge^{2}T_{X})$ onto $H^{0,2}(X)$, as in the proof of
\Cref{thm:factor}(iii). Hence $c_{\pi}$ runs over all of $H^{2}(\cO_{X})$ as
$\pi$ varies, and $(c_{\pi}\cdot\id_{L})\circ p=c_{\pi}\cdot p$. The
statement for $q$ is the same with the square
$\operatorname{ev}_{F}(x_{\pi})\circ q=q\circ\operatorname{ev}_{L[m]}(x_{\pi})$.
For the last assertion apply the first to the inclusion
$p\colon I_{Z}\otimes L'\to L'$: the sequence
$0\to I_{Z}\to\cO_{X}\to\cO_{Z}\to0$ gives an exact sequence
\[
  \Ext^{2}(\cO_{Z},\cO_{X})\longrightarrow H^{2}(\cO_{X})
  \longrightarrow\Ext^{2}(I_{Z},\cO_{X}),
\]
in which the second map is $z\mapsto z\cdot p$; so $z\cdot p=0$ for all $z$
means that the first map is surjective, and by Serre duality with $\omega_{X}$
trivial it is the dual of the restriction
$H^{n-2}(\cO_{X})\to H^{n-2}(\cO_{Z})$.
\end{proof}

\begin{theorem}[No two-term complex of powers of the polarisation runs the
criterion]\label{thm:nosplit}
Let $(X,t)$ be a polarised abelian variety of dimension $n\ge3$ and $M$ an
ample line bundle with $c_{1}(M)\in\QQ t$. Let $A$ and $B$ be direct sums of
tensor powers $M^{\otimes e}$, $e\in\ZZ$, let $\varphi\colon A\to B$ be any
morphism, and let $F$ be the cone of $\varphi$, or any shift of it. If $F\ne0$
and $\ch(F)=au_{t}+bv_{t}$ with $(a,b)\ne(0,0)$, then $F$ does not satisfy
the weakened criterion. In particular no object with a resolution
$0\to E_{1}\to E_{0}\to F\to0$ by sums of powers of $M$, whatever the rank,
does.
\end{theorem}

\begin{proof}
A pair of summands $M^{\otimes e}$ of $A$ and $M^{\otimes e}$ of $B$ joined
by a nonzero entry of $\varphi$, which is then a nonzero scalar, can be
cancelled without changing the cone up to isomorphism; do this until no such
pair remains. If $A=0$ or $B=0$ after cancelling, $F$ is a shifted sum of
line bundles, $\operatorname{ev}_{F}(x_{\pi})$ is diagonal with the nonzero
entries $c_{\pi}\cdot\id$ of the proof of \Cref{lem:annihilation}, and the
criterion fails as in \Cref{cor:noline}; so suppose both are nonzero and
consider the triangle $A\xrightarrow{\varphi}B\xrightarrow{\;q\;}F\to A[1]$.
Let $M^{\otimes b_{j}}$ be a summand of $B$, $\iota_{j}$ its inclusion into
$B$, and $q_{j}=q\circ\iota_{j}\colon M^{\otimes b_{j}}\to F$. Applying
$\Hom(M^{\otimes b_{j}},-)$ to the triangle gives the exact sequence
\[
  \Hom(M^{\otimes b_{j}},A[2])\xrightarrow{\;\varphi_{*}\;}
  \Hom(M^{\otimes b_{j}},B[2])\xrightarrow{\;q_{*}\;}
  \Hom(M^{\otimes b_{j}},F[2]),
\]
and $z\cdot q_{j}=q_{*}(z\cdot\iota_{j})$ for $z\in H^{2}(\cO_{X})$. Now
$\Hom(M^{\otimes b_{j}},B[2])=\bigoplus_{k}H^{2}(M^{\otimes(b_{k}-b_{j})})$,
and $H^{2}(M^{\otimes m})=0$ for $m\ne0$: for $m>0$ by Kodaira vanishing,
$\omega_{X}$ being trivial, and for $m<0$ because a negative line bundle on
an abelian variety has cohomology only in degree $n\ge3$. So $z\cdot\iota_{j}$
is the element with component $z$ in the summand $k=j$ and zero elsewhere. It
lies in the image of $\varphi_{*}$ only if there is
$\psi\in\Hom(M^{\otimes b_{j}},A[2])=\bigoplus_{i}H^{2}(M^{\otimes(a_{i}-b_{j})})$
with $(\varphi\circ\psi)_{j}=\sum_{i}\varphi_{ji}\psi_{i}=z$; but $\psi_{i}\ne0$
forces $a_{i}=b_{j}$, and for such $i$ the entry $\varphi_{ji}$ is a scalar,
zero after the cancellation. So $z\cdot q_{j}\ne0$ for every $z\ne0$, which
contradicts \Cref{lem:annihilation}. The last assertion is the case of a
resolution, where $F\cong\coker\varphi$ is the cone of $\varphi\colon E_{1}\to E_{0}$.
\end{proof}

The powers of the polarisation are not the only bundles on which the
compensated Poisson classes act by scalars. So does every bundle whose Atiyah
class is a multiple of the identity, and on an abelian variety these are the
semi-homogeneous bundles of Mukai \cite{Muk78}; the argument of
\Cref{thm:nosplit} goes through for them, with the cohomology of the fibre
product of two isogenies in place of the cohomology of a power of $\Theta$.
We state the result in a form that is self-contained.

\begin{definition}[Isogeny blocks]\label{def:block}
Let $(X,\Theta)$ be a polarised abelian variety. An \emph{isogeny block} is a
vector bundle $E=\pi_{*}L$, where $\pi\colon Y\to X$ is an isogeny and $L$ is
a line bundle on $Y$ with $c_{1}(L)=\nu\,\pi^{*}\Theta$ in $H^{2}(Y,\QQ)$,
$\nu\in\QQ$; its \emph{slope} is $\nu$, so that $c_{1}(E)=\nu\deg(\pi)\,\Theta$.
The block is \emph{simple} if $\tau_{g}^{*}L\not\cong L$ for every nonzero
$g\in\ker\pi$. The powers $\cO(e\Theta)$ are the blocks with $\pi=\id$.
\end{definition}

\begin{lemma}[Three properties of a block]\label{lem:blocks}
Let $E=\pi_{*}L$ be an isogeny block of slope $\nu$ on an abelian $n$-fold.
\begin{enumerate}[label=\textup{(\roman*)},nosep]
\item $\mathrm{At}(E)=\nu\Theta\cdot\id_{E}$; hence for a constant bivector
  $\pi$ the compensated class $x_{\pi}$ of \Cref{thm:factor}(ii) acts by
  $\operatorname{ev}_{E}(x_{\pi})=\tfrac{\nu^{2}+d}{2}\,(\pi\lrcorner\,\Theta^{2})\cdot\id_{E}$,
  and \Cref{lem:annihilation} holds with $E$ in place of the line bundle $L$.
\item For a second block $E'=\pi'_{*}L'$ of slope $\nu'$,
  $\Ext^{q}(E',E)=H^{q}(Z,M)$ with $Z=Y'\times_{X}Y$, a disjoint union of
  abelian varieties isogenous to $X$, and $M$ a line bundle whose class on
  each component is $(\nu-\nu')$ times the pullback of $\Theta$. If
  $\nu\ne\nu'$ and $n\ge3$ then $\Ext^{2}(E',E)=0$. If $E$ is simple then
  $\Ext^{2}(E,E)=H^{2}(\cO_{X})\cdot\id_{E}$.
\item A simple block is $\mu$-stable with respect to $\Theta$; hence a
  nonzero morphism between simple blocks of the same slope is an
  isomorphism.
\end{enumerate}
\end{lemma}

\begin{proof}
(i) Since $\pi$ is \'etale, $\pi^{*}\Omega^{1}_{X}=\Omega^{1}_{Y}$ and $\pi_{*}$
is exact, so the first jet sequence of $\pi_{*}L$ is the direct image of
that of $L$, and $\mathrm{At}(\pi_{*}L)$ is the image of $\mathrm{At}(L)$
under $\pi_{*}\colon\Ext^{1}_{Y}(L,L\otimes\pi^{*}\Omega^{1}_{X})\to
\Ext^{1}_{X}(\pi_{*}L,\pi_{*}L\otimes\Omega^{1}_{X})$. As $\mathrm{At}(L)=c_{1}(L)\cdot\id_{L}
=\nu\,\pi^{*}\Theta\cdot\id_{L}$, the projection formula gives
$\nu\Theta\cdot\id_{\pi_{*}L}$. Then \eqref{eq:evgeneral} with a scalar Atiyah
class gives the value of $\operatorname{ev}_{E}(x_{\pi})$, and the proof of
\Cref{lem:annihilation} used nothing else about $L$.

(ii) Adjunction for the finite \'etale $\pi'$, flat base change along the
cartesian square with vertices $Z$, $Y'$, $Y$, $X$, and the projection formula
give $\Ext^{q}_{X}(\pi'_{*}L',\pi_{*}L)=H^{q}(Y',L'^{-1}\otimes\pi'^{*}\pi_{*}L)
=H^{q}(Z,p'^{*}L'^{-1}\otimes p^{*}L)$, with $p,p'$ the two projections; the
class of $M=p'^{*}L'^{-1}\otimes p^{*}L$ is $(\nu-\nu')$ times the pullback of
$\Theta$ on every component, which is an abelian variety finite \'etale over
$Y'$. A line bundle on an abelian $n$-fold whose class is a positive
multiple of an ample class has cohomology only in degree $0$, and one whose
class is a negative multiple only in degree $n$, by Kodaira vanishing and
Serre duality with trivial canonical bundle; so $H^{2}(Z,M)=0$ when
$\nu\ne\nu'$ and $n\ge3$. When $E'=E$, the components of $Y\times_{X}Y$ are
the graphs of the translations by $g\in\ker\pi$, on which $M$ restricts to
$L^{-1}\otimes\tau_{g}^{*}L\in\Pic^{0}(Y)$; a nontrivial element of $\Pic^{0}$
has no cohomology, so $\Ext^{q}(E,E)=\bigoplus_{g}H^{q}(\cO_{Y})$ over the
$g$ with $\tau_{g}^{*}L\cong L$, which for a simple block is $g=0$ alone, and
$H^{q}(\cO_{Y})=H^{q}(\cO_{X})$ through $\pi$.

(iii) The pullback $\pi^{*}E=\bigoplus_{g\in\ker\pi}\tau_{g}^{*}L$ is a direct
sum of line bundles of one slope, hence $\mu$-semistable, and semistability
descends along the finite \'etale $\pi$; so $E$ is semistable. If $F\subset E$
is a saturated subsheaf of the same slope, then $\pi^{*}F\subset\pi^{*}E$ is
a saturated subsheaf of the same slope of a direct sum of pairwise
non-isomorphic line bundles, the $\tau_{g}^{*}L$ being distinct because the
block is simple, hence a partial sum $\bigoplus_{g\in S}\tau_{g}^{*}L$; it is
invariant under the translations by $\ker\pi$, which permute the summands
transitively, so $S$ is empty or everything and $F=0$ or $F=E$. A nonzero
morphism between stable bundles of the same slope is an isomorphism.
\end{proof}

\begin{theorem}[No two-term complex of simple blocks runs the criterion]
\label{thm:noblocks}
Let $(X,\Theta)$ be a polarised abelian variety of dimension $n\ge3$, let $A$
and $B$ be direct sums of simple isogeny blocks, let $\varphi\colon A\to B$
be any morphism and let $F$ be the cone of $\varphi$ or a shift of it. If
$F\ne0$ and $\ch(F)=au_{\Theta}+bv_{\Theta}$ with $(a,b)\ne(0,0)$, then $F$
does not satisfy the weakened criterion. In particular no object with a
resolution $0\to E_{1}\to E_{0}\to F\to0$ by sums of simple isogeny blocks
does, whatever the rank.
\end{theorem}

\begin{proof}
Cancel, by Gaussian elimination in the two-term complex, every pair of
isomorphic blocks of $A$ and $B$ joined by an entry of $\varphi$ that is an
isomorphism; this changes neither the cone nor the block structure of the
remaining terms. By \Cref{lem:blocks}(iii), every remaining entry between
blocks of the same slope is zero. The rest is the proof of
\Cref{thm:nosplit} with \Cref{lem:blocks} in place of the vanishing of
$H^{2}(M^{\otimes m})$: for a block $E'$ of $B$ with inclusion $\iota$ and
$q'=q\circ\iota\colon E'\to F$, the class $z\cdot\iota\in\Hom(E',B[2])$ has
component $z\cdot\id_{E'}\ne0$ in $\Ext^{2}(E',E')=H^{2}(\cO_{X})\cdot\id_{E'}$
by (ii); it lies in the image of $\varphi_{*}$ only if
$\sum_{i}\varphi_{ji}\psi_{i}=z\cdot\id_{E'}$ for some
$\psi_{i}\in\Ext^{2}(E',A_{i})$, and $\psi_{i}\ne0$ forces $A_{i}$ to have
the slope of $E'$ by (ii), for which $\varphi_{ji}=0$. So $z\cdot q'\ne0$,
against \Cref{lem:annihilation} in the form of (i). If $A$ or $B$ is empty
after cancelling, $F$ is a shifted sum of blocks, on each of which
$\operatorname{ev}(x_{\pi})$ is a nonzero scalar by (i).
\end{proof}

\begin{remark}[What the blocks are]\label{rem:blocks}
By Mukai \cite{Muk78} the simple isogeny blocks are exactly the simple
semi-homogeneous bundles, the bundles $E$ with $\tau_{x}^{*}E\cong E\otimes P$
for every $x$ and some $P\in\Pic^{0}(X)$, and every semi-homogeneous bundle
on a variety with $\NS(X)_{\QQ}=\QQ\Theta$ is an iterated extension of such
blocks and their twists by $\Pic^{0}$; we do not use this. What
\Cref{thm:noblocks} says is that the bundles in a Hilbert--Burch resolution
of a candidate cannot all have scalar Atiyah class. The obstruction is the
same one every time: the compensated Poisson classes act on a block by the
scalar $\tfrac{\nu^{2}+d}{2}\,\pi\lrcorner\,\Theta^{2}$, which would vanish only
for $\nu^{2}=-d$, and cones of maps between blocks inherit the scalar through
the summand inclusions because the $\Ext^{2}$ between blocks of different
slopes vanishes. A candidate must therefore be built from a bundle whose
Atiyah class is not a scalar, that is, one which is not projectively flat.
\end{remark}

\begin{corollary}[The split shape is closed, and what the two routes must be]
\label{cor:disjointroutes}
Let $(X,\Theta)$ be a principally polarised abelian fourfold.
\begin{enumerate}[label=\textup{(\roman*)},nosep]
\item No secant object of rank one whose support has a resolution
  \eqref{eq:burchres} by sums of powers of $\Theta$ satisfies the weakened
  criterion, for any $r$, any twists and any $d$; in particular every
  candidate of \Cref{tab:burch} is excluded, whether or not its matrix
  degenerates in codimension two.
\item If a secant object $\cI_{Z}(b\Theta)$ satisfies the criterion, then
  $Z$ is Cohen--Macaulay, $\cI_{Z}$ has projective dimension one, its
  Hilbert--Burch resolution does not split into powers of $\Theta$, nor into
  simple semi-homogeneous bundles, and the restriction
  $H^{2}(\cO_{X})\to H^{2}(\cO_{Z})$ is injective. If $Z$ is smooth then
  moreover $d\in\{1,3,5,7\}$.
\item In the range of \Cref{thm:burchsearch} the two shapes do not meet:
  a smooth support forces $d\le9$ and a split resolution forces $d\ge15$,
  so even before (i) no support could be both.
\end{enumerate}
\end{corollary}

\begin{proof}
(i) is \Cref{thm:nosplit} with $M=\cO_{X}(\Theta)$, the resolution
\eqref{eq:burchres} presenting $\cI_{Z}(b\Theta)$ as the cone of
$\varphi\otimes\cO(b\Theta)$. (ii): Cohen--Macaulay and projective
dimension one are \Cref{cor:cmdichotomy}(ii); non-splitting is (i) and
\Cref{thm:noblocks};
injectivity is the last assertion of \Cref{lem:annihilation}; the four
discriminants are \Cref{cor:fourdiscriminants}. (iii) is
\Cref{cor:fourdiscriminants} against \Cref{thm:burchsearch}(iii), the sets
$\{1,3,5,7\}$ and $\{15,23\}$ being disjoint.
\end{proof}

\begin{remark}[The route map]\label{rem:routemap}
Collecting what is now known about the first of the six statements of
\Cref{ssec:closuremap}, the position is this. The local part of the criterion
is finished: it is a condition on depth, nothing else
(\Cref{thm:cmvanishing}, \Cref{cor:cmdichotomy}). The support must therefore
be Cohen--Macaulay, its ideal has projective dimension one, and there are
exactly two shapes left.

The first is a smooth support. Then every invariant is forced, the
discriminant is one of four (\Cref{cor:fourdiscriminants}), and the resolution
cannot split, now for every $d$ and every rank (\Cref{thm:nosplit}). What is
needed is the existence of
such a surface on a principally polarised abelian fourfold, with a
Hilbert--Burch matrix over a bundle that is not a sum of line bundles, and
then the vanishing of the image of $A_{\theta}A_{\theta'}$ in
$H^{1}(S,N_{S/X})$. The first half is a question about surfaces with
prescribed Chern numbers, of the kind settled by construction or by
classification; the second is a question about one cohomology group.

The second is a singular Cohen--Macaulay support. Then the invariants are
still forced by \eqref{eq:smoothinvariants} with $c_{1}(N_{Z/X})$ in place of
$K_{Z}$, and by \Cref{cor:disjointroutes}(i) the Hilbert--Burch resolution
must not split into powers of $\Theta$ or into simple semi-homogeneous
bundles (\Cref{thm:noblocks}): the explicit candidates of \Cref{tab:burch}
are excluded, and what is needed is a resolution over a bundle that is not
projectively flat, together with the vanishing in
$H^{1}(Z,\cE xt^{1}(\cI_{Z},\cI_{Z}))$. In both shapes the restriction
$H^{2}(\cO_{X})\to H^{2}(\cO_{Z})$ must be injective
(\Cref{lem:annihilation}), which for a smooth $S$ says that the six classes
of $\bigwedge^{2}H^{1}(\cO_{X})$ remain independent in $H^{2}(\cO_{S})$.

Both are finite questions in the sense that they are questions about one
variety and one sheaf, not about a family or a limit. Neither is settled here,
and settling both would give $\mathbf{W}(K,4,\delta)$ for the discriminants
reached and nothing beyond them.
\end{remark}
